% Generated mechanically from the frozen coherent.tex target.
% Do not edit this generated file; edit the bound locale source instead.

% OK, start here.
%

\setcounter{chapter}{29}
\chapter{概形的上同调}



\phantomsection
\label{coherent-section-phantom}


\section{引言}
\label{coherent-section-introduction}

\noindent
本章首先证明关于准凝聚层上同调的一系列结果。
一个基础参考文献是 \cite{EGA}。
在此基础上，我们将进一步讨论 Noether 情形下凝聚层的上同调。
见 \cite{FAC}。







\section{准凝聚层的 {\v C}ech 上同调}
\label{coherent-section-cech-quasi-coherent}

\noindent
设 $X$ 为概形。
设 $U \subset X$ 为仿射开集。
回忆，$U$ 的一个{\it 标准开覆盖}是形如
$\mathcal{U} : U = \bigcup_{i = 1}^n D(f_i)$ 的覆盖，其中
$f_1, \ldots, f_n \in \Gamma(U, \mathcal{O}_X)$ 生成单位理想；见
《概形》定义 \ref{schemes-definition-standard-covering}。

\begin{lemma}
\label{coherent-lemma-cech-cohomology-quasi-coherent-trivial}
设 $X$ 为概形。
设 $\mathcal{F}$ 为准凝聚 $\mathcal{O}_X$-模。
设 $\mathcal{U} : U = \bigcup_{i = 1}^n D(f_i)$ 为 $X$ 的一个仿射开集的
标准开覆盖。
则 $\check{H}^p(\mathcal{U}, \mathcal{F}) = 0$ 对所有 $p > 0$ 均成立。
\end{lemma}


\begin{proof}
写成 $U = \Spec(A)$，其中 $A$ 为某个环。
换言之，$f_1, \ldots, f_n$ 是 $A$ 中生成 $A$ 的单位理想的元素。
写成 $\mathcal{F}|_U = \widetilde{M}$，其中这个 $A$-模记作 $M$。
显然，{\v C}ech 复形
$\check{\mathcal{C}}^\bullet(\mathcal{U}, \mathcal{F})$
可与下列复形等同：
$$
\prod\nolimits_{i_0} M_{f_{i_0}} \to
\prod\nolimits_{i_0i_1} M_{f_{i_0}f_{i_1}} \to
\prod\nolimits_{i_0i_1i_2} M_{f_{i_0}f_{i_1}f_{i_2}} \to
\ldots
$$
我们需要证明扩充复形
\begin{equation}
\label{coherent-equation-extended}
0 \to
M \to
\prod\nolimits_{i_0} M_{f_{i_0}} \to
\prod\nolimits_{i_0i_1} M_{f_{i_0}f_{i_1}} \to
\prod\nolimits_{i_0i_1i_2} M_{f_{i_0}f_{i_1}f_{i_2}} \to
\ldots
\end{equation}

正合；我们已在《代数》引理 \ref{algebra-lemma-cover-module} 中研究过
它的截断。
只需证明复形 (\ref{coherent-equation-extended}) 在任意素理想 $\mathfrak p$
处局部化后正合；见《代数》引理
\ref{algebra-lemma-characterize-zero-local}。
事实上，我们将证明该扩充复形在 $\mathfrak p$ 处局部化后零伦。

\medskip\noindent
存在指标 $i$，使得 $f_i \not \in \mathfrak p$。
选定这样的指标，并记为 $i_{\text{fix}}$。注意
$M_{f_{i_{\text{fix}}}, \mathfrak p} = M_{\mathfrak p}$。类似地，
对乘积 $f_{i_0} \ldots f_{i_p}$ 再在 $\mathfrak p$ 处局部化时，
其中任一因子 $f_{i_j}$，只要 $i_j = i_{\text{fix}}$，均可略去。
定义同伦
$$
h :
\prod\nolimits_{i_0 \ldots i_{p + 1}}
M_{f_{i_0} \ldots f_{i_{p + 1}}, \mathfrak p}
\longrightarrow
\prod\nolimits_{i_0 \ldots i_p}
M_{f_{i_0} \ldots f_{i_p}, \mathfrak p}
$$
如下：
$$
h(s)_{i_0 \ldots i_p} = s_{i_{\text{fix}} i_0 \ldots i_p}
$$
（这与《上同调》引理
\ref{cohomology-lemma-homology-complex} 证明中的同伦“对偶”。）
换言之，$h : \prod_{i_0} M_{f_{i_0}, \mathfrak p} \to M_\mathfrak p$
是到因子
$M_{f_{i_{\text{fix}}}, \mathfrak p} = M_{\mathfrak p}$ 的投影；一般地，
映射 $h$ 是到诸因子
$M_{f_{i_{\text{fix}}} f_{i_1} \ldots f_{i_{p + 1}}, \mathfrak p}
= M_{f_{i_1} \ldots f_{i_{p + 1}}, \mathfrak p}$ 的投影。计算得
\begin{align*}
(dh + hd)(s)_{i_0 \ldots i_p}
& =
\sum\nolimits_{j = 0}^p
(-1)^j
h(s)_{i_0 \ldots \hat i_j \ldots i_p}
+
d(s)_{i_{\text{fix}} i_0 \ldots i_p}\\
& =
\sum\nolimits_{j = 0}^p
(-1)^j
s_{i_{\text{fix}} i_0 \ldots \hat i_j \ldots i_p}
+
s_{i_0 \ldots i_p}
+
\sum\nolimits_{j = 0}^p
(-1)^{j + 1}
s_{i_{\text{fix}} i_0 \ldots \hat i_j \ldots i_p} \\
& =
s_{i_0 \ldots i_p}
\end{align*}
这就按要求证明了恒等映射零伦。
\end{proof}


\noindent
下一个引理特别表明：对任意仿射概形 $X$ 以及任意准凝聚层
$\mathcal{F}$（定义在 $X$ 上），均有
$$
H^p(X, \mathcal{F}) = 0
$$
其中 $p > 0$。

\begin{lemma}
\label{coherent-lemma-quasi-coherent-affine-cohomology-zero}
\begin{slogan}
Serre 消失：仿射概形上准凝聚模的高阶上同调消失。
\end{slogan}
设 $X$ 为概形。
设 $\mathcal{F}$ 为准凝聚 $\mathcal{O}_X$-模。
对任意仿射开集 $U \subset X$，均有
$H^p(U, \mathcal{F}) = 0$，其中 $p > 0$。
\end{lemma}

\begin{proof}
我们应用《上同调》引理
\ref{cohomology-lemma-cech-vanish-basis}。
取拓扑基 $\mathcal{B}$；这里考虑 $X$ 的拓扑，并使用 $X$ 的全部仿射开集。
取开覆盖集 $\text{Cov}$ 为 $X$ 的仿射开集的全部标准开覆盖。
下面验证《上同调》引理
\ref{cohomology-lemma-cech-vanish-basis} 的条件 (1)、(2) 和 (3)。
注意，仿射概形中两个标准开集的交仍为标准开集，故性质 (1) 成立。
这些覆盖构成 $\mathcal{B}$ 中任一元素的开覆盖的共尾系；见
《概形》引理 \ref{schemes-lemma-standard-open}。
故 (2) 成立。
最后，条件 (3) 来自引理
\ref{coherent-lemma-cech-cohomology-quasi-coherent-trivial}。
\end{proof}

\noindent
下面给出仿射概形上准凝聚层上同调消失的相对版本。

\begin{lemma}
\label{coherent-lemma-relative-affine-vanishing}
设 $f : X \to S$ 为概形态射。
设 $\mathcal{F}$ 为准凝聚 $\mathcal{O}_X$-模。
若 $f$ 仿射，则 $R^if_*\mathcal{F} = 0$ 对所有 $i > 0$ 均成立。
\end{lemma}

\begin{proof}
按照《上同调》引理
\ref{cohomology-lemma-describe-higher-direct-images}，层
$R^if_*\mathcal{F}$ 是预层
$V \mapsto H^i(f^{-1}(V), \mathcal{F}|_{f^{-1}(V)})$ 的层化。
由假设，只要 $V$ 仿射，$f^{-1}(V)$ 就仿射；见《态射》定义
\ref{morphisms-definition-affine}。
由引理 \ref{coherent-lemma-quasi-coherent-affine-cohomology-zero} 可知，
$H^i(f^{-1}(V), \mathcal{F}|_{f^{-1}(V)}) = 0$
只要 $V$ 仿射就成立。
由于 $S$ 有一个由仿射开集构成的基，结论成立。
\end{proof}

\begin{lemma}
\label{coherent-lemma-relative-affine-cohomology}
设 $f : X \to S$ 为概形的仿射态射。
设 $\mathcal{F}$ 为准凝聚 $\mathcal{O}_X$-模。
则 $H^i(X, \mathcal{F}) = H^i(S, f_*\mathcal{F})$
对所有 $i \geq 0$ 均成立。
\end{lemma}

\begin{proof}
这来自引理 \ref{coherent-lemma-relative-affine-vanishing} 以及 Leray 谱序列。
见《上同调》引理 \ref{cohomology-lemma-apply-Leray}。
\end{proof}

\noindent
下面两个引理说明何时可用 {\v C}ech 上同调计算准凝聚模的上同调。

\begin{lemma}
\label{coherent-lemma-affine-diagonal}
设 $X$ 为概形。下列条件等价：
\begin{enumerate}
\item $X$ 的对角态射 $\Delta : X \to X \times X$ 仿射；
\item 对任意仿射开集 $U, V \subset X$，交 $U \cap V$ 仿射；以及
\item 存在开覆盖 $\mathcal{U} : X = \bigcup_{i \in I} U_i$，使得
$U_{i_0 \ldots i_p}$ 对所有 $p \ge 0$ 以及所有
$i_0, \ldots, i_p \in I$ 都是仿射开集。
\end{enumerate}
特别地，若 $X$ 分离，则这些条件成立。
\end{lemma}

\begin{proof}
假设 $X$ 的对角态射仿射。设 $U, V \subset X$ 为仿射开集。
则 $U \cap V = \Delta^{-1}(U \times V)$ 仿射，故 (2) 成立。
(2) 蕴含 (3) 是立即的。反之，若 $X$ 有如 (3) 所述的覆盖，则
$X \times X = \bigcup U_i \times U_{i'}$ 是仿射开覆盖，并且
$\Delta^{-1}(U_i \times U_{i'}) = U_i \cap U_{i'}$
仿射。于是由《态射》引理
\ref{morphisms-lemma-characterize-affine}，$\Delta$ 是仿射态射。
最后一个断言来自《概形》引理
\ref{schemes-lemma-characterize-separated}。
\end{proof}

\begin{lemma}
\label{coherent-lemma-cech-cohomology-quasi-coherent}
设 $X$ 为概形。
设 $\mathcal{U} : X = \bigcup_{i \in I} U_i$ 为开覆盖，并且
$U_{i_0 \ldots i_p}$ 对所有 $p \ge 0$ 以及所有
$i_0, \ldots, i_p \in I$ 都是仿射开集。
在这种情况下，对任意准凝聚层 $\mathcal{F}$，均有
$$
\check{H}^p(\mathcal{U}, \mathcal{F}) = H^p(X, \mathcal{F})
$$
作为 $\Gamma(X, \mathcal{O}_X)$-模的等式，并对所有 $p$ 成立。
\end{lemma}

\begin{proof}
由引理 \ref{coherent-lemma-quasi-coherent-affine-cohomology-zero}，这是
《上同调》引理
\ref{cohomology-lemma-cech-spectral-sequence-application} 的特殊情形。
\end{proof}







\section{上同调的消失}
\label{coherent-section-vanishing}

\noindent
我们已经看到，在仿射概形上，任意准凝聚层的高阶上同调群都消失
（引理 \ref{coherent-lemma-quasi-coherent-affine-cohomology-zero}）。
事实证明，这一性质也刻画仿射概形。我们给出两个版本。

\begin{lemma}
\label{coherent-lemma-quasi-compact-h1-zero-covering}
\begin{reference}
\cite{Serre-criterion}, \cite[II, Theorem 5.2.1 (d') and IV (1.7.17)]{EGA}
\end{reference}
\begin{slogan}
Serre 仿射性判别准则。
\end{slogan}
设 $X$ 为概形。假设：
\begin{enumerate}
\item $X$ 拟紧；
\item 对每个准凝聚理想层
$\mathcal{I} \subset \mathcal{O}_X$，均有 $H^1(X, \mathcal{I}) = 0$。
\end{enumerate}
则 $X$ 仿射。
\end{lemma}

\begin{proof}
设 $x \in X$ 为闭点。设 $U \subset X$ 为 $x$ 的一个仿射开邻域。
写成 $U = \Spec(A)$，并令 $\mathfrak m \subset A$ 为对应于 $x$ 的极大理想。
置 $Z = X \setminus U$ 以及 $Z' = Z \cup \{x\}$。
由《概形》引理 \ref{schemes-lemma-reduced-closed-subscheme}，存在准凝聚理想层
$\mathcal{I}$ 和 $\mathcal{I}'$，它们分别截出约化闭子概形
$Z$ 和 $Z'$。
考虑短正合列
$$
0 \to \mathcal{I}' \to \mathcal{I} \to \mathcal{I}/\mathcal{I}' \to 0.
$$
由于 $x$ 是 $X$ 的闭点且 $x \not \in Z$，
$\mathcal{I}/\mathcal{I}'$ 支持在 $x$ 上。事实上，
$\mathcal{I}/\mathcal{I'}$ 到 $U$ 的限制对应于 $A$-模
$A/\mathfrak m$。因此
$\Gamma(X, \mathcal{I}/\mathcal{I'})
= A/\mathfrak m$。
由于依假设 $H^1(X, \mathcal{I}') = 0$，存在整体截面
$f \in \Gamma(X, \mathcal{I})$，它映到元素 $1 \in A/\mathfrak m$；
这里把它视为 $\mathcal{I}/\mathcal{I'}$ 的截面。
显然 $x \in X_f \subset U$。这意味着
$X_f = D(f_A)$，其中 $f_A$ 是 $f$ 在
$A = \Gamma(U, \mathcal{O}_X)$ 中的像。特别地，$X_f$ 仿射。

\medskip\noindent
考虑并集 $W = \bigcup X_f$，其中遍历所有
$f \in \Gamma(X, \mathcal{O}_X)$ 且 $X_f$ 仿射。显然 $W$ 是 $X$ 的开集。
由上述论证，$X$ 的每个闭点都包含在 $W$ 中。
闭子集 $X \setminus W$（它是 $X$ 中的闭子集）也拟紧
（见《拓扑学》引理 \ref{topology-lemma-closed-in-quasi-compact}）。
因此，如果它非空，就有闭点（见《拓扑学》引理
\ref{topology-lemma-quasi-compact-closed-point}）。
这将与所有闭点都在 $W$ 中矛盾。故 $X = W$。

\medskip\noindent
选择有限多个 $f_1, \ldots, f_n \in \Gamma(X, \mathcal{O}_X)$，
使得 $X = X_{f_1} \cup \ldots \cup X_{f_n}$，并且每个
$X_{f_i}$ 都仿射。由上述论证，这样的选择是可能的。
由《性质》引理 \ref{properties-lemma-characterize-affine}，
为完成证明，只需证明 $f_1, \ldots, f_n$ 在
$\Gamma(X, \mathcal{O}_X)$ 中生成单位理想。考虑短正合列
$$
\xymatrix{
0 \ar[r] &
\mathcal{F} \ar[r] &
\mathcal{O}_X^{\oplus n} \ar[rr]^{f_1, \ldots, f_n} & &
\mathcal{O}_X \ar[r] &
0
}
$$
由 $f_1, \ldots, f_n$ 定义的箭头满射，因为诸开集
$X_{f_i}$ 覆盖 $X$。令 $\mathcal{F}$ 为该满射的核。
注意，$\mathcal{F}$ 有滤过
$$
0 = \mathcal{F}_0 \subset \mathcal{F}_1 \subset
\ldots \subset \mathcal{F}_n = \mathcal{F}
$$
使得每个子商 $\mathcal{F}_i/\mathcal{F}_{i - 1}$
都同构于一个准凝聚理想层。
具体而言，可取 $\mathcal{F}_i$ 为 $\mathcal{F}$ 与前 $i$ 个直和因子的交；
这些因子属于 $\mathcal{O}_X^{\oplus n}$。
引理的假设蕴含
$H^1(X, \mathcal{F}_i/\mathcal{F}_{i - 1}) = 0$
对所有 $i$ 都成立。于是 $H^1(X, \mathcal{F}_2) = 0$，
因为它夹在 $H^1(X, \mathcal{F}_1)$ 与
$H^1(X, \mathcal{F}_2/\mathcal{F}_1)$ 之间。
如此继续，可得 $H^1(X, \mathcal{F}) = 0$。
因此，映射
$$
\xymatrix{
\bigoplus\nolimits_{i = 1, \ldots, n} \Gamma(X, \mathcal{O}_X)
\ar[rr]^{f_1, \ldots, f_n} & &
\Gamma(X, \mathcal{O}_X)
}
$$
如所需为满射。
\end{proof}

\noindent
注意，若 $X$ 是 Noether 概形，则每个准凝聚理想层自动是凝聚理想层，
也是有限型准凝聚理想层。因此在这种情况下，前一引理与下一引理均适用。

\begin{lemma}
\label{coherent-lemma-quasi-separated-h1-zero-covering}
\begin{reference}
\cite{Serre-criterion}, \cite[II, Theorem 5.2.1]{EGA}
\end{reference}
\begin{slogan}
Serre 仿射性判别准则。
\end{slogan}
设 $X$ 为概形。假设：
\begin{enumerate}
\item $X$ 拟紧；
\item $X$ 拟分离；并且
\item $H^1(X, \mathcal{I}) = 0$ 对每个有限型准凝聚理想层
$\mathcal{I}$ 均成立。
\end{enumerate}
则 $X$ 仿射。
\end{lemma}

\begin{proof}
由《性质》引理
\ref{properties-lemma-quasi-coherent-colimit-finite-type}，
每个准凝聚理想层都是有限型准凝聚理想层的有向余极限。
由《上同调》引理
\ref{cohomology-lemma-quasi-separated-cohomology-colimit}，
在 $X$ 上取上同调与有向余极限可交换。
因此，$H^1(X, \mathcal{I}) = 0$ 对 $X$ 上的每个准凝聚理想层均成立。
换言之，引理
\ref{coherent-lemma-quasi-compact-h1-zero-covering} 适用。
\end{proof}

\noindent
可以利用上述论证给出一个判定可逆层何时丰沛的充分条件。
不过我们提醒读者，这个条件并非必要。

\begin{lemma}
\label{coherent-lemma-quasi-compact-h1-zero-invertible}
设 $X$ 为概形。设 $\mathcal{L}$ 为可逆 $\mathcal{O}_X$-模。
假设：
\begin{enumerate}
\item $X$ 拟紧；
\item 对每个准凝聚理想层
$\mathcal{I} \subset \mathcal{O}_X$，存在 $n \geq 1$，使得
$H^1(X, \mathcal{I} \otimes_{\mathcal{O}_X} \mathcal{L}^{\otimes n}) = 0$。
\end{enumerate}
则 $\mathcal{L}$ 丰沛。
\end{lemma}

\begin{proof}
证明与引理 \ref{coherent-lemma-quasi-compact-h1-zero-covering} 完全相同。
设 $x \in X$ 为闭点。设 $U \subset X$ 为 $x$ 的仿射开邻域，并且
$\mathcal{L}|_U \cong \mathcal{O}_U$。
写成 $U = \Spec(A)$，并令 $\mathfrak m \subset A$
为对应于 $x$ 的极大理想。
置 $Z = X \setminus U$ 以及 $Z' = Z \cup \{x\}$。
由《概形》引理 \ref{schemes-lemma-reduced-closed-subscheme}，存在准凝聚理想层
$\mathcal{I}$ 和 $\mathcal{I}'$，它们分别截出约化闭子概形
$Z$ 和 $Z'$。考虑短正合列
$$
0 \to \mathcal{I}' \to \mathcal{I} \to \mathcal{I}/\mathcal{I}' \to 0.
$$
对每个 $n \geq 1$，得到短正合列
$$
0 \to \mathcal{I}' \otimes_{\mathcal{O}_X} \mathcal{L}^{\otimes n}
\to \mathcal{I} \otimes_{\mathcal{O}_X} \mathcal{L}^{\otimes n} \to
\mathcal{I}/\mathcal{I}' \otimes_{\mathcal{O}_X} \mathcal{L}^{\otimes n} \to 0.
$$
由假设，可选择 $n$，使得
$H^1(X, \mathcal{I}' \otimes_{\mathcal{O}_X} \mathcal{L}^{\otimes n}) = 0$。
由于 $x$ 是 $X$ 的闭点且 $x \not \in Z$，
$\mathcal{I}/\mathcal{I}'$ 支持在 $x$ 上。事实上，
$\mathcal{I}/\mathcal{I'}$ 到 $U$ 的限制对应于 $A$-模
$A/\mathfrak m$。由于 $\mathcal{L}$ 在 $U$ 上平凡，
$\mathcal{I}/\mathcal{I}' \otimes_{\mathcal{O}_X} \mathcal{L}^{\otimes n}$
到 $U$ 的限制也对应于 $A$-模 $A/\mathfrak m$。
因此
$\Gamma(X, \mathcal{I}/\mathcal{I'} \otimes_{\mathcal{O}_X}
\mathcal{L}^{\otimes n}) = A/\mathfrak m$。
根据所选的 $n$，存在整体截面
$s \in \Gamma(X, \mathcal{I} \otimes_{\mathcal{O}_X} \mathcal{L}^{\otimes n})$，
它映到元素 $1 \in A/\mathfrak m$。显然
$x \in X_s \subset U$，因为 $s$ 在 $Z$ 的点处消失。
这意味着 $X_s = D(f)$，其中 $f \in A$ 是 $s$ 在
$A \cong \Gamma(U, \mathcal{L}^{\otimes n})$ 中的像。
特别地，$X_s$ 仿射。

\medskip\noindent
考虑并集 $W = \bigcup X_s$，其中遍历所有
$s \in \Gamma(X, \mathcal{L}^{\otimes n})$，取 $n \geq 1$，并要求
$X_s$ 仿射。显然 $W$ 是 $X$ 的开集。
由上述论证，$X$ 的每个闭点都包含在 $W$ 中。
闭子集 $X \setminus W$（它是 $X$ 中的闭子集）也拟紧
（见《拓扑学》引理 \ref{topology-lemma-closed-in-quasi-compact}）。
因此，如果它非空，就有闭点（见《拓扑学》引理
\ref{topology-lemma-quasi-compact-closed-point}）。
这将与所有闭点都在 $W$ 中矛盾。故 $X = W$。
由《性质》定义 \ref{properties-definition-ample}，
这意味着 $\mathcal{L}$ 丰沛。
\end{proof}

\noindent
引理 \ref{coherent-lemma-quasi-compact-h1-zero-invertible} 有一个使用有限型理想层的
变体；若以后需要，我们将在此陈述并证明。

\begin{lemma}
\label{coherent-lemma-criterion-affine-morphism}
设 $f : X \to Y$ 为拟紧态射，并且 $X$ 和 $Y$ 均拟分离。
若 $R^1f_*\mathcal{I} = 0$ 对每个准凝聚理想层
$\mathcal{I}$（定义在 $X$ 上）均成立，则 $f$ 仿射。
\end{lemma}

\begin{proof}
设 $V \subset Y$ 为仿射开子概形。我们需要证明
$U = f^{-1}(V)$ 仿射。由《概形》引理
\ref{schemes-lemma-quasi-compact-permanence}，包含态射 $V \to Y$ 拟紧。
因此基变换 $U \to X$ 拟紧；见《概形》引理
\ref{schemes-lemma-quasi-compact-preserved-base-change}。
于是，任意准凝聚理想层 $\mathcal{I}$（定义在 $U$ 上）
都可延拓为 $X$ 上的准凝聚理想层；见《性质》引理
\ref{properties-lemma-extend-trivial}。
由于 $R^1f_*$ 的构造在 $Y$ 上是局部的
（《上同调》节 \ref{cohomology-section-locality}），
由引理中的假设可得 $R^1(U \to V)_*\mathcal{I} = 0$。
因此由 Leray 谱序列
（《上同调》引理 \ref{cohomology-lemma-Leray}），可得
$H^1(U, \mathcal{I}) = H^1(V, (U \to V)_*\mathcal{I})$。
由《概形》引理
\ref{schemes-lemma-push-forward-quasi-coherent}，
$(U \to V)_*\mathcal{I}$ 准凝聚，故由引理
\ref{coherent-lemma-quasi-coherent-affine-cohomology-zero}，
$H^1(V, (U \to V)_*\mathcal{I}) = 0$。
于是由引理 \ref{coherent-lemma-quasi-compact-h1-zero-covering}，
$U$ 仿射。
\end{proof}















\section{高阶直像的准凝聚性}
\label{coherent-section-quasi-coherence}

\noindent
我们已经看到，仿射概形上一个准凝聚模的高阶上同调群为零。
对拟分离拟紧概形 $X$，这蕴含准凝聚层的上同调群在某一次数以上消失。
然而，$X$ 未必具有有限上同调维数，因为上同调维数的定义要求
{\it 所有} $\mathcal{O}_X$-模的上同调都消失。

\begin{lemma}[归纳原理]
\label{coherent-lemma-induction-principle}
\begin{reference}
\cite[Proposition 3.3.1]{BvdB}
\end{reference}
设 $X$ 为拟紧拟分离概形。设 $P$ 是 $X$ 的拟紧开集所具有的一个性质。
假设：
\begin{enumerate}
\item $P$ 对 $X$ 的每个仿射开集成立；
\item 若 $U$ 是拟紧开集，$V$ 是仿射开集，且 $P$ 对
$U$、$V$ 和 $U \cap V$ 成立，则 $P$ 对 $U \cup V$ 成立。
\end{enumerate}
则 $P$ 对 $X$ 的每个拟紧开集成立，特别地对 $X$ 成立。
\end{lemma}

\begin{proof}
首先用归纳法证明 $P$ 对{\it 分离的}拟紧开集 $W \subset X$ 成立。
事实上，可将这样的开集写成
$W = U_1 \cup \ldots \cup U_n$，并对 $n$ 归纳；
在性质 (2) 中取 $U = U_1 \cup \ldots \cup U_{n - 1}$ 和 $V = U_n$。
这是允许的，因为
$U \cap V = (U_1 \cap U_n) \cup \ldots \cup (U_{n - 1} \cap U_n)$
也是 $n - 1$ 个仿射开子概形之并；这里把《概形》引理
\ref{schemes-lemma-characterize-separated} 应用于仿射开集 $U_i$ 和 $U_n$，
它们均为 $W$ 的仿射开集。
由此，对任意拟紧开集 $W \subset X$，可对覆盖 $W$ 所需的仿射开集数目
归纳，使用与前面相同的办法，并利用如下事实：拟紧开集
$U_i \cap U_n$ 作为仿射概形 $U_n$ 的开子概形是分离的。
\end{proof}

\begin{lemma}
\label{coherent-lemma-vanishing-nr-affines}
\begin{slogan}
对具有仿射对角态射的概形，准凝聚模的上同调在次数大于某个覆盖所需的
仿射开集数目时消失。
\end{slogan}
设 $X$ 为具有仿射对角态射的拟紧概形（例如，若 $X$ 分离）。
令 $t = t(X)$ 为覆盖 $X$ 所需的仿射开集的最小数目。
则均有 $H^n(X, \mathcal{F}) = 0$，其中 $n \geq t$ 任意，
$\mathcal{F}$ 为任意准凝聚层。
\end{lemma}

\begin{proof}
第一种证明。
对 $t$ 归纳。
若 $t = 1$，结论来自引理
\ref{coherent-lemma-quasi-coherent-affine-cohomology-zero}。
若 $t > 1$，写成 $X = U \cup V$，其中 $V$ 为仿射开集，而
$U = U_1 \cup \ldots \cup U_{t - 1}$ 是 $t - 1$ 个仿射开集之并。
注意，此时
$U \cap V =  (U_1 \cap V) \cup \ldots (U_{t - 1} \cap V)$
也是 $t - 1$ 个仿射开子概形之并。
事实上，因为对角态射仿射，两个仿射开集的交仿射；见引理
\ref{coherent-lemma-affine-diagonal}。
应用 Mayer--Vietoris 长正合列
$$
0 \to
H^0(X, \mathcal{F}) \to
H^0(U, \mathcal{F}) \oplus H^0(V, \mathcal{F}) \to
H^0(U \cap V, \mathcal{F}) \to
H^1(X, \mathcal{F}) \to \ldots
$$
见《上同调》引理 \ref{cohomology-lemma-mayer-vietoris}。
由归纳可知，群 $H^i(U, \mathcal{F})$、
$H^i(V, \mathcal{F})$ 和 $H^i(U \cap V, \mathcal{F})$
在 $i \geq t - 1$ 时为零。因此立即得到
$H^i(X, \mathcal{F})$ 在 $i \geq t$ 时为零。

\medskip\noindent
第二种证明。
设 $\mathcal{U} : X = \bigcup_{i = 1}^t U_i$ 为有限仿射开覆盖。
由于 $X$ 具有仿射对角态射，多重交
$U_{i_0 \ldots i_p}$ 都仿射；见引理 \ref{coherent-lemma-affine-diagonal}。
由引理 \ref{coherent-lemma-cech-cohomology-quasi-coherent}，{\v C}ech 上同调群
$\check{H}^p(\mathcal{U}, \mathcal{F})$ 与上同调群相同。
由《上同调》引理 \ref{cohomology-lemma-alternating-usual}，
可用交错 {\v C}ech 复形
$\check{\mathcal{C}}_{alt}^\bullet(\mathcal{U}, \mathcal{F})$
计算 {\v C}ech 上同调群。由于覆盖由 $t$ 个元素组成，立即得到
$\check{\mathcal{C}}_{alt}^p(\mathcal{U}, \mathcal{F}) = 0$
对所有 $p \geq t$ 成立。因此结论成立。
\end{proof}

\begin{lemma}
\label{coherent-lemma-affine-diagonal-universal-delta-functor}
设 $X$ 为具有仿射对角态射的拟紧概形
（例如，若 $X$ 分离）。则：
\begin{enumerate}
\item 给定准凝聚 $\mathcal{O}_X$-模 $\mathcal{F}$，存在嵌入
$\mathcal{F} \to \mathcal{F}'$，其中后者是准凝聚
$\mathcal{O}_X$-模，
使得 $H^p(X, \mathcal{F}') = 0$ 对所有 $p \geq 1$ 成立；并且
\item $\{H^n(X, -)\}_{n \geq 0}$ 是泛 $\delta$-函子，
其源为 $\QCoh(\mathcal{O}_X)$，靶为 $\textit{Ab}$。
\end{enumerate}
\end{lemma}

\begin{proof}
设 $X = \bigcup U_i$ 为有限仿射开覆盖。
置 $U = \coprod U_i$，并以 $j : U \to X$ 表示诱导给定开浸入
$U_i \to X$ 的态射。
由于 $U$ 是仿射概形，且 $X$ 有仿射对角态射，态射 $j$ 仿射；见《态射》引理
\ref{morphisms-lemma-affine-permanence}。
对每个 $\mathcal{O}_X$-模 $\mathcal{F}$，存在典范映射
$\mathcal{F} \to j_*j^*\mathcal{F}$。
在茎上检验可知该映射单射：若 $x \in U_i$，则有分解
$$
\mathcal{F}_x \to (j_*j^*\mathcal{F})_x
\to (j^*\mathcal{F})_{x'} = \mathcal{F}_x
$$
其中 $x' \in U$ 是把点 $x$ 看成 $U_i \subset U$ 的点。
现在，若 $\mathcal{F}$ 准凝聚，则 $j^*\mathcal{F}$
在仿射概形 $U$ 上准凝聚，故由引理
\ref{coherent-lemma-quasi-coherent-affine-cohomology-zero}，其高阶上同调消失。
又有 $H^p(X, j_*j^*\mathcal{F}) = 0$ 对 $p > 0$ 成立；
这是由引理 \ref{coherent-lemma-relative-affine-cohomology} 以及 $j$ 仿射得到的。
这证明 (1)。
最后，映射
$H^p(X, \mathcal{F}) \to H^p(X, j_*j^*\mathcal{F})$
为零，而 (2) 来自《同调代数》引理
\ref{homology-lemma-efface-implies-universal}。
\end{proof}

\begin{lemma}
\label{coherent-lemma-vanishing-nr-affines-quasi-separated}
设 $X$ 为拟紧拟分离概形。
设 $X = U_1 \cup \ldots \cup U_n$ 为开覆盖，其中每个 $U_i$
都拟紧且分离（例如仿射）。置
$$
d = \max\nolimits_{I \subset \{1, \ldots, n\}}
\left(|I| + t(\bigcap\nolimits_{i \in I} U_i) - 1\right)
$$
其中 $t(U)$ 是覆盖概形 $U$ 所需的仿射开集的最小数目。
则均有 $H^p(X, \mathcal{F}) = 0$，其中 $p \geq d$ 任意，
$\mathcal{F}$ 为任意准凝聚层。
\end{lemma}

\begin{proof}
注意，由于 $X$ 拟分离且 $U_i$ 拟紧，数
$t(\bigcap_{i \in I} U_i)$ 有限。对 $n$ 归纳证明。
若 $n = 1$，结论来自引理 \ref{coherent-lemma-vanishing-nr-affines}。
若 $n > 1$，写成 $X = U \cup V$，其中
$U = U_1 \cup \ldots \cup U_{n - 1}$ 且 $V = U_n$。
应用 Mayer--Vietoris 长正合列
$$
0 \to
H^0(X, \mathcal{F}) \to
H^0(U, \mathcal{F}) \oplus H^0(V, \mathcal{F}) \to
H^0(U \cap V, \mathcal{F}) \to
H^1(X, \mathcal{F}) \to \ldots
$$
见《上同调》引理 \ref{cohomology-lemma-mayer-vietoris}。
为对 $q \geq d$ 完成证明，我们将证明
$H^q(V, \mathcal{F})$、$H^q(U, \mathcal{F})$ 和
$H^{q - 1}(U \cap V, \mathcal{F})$ 消失。
由 $n = 1$ 的情形，$H^q(V, \mathcal{F}) = 0$ 对
$q \geq t(V) = t(U_n)$ 成立。由于 $t(V) \leq d$，这就证明了所需断言。
由归纳假设，$H^q(U, \mathcal{F}) = 0$ 在下列范围成立：
$$
q \geq \max\nolimits_{I \subset \{1, \ldots, n - 1\}}
\left(|I| + t(\bigcap\nolimits_{i \in I} U_i) - 1\right)
$$
由于右端整数小于或等于 $d$，这就证明了所需断言。
最后，可对开集
$U \cap V = (U_1 \cap U_n) \cup \ldots \cup (U_{n - 1} \cap U_n)$
使用归纳假设，得到 $H^q(U \cap V, \mathcal{F}) = 0$ 在下列范围成立：
$$
q \geq \max\nolimits_{I \subset \{1, \ldots, n - 1\}}
\left(|I| + t(U_n \cap \bigcap\nolimits_{i \in I} U_i) - 1\right)
$$
由于右端整数严格小于 $d$，引理得证。
\end{proof}

\begin{lemma}
\label{coherent-lemma-quasi-coherence-higher-direct-images}
设 $f : X \to S$ 为概形态射。
假设 $f$ 拟分离且拟紧。
\begin{enumerate}
\item 对任意准凝聚 $\mathcal{O}_X$-模 $\mathcal{F}$，
高阶直像 $R^pf_*\mathcal{F}$ 在 $S$ 上准凝聚。
\item 若 $S$ 拟紧，则存在整数 $n = n(X, S, f)$，使得
$R^pf_*\mathcal{F} = 0$；这里 $p \geq n$ 任意，$\mathcal{F}$
为定义在 $X$ 上的任意准凝聚层。
\item 事实上，若 $S$ 拟紧，则可找到 $n = n(X, S, f)$，使得对每个
概形态射 $S' \to S$，都有 $R^p(f')_*\mathcal{F}' = 0$
对 $p \geq n$ 以及任意准凝聚层 $\mathcal{F}'$ 成立，后者定义在 $X'$ 上。
这里 $f' : X' = S' \times_S X \to S'$ 是 $f$ 的基变换。
\end{enumerate}
\end{lemma}

\begin{proof}
先证明 (1)。注意，在引理的假设下，层
$R^0f_*\mathcal{F} = f_*\mathcal{F}$ 由《概形》引理
\ref{schemes-lemma-push-forward-quasi-coherent} 是准凝聚的。
利用《上同调》引理
\ref{cohomology-lemma-localize-higher-direct-images}，
可知形成高阶直像与限制到开子概形可交换。
由于准凝聚性在 $S$ 上是局部性质，问题归约为下一段所讨论的情形。

\medskip\noindent
在 $S$ 仿射的情形下证明 (1)。我们使用归纳原理。
由于 $f$ 拟紧且拟分离，$X$ 拟紧且拟分离。
对拟紧开集 $U \subset X$，令 $a = f|_U$，并令 $P(U)$ 表示如下性质：
$R^pa_*\mathcal{F}$ 在 $S$ 上准凝聚；这对所有准凝聚模
$\mathcal{F}$（定义在 $U$ 上）以及所有 $p \geq 0$ 成立。
由于 $P(X)$ 就是 (1)，只需证明引理
\ref{coherent-lemma-induction-principle} 的条件 (1) 和 (2) 成立。
若 $U$ 仿射，则 $P(U)$ 成立，因为
$R^pa_*\mathcal{F} = 0$ 对 $p \geq 1$ 成立
（由引理 \ref{coherent-lemma-relative-affine-vanishing} 以及《态射》引理
\ref{morphisms-lemma-morphism-affines-affine}），并且我们在第一段已经观察到
$p = 0$ 时结论成立。
接下来，设 $U \subset X$ 为拟紧开集，$V \subset X$ 为仿射开集，并假设
$P(U)$、$P(V)$ 和 $P(U \cap V)$ 成立。
令 $a = f|_U$、$b = f|_V$、$c = f|_{U \cap V}$ 以及
$g = f|_{U \cup V}$。则对任意准凝聚
$\mathcal{O}_{U \cup V}$-模 $\mathcal{F}$，有相对
Mayer--Vietoris 序列
$$
0 \to
g_*\mathcal{F} \to
a_*(\mathcal{F}|_U) \oplus b_*(\mathcal{F}|_V) \to
c_*(\mathcal{F}|_{U \cap V}) \to
R^1g_*\mathcal{F} \to \ldots
$$
见《上同调》引理 \ref{cohomology-lemma-relative-mayer-vietoris}。
由 $P(U)$、$P(V)$ 和 $P(U \cap V)$ 可知，
$R^pa_*(\mathcal{F}|_U)$、$R^pb_*(\mathcal{F}|_V)$ 以及
$R^pc_*(\mathcal{F}|_{U \cap V})$ 都准凝聚。
使用《概形》节 \ref{schemes-section-quasi-coherent} 中关于准凝聚层的结果，
可知每个层 $R^pg_*\mathcal{F}$ 都准凝聚，因为它位于一个短正合列中间：
其左侧是准凝聚层之间一个映射的余核，右侧是准凝聚层之间一个映射的核。
故 $P(U \cup V)$ 成立，(1) 得证。

\medskip\noindent
下面证明 (3)，从而更证明 (2)。选取有限仿射开覆盖
$S = \bigcup_{j = 1, \ldots m} S_j$。对每个 $j$，选取有限仿射开覆盖
$f^{-1}(S_j) = \bigcup_{i = 1, \ldots t_j} U_{ji} $。
令
$$
d_j = \max\nolimits_{I \subset \{1, \ldots, t_j\}}
\left(|I| + t(\bigcap\nolimits_{i \in I} U_{ji})\right)
$$
为引理 \ref{coherent-lemma-vanishing-nr-affines-quasi-separated} 中所得到的整数。
我们断言 $n(X, S, f) = \max d_j$ 可用。

\medskip\noindent
事实上，设 $S' \to S$ 为概形态射，并设 $\mathcal{F}'$ 为
$X' = S' \times_S X$ 上的准凝聚层。
我们要证明 $R^pf'_*\mathcal{F}' = 0$ 对
$p \geq n(X, S, f)$ 成立。
由于该问题在 $S'$ 上是局部的，可假设 $S'$ 仿射，并映入某个
$S_j$，其中 $j$ 为相应指标。于是 $X' = S' \times_{S_j} f^{-1}(S_j)$
由仿射开集 $S' \times_{S_j} U_{ji}$（$i = 1, \ldots t_j$）覆盖，且诸交
$$
\bigcap\nolimits_{i \in I} S' \times_{S_j} U_{ji} =
S' \times_{S_j} \bigcap\nolimits_{i \in I} U_{ji}
$$
可由与基变换前相同数目的仿射开集覆盖。
应用引理 \ref{coherent-lemma-vanishing-nr-affines-quasi-separated}，得到
$H^p(X', \mathcal{F}') = 0$。
由证明的第一部分，已经知道每个 $R^qf'_*\mathcal{F}'$ 都准凝聚，
因而在仿射概形 $S'$ 上具有消失的高阶上同调群。
于是由《上同调》引理 \ref{cohomology-lemma-apply-Leray} 可得
$H^0(S', R^pf'_*\mathcal{F}') = H^p(X', \mathcal{F}') = 0$。
由于 $R^pf'_*\mathcal{F}'$ 准凝聚，结论是
$R^pf'_*\mathcal{F}' = 0$。
\end{proof}

\begin{lemma}
\label{coherent-lemma-quasi-coherence-higher-direct-images-application}
设 $f : X \to S$ 为概形态射。
假设 $f$ 拟分离且拟紧。
假设 $S$ 仿射。
对任意准凝聚 $\mathcal{O}_X$-模 $\mathcal{F}$，均有
$$
H^q(X, \mathcal{F}) = H^0(S, R^qf_*\mathcal{F})
$$
其中 $q \in \mathbf{Z}$。
\end{lemma}

\begin{proof}
考虑 Leray 谱序列 $E_2^{p, q} = H^p(S, R^qf_*\mathcal{F})$，
它收敛到 $H^{p + q}(X, \mathcal{F})$；见《上同调》引理
\ref{cohomology-lemma-Leray}。
由引理 \ref{coherent-lemma-quasi-coherence-higher-direct-images}，
层 $R^qf_*\mathcal{F}$ 准凝聚。
由引理 \ref{coherent-lemma-quasi-coherent-affine-cohomology-zero}，
有 $E_2^{p, q} = 0$，只要 $p > 0$。
因此谱序列在 $E_2$ 处退化，结论成立。
一般原理另见《上同调》引理
\ref{cohomology-lemma-apply-Leray} (2)。
\end{proof}


\section{上同调与基变换（一）}
\label{coherent-section-cohomology-and-base-change}

\noindent
设 $f : X \to S$ 为概形态射。
设 $\mathcal{F}$ 为 $X$ 上的准凝聚层。
再设 $g : S' \to S$ 为任意概形态射。记
$X' = X_{S'} = S' \times_S X$ 为 $X$ 的基变换，并记
$f' : X' \to S'$ 为 $f$ 的基变换。
还以 $g' : X' \to X$ 表示投影，并置
$\mathcal{F}' = (g')^*\mathcal{F}$。
下图表示这一情形：
\begin{equation}
\label{coherent-equation-base-change-diagram}
\vcenter{
\xymatrix{
\mathcal{F}' = (g')^*\mathcal{F} &
X' \ar[r]_{g'} \ar[d]_{f'} &
X \ar[d]^f &
\mathcal{F} \\
Rf'_*\mathcal{F}' &
S' \ar[r]^g &
S &
Rf_*\mathcal{F}
}
}
\end{equation}
下面是我们所考虑的基变换性质最简单的情形。

\begin{lemma}
\label{coherent-lemma-affine-base-change}
设 $f : X \to S$ 为概形态射。
设 $\mathcal{F}$ 为准凝聚 $\mathcal{O}_X$-模。
假设 $f$ 仿射。
此时 $f_*\mathcal{F} \cong Rf_*\mathcal{F}$ 是准凝聚层，并且对每个基变换图
(\ref{coherent-equation-base-change-diagram}) 都有
$$
g^*f_*\mathcal{F} = f'_*(g')^*\mathcal{F}.
$$
\end{lemma}

\begin{proof}
高阶直像的消失即引理 \ref{coherent-lemma-relative-affine-vanishing}。
该断言在 $S$ 和 $S'$ 上都是局部的。因此可假设
$X = \Spec(A)$、$S = \Spec(R)$、
$S' = \Spec(R')$ 且 $\mathcal{F} = \widetilde{M}$，
其中后者由某个 $A$-模 $M$ 给出。
我们用《概形》引理 \ref{schemes-lemma-widetilde-pullback}
描述 $\mathcal{F}$ 的拉回与直像。
具体地，$X' = \Spec(R' \otimes_R A)$，而
$\mathcal{F}'$ 是由 $(R' \otimes_R A) \otimes_A M$
所对应的准凝聚层。因此引理归结为 $R'$-模的等式
$$
(R' \otimes_R A) \otimes_A M = R' \otimes_R M
$$
。
\end{proof}

\noindent
在许多情形下，只需知道下面这个上同调与基变换的特殊情形。
它立即推出自节 \ref{coherent-section-cohomology-and-base-change-derived}
中的更强结果；但由于它如此重要，值得给出独立证明。

\begin{lemma}[平坦基变换]
\label{coherent-lemma-flat-base-change-cohomology}
考虑概形的笛卡尔图
$$
\xymatrix{
X' \ar[d]_{f'} \ar[r]_{g'} & X \ar[d]^f \\
S' \ar[r]^g & S
}
$$
设 $\mathcal{F}$ 为准凝聚 $\mathcal{O}_X$-模，其拉回为
$\mathcal{F}' = (g')^*\mathcal{F}$。
假设 $g$ 平坦，且 $f$ 拟紧拟分离。
对任意 $i \geq 0$，有：
\begin{enumerate}
\item 《上同调》引理 \ref{cohomology-lemma-base-change-map-flat-case}
中的基变换映射是同构
$$
g^*R^if_*\mathcal{F} \longrightarrow R^if'_*\mathcal{F}',
$$
\item 若 $S = \Spec(A)$ 且 $S' = \Spec(B)$，则
$H^i(X, \mathcal{F}) \otimes_A B = H^i(X', \mathcal{F}')$。
\end{enumerate}
\end{lemma}

\begin{proof}
在 (1) 中可以使用《上同调》引理
\ref{cohomology-lemma-base-change-map-flat-case}，因为由《态射》引理
\ref{morphisms-lemma-base-change-flat}，$g'$ 平坦。
说明这一点后，(1) 推出自 (2)。事实上，(1) 在 $S'$ 上是局部的，
故可假设 $S$ 和 $S'$ 都仿射。换言之，如 (2) 中那样，有
$S = \Spec(A)$ 和 $S' = \Spec(B)$。
于是，由于 $R^if_*\mathcal{F}$ 准凝聚
（引理 \ref{coherent-lemma-quasi-coherence-higher-direct-images}），
它是一个准凝聚 $\mathcal{O}_S$-模，该模与 $A$-模
$H^0(S, R^if_*\mathcal{F}) = H^i(X, \mathcal{F})$ 相对应
（等式来自引理
\ref{coherent-lemma-quasi-coherence-higher-direct-images-application}）。
类似地，$R^if'_*\mathcal{F}'$ 是一个准凝聚
$\mathcal{O}_{S'}$-模，该模与 $B$-模
$H^i(X', \mathcal{F}')$ 相对应。由于沿 $g$ 拉回在模上对应于
$- \otimes_A B$
（《概形》引理 \ref{schemes-lemma-widetilde-pullback}），
只需证明 (2)。

\medskip\noindent
设 $A \to B$ 为平坦环同态。
设 $X$ 为 $A$ 上的拟紧拟分离概形。
设 $\mathcal{F}$ 为准凝聚 $\mathcal{O}_X$-模。
置 $X_B = X \times_{\Spec(A)} \Spec(B)$，并以
$\mathcal{F}_B$ 表示 $\mathcal{F}$ 的拉回。
我们要证明映射
$$
H^i(X, \mathcal{F}) \otimes_A B \longrightarrow H^i(X_B, \mathcal{F}_B)
$$
（由引理陈述中的参考结果给出）是同构。

\medskip\noindent
若 $X$ 分离，选取仿射开覆盖
$\mathcal{U} : X = U_1 \cup \ldots \cup U_t$，并回忆
$$
\check{H}^p(\mathcal{U}, \mathcal{F}) = H^p(X, \mathcal{F}),
$$
见引理 \ref{coherent-lemma-cech-cohomology-quasi-coherent}。
若 $\mathcal{U}_B : X_B = (U_1)_B \cup \ldots \cup (U_t)_B$
是由基变换所得的覆盖，则每个 $(U_i)_B$ 仍仿射，且 $X_B$ 分离。
因此得到
$$
\check{H}^p(\mathcal{U}_B, \mathcal{F}_B) = H^p(X_B, \mathcal{F}_B).
$$
两个 {\v C}ech 复形之间有关系
$$
\check{\mathcal{C}}^\bullet(\mathcal{U}_B, \mathcal{F}_B) =
\check{\mathcal{C}}^\bullet(\mathcal{U}, \mathcal{F}) \otimes_A B
$$
这是引理 \ref{coherent-lemma-affine-base-change} 的推论。
由于 $A \to B$ 平坦，取上同调后同一结论仍成立。

\medskip\noindent
若 $X$ 拟分离，选取仿射开覆盖
$\mathcal{U} : X = U_1 \cup \ldots \cup U_t$。
我们将使用《上同调》引理
\ref{cohomology-lemma-cech-spectral-sequence} 的
{\v C}ech 到上同调谱序列。
希望避开这一谱序列的读者，可如引理
\ref{coherent-lemma-quasi-coherence-higher-direct-images} 的证明中那样，
使用 Mayer--Vietoris 并对 $t$ 归纳。
该谱序列的 $E_2$ 页为
$E_2^{p, q} = \check{H}^p(\mathcal{U}, \underline{H}^q(\mathcal{F}))$，
并收敛到 $H^{p + q}(X, \mathcal{F})$。
类似地，还有一个 $E_2$ 页为
$E_2^{p, q} = \check{H}^p(\mathcal{U}_B, \underline{H}^q(\mathcal{F}_B))$
的谱序列，它收敛到 $H^{p + q}(X_B, \mathcal{F}_B)$。
由于诸交 $U_{i_0 \ldots i_p}$ 拟紧且分离，
证明第二段的结果给出
$\check{H}^p(\mathcal{U}_B, \underline{H}^q(\mathcal{F}_B)) =
\check{H}^p(\mathcal{U}, \underline{H}^q(\mathcal{F})) \otimes_A B$。
利用 $A \to B$ 平坦，得到
$H^i(X, \mathcal{F}) \otimes_A B \to H^i(X_B, \mathcal{F}_B)$
对所有 $i$ 都是同构，证毕。
\end{proof}

\begin{lemma}[有限局部自由基变换]
\label{coherent-lemma-finite-locally-free-base-change-cohomology}
考虑概形的笛卡尔图
$$
\xymatrix{
Y \ar[d]_{g} \ar[r]_h & X \ar[d]^f \\
\Spec(B) \ar[r] & \Spec(A)
}
$$
设 $\mathcal{F}$ 为准凝聚 $\mathcal{O}_X$-模，其拉回为
$\mathcal{G} = h^*\mathcal{F}$。
若 $B$ 是有限局部自由 $A$-模，则
$H^i(X, \mathcal{F}) \otimes_A B = H^i(Y, \mathcal{G})$。
\end{lemma}

\noindent
{\bf 警告}：除非理解本引理与引理
\ref{coherent-lemma-flat-base-change-cohomology} 的差别，否则不要使用本引理。

\begin{proof}
若 $X$ 分离，选取仿射开覆盖
$\mathcal{U} : X = \bigcup_{i \in I} U_i$，并回忆
$$
\check{H}^p(\mathcal{U}, \mathcal{F}) = H^p(X, \mathcal{F}),
$$
见引理 \ref{coherent-lemma-cech-cohomology-quasi-coherent}。
令 $\mathcal{V} : Y = \bigcup_{i \in I} g^{-1}(U_i)$
为 $Y$ 的相应仿射开覆盖。
开集 $V_i = g^{-1}(U_i) = U_i \times_{\Spec(A)} \Spec(B)$
均仿射，且 $Y$ 分离。因此得到
$$
\check{H}^p(\mathcal{V}, \mathcal{G}) = H^p(Y, \mathcal{G}).
$$
我们断言 {\v C}ech 复形的映射
$$
\check{\mathcal{C}}^\bullet(\mathcal{U}, \mathcal{F}) \otimes_A B
\longrightarrow
\check{\mathcal{C}}^\bullet(\mathcal{V}, \mathcal{G})
$$
是同构。事实上，由于 $B$ 作为 $A$-模有限表示，
和 $B$ 张量在 $A$ 上与乘积可交换；见《代数》命题
\ref{algebra-proposition-fp-tensor}。
因此只需证明映射
$\Gamma(U_{i_0 \ldots i_p}, \mathcal{F}) \otimes_A B \to
\Gamma(V_{i_0 \ldots i_p}, \mathcal{G})$
是同构，而这来自引理 \ref{coherent-lemma-affine-base-change}。
由于 $A \to B$ 平坦，取上同调后同一结论仍成立。

\medskip\noindent
在一般情形下，论证完全相同：使用仿射开覆盖
$\mathcal{U} : X = \bigcup_{i \in I} U_i$ 及其相应覆盖
$\mathcal{V} : Y = \bigcup_{i \in I} V_i$，
其中 $V_i = g^{-1}(U_i)$ 如上。
我们使用《上同调》引理
\ref{cohomology-lemma-cech-spectral-sequence} 的
{\v C}ech 到上同调谱序列。
该谱序列的 $E_2$ 页为
$E_2^{p, q} = \check{H}^p(\mathcal{U}, \underline{H}^q(\mathcal{F}))$，
并收敛到 $H^{p + q}(X, \mathcal{F})$。
类似地，还有一个 $E_2$ 页为
$E_2^{p, q} = \check{H}^p(\mathcal{V}, \underline{H}^q(\mathcal{G}))$
的谱序列，它收敛到 $H^{p + q}(Y, \mathcal{G})$。
由于诸交 $U_{i_0 \ldots i_p}$ 分离，上一段的结果给出同构
$\Gamma(U_{i_0 \ldots i_p}, \underline{H}^q(\mathcal{F})) \otimes_A B
\to \Gamma(V_{i_0 \ldots i_p}, \underline{H}^q(\mathcal{G}))$。
利用 $- \otimes_A B$ 与乘积可交换且正合，得到
$\check{H}^p(\mathcal{U}, \underline{H}^q(\mathcal{F})) \otimes_A B
\to \check{H}^p(\mathcal{V}, \underline{H}^q(\mathcal{G}))$
是同构。再利用 $A \to B$ 平坦，得到
$H^i(X, \mathcal{F}) \otimes_A B \to H^i(Y, \mathcal{G})$
对所有 $i$ 都是同构，证毕。
\end{proof}


\section{余极限与高阶直像}
\label{coherent-section-colimits}

\noindent
这一性质的一般性结果见《上同调》节
\ref{cohomology-section-limits}、《层》引理
\ref{sheaves-lemma-directed-colimits-sections} 以及《模》引理
\ref{modules-lemma-finite-presentation-quasi-compact-colimit}。

\begin{lemma}
\label{coherent-lemma-colimit-cohomology}
设 $f : X \to S$ 为概形之间的拟紧拟分离态射。
设 $\mathcal{F} = \colim \mathcal{F}_i$ 为 $X$ 上阿贝尔层的滤过余极限。
则对任意 $p \geq 0$，有
$$
R^pf_*\mathcal{F} = \colim R^pf_*\mathcal{F}_i.
$$
\end{lemma}

\begin{proof}
应用《上同调》引理
\ref{cohomology-lemma-higher-direct-image-colimit}。
由于仿射开集构成 $S$ 的拓扑基，只需证明：对仿射开集
$U \subset S$，有
$H^p(f^{-1}U, \mathcal{F}) = \colim H^p(f^{-1}U, \mathcal{F}_i)$。
因为 $f^{-1}U$ 拟紧拟分离，结论来自《上同调》引理
\ref{cohomology-lemma-quasi-separated-cohomology-colimit}。
（这是因为 $f^{-1}U$ 中仿射开集所成的基满足该引理的假设。）
\end{proof}


\section{上同调与基变换（二）}
\label{coherent-section-cohomology-and-base-change-derived}

\noindent
设 $f : X \to S$ 为概形态射，并设 $\mathcal{F}$
为准凝聚 $\mathcal{O}_X$-模。若 $f$ 拟紧拟分离，我们希望将
$Rf_*\mathcal{F}$ 表示为 $S$ 上的准凝聚层复形。这由如下事实推出：
层 $R^if_*\mathcal{F}$ 在 $S$ 拟紧且具有仿射对角态射时准凝聚；
再利用 $D_\QCoh(S)$ 与 $D(\QCoh(\mathcal{O}_S))$ 的等价，
见《概形的导出范畴》命题
\ref{perfect-proposition-quasi-compact-affine-diagonal}。

\medskip\noindent
本节采用另一种方法，构造一个具有良好基变换性质的显式复形。
当 $f$ 与 $S$ 分离，或者更一般地具有仿射对角态射时，
这一构造尤其简单。由于这是迄今最常用的情形，我们将其单独处理。

\begin{lemma}
\label{coherent-lemma-separated-case-relative-cech}
设 $f : X \to S$ 为概形态射。
设 $\mathcal{F}$ 为准凝聚 $\mathcal{O}_X$-模。
假设 $X$ 拟紧，且 $X$ 与 $S$ 都具有仿射对角态射
（例如，当 $X$ 与 $S$ 都分离时）。
在此情形下，可按如下方式计算 $Rf_*\mathcal{F}$：
\begin{enumerate}
\item 选取有限仿射开覆盖
$\mathcal{U} : X = \bigcup_{i = 1, \ldots, n} U_i$。
\item 对 $i_0, \ldots, i_p \in \{1, \ldots, n\}$，记
$f_{i_0 \ldots i_p} : U_{i_0 \ldots i_p} \to S$ 为 $f$
在交集 $U_{i_0 \ldots i_p} = U_{i_0} \cap \ldots \cap U_{i_p}$
上的限制。
\item 令 $\mathcal{F}_{i_0 \ldots i_p}$ 等于 $\mathcal{F}$
在 $U_{i_0 \ldots i_p}$ 上的限制。
\item 置
$$
\check{\mathcal{C}}^p(\mathcal{U}, f, \mathcal{F}) =
\bigoplus\nolimits_{i_0 \ldots i_p}
f_{i_0 \ldots i_p *} \mathcal{F}_{i_0 \ldots i_p}
$$
并按《上同调》公式 (\ref{cohomology-equation-d-cech}) 定义微分
$d : \check{\mathcal{C}}^p(\mathcal{U}, f, \mathcal{F})
\to \check{\mathcal{C}}^{p + 1}(\mathcal{U}, f, \mathcal{F})$。
\end{enumerate}
则复形 $\check{\mathcal{C}}^\bullet(\mathcal{U}, f, \mathcal{F})$
是 $S$ 上的准凝聚层复形，并带有同构
$$
\check{\mathcal{C}}^\bullet(\mathcal{U}, f, \mathcal{F})
\longrightarrow
Rf_*\mathcal{F}
$$
它位于 $D^{+}(S)$ 中。该同构关于准凝聚层 $\mathcal{F}$ 具有函子性。
\end{lemma}

\begin{proof}
考虑解消
$\mathcal{F} \to {\mathfrak C}^\bullet(\mathcal{U}, \mathcal{F})$，
即《上同调》引理 \ref{cohomology-lemma-covering-resolution} 中的解消。
我们有复形等式
$\check{\mathcal{C}}^\bullet(\mathcal{U}, f, \mathcal{F}) =
f_*{\mathfrak C}^\bullet(\mathcal{U}, \mathcal{F})$；
二者都是准凝聚 $\mathcal{O}_S$-模的复形。
态射 $j_{i_0 \ldots i_p} : U_{i_0 \ldots i_p} \to X$
与态射 $f_{i_0 \ldots i_p} : U_{i_0 \ldots i_p} \to S$
由《态射》引理 \ref{morphisms-lemma-affine-permanence}
及引理 \ref{coherent-lemma-affine-diagonal} 都是仿射的。
故 $R^qj_{i_0 \ldots i_p *}\mathcal{F}_{i_0 \ldots i_p}$
以及 $R^qf_{i_0 \ldots i_p *}\mathcal{F}_{i_0 \ldots i_p}$
在 $q > 0$ 时均为零（引理 \ref{coherent-lemma-relative-affine-vanishing}）。
利用 $f \circ j_{i_0 \ldots i_p} = f_{i_0 \ldots i_p}$ 及
《上同调》引理 \ref{cohomology-lemma-relative-Leray} 中的谱序列，
可知
$R^qf_*(j_{i_0 \ldots i_p *}\mathcal{F}_{i_0 \ldots i_p}) = 0$
对 $q > 0$ 成立。
由于复形 ${\mathfrak C}^\bullet(\mathcal{U}, \mathcal{F})$ 的各项
是层 $j_{i_0 \ldots i_p *}\mathcal{F}_{i_0 \ldots i_p}$ 的有限直和，
由 Leray 无环性引理
（《导出范畴》引理 \ref{derived-lemma-leray-acyclicity}）得到
$$
Rf_* \mathcal{F} = f_*{\mathfrak C}^\bullet(\mathcal{U}, \mathcal{F}) =
\check{\mathcal{C}}^\bullet(\mathcal{U}, f, \mathcal{F})
$$
如所欲。
\end{proof}

\noindent
下面考察进行基变换时会发生什么。

\begin{lemma}
\label{coherent-lemma-base-change-complex}
采用图 (\ref{coherent-equation-base-change-diagram}) 中的记号。
假设 $f : X \to S$ 与 $\mathcal{F}$ 满足引理
\ref{coherent-lemma-separated-case-relative-cech} 的假设。选取有限仿射开覆盖
$\mathcal{U} : X = \bigcup U_i$ 作为 $X$ 的覆盖。
存在典范同构
$$
g^*\check{\mathcal{C}}^\bullet(\mathcal{U}, f, \mathcal{F})
\longrightarrow
Rf'_*\mathcal{F}'
$$
它位于 $D^{+}(S')$ 中。此外，若 $S' \to S$ 仿射，则实际上
$$
g^*\check{\mathcal{C}}^\bullet(\mathcal{U}, f, \mathcal{F})
=
\check{\mathcal{C}}^\bullet(\mathcal{U}', f', \mathcal{F}')
$$
其中 $\mathcal{U}' : X' = \bigcup U_i'$，而
$U_i' = (g')^{-1}(U_i) = U_{i, S'}$ 也仿射。
\end{lemma}

\begin{proof}
可定义 $U_i' = (g')^{-1}(U_i) = U_{i, S'}$，
而不论 $S'$ 在 $S$ 上是否仿射。
令 $\mathcal{U}' : X' = \bigcup U_i'$ 为由此得到的 $X'$ 的覆盖。
在此情形下，我们断言
$$
g^*\check{\mathcal{C}}^\bullet(\mathcal{U}, f, \mathcal{F})
=
\check{\mathcal{C}}^\bullet(\mathcal{U}', f', \mathcal{F}')
$$
其中 $\check{\mathcal{C}}^\bullet(\mathcal{U}', f', \mathcal{F}')$
的定义与引理 \ref{coherent-lemma-separated-case-relative-cech} 中完全相同。
这可通过在 $S'$ 上局部工作，由仿射态射的情形
（引理 \ref{coherent-lemma-affine-base-change}）立即得到。
此外，完全如引理 \ref{coherent-lemma-separated-case-relative-cech} 的证明，
可见存在同构
$$
\check{\mathcal{C}}^\bullet(\mathcal{U}', f', \mathcal{F}')
\longrightarrow
Rf'_*\mathcal{F}'
$$
它位于 $D^{+}(S')$ 中，因为态射 $U_i' \to X'$ 与 $U_i' \to S'$
仍然仿射（它们是仿射态射的基变换）。细节从略。
\end{proof}

\noindent
上述引理说明，复形
$$
\mathcal{K}^\bullet = \check{\mathcal{C}}^\bullet(\mathcal{U}, f, \mathcal{F})
$$
是 $S$ 上准凝聚层的下有界复形，它\emph{泛地}计算
$f : X \to S$ 的高阶直像。
这是该特定复形所具有的性质；一般而言，将
$\check{\mathcal{C}}^\bullet(\mathcal{U}, f, \mathcal{F})$
换成拟同构复形并不能保持这一性质！换言之，
这不是一个在导出范畴中有意义的陈述。
原因在于，拉回 $g^*\mathcal{K}^\bullet$ 一般\emph{不}等于
$Lg^*\mathcal{K}^\bullet$，后者是 $\mathcal{K}^\bullet$ 的导出拉回！

\medskip\noindent
下面给出可以证明该陈述的一个更一般情形。
我们指出，在大多数应用中，$S$ 分离这一条件无伤大雅，
因为这里通常是用它证明全导出像的某个局部性质。
证明要复杂得多，并使用超覆盖；这也是有时如何使用超覆盖的一个好例子。

\begin{lemma}
\label{coherent-lemma-hypercoverings}
设 $f : X \to S$ 为概形态射。
设 $\mathcal{F}$ 为 $X$ 上的准凝聚层。
假设 $f$ 拟紧拟分离，且 $S$ 拟紧并分离。
则存在下有界复形 $\mathcal{K}^\bullet$，
它由准凝聚 $\mathcal{O}_S$-模组成，并具有如下性质：对每个态射
$g : S' \to S$，复形 $g^*\mathcal{K}^\bullet$ 是
$Rf'_*\mathcal{F}'$ 的一个代表元，其中记号如图
(\ref{coherent-equation-base-change-diagram})。
\end{lemma}

\begin{proof}
（若 $f$ 也分离，请见引理 \ref{coherent-lemma-base-change-complex}。）
这些假设特别推出 $X$ 作为概形拟紧拟分离。
令 $\mathcal{B}$ 为 $X$ 的仿射开集所成的集合。由《超覆盖》引理
\ref{hypercovering-lemma-quasi-separated-quasi-compact-hypercovering}，
可找到超覆盖 $K = (I, \{U_i\})$，使每个 $I_n$ 有限，
且每个 $U_i$ 都是 $X$ 的仿射开集。由《超覆盖》引理
\ref{hypercovering-lemma-cech-spectral-sequence}，
存在一个 $E_2$ 页为
$$
E_2^{p, q} = \check{H}^p(K, \underline{H}^q(\mathcal{F}))
$$
且收敛到 $H^{p + q}(X, \mathcal{F})$ 的谱序列。注意，
$\check{H}^p(K, \underline{H}^q(\mathcal{F}))$ 是下列复形的第 $p$ 个上同调群：
$$
\prod\nolimits_{i \in I_0} H^q(U_i, \mathcal{F})
\to
\prod\nolimits_{i \in I_1} H^q(U_i, \mathcal{F})
\to
\prod\nolimits_{i \in I_2} H^q(U_i, \mathcal{F})
\to \ldots
$$
由于每个 $U_i$ 仿射，可见除非 $q = 0$，否则它为零；
而在后一情形得到
$$
\prod\nolimits_{i \in I_0} \mathcal{F}(U_i)
\to
\prod\nolimits_{i \in I_1} \mathcal{F}(U_i)
\to
\prod\nolimits_{i \in I_2} \mathcal{F}(U_i)
\to \ldots
$$
因此可知 $R\Gamma(X, \mathcal{F})$ 由这个复形计算。

\medskip\noindent
对任意 $n$ 与 $i \in I_n$，记 $f_i : U_i \to S$ 为 $f$
在 $U_i$ 上的限制。由于 $S$ 分离且 $U_i$ 仿射，该态射仿射。
考虑 $S$ 上的准凝聚层复形
$$
\mathcal{K}^\bullet = (
\prod\nolimits_{i \in I_0} f_{i, *}\mathcal{F}|_{U_i}
\to
\prod\nolimits_{i \in I_1} f_{i, *}\mathcal{F}|_{U_i}
\to
\prod\nolimits_{i \in I_2} f_{i, *}\mathcal{F}|_{U_i}
\to \ldots )
$$
与《超覆盖》引理 \ref{hypercovering-lemma-cech-spectral-sequence}
中一样，我们得到映射 $\mathcal{K}^\bullet \to Rf_*\mathcal{F}$，
它位于 $D(\mathcal{O}_S)$ 中；其构造是选取 $\mathcal{F}$ 的单射解消
（细节从略）。
考虑任意仿射概形 $V$ 及态射 $g : V \to S$。
则基变换 $X_V$ 具有由基变换得到的超覆盖
$K_V = (I, \{U_{i, V}\})$。此外，
$g^*f_{i, *}\mathcal{F} = f_{i, V, *}(g')^*\mathcal{F}|_{U_{i, V}}$。
因此上述论证证明，$\Gamma(V, g^*\mathcal{K}^\bullet)$
计算 $R\Gamma(X_V, (g')^*\mathcal{F})$。
要证明复形相等，只需在 $S'$ 上 Zariski 局部证明，故引理得证。
\end{proof}

\noindent
下面的引理是平坦基变换的一种变式。

\begin{lemma}
\label{coherent-lemma-base-change-tensor-with-flat}
考虑概形的笛卡尔图
$$
\xymatrix{
X' \ar[d]_{f'} \ar[r]_{g'} & X \ar[d]^f \\
S' \ar[r]^g & S
}
$$
设 $\mathcal{F}$ 为准凝聚 $\mathcal{O}_X$-模。
设 $\mathcal{G}$ 为准凝聚 $\mathcal{O}_{S'}$-模，
并且在 $S$ 上平坦。假设 $f$ 拟紧拟分离。
对任意 $i \geq 0$，有等同
$$
\mathcal{G} \otimes_{\mathcal{O}_{S'}} g^*R^if_*\mathcal{F} =
R^if'_*\left((f')^*\mathcal{G}
\otimes_{\mathcal{O}_{X'}} (g')^*\mathcal{F}\right)
$$
\end{lemma}

\begin{proof}
先构造从左到右的映射。首先，有基变换映射
$Lg^*Rf_*\mathcal{F} \to Rf'_*L(g')^*\mathcal{F}$。另有伴随映射
$\mathcal{G} \to Rf'_*L(f')^*\mathcal{G}$。利用相对杯积，得到
\begin{align*}
\mathcal{G}
\otimes_{\mathcal{O}_{S'}}^\mathbf{L}
Lg^*Rf_*\mathcal{F}
& \to
Rf'_*L(f')^*\mathcal{G}
\otimes_{\mathcal{O}_{S'}}^\mathbf{L}
Rf'_*L(g')^*\mathcal{F} \\
& \to
Rf'_*\left(L(f')^*\mathcal{G}
\otimes_{\mathcal{O}_{X'}}^\mathbf{L}
L(g')^*\mathcal{F}\right) \\
& \to
Rf'_*\left((f')^*\mathcal{G}
\otimes_{\mathcal{O}_{X'}} (g')^*\mathcal{F}\right)
\end{align*}
其中中间的箭头使用了相对杯积，见《上同调》备注
\ref{cohomology-remark-cup-product}。
该复合的源为
$$
\mathcal{G} \otimes_{\mathcal{O}_{S'}}^\mathbf{L} Lg^*Rf_*\mathcal{F} =
\mathcal{G} \otimes_{g^{-1}\mathcal{O}_S}^\mathbf{L} g^{-1}Rf_*\mathcal{F}
$$
这是《上同调》引理 \ref{cohomology-lemma-variant-derived-pullback}。
由于 $\mathcal{G}$ 作为 $g^{-1}\mathcal{O}_S$-模层平坦，
且 $g^{-1}$ 是正合函子，这是一个复形，其第 $i$ 个上同调层为
$\mathcal{G} \otimes_{g^{-1}\mathcal{O}_S} g^{-1}R^if_*\mathcal{F} =
\mathcal{G} \otimes_{\mathcal{O}_{S'}} g^*R^if_*\mathcal{F}$。
这样便得到引理等式中从左到右的全局映射。
为证明此映射是同构，可在 $S'$ 上局部工作。
因此可以并且确实假设 $S$ 与 $S'$ 都是仿射概形。

\medskip\noindent
现在证明 $S$ 与 $S'$ 仿射的情形。设 $S = \Spec(A)$ 且 $S' = \Spec(B)$，
并设 $\mathcal{G}$ 对应于 $B$-模 $N$，后者按假设在 $A$ 上平坦。
由于 $S$ 仿射，$X$ 拟紧拟分离。
我们用超覆盖论证完成证明；若 $X$ 分离或具有仿射对角态射，
则可使用一个 {\v C}ech 覆盖。
令 $\mathcal{B}$ 为 $X$ 的仿射开集所成的集合。由《超覆盖》引理
\ref{hypercovering-lemma-quasi-separated-quasi-compact-hypercovering}，
可找到超覆盖 $K = (I, \{U_i\})$，它是 $X$ 的超覆盖，且每个
$I_n$ 有限、每个 $U_i$ 都是 $X$ 的仿射开集。由《超覆盖》引理
\ref{hypercovering-lemma-cech-spectral-sequence}，
存在一个 $E_2$ 页为
$$
E_2^{p, q} = \check{H}^p(K, \underline{H}^q(\mathcal{F}))
$$
且收敛到 $H^{p + q}(X, \mathcal{F})$ 的谱序列。由于每个 $U_i$ 仿射，
且 $\mathcal{F}$ 准凝聚，$\underline{H}^q(\mathcal{F})$
在 $U_i$ 上的值对 $q > 0$ 为零。因此谱序列退化，
从而上同调模 $H^q(X, \mathcal{F})$ 由下列复形计算：
$$
\prod\nolimits_{i \in I_0} \mathcal{F}(U_i)
\to
\prod\nolimits_{i \in I_1} \mathcal{F}(U_i)
\to
\prod\nolimits_{i \in I_2} \mathcal{F}(U_i)
\to \ldots
$$
接着注意，我们的超覆盖在 $S'$ 上的基变换是
$X' = S' \times_S X$ 的一个超覆盖。
概形 $S' \times_S U_i$ 也都仿射，并且有
$$
\left((f')^*\mathcal{G} \otimes_{\mathcal{O}_{S'}} (g')^*\mathcal{F}\right)
(S' \times_S U_i) =
N \otimes_A \mathcal{F}(U_i)
$$
由此可知，上同调模
$H^q(X', (f')^*\mathcal{G} \otimes_{\mathcal{O}_{S'}} (g')^*\mathcal{F})$
由下列复形计算：
$$
N \otimes_A \left(
\prod\nolimits_{i \in I_0} \mathcal{F}(U_i)
\to
\prod\nolimits_{i \in I_1} \mathcal{F}(U_i)
\to
\prod\nolimits_{i \in I_2} \mathcal{F}(U_i)
\to \ldots
\right)
$$
由于 $N$ 在 $A$ 上平坦，可知
$$
H^q(X', (f')^*\mathcal{G} \otimes_{\mathcal{O}_{S'}} (g')^*\mathcal{F})
=
N \otimes_A H^q(X, \mathcal{F})
$$
这正是待证陈述的代数表述，故证明完成。
\end{proof}







\section{射影空间的上同调}
\label{coherent-section-cohomology-projective-space}

\noindent
本节计算 $\mathbf{P}^n_S$ 的结构层各扭转的上同调，其中基概形为 $S$。
回忆，$\mathbf{P}^n_S$ 定义为纤维积
$
\mathbf{P}^n_S = S \times_{\Spec(\mathbf{Z})} \mathbf{P}^n_{\mathbf{Z}}
$
见《构造》定义 \ref{constructions-definition-projective-space}。
《构造》引理 \ref{constructions-lemma-projective-space-bundle} 已证明它等于
$$
\mathbf{P}^n_S = \underline{\text{Proj}}_S(\mathcal{O}_S[T_0, \ldots, T_n])
$$
特别地，射影空间是射影丛的一个特殊情形。
若 $S = \Spec(R)$ 仿射，则有
$$
\mathbf{P}^n_S = \mathbf{P}^n_R = \text{Proj}(R[T_0, \ldots, T_n]).
$$
所有这些等同彼此相容，也与扭转结构层
$\mathcal{O}_{\mathbf{P}^n_S}(d)$ 的构造相容。

\medskip\noindent
在陈述结果前，需要引入一些记号。
设 $R$ 为环。
回忆，$R[T_0, \ldots, T_n]$ 是分次 $R$-代数，
其中每个 $T_i$ 都是次数为 $1$ 的齐次元。
以 $(R[T_0, \ldots, T_n])_d$ 表示次数为 $d$ 的分量。
它是有限自由 $R$-模，秩为 $\binom{n + d}{d}$，
当 $d \geq 0$；否则为零。它的一组基由单项式
$T_0^{e_0} \ldots T_n^{e_n}$ 组成，其中 $\sum e_i = d$。
还要使用如下记号：
$R[\frac{1}{T_0}, \ldots, \frac{1}{T_n}]$ 表示 $\mathbf{Z}$-分次环，
其中 $\frac{1}{T_i}$ 的次数为 $-1$。特别地，下述陈述中出现的
$\mathbf{Z}$-分次 $R[\frac{1}{T_0}, \ldots, \frac{1}{T_n}]$-模
$$
\frac{1}{T_0 \ldots T_n} R[\frac{1}{T_0}, \ldots, \frac{1}{T_n}]
$$
在次数 $\geq -n$ 时为零，以生成元 $\frac{1}{T_0 \ldots T_n}$
在次数 $-n - 1$ 时自由生成，并且是秩为
$(-1)^n\binom{n + d}{d}$ 的自由模，当 $d \leq -n - 1$。

\begin{lemma}
\label{coherent-lemma-cohomology-projective-space-over-ring}
\begin{reference}
\cite[III Proposition 2.1.12]{EGA}
\end{reference}
设 $R$ 为环。
设 $n \geq 0$ 为整数。
则
$$
H^q(\mathbf{P}^n, \mathcal{O}_{\mathbf{P}^n_R}(d)) =
\left\{
\begin{matrix}
(R[T_0, \ldots, T_n])_d & \text{若} & q = 0,\ d \geq 0 \\
\left(\frac{1}{T_0 \ldots T_n} R[\frac{1}{T_0}, \ldots, \frac{1}{T_n}]\right)_d
& \text{若} & q = n,\ d < 0 \\
0 & \text{若} & q, n, d\text{ not as above}
\end{matrix}
\right.
$$
这是 $R$-模的等式。
\end{lemma}

\begin{proof}
使用标准仿射开覆盖
$$
\mathcal{U} : \mathbf{P}^n_R = \bigcup\nolimits_{i = 0}^n D_{+}(T_i)
$$
并利用 {\v C}ech 复形计算上同调。
这是允许的，因为由引理 \ref{coherent-lemma-cech-cohomology-quasi-coherent}，
有限多个仿射开集 $D_{+}(T_i)$ 的任意交仍是标准仿射开集
（见《构造》节 \ref{constructions-section-proj}）。
事实上，根据《上同调》引理 \ref{cohomology-lemma-ordered-alternating}
及 \ref{cohomology-lemma-alternating-usual}，
可使用交错或有序 {\v C}ech 复形。

\medskip\noindent
在 $\{0, \ldots, n\}$ 上采用通常次序。
于是所考察复形的各项为
$$
\check{\mathcal{C}}_{ord}^p(\mathcal{U}, \mathcal{O}_{\mathbf{P}_R}(d))
=
\bigoplus\nolimits_{i_0 < \ldots < i_p}
(R[T_0, \ldots, T_n, \frac{1}{T_{i_0} \ldots T_{i_p}}])_d
$$
此外，映射由通常公式给出：
$$
d(s)_{i_0 \ldots i_{p + 1}} =
\sum\nolimits_{j = 0}^{p + 1} (-1)^j s_{i_0 \ldots \hat i_j \ldots i_{p + 1}}
$$
见《上同调》节 \ref{cohomology-section-alternating-cech}。
注意，此复形的每一项都有自然的 $\mathbf{Z}^{n + 1}$-分次。
具体而言，规定单项式 $T_0^{e_0} \ldots T_n^{e_n}$ 为权
$(e_0, \ldots, e_n) \in \mathbf{Z}^{n + 1}$ 的齐次元。
显然，上述微分保持该分次。用公式表示即
$$
\check{\mathcal{C}}_{ord}^\bullet(\mathcal{U}, \mathcal{O}_{\mathbf{P}_R}(d))
=
\bigoplus\nolimits_{\vec{e} \in \mathbf{Z}^{n + 1}}
\check{\mathcal{C}}^\bullet(\vec{e})
$$
其中右侧并非所有直和项都会出现（见下文）。
因此，为计算该复形的上同调模，只需计算各分次部分的上同调，
最后取直和。

\medskip\noindent
固定 $\vec{e} = (e_0, \ldots, e_n) \in \mathbf{Z}^{n + 1}$。
要使这个权出现在上述复形中，需要假设
$e_0 + \ldots + e_n = d$（否则它当然会出现在结构层的另一扭转中）。
在此假设下，置
$$
NEG(\vec{e}) = \{i \in \{0, \ldots, n\} \mid e_i < 0\}.
$$
采用此记号，上述 {\v C}ech 复形的权 $\vec{e}$ 直和项
$\check{\mathcal{C}}^\bullet(\vec{e})$ 的各项为
$$
\check{\mathcal{C}}^p(\vec{e})
=
\bigoplus\nolimits_{i_0 < \ldots < i_p,
\ NEG(\vec{e}) \subset \{i_0, \ldots, i_p\}}
R \cdot T_0^{e_0} \ldots T_n^{e_n}
$$
换言之，对满足 $i_0 < \ldots < i_p$ 但
$NEG(\vec{e})$ 不包含于 $\{i_0 \ldots i_p\}$ 的指标，
相应项为零。复形 $\check{\mathcal{C}}^\bullet(\vec{e})$ 的微分
仍由上述完全相同的公式给出。

\medskip\noindent
假设 $NEG(\vec{e}) = \{0, \ldots, n\}$，即所有指数
$e_i$ 都为负。
此时复形 $\check{\mathcal{C}}^\bullet(\vec{e})$ 只有一项，即
$\check{\mathcal{C}}^n(\vec{e}) =
R \cdot \frac{1}{T_0^{-e_0} \ldots T_n^{-e_n}}$。因此在此情形下
$$
H^q(\check{\mathcal{C}}^\bullet(\vec{e})) =
\left\{
\begin{matrix}
R \cdot \frac{1}{T_0^{-e_0} \ldots T_n^{-e_n}} & \text{若} & q = n \\
0 & \text{若} & \text{否则}
\end{matrix}
\right.
$$
所有这些项的直和显然给出
$$
\left(\frac{1}{T_0 \ldots T_n} R[\frac{1}{T_0}, \ldots, \frac{1}{T_n}]\right)_d
$$
在次数 $n$ 的值，如引理所述。此外，这些项对其他次数的上同调
没有贡献（也与引理的陈述一致）。

\medskip\noindent
假设 $NEG(\vec{e}) = \emptyset$。此时复形
$\check{\mathcal{C}}^\bullet(\vec{e})$ 有一个直和项 $R$，
对应于所有 $i_0 < \ldots < i_p$。
将复形 $\check{\mathcal{C}}^\bullet(\vec{e})$ 与另一个复形比较。
具体地，考虑仿射开覆盖
$$
\mathcal{V} : \Spec(R) = \bigcup\nolimits_{i \in \{0, \ldots, n\}} V_i
$$
其中 $V_i = \Spec(R)$ 对所有 $i$ 都成立。考虑交错 {\v C}ech 复形
$$
\check{\mathcal{C}}_{ord}^\bullet(\mathcal{V}, \mathcal{O}_{\Spec(R)})
$$
由与上面相同的论证，它计算 $\Spec(R)$ 上结构层的上同调。
因此可见\par\noindent
$H^p(
\check{\mathcal{C}}_{ord}^\bullet(\mathcal{V}, \mathcal{O}_{\Spec(R)})
) = R$\hfill\break
当 $p = 0$，而为 $0$ 当 $p > 0$。\par
这些事实见引理 \ref{coherent-lemma-cech-cohomology-quasi-coherent-trivial} 及其证明。
还要注意，
$\check{\mathcal{C}}_{ord}^\bullet(\mathcal{V}, \mathcal{O}_{\Spec(R)})$
也有一个直和项 $R$ 对应于每个 $i_0 < \ldots < i_p$，
并且其微分与 $\check{\mathcal{C}}^\bullet(\vec{e})$ 完全相同。
换言之，这些复形彼此同构，因而有相同的上同调。由此得到
$$
H^q(\check{\mathcal{C}}^\bullet(\vec{e})) =
\left\{
\begin{matrix}
R \cdot T_0^{e_0} \ldots T_n^{e_n} & \text{若} & q = 0 \\
0 & \text{若} & \text{否则}
\end{matrix}
\right.
$$
这里假设 $NEG(\vec{e}) = \emptyset$。
所有这些项的直和显然给出
$$
(R[T_0, \ldots, T_n])_d
$$
在次数 $0$ 的值，如引理所述。此外，这些项对其他次数的上同调
没有贡献（也与引理的陈述一致）。

\medskip\noindent
为完成引理的证明，还须说明：当
$\check{\mathcal{C}}^\bullet(\vec{e})$ 满足
$NEG(\vec{e})$ 既非空也不等于 $\{0, \ldots, n\}$ 时，它是无环的。
选取指标 $i_{\text{fix}} \not \in NEG(\vec{e})$（这样的指标存在）。
考虑映射
$$
h :
\check{\mathcal{C}}^{p + 1}(\vec{e})
\to
\check{\mathcal{C}}^p(\vec{e})
$$
其规则为：对 $i_0 < \ldots < i_p$，有
$$
h(s)_{i_0 \ldots i_p} =
\left\{
\begin{matrix}
0 & \text{若} & p \not \in \{0, \ldots, n - 1\} \\
0 & \text{若} & i_{\text{fix}} \in \{i_0, \ldots, i_p\} \\
s_{i_{\text{fix}} i_0 \ldots i_p} & \text{若} & i_{\text{fix}} < i_0 \\
(-1)^a s_{i_0 \ldots i_{a - 1} i_{\text{fix}} i_a \ldots i_p} &
\text{若} & i_{a - 1} < i_{\text{fix}} < i_a \\
(-1)^{p + 1} s_{i_0 \ldots i_p i_{\text{fix}}} &
\text{若} & i_p < i_{\text{fix}}
\end{matrix}
\right.
$$
请与引理 \ref{coherent-lemma-cech-cohomology-quasi-coherent-trivial} 的证明比较。
该定义有意义，因为
$$
NEG(\vec{e}) \subset \{i_0, \ldots, i_p\}
\Leftrightarrow
NEG(\vec{e}) \subset \{i_{\text{fix}}, i_0, \ldots, i_p\}
$$
与该引理证明中完全相同的组合计算\footnote{
例如，假设 $i_0 < \ldots < i_p$ 满足
$i_{\text{fix}} \not \in \{i_0, \ldots, i_p\}$，并且有
$i_{a - 1} < i_{\text{fix}} < i_a$，其中 $1 \leq a \leq p$。则
\begin{align*}
&
(dh + hd)(s)_{i_0 \ldots i_p} \\
& =
\sum\nolimits_{j = 0}^p
(-1)^j
h(s)_{i_0 \ldots \hat i_j \ldots i_p}
+
(-1)^a d(s)_{i_0 \ldots i_{a - 1} i_{\text{fix}} i_a \ldots i_p}\\
& =
\sum\nolimits_{j = 0}^{a - 1}
(-1)^{j + a - 1}
s_{i_0 \ldots \hat i_j \ldots i_{a - 1} i_{\text{fix}} i_a \ldots i_p} +
\sum\nolimits_{j = a}^p
(-1)^{j + a}
s_{i_0 \ldots i_{a - 1} i_{\text{fix}} i_a \ldots \hat i_j \ldots i_p}
+ \\
&
\sum\nolimits_{j = 0}^{a - 1}
(-1)^{a + j}
s_{i_0 \ldots \hat i_j \ldots i_{a - 1} i_{\text{fix}} i_a \ldots i_p} +
(-1)^{2a} s_{i_0 \ldots i_p} +
\sum\nolimits_{j = a}^p
(-1)^{a + j + 1}
s_{i_0 \ldots i_{a - 1} i_{\text{fix}} i_a \ldots \hat i_j \ldots i_p}
\\
& =
s_{i_0 \ldots i_p}
\end{align*}
如所欲。其他情形类似。}
与引理 \ref{coherent-lemma-cech-cohomology-quasi-coherent-trivial} 的证明一样，这说明
$$
(hd + dh)(s)_{i_0 \ldots i_p}
=
s_{i_0 \ldots i_p}
$$
因此复形 $\check{\mathcal{C}}^\bullet(\vec{e})$ 的恒等映射
零伦，从而该复形无环。
\end{proof}

\noindent
在下面的引理中，将使用自由 $R$-模之间的配对
$$
R[T_0, \ldots, T_n]
\times
\frac{1}{T_0 \ldots T_n} R[\frac{1}{T_0}, \ldots, \frac{1}{T_n}]
\longrightarrow
R
$$
它由如下规则定义：
$$
(f, g)
\longmapsto
\text{coefficient of }
\frac{1}{T_0 \ldots T_n}
\text{ 于 }fg.
$$
换言之，基元素 $T_0^{e_0} \ldots T_n^{e_n}$ 与基元素
$T_0^{d_0} \ldots T_n^{d_n}$ 配对得到 $1$，当且仅当
$e_i + d_i = -1$ 对所有 $i$ 成立；在其余所有情形下配对为零。
利用这个配对，得到等同
$$
\left(\frac{1}{T_0 \ldots T_n} R[\frac{1}{T_0}, \ldots, \frac{1}{T_n}]\right)_d
=
\Hom_R((R[T_0, \ldots, T_n])_{-n - 1 - d}, R)
$$
因此可将引理 \ref{coherent-lemma-cohomology-projective-space-over-ring} 的结果
重新表述为
\begin{equation}
\label{coherent-equation-identify}
H^q(\mathbf{P}^n, \mathcal{O}_{\mathbf{P}^n_R}(d)) =
\left\{
\begin{matrix}
(R[T_0, \ldots, T_n])_d & \text{若} & q = 0 \\
0 & \text{若} & q \not = 0, n \\
\Hom_R((R[T_0, \ldots, T_n])_{-n - 1 - d}, R)
& \text{若} & q = n
\end{matrix}
\right.
\end{equation}

\begin{lemma}
\label{coherent-lemma-identify-functorially}
公式 (\ref{coherent-equation-identify}) 中的等同与环同态 $R \to R'$ 的基变换相容。
此外，对任意齐次元 $f \in R[T_0, \ldots, T_n]$，若其次数为 $m$，
则乘以 $f$ 的映射
$$
\mathcal{O}_{\mathbf{P}^n_R}(d)
\longrightarrow
\mathcal{O}_{\mathbf{P}^n_R}(d + m)
$$
在公式 (\ref{coherent-equation-identify}) 的等同下，于上同调群上诱导的映射
为乘以 $f$（在 $H^0$ 上），以及乘以 $f$ 的对偶映射
$$
(R[T_0, \ldots, T_n])_{-n - 1 - (d + m)}
\longrightarrow
(R[T_0, \ldots, T_n])_{-n - 1 - d}
$$
（在 $H^n$ 上）。
\end{lemma}

\begin{proof}
假设 $R \to R'$ 为环同态。
令 $\mathcal{U}$ 为 $\mathbf{P}^n_R$ 的标准仿射开覆盖，
令 $\mathcal{U}'$ 为 $\mathbf{P}^n_{R'}$ 的标准仿射开覆盖。
注意，$\mathcal{U}'$ 是覆盖 $\mathcal{U}$ 沿典范态射
$\mathbf{P}^n_{R'} \to \mathbf{P}^n_R$ 的拉回。因此有 {\v C}ech 复形的映射
$$
\gamma :
\check{\mathcal{C}}_{ord}^\bullet(\mathcal{U},
\mathcal{O}_{\mathbf{P}_R}(d))
\longrightarrow
\check{\mathcal{C}}_{ord}^\bullet(\mathcal{U}',
\mathcal{O}_{\mathbf{P}_{R'}}(d))
$$
由《上同调》引理 \ref{cohomology-lemma-functoriality-cech}，
它与上同调上的映射相容。
从引理 \ref{coherent-lemma-cohomology-projective-space-over-ring} 证明中的计算可立即看出，
这个 {\v C}ech 复形映射与所论上同调群的等同相容。
（也就是说，$R$ 上 {\v C}ech 复形的基元素恰好映到
$R'$ 上 {\v C}ech 复形的相应基元素。）这证明了引理的第一项陈述。

\medskip\noindent
现在固定环 $R$，并考虑两个齐次多项式
$f, g \in R[T_0, \ldots, T_n]$，二者次数同为 $m$。
由于上同调是加性函子，乘以 $f + g$ 所诱导的映射显然等于
乘以 $f$ 所诱导的映射与乘以 $g$ 所诱导的映射之和。
此外，由于上同调是函子，乘以乘积 $fg$ 也有类似结论，
其中 $f, g$ 都齐次（但次数不必相同）。
因此，为验证引理的第二项陈述，只需证明 $f = x \in R$ 或 $f = T_i$ 的情形。
在乘以元素 $x \in R$ 的情形，结论成立，因为所涉及的每个上同调群或复形
都具有 $R$-模或 $R$-模复形的结构。
最后，考虑乘以 $T_i$，将其视为 $\mathcal{O}_{\mathbf{P}^n_R}$-线性映射
$$
\mathcal{O}_{\mathbf{P}^n_R}(d)
\longrightarrow
\mathcal{O}_{\mathbf{P}^n_R}(d + 1)
$$
关于 $H^0$ 的陈述显然。关于 $H^n$ 的陈述，
将乘以 $T_i$ 视为 {\v C}ech 复形上的映射
$$
\check{\mathcal{C}}_{ord}^\bullet(\mathcal{U},
\mathcal{O}_{\mathbf{P}_R}(d))
\longrightarrow
\check{\mathcal{C}}_{ord}^\bullet(\mathcal{U},
\mathcal{O}_{\mathbf{P}_R}(d + 1))
$$
将使用引理 \ref{coherent-lemma-cohomology-projective-space-over-ring} 的证明中引入的记号。
依据分解
$$
\check{\mathcal{C}}_{ord}^\bullet(\mathcal{U}, \mathcal{O}_{\mathbf{P}_R}(d))
=
\bigoplus\nolimits_{\vec{e} \in \mathbf{Z}^{n + 1}, \ \sum e_i = d}
\check{\mathcal{C}}^\bullet(\vec{e})
$$
以及
$$
\check{\mathcal{C}}_{ord}^\bullet(\mathcal{U},
\mathcal{O}_{\mathbf{P}_R}(d + 1))
=
\bigoplus\nolimits_{\vec{e} \in \mathbf{Z}^{n + 1}, \ \sum e_i = d + 1}
\check{\mathcal{C}}^\bullet(\vec{e})
$$
考察乘以 $T_i$ 的作用。
显然，它将子复形 $\check{\mathcal{C}}^\bullet(\vec{e})$ 映到子复形
$\check{\mathcal{C}}^\bullet(\vec{e} + \vec{b}_i)$，其中
$\vec{b}_i = (0, \ldots, 0, 1, 0, \ldots, 0))$ 是第 $i$ 个基向量。
换言之，它将 $H^n$ 中对应于 $\vec{e}$、满足 $e_i < 0$ 且
$\sum e_i = d$ 的直和项映到 $H^n$ 中对应于
$\vec{e} + \vec{b}_i$ 的直和项（若 $e_i + b_i \geq 0$，后者为零）。
容易看出，这恰好对应于乘以 $T_i$ 的对偶映射的作用，即映射
$$
(R[T_0, \ldots, T_n])_{-n - 1 - (d + 1)}
\longrightarrow
(R[T_0, \ldots, T_n])_{-n - 1 - d}
$$
引理得证。
\end{proof}

\noindent
在陈述相对版本前，需要引入一些记号。
回忆，$\mathcal{O}_S[T_0, \ldots, T_n]$ 是分次
$\mathcal{O}_S$-模，其中每个 $T_i$ 都是次数为 $1$ 的齐次元。
以 $(\mathcal{O}_S[T_0, \ldots, T_n])_d$ 表示次数为 $d$ 的分量。
它是秩为 $\binom{n + d}{d}$ 的有限局部自由层，定义在 $S$ 上。

\begin{lemma}
\label{coherent-lemma-cohomology-projective-space-over-base}
设 $S$ 为概形。
设 $n \geq 0$ 为整数。
考虑结构态射
$$
f : \mathbf{P}^n_S \longrightarrow S.
$$
则
$$
R^qf_*(\mathcal{O}_{\mathbf{P}^n_S}(d)) =
\left\{
\begin{matrix}
(\mathcal{O}_S[T_0, \ldots, T_n])_d & \text{若} & q = 0 \\
0 & \text{若} & q \not = 0, n \\
\SheafHom_{\mathcal{O}_S}(
(\mathcal{O}_S[T_0, \ldots, T_n])_{- n - 1 - d}, \mathcal{O}_S)
& \text{若} & q = n
\end{matrix}
\right.
$$
\end{lemma}

\begin{proof}
从略。提示：这由公式 (\ref{coherent-equation-identify}) 中的等同
与仿射基变换相容而得，见引理 \ref{coherent-lemma-identify-functorially}。
\end{proof}

\noindent
下面陈述与有限局部自由层相伴的射影丛版本。
设 $S$ 为概形。设 $\mathcal{E}$ 为有限局部自由 $\mathcal{O}_S$-模，
其秩恒为 $n + 1$；见《模》节 \ref{modules-section-locally-free}。
此时把 $\text{Sym}(\mathcal{E})$ 看作分次 $\mathcal{O}_S$-模，
其中 $\mathcal{E}$ 是次数为 $1$ 的分次部分。
而 $\text{Sym}^d(\mathcal{E})$ 是次数为 $d$ 的分量。
它是秩为 $\binom{n + d}{d}$ 的有限局部自由层，定义在 $S$ 上。
回忆，我们的规范化约定为
$$
\pi :
\mathbf{P}(\mathcal{E})
=
\underline{\text{Proj}}_S(\text{Sym}(\mathcal{E}))
\longrightarrow
S
$$
并且有自然映射
$\text{Sym}^d(\mathcal{E}) \to \pi_*\mathcal{O}_{\mathbf{P}(\mathcal{E})}(d)$。

\begin{lemma}
\label{coherent-lemma-cohomology-projective-bundle}
设 $S$ 为概形。设 $n \geq 1$。
设 $\mathcal{E}$ 为有限局部自由 $\mathcal{O}_S$-模，
其秩恒为 $n + 1$。
考虑结构态射
$$
\pi : \mathbf{P}(\mathcal{E}) \longrightarrow S.
$$
则
$$
R^q\pi_*(\mathcal{O}_{\mathbf{P}(\mathcal{E})}(d)) =
\left\{
\begin{matrix}
\text{Sym}^d(\mathcal{E}) & \text{若} & q = 0 \\
0 & \text{若} & q \not = 0, n \\
\SheafHom_{\mathcal{O}_S}(
\text{Sym}^{- n - 1 - d}(\mathcal{E})
\otimes_{\mathcal{O}_S}
\wedge^{n + 1}\mathcal{E},
\mathcal{O}_S)
& \text{若} & q = n
\end{matrix}
\right.
$$
这些等同与基变换以及局部自由层之间的同构相容。
\end{lemma}

\begin{proof}
考虑典范映射
$$
\pi^*\mathcal{E} \longrightarrow \mathcal{O}_{\mathbf{P}(\mathcal{E})}(1)
$$
向下扭转 $1$，得到
$$
\pi^*(\mathcal{E})(-1) \longrightarrow \mathcal{O}_{\mathbf{P}(\mathcal{E})}
$$
这是从秩为 $n + 1$ 的局部自由层到结构层的满射。
因此相应的 Koszul 复形正合（《更多代数》引理
\ref{more-algebra-lemma-homotopy-koszul-abstract}）。
换言之，存在正合复形
$$
0 \to
\pi^*(\wedge^{n + 1}\mathcal{E})(-n - 1) \to
\ldots \to
\pi^*(\wedge^i\mathcal{E})(-i) \to
\ldots \to
\pi^*\mathcal{E}(-1) \to
\mathcal{O}_{\mathbf{P}(\mathcal{E})} \to 0
$$
将项 $\pi^*(\wedge^i\mathcal{E})(-i)$ 视为位于次数 $-i$。
使用《导出范畴》引理 \ref{derived-lemma-two-ss-complex-functor}
中的第一谱序列，计算这个无环复形的高阶直像。
具体而言，存在各项为
$$
E_1^{p, q} = R^q\pi_*\left(\pi^*(\wedge^{-p}\mathcal{E})(p)\right)
$$
且收敛到零的谱序列！由投影公式
（《上同调》引理 \ref{cohomology-lemma-projection-formula}），有
$$
E_1^{p, q} = \wedge^{-p} \mathcal{E} \otimes_{\mathcal{O}_S}
R^q\pi_*\left(\mathcal{O}_{\mathbf{P}(\mathcal{E})}(p)\right).
$$
注意，在 $S$ 上局部，层 $\mathcal{E}$ 平凡，
即同构于 $\mathcal{O}_S^{\oplus n + 1}$；故在 $S$ 上局部，
态射 $\mathbf{P}(\mathcal{E}) \to S$ 可与
$\mathbf{P}^n_S \to S$ 等同。因此在 $S$ 上局部，
可使用引理 \ref{coherent-lemma-cohomology-projective-space-over-ring}、
\ref{coherent-lemma-identify-functorially} 或
\ref{coherent-lemma-cohomology-projective-space-over-base} 的结果。
由此，$E_1^{p, q} = 0$，除非 $(p, q)$ 为 $(0, 0)$
或 $(-n - 1, n)$。非零项为
\begin{align*}
E_1^{0, 0} & = \pi_*\mathcal{O}_{\mathbf{P}(\mathcal{E})} = \mathcal{O}_S \\
E_1^{-n - 1, n} & =
R^n\pi_*\left(\pi^*(\wedge^{n + 1}\mathcal{E})(-n - 1)\right) =
\wedge^{n + 1}\mathcal{E} \otimes_{\mathcal{O}_S}
R^n\pi_*\left(\mathcal{O}_{\mathbf{P}(\mathcal{E})}(-n - 1)\right)
\end{align*}
因此谱序列中只能有一个非零微分，即映射
$d_{n + 1}^{-n - 1, n} : E_{n + 1}^{-n - 1, n} \to E_{n + 1}^{0, 0}$；
它必为同构（因为谱序列收敛到 $0$ 层）。
于是 $E_1^{p, q} = E_{n + 1}^{p, q}$，并得到典范同构
$$
\wedge^{n + 1}\mathcal{E} \otimes_{\mathcal{O}_S}
R^n\pi_*\left(\mathcal{O}_{\mathbf{P}(\mathcal{E})}(-n - 1)\right) =
R^n\pi_*\left(\pi^*(\wedge^{n + 1}\mathcal{E})(-n - 1)\right)
\xrightarrow{d_{n + 1}^{-n - 1, n}}
\mathcal{O}_S
$$
由于 $\wedge^{n + 1}\mathcal{E}$ 是可逆层，这推出
$R^n\pi_*\mathcal{O}_{\mathbf{P}(\mathcal{E})}(-n - 1)$ 也可逆，
并且典范同构于 $\wedge^{n + 1}\mathcal{E}$ 的逆。
换言之，已经证明引理中 $d = - n - 1$ 的情形。

\medskip\noindent
在 $S$ 上局部工作，由引理
\ref{coherent-lemma-cohomology-projective-space-over-ring}、
\ref{coherent-lemma-identify-functorially} 或
\ref{coherent-lemma-cohomology-projective-space-over-base} 的上同调计算，
立即得到引理中的消失陈述。此外，关于 $R^0\pi_*$ 的结果也很清楚，
因为存在典范映射
$\text{Sym}^d(\mathcal{E}) \to \pi_* \mathcal{O}_{\mathbf{P}(\mathcal{E})}(d)$
对所有 $d$ 都成立。
还须证明 $R^n\pi_*\mathcal{O}_{\mathbf{P}(\mathcal{E})}(d)$ 的描述
对 $d < -n - 1$ 正确。为此考虑映射
$$
\pi^*(\text{Sym}^{-d - n - 1}(\mathcal{E}))
\otimes_{\mathcal{O}_{\mathbf{P}(\mathcal{E})}}
\mathcal{O}_{\mathbf{P}(\mathcal{E})}(d)
\longrightarrow
\mathcal{O}_{\mathbf{P}(\mathcal{E})}(-n - 1)
$$
应用 $R^n\pi_*$ 及投影公式（见上文），得到映射
$$
\text{Sym}^{-d - n - 1}(\mathcal{E})
\otimes_{\mathcal{O}_S}
R^n\pi_*(\mathcal{O}_{\mathbf{P}(\mathcal{E})}(d))
\longrightarrow
R^n\pi_*\mathcal{O}_{\mathbf{P}(\mathcal{E})}(-n - 1) =
(\wedge^{n + 1}\mathcal{E})^{\otimes -1}
$$
（最后一个等式已在上面证明）。
再次利用引理
\ref{coherent-lemma-cohomology-projective-space-over-ring}、
\ref{coherent-lemma-identify-functorially} 或
\ref{coherent-lemma-cohomology-projective-space-over-base} 的局部计算，
可知此映射在
$R^n\pi_*(\mathcal{O}_{\mathbf{P}(\mathcal{E})}(d))$ 与
$\text{Sym}^{-d - n - 1}(\mathcal{E}) \otimes \wedge^{n + 1}(\mathcal{E})$
之间诱导完美配对，如所欲。
\end{proof}

















\section{局部 Noether 概形上的凝聚层}
\label{coherent-section-coherent-sheaves}

\noindent
我们已在《模》节 \ref{modules-section-coherent} 中定义了
任意环层空间上凝聚模的概念。
尽管可以考虑非 Noether 概形上的凝聚层，
在考察凝聚层时，我们总是假设基概形局部 Noether。
下面给出局部 Noether 概形上凝聚层的一个刻画。

\begin{lemma}
\label{coherent-lemma-coherent-Noetherian}
设 $X$ 为局部 Noether 概形。
设 $\mathcal{F}$ 为 $\mathcal{O}_X$-模。
下列条件等价：
\begin{enumerate}
\item $\mathcal{F}$ 凝聚；
\item $\mathcal{F}$ 是拟凝聚有限型 $\mathcal{O}_X$-模；
\item $\mathcal{F}$ 是有限表示 $\mathcal{O}_X$-模；
\item 对任意仿射开集 $\Spec(A) = U \subset X$，都有
$\mathcal{F}|_U = \widetilde M$，其中 $M$ 为有限 $A$-模；
\item 存在仿射开覆盖 $X = \bigcup U_i$，
$U_i = \Spec(A_i)$，使得每个
$\mathcal{F}|_{U_i} = \widetilde M_i$，其中 $M_i$ 为有限 $A_i$-模。
\end{enumerate}
特别地，$\mathcal{O}_X$ 凝聚；任意可逆 $\mathcal{O}_X$-模凝聚；
更一般地，任意有限局部自由 $\mathcal{O}_X$-模凝聚。
\end{lemma}

\begin{proof}
蕴含 (1) $\Rightarrow$ (2) 与 (1) $\Rightarrow$ (3) 一般成立，见
《模》引理 \ref{modules-lemma-coherent-finite-presentation}。
若 $\mathcal{F}$ 有限表示，则 $\mathcal{F}$ 拟凝聚；见
《模》引理 \ref{modules-lemma-finite-presentation-quasi-coherent}。
故 (3) $\Rightarrow$ (2) 也成立。

\medskip\noindent
假设 $\mathcal{F}$ 是拟凝聚有限型 $\mathcal{O}_X$-模。
由《性质》引理 \ref{properties-lemma-finite-type-module}，
在任意仿射开集 $\Spec(A) = U \subset X$ 上都有
$\mathcal{F}|_U = \widetilde M$，其中 $M$ 为有限 $A$-模。
由于 $A$ Noether，$M$ 有有限解消
$$
A^{\oplus m} \to A^{\oplus n} \to M \to 0.
$$
因此由《性质》引理 \ref{properties-lemma-finite-presentation-module}，
$\mathcal{F}$ 有限表示。换言之，(2) $\Rightarrow$ (3)。

\medskip\noindent
由《模》引理 \ref{modules-lemma-coherent-structure-sheaf}，
为证明 (3) 推出 (1)，只需证明 $\mathcal{O}_X$ 凝聚。
因此须证明：给定任意开集 $U \subset X$ 与任意有限组截面
$f_i \in \mathcal{O}_X(U)$，$i = 1, \ldots, n$，映射
$\bigoplus_{i = 1, \ldots, n} \mathcal{O}_U \to \mathcal{O}_U$
的核为有限型。有限型是局部性质，故只需在任意 $x \in U$ 的邻域中检验。
于是可假设 $U = \Spec(A)$ 仿射。此时
$f_1, \ldots, f_n \in A$ 是 $A$ 的元素。由于 $A$ Noether，见
《性质》引理 \ref{properties-lemma-locally-Noetherian}，其核 $K$，即映射
$\bigoplus_{i = 1, \ldots, n} A \to A$ 的核，是有限 $A$-模；
例如见《代数》引理
\ref{algebra-lemma-Noetherian-basic}。
由于函子\ $\widetilde{ }$\ 正合，见《概形》引理
\ref{schemes-lemma-spec-sheaves}，得到正合序列
$$
\widetilde K \to
\bigoplus\nolimits_{i = 1, \ldots, n} \mathcal{O}_U \to
\mathcal{O}_U
$$
再由《性质》引理 \ref{properties-lemma-finite-type-module}，
可见 $\widetilde K$ 为有限型。因此 (1)、(2)、(3) 全都等价。

\medskip\noindent
由《性质》引理 \ref{properties-lemma-finite-type-module}，
(2) 推出 (4)；(4) 推出 (5) 是平凡的。
《概形》节 \ref{schemes-section-quasi-coherent} 中的讨论表明，
(5) 推出 $\mathcal{F}$ 拟凝聚；显然，(5) 还推出
$\mathcal{F}$ 为有限型。因此 (5) 推出 (2)，证毕。
\end{proof}

\begin{lemma}
\label{coherent-lemma-coherent-abelian-Noetherian}
设 $X$ 为局部 Noether 概形。
凝聚 $\mathcal{O}_X$-模所成的范畴是阿贝尔范畴。
更确切地，凝聚 $\mathcal{O}_X$-模之间映射的核与余核均凝聚。
凝聚层的任意扩张均凝聚。
\end{lemma}

\begin{proof}
这是《模》引理 \ref{modules-lemma-coherent-abelian}
在这一特殊情形下的重述。
\end{proof}

\noindent
下列引理对于凝聚 $\mathcal{O}_X$-模的范畴，
在一般环层空间 $X$ 上并不总是成立。

\begin{lemma}
\label{coherent-lemma-coherent-Noetherian-quasi-coherent-sub-quotient}
设 $X$ 为局部 Noether 概形。
设 $\mathcal{F}$ 为凝聚 $\mathcal{O}_X$-模。
$\mathcal{F}$ 的任意拟凝聚子模均凝聚。
$\mathcal{F}$ 的任意拟凝聚商模均凝聚。
\end{lemma}

\begin{proof}
可假设 $X$ 仿射，设 $X = \Spec(A)$。
《性质》引理 \ref{properties-lemma-locally-Noetherian} 推出 $A$ Noether。
引理 \ref{coherent-lemma-coherent-Noetherian} 将问题化为代数。
该引理的代数对应是：有限 $A$-模的商或子模仍为有限 $A$-模；
例如见《代数》引理 \ref{algebra-lemma-Noetherian-basic}。
\end{proof}

\begin{lemma}
\label{coherent-lemma-tensor-hom-coherent}
设 $X$ 为局部 Noether 概形。
设 $\mathcal{F}$、$\mathcal{G}$ 为凝聚 $\mathcal{O}_X$-模。
则 $\mathcal{O}_X$-模 $\mathcal{F} \otimes_{\mathcal{O}_X} \mathcal{G}$
与 $\SheafHom_{\mathcal{O}_X}(\mathcal{F}, \mathcal{G})$ 都凝聚。
\end{lemma}

\begin{proof}
《模》引理 \ref{modules-lemma-internal-hom-locally-kernel-direct-sum} 已证明
$\SheafHom_{\mathcal{O}_X}(\mathcal{F}, \mathcal{G})$ 凝聚。
张量积的结论是《模》引理 \ref{modules-lemma-tensor-product-permanence}。
\end{proof}

\begin{lemma}
\label{coherent-lemma-local-isomorphism}
设 $X$ 为局部 Noether 概形。
设 $\mathcal{F}$、$\mathcal{G}$ 为凝聚 $\mathcal{O}_X$-模。
设 $\varphi : \mathcal{G} \to \mathcal{F}$ 为 $\mathcal{O}_X$-模同态。
设 $x \in X$。
\begin{enumerate}
\item 若 $\mathcal{F}_x = 0$，则存在开邻域 $U \subset X$，
它是 $x$ 的邻域，且 $\mathcal{F}|_U = 0$。
\item 若 $\varphi_x : \mathcal{G}_x \to \mathcal{F}_x$ 单射，
则存在开邻域 $U \subset X$，它是 $x$ 的邻域，且 $\varphi|_U$ 单射。
\item 若 $\varphi_x : \mathcal{G}_x \to \mathcal{F}_x$ 满射，
则存在开邻域 $U \subset X$，它是 $x$ 的邻域，且 $\varphi|_U$ 满射。
\item 若 $\varphi_x : \mathcal{G}_x \to \mathcal{F}_x$ 双射，
则存在开邻域 $U \subset X$，它是 $x$ 的邻域，且 $\varphi|_U$ 为同构。
\end{enumerate}
\end{lemma}

\begin{proof}
见《模》引理
\ref{modules-lemma-finite-type-surjective-on-stalk}、
\ref{modules-lemma-finite-type-stalk-zero} 及
\ref{modules-lemma-finite-type-to-coherent-injective-on-stalk}。
\end{proof}

\begin{lemma}
\label{coherent-lemma-map-stalks-local-map}
设 $X$ 为局部 Noether 概形。
设 $\mathcal{F}$、$\mathcal{G}$ 为凝聚 $\mathcal{O}_X$-模。
设 $x \in X$。
假设 $\psi : \mathcal{G}_x \to \mathcal{F}_x$ 是
$\mathcal{O}_{X, x}$-模映射。
则存在开邻域 $U \subset X$（它是 $x$ 的邻域）及映射
$\varphi : \mathcal{G}|_U \to \mathcal{F}|_U$，使 $\varphi_x = \psi$。
\end{lemma}

\begin{proof}
根据引理 \ref{coherent-lemma-coherent-Noetherian}，
这是《模》引理 \ref{modules-lemma-stalk-internal-hom} 的重述。
\end{proof}

\begin{lemma}
\label{coherent-lemma-coherent-support-closed}
设 $X$ 为局部 Noether 概形。设 $\mathcal{F}$ 为凝聚
$\mathcal{O}_X$-模。则 $\text{Supp}(\mathcal{F})$ 闭，且
$\mathcal{F}$ 来自 $\mathcal{F}$ 的概形论支集上的凝聚层；见
《态射》定义 \ref{morphisms-definition-scheme-theoretic-support}。
\end{lemma}

\begin{proof}
设 $i : Z \to X$ 为 $\mathcal{F}$ 的概形论支集，
并设 $\mathcal{G}$ 为 $Z$ 上的有限型拟凝聚层，使
$i_*\mathcal{G} \cong \mathcal{F}$。
由于 $Z = \text{Supp}(\mathcal{F})$，支集闭。
由《态射》引理 \ref{morphisms-lemma-immersion-locally-finite-type}
及 \ref{morphisms-lemma-finite-type-noetherian}，概形 $Z$ 局部 Noether。
最后，由引理 \ref{coherent-lemma-coherent-Noetherian}，
$\mathcal{G}$ 是凝聚 $\mathcal{O}_Z$-模。
\end{proof}

\begin{lemma}
\label{coherent-lemma-i-star-equivalence}
设 $i : Z \to X$ 为局部 Noether 概形的闭浸入。
设 $\mathcal{I} \subset \mathcal{O}_X$ 为切出 $Z$ 的拟凝聚理想层。
函子 $i_*$ 在凝聚 $\mathcal{O}_X$-模（由 $\mathcal{I}$ 零化）所成的范畴与
凝聚 $\mathcal{O}_Z$-模范畴之间诱导等价。
\end{lemma}

\begin{proof}
由《态射》引理 \ref{morphisms-lemma-i-star-equivalence}，该函子全忠实。
设 $\mathcal{F}$ 为凝聚 $\mathcal{O}_X$-模，并由 $\mathcal{I}$ 零化。
由《态射》引理 \ref{morphisms-lemma-i-star-equivalence}，
可写成 $\mathcal{F} = i_*\mathcal{G}$，其中 $\mathcal{G}$
为 $Z$ 上的某个拟凝聚层。由《模》引理
\ref{modules-lemma-i-star-reflects-finite-type}，可见 $\mathcal{G}$ 有限型。
因此由引理 \ref{coherent-lemma-coherent-Noetherian}，$\mathcal{G}$ 凝聚。
故该函子也本质满，如所欲。
\end{proof}

\begin{lemma}
\label{coherent-lemma-finite-pushforward-coherent}
设 $f : X \to Y$ 为概形态射。
设 $\mathcal{F}$ 为拟凝聚 $\mathcal{O}_X$-模。
假设 $f$ 有限且 $Y$ 局部 Noether。
则 $R^pf_*\mathcal{F} = 0$ 对 $p > 0$ 成立，并且 $f_*\mathcal{F}$ 凝聚，
如果 $\mathcal{F}$ 凝聚。
\end{lemma}

\begin{proof}
高阶直像由引理 \ref{coherent-lemma-relative-affine-vanishing} 消失，
且有限态射按定义仿射。
注意，这些假设还推出 $X$ 局部 Noether
（见《态射》引理 \ref{morphisms-lemma-finite-type-noetherian}），
故陈述有意义。
设 $\Spec(A) = V \subset Y$ 为仿射开子集。
由《态射》定义 \ref{morphisms-definition-integral}，
有 $f^{-1}(V) = \Spec(B)$，其中 $A \to B$ 有限。
引理 \ref{coherent-lemma-coherent-Noetherian} 将引理的陈述化为如下代数事实：
若 $M$ 是有限 $B$-模，则把 $M$ 看作 $A$-模时也有限；
见《代数》引理 \ref{algebra-lemma-finite-module-over-finite-extension}。
\end{proof}

\noindent
在该引理的情形下，高阶直像也凝聚，因为它们消失。
我们将证明，对局部 Noether 概形之间的任意固有态射，
这一结论总是成立（命题 \ref{coherent-proposition-proper-pushforward-coherent}）。

\begin{lemma}
\label{coherent-lemma-coherent-support-dimension-0}
设 $X$ 为局部 Noether 概形。设 $\mathcal{F}$ 为凝聚层，且
$\dim(\text{Supp}(\mathcal{F})) \leq 0$。
则 $\mathcal{F}$ 由全局截面生成，并且
$H^i(X, \mathcal{F}) = 0$ 对 $i > 0$ 成立。
\end{lemma}

\begin{proof}
由引理 \ref{coherent-lemma-coherent-support-closed}，可写
$\mathcal{F} = i_*\mathcal{G}$，其中 $i : Z \to X$ 是
$\mathcal{F}$ 的概形论支集的包含，而 $\mathcal{G}$ 是凝聚
$\mathcal{O}_Z$-模。由于 $Z$ 的维数为 $0$，
可知 $Z$ 是仿射概形的不交并（《性质》引理
\ref{properties-lemma-locally-Noetherian-dimension-0}）。
因此 $\mathcal{G}$ 全局生成，且 $\mathcal{G}$ 的高阶上同调群为零
（引理 \ref{coherent-lemma-quasi-coherent-affine-cohomology-zero}）。
故 $\mathcal{F} = i_*\mathcal{G}$ 全局生成。
由于 $\mathcal{F}$ 与 $\mathcal{G}$ 的上同调相同
（闭浸入仿射，故可应用引理 \ref{coherent-lemma-relative-affine-cohomology}），
可知 $\mathcal{F}$ 的高阶上同调群为零。
\end{proof}

\begin{lemma}
\label{coherent-lemma-pushforward-coherent-on-open}
设 $X$ 为概形。设 $j : U \to X$ 为开集的包含。
设 $T \subset X$ 为包含于 $U$ 的闭子集。
若 $\mathcal{F}$ 是凝聚 $\mathcal{O}_U$-模，且
$\text{Supp}(\mathcal{F}) \subset T$，则
$j_*\mathcal{F}$ 是凝聚 $\mathcal{O}_X$-模。
\end{lemma}

\begin{proof}
考虑开覆盖 $X = U \cup (X \setminus T)$。
则 $j_*\mathcal{F}|_U = \mathcal{F}$ 凝聚，
且 $j_*\mathcal{F}|_{X \setminus T} = 0$ 也凝聚。
故 $j_*\mathcal{F}$ 凝聚。
\end{proof}

\section{Noether 概形上的凝聚层}
\label{coherent-section-coherent-quasi-compact}

\noindent
本节列举 Noether 概形上凝聚层的一些性质。

\begin{lemma}
\label{coherent-lemma-acc-coherent}
设 $X$ 为 Noether 概形。
设 $\mathcal{F}$ 为凝聚 $\mathcal{O}_X$-模。
$\mathcal{F}$ 的准凝聚子模满足升链条件。换言之，给定准凝聚子模的任意序列
$$
\mathcal{F}_1 \subset \mathcal{F}_2 \subset \ldots \subset \mathcal{F}
$$
则有 $\mathcal{F}_n = \mathcal{F}_{n + 1} = \ldots $
对某个 $n \geq 0$ 成立。
\end{lemma}

\begin{proof}
选取一个有限仿射开覆盖。
在覆盖的每个成员上，由《代数》引理
\ref{algebra-lemma-Noetherian-basic} 可知该链稳定。
故引理成立。
\end{proof}

\begin{lemma}
\label{coherent-lemma-power-ideal-kills-sheaf}
设 $X$ 为 Noether 概形。
设 $\mathcal{F}$ 为 $X$ 上的凝聚层。
设 $\mathcal{I} \subset \mathcal{O}_X$ 为对应于闭子概形
$Z \subset X$ 的准凝聚理想层。
则存在某个 $n \geq 0$ 使得 $\mathcal{I}^n\mathcal{F} = 0$，
当且仅当 $\text{Supp}(\mathcal{F}) \subset Z$（在集合论意义下）。
\end{lemma}

\begin{proof}
由于 $X$ 可由有限多个 Noether 环的谱覆盖，结论立即由
《代数》引理 \ref{algebra-lemma-Noetherian-power-ideal-kills-module}
推出。
\end{proof}

\begin{lemma}[Artin--Rees]
\label{coherent-lemma-Artin-Rees}
设 $X$ 为 Noether 概形。
设 $\mathcal{F}$ 为 $X$ 上的凝聚层。
设 $\mathcal{G} \subset \mathcal{F}$ 为准凝聚子层。
设 $\mathcal{I} \subset \mathcal{O}_X$ 为准凝聚理想层。
则存在 $c \geq 0$，使得对所有 $n \geq c$ 都有
$$
\mathcal{I}^{n - c}(\mathcal{I}^c\mathcal{F} \cap \mathcal{G})
=
\mathcal{I}^n\mathcal{F} \cap \mathcal{G}
$$
\end{lemma}

\begin{proof}
由于 $X$ 可由有限多个 Noether 环的谱覆盖，结论立即由
《代数》引理 \ref{algebra-lemma-Artin-Rees} 推出。
\end{proof}

\begin{lemma}
\label{coherent-lemma-directed-colimit-coherent}
设 $X$ 为 Noether 概形。每个准凝聚 $\mathcal{O}_X$-模
都是其凝聚子模的滤过余极限。
\end{lemma}

\begin{proof}
本结论只是《性质》引理
\ref{properties-lemma-quasi-coherent-colimit-finite-type}
的另一种表述；这里用到由
引理 \ref{coherent-lemma-coherent-Noetherian} 可知有限型准凝聚
$\mathcal{O}_X$-模是凝聚的这一事实。
\end{proof}

\begin{lemma}
\label{coherent-lemma-homs-over-open}
设 $X$ 为 Noether 概形。
设 $\mathcal{F}$ 为准凝聚 $\mathcal{O}_X$-模。
设 $\mathcal{G}$ 为凝聚 $\mathcal{O}_X$-模。
设 $\mathcal{I} \subset \mathcal{O}_X$ 为准凝聚理想层。
以 $Z \subset X$ 表示相应的闭子概形，并令
$U = X \setminus Z$。存在典范同构
$$
\colim_n \Hom_{\mathcal{O}_X}(\mathcal{I}^n\mathcal{G}, \mathcal{F})
\longrightarrow
\Hom_{\mathcal{O}_U}(\mathcal{G}|_U, \mathcal{F}|_U).
$$
特别地，存在同构
$$
\colim_n \Hom_{\mathcal{O}_X}(
\mathcal{I}^n, \mathcal{F})
\longrightarrow
\Gamma(U, \mathcal{F}).
$$
\end{lemma}

\begin{proof}
先证明第二个映射是同构。由《性质》引理
\ref{properties-lemma-sections-over-quasi-compact-open}，
它是单射。由于 $\mathcal{F}$ 是其凝聚子模的并，参见
《性质》引理 \ref{properties-lemma-quasi-coherent-colimit-finite-type}
（以及引理 \ref{coherent-lemma-coherent-Noetherian}），
为证明满射性，可以并且确实假设 $\mathcal{F}$ 凝聚。
以 $\mathcal{F}_n$ 表示 $\mathcal{F}$ 中由
$\mathcal{I}^n$ 零化的截面组成的准凝聚子层，参见
《性质》引理 \ref{properties-lemma-sections-over-quasi-compact-open}。
由于 $\mathcal{F}_1 \subset \mathcal{F}_2 \subset \ldots$，
可知 $\mathcal{F}_n = \mathcal{F}_{n + 1} = \ldots $
对某个 $n \geq 0$ 成立，这由引理
\ref{coherent-lemma-acc-coherent} 推出。
令 $\mathcal{H} = \mathcal{F}_n$，其中 $n$ 为上述整数。由 Artin--Rees
（引理 \ref{coherent-lemma-Artin-Rees}），存在 $c \geq 0$ 使得
$\mathcal{I}^m\mathcal{F} \cap \mathcal{H}
\subset \mathcal{I}^{m - c}\mathcal{H}$。
取 $m = n + c$，得到
$\mathcal{I}^m\mathcal{F} \cap \mathcal{H} \subset \mathcal{I}^n\mathcal{H}
= 0$。因此，若令 $\mathcal{F}' = \mathcal{I}^m\mathcal{F}$，
则有 $\mathcal{F}' \cap \mathcal{F}_n = 0$ 且
$\mathcal{F}'|_U = \mathcal{F}|_U$。特别注意到，子层
$(\mathcal{F}')_N$，即由 $\mathcal{I}^N$ 零化的截面所成子层，
对所有 $N \geq 0$ 均为零。因此由《性质》引理
\ref{properties-lemma-sections-over-quasi-compact-open}
可知，下列交换图中的上方水平箭头是双射：
$$
\xymatrix{
\colim_n \Hom_{\mathcal{O}_X}(
\mathcal{I}^n, \mathcal{F}')
\ar[r] \ar[d] &
\Gamma(U, \mathcal{F}') \ar[d] \\
\colim_n \Hom_{\mathcal{O}_X}(
\mathcal{I}^n, \mathcal{F})
\ar[r] &
\Gamma(U, \mathcal{F})
}
$$
右侧竖直箭头也是双射，故下方水平箭头如所要求是满射。

\medskip\noindent
接下来证明引理中的第一个箭头是双射。
由引理 \ref{coherent-lemma-coherent-Noetherian}，层 $\mathcal{G}$
是有限表示的，因而层
$\mathcal{H} = \SheafHom_{\mathcal{O}_X}(\mathcal{G}, \mathcal{F})$
是准凝聚的，参见《概形》第
\ref{schemes-section-quasi-coherent} 节。按定义有
$$
\mathcal{H}(U)
=
\Hom_{\mathcal{O}_U}(\mathcal{G}|_U, \mathcal{F}|_U)
$$
在引理第一个箭头的右端选取 $\psi$，即
$\psi \in \mathcal{H}(U)$。刚刚证明的结果适用于
$\mathcal{H}$，故存在 $n \geq 0$ 以及
$\varphi : \mathcal{I}^n \to \mathcal{H}$，它恢复
$\psi$（限制到 $U$ 后）。由《模》引理
\ref{modules-lemma-internal-hom}，$\varphi$ 对应于映射
$$
\varphi :
\mathcal{I}^n \otimes_{\mathcal{O}_X} \mathcal{G}
\longrightarrow
\mathcal{F}.
$$
这几乎就是所需结论，唯一差别是箭头的源是
$\mathcal{I}^n$ 与 $\mathcal{G}$ 的张量积，而非它们的积。
我们将说明，以增大 $n$ 为代价，这一区别无关紧要。
考虑短正合列
$$
0 \to \mathcal{K} \to
\mathcal{I}^n \otimes_{\mathcal{O}_X} \mathcal{G} \to
\mathcal{I}^n\mathcal{G} \to 0
$$
其中 $\mathcal{K}$ 定义为核。注意
$\mathcal{I}^n\mathcal{K} = 0$（证明略）。
再次应用 Artin--Rees 可知，对某个充分大的 $m$ 有
$$
\mathcal{K}
\cap
\mathcal{I}^m(\mathcal{I}^n \otimes_{\mathcal{O}_X} \mathcal{G})
=
0
$$
换言之，映射
$$
\mathcal{I}^m(\mathcal{I}^n \otimes_{\mathcal{O}_X} \mathcal{G})
\longrightarrow
\mathcal{I}^{n + m}\mathcal{G}
$$
是同构。设 $\varphi'$ 为 $\varphi$ 在该子模上的限制，
将它视为映射
$\mathcal{I}^{m + n}\mathcal{G} \to \mathcal{F}$。
于是 $\varphi'$ 给出引理第一个箭头左端的一个元素，
该元素经此箭头映为 $\psi$。换言之，我们已经证明该箭头满射。
其单射性的证明从略。
\end{proof}

\begin{lemma}
\label{coherent-lemma-extend-coherent}
设 $X$ 为局部 Noether 概形。设 $\mathcal{F}$、$\mathcal{G}$
为凝聚 $\mathcal{O}_X$-模。设 $U \subset X$ 为开集，并设
$\varphi : \mathcal{F}|_U \to \mathcal{G}|_U$ 为
$\mathcal{O}_U$-模映射。则存在凝聚子模
$\mathcal{F}' \subset \mathcal{F}$，它与
$\mathcal{F}$ 在 $U$ 上相同，且 $\varphi$ 可延拓为
$\varphi' : \mathcal{F}' \to \mathcal{G}$。
\end{lemma}

\begin{proof}
设 $\mathcal{I} \subset \mathcal{O}_X$ 为定义
$X \setminus U$ 上既约诱导概形结构的凝聚理想层。
若 $X$ 是 Noether 的，则引理 \ref{coherent-lemma-homs-over-open}
表明可取 $\mathcal{F}' = \mathcal{I}^n\mathcal{F}$，
其中 $n$ 为某个整数。
一般情形将借助 Zorn 引理由此推出。

\medskip\noindent
考虑所有三元组 $(U', \mathcal{F}', \varphi')$ 所成的集合，其中
$U \subset U' \subset X$ 为开集，
$\mathcal{F}' \subset \mathcal{F}|_{U'}$
是一个与 $\mathcal{F}$ 在 $U$ 上相同的凝聚子层，而
$\varphi' : \mathcal{F}' \to \mathcal{G}|_{U'}$
限制为 $\varphi$（在 $U$ 上）。我们约定
$(U'', \mathcal{F}'', \varphi'') \geq (U', \mathcal{F}', \varphi')$
当且仅当 $U'' \supset U'$、
$\mathcal{F}''|_{U'} = \mathcal{F}'$ 且
$\varphi''|_{U'} = \varphi'$。
显然，若有三元组的一个全序族
$(U_i, \mathcal{F}_i, \varphi_i)$，则可将
$\mathcal{F}_i$ 粘合成子层 $\mathcal{F}'$，它是
$\mathcal{F}$ 在 $U' = \bigcup U_i$ 上的子层，
并将 $\varphi$ 延拓为映射
$\varphi' : \mathcal{F}' \to \mathcal{G}|_{U'}$。
因此三元组的每个全序子集都有上界。
最后，设 $(U', \mathcal{F}', \varphi')$ 为任意三元组，
但 $U' \not = X$。则可选取仿射开集
$W \subset X$，它不包含于 $U'$。由第一段的结果，可以把子层
$\mathcal{F}'|_{W \cap U'}$ 及限制
$\varphi'|_{W \cap U'}$ 延拓为某个子层
$\mathcal{F}'' \subset \mathcal{F}|_W$ 及映射
$\varphi'' : \mathcal{F}'' \to \mathcal{G}|_W$。
当然，$(\mathcal{F}', \varphi')$ 与
$(\mathcal{F}'', \varphi'')$ 在 $W \cap U'$ 上相合，
恰好意味着可以将其延拓成三元组
$(U' \cup W, \mathcal{F}''', \varphi''')$。
因此任意极大三元组 $(U', \mathcal{F}', \varphi')$
（其存在性由 Zorn 引理保证）必有 $U' = X$，证明完毕。
\end{proof}

\section{深度}
\label{coherent-section-depth}

\noindent
本节简要讨论局部 Noether 概形上凝聚模的深度及性质
$(S_k)$。注意，我们已经在《性质》第
\ref{properties-section-Rk} 节讨论过局部 Noether 概形的这一概念。

\begin{definition}
\label{coherent-definition-depth}
设 $X$ 为局部 Noether 概形。
设 $\mathcal{F}$ 为凝聚 $\mathcal{O}_X$-模。
设 $k \geq 0$ 为整数。
\begin{enumerate}
\item 称 $\mathcal{F}$ 在一点处有{\it 深度 $k$}，若该点
$x$ 属于 $X$ 且
$\text{depth}_{\mathcal{O}_{X, x}}(\mathcal{F}_x) = k$。
\item 称 $X$ 在一点处有{\it 深度 $k$}，若该点
$x$ 属于 $X$ 且
$\text{depth}(\mathcal{O}_{X, x}) = k$。
\item 若 $\mathcal{F}$ 满足性质 {\it $(S_k)$}，是指
$$
\text{depth}_{\mathcal{O}_{X, x}}(\mathcal{F}_x)
\geq \min(k, \dim(\text{Supp}(\mathcal{F}_x)))
$$
对所有 $x \in X$ 成立。
\item 称 $X$ 满足性质 {\it $(S_k)$}，若 $\mathcal{O}_X$
满足性质 $(S_k)$。
\end{enumerate}
\end{definition}

\noindent
任意凝聚层都满足条件 $(S_0)$。
条件 $(S_1)$ 等价于没有嵌入伴随点，参见《除子》引理
\ref{divisors-lemma-S1-no-embedded}。

\begin{lemma}
\label{coherent-lemma-hom-into-depth}
设 $X$ 为局部 Noether 概形。设 $\mathcal{F}$、$\mathcal{G}$
为凝聚 $\mathcal{O}_X$-模，且 $x \in X$。
\begin{enumerate}
\item 若 $\mathcal{G}_x$ 的深度 $\geq 1$，则
$\SheafHom_{\mathcal{O}_X}(\mathcal{F}, \mathcal{G})_x$
的深度 $\geq 1$。
\item 若 $\mathcal{G}_x$ 的深度 $\geq 2$，则
$\Hom_{\mathcal{O}_X}(\mathcal{F}, \mathcal{G})_x$ 的深度 $\geq 2$。
\end{enumerate}
\end{lemma}

\begin{proof}
由引理 \ref{coherent-lemma-tensor-hom-coherent}，
$\SheafHom_{\mathcal{O}_X}(\mathcal{F}, \mathcal{G})$
是凝聚 $\mathcal{O}_X$-模。
凝聚模是有限表示的（引理 \ref{coherent-lemma-coherent-Noetherian}），
故取层茎与取 $\SheafHom$ 及 $\Hom$ 可交换，参见《模》引理
\ref{modules-lemma-stalk-internal-hom}。
因而可归结为局部环上的有限模情形，即《代数进阶》引理
\ref{more-algebra-lemma-hom-into-depth}。
\end{proof}

\begin{lemma}
\label{coherent-lemma-hom-into-S2}
设 $X$ 为局部 Noether 概形。设 $\mathcal{F}$、$\mathcal{G}$
为凝聚 $\mathcal{O}_X$-模。
\begin{enumerate}
\item 若 $\mathcal{G}$ 满足性质 $(S_1)$，则
$\SheafHom_{\mathcal{O}_X}(\mathcal{F}, \mathcal{G})$ 满足性质 $(S_1)$。
\item 若 $\mathcal{G}$ 满足性质 $(S_2)$，则
$\SheafHom_{\mathcal{O}_X}(\mathcal{F}, \mathcal{G})$ 满足性质 $(S_2)$。
\end{enumerate}
\end{lemma}

\begin{proof}
立即由引理 \ref{coherent-lemma-hom-into-depth} 和定义推出。
\end{proof}

\noindent
我们已在《性质》引理
\ref{properties-lemma-scheme-CM-iff-all-Sk} 中看到，
局部 Noether 概形是 Cohen--Macaulay 的，当且仅当
$(S_k)$ 对所有 $k$ 均成立。
因此引入下述定义是合理的；它等价于所有层茎均为
Cohen--Macaulay 模这一条件。

\begin{definition}
\label{coherent-definition-Cohen-Macaulay}
设 $X$ 为局部 Noether 概形。
设 $\mathcal{F}$ 为凝聚 $\mathcal{O}_X$-模。
称 $\mathcal{F}$ 是 {\it Cohen--Macaulay 的}，当且仅当
$(S_k)$ 对所有 $k \geq 0$ 均成立。
\end{definition}

\begin{lemma}
\label{coherent-lemma-Cohen-Macaulay-over-regular}
设 $X$ 为正则概形。设 $\mathcal{F}$ 为凝聚
$\mathcal{O}_X$-模。下列条件等价：
\begin{enumerate}
\item $\mathcal{F}$ 是 Cohen--Macaulay 的，且
$\text{Supp}(\mathcal{F}) = X$；
\item $\mathcal{F}$ 是秩 $> 0$ 的有限局部自由模。
\end{enumerate}
\end{lemma}

\begin{proof}
设 $x \in X$。若 (2) 成立，则 $\mathcal{F}_x$ 是自由
$\mathcal{O}_{X, x}$-模，且其秩 $> 0$。因此
$\text{depth}(\mathcal{F}_x) = \dim(\mathcal{O}_{X, x})$，
因为正则局部环是 Cohen--Macaulay 的（《代数》引理
\ref{algebra-lemma-regular-ring-CM}）。
反之，若 (1) 成立，则 $\mathcal{F}_x$ 是
$\mathcal{O}_{X, x}$ 上的极大 Cohen--Macaulay 模
（《代数》定义 \ref{algebra-definition-maximal-CM}）。
故由《代数》引理 \ref{algebra-lemma-regular-mcm-free}，
$\mathcal{F}_x$ 是自由的。
\end{proof}

\section{凝聚层的逐层剥离}
\label{coherent-section-devissage}

\noindent
设 $X$ 为 Noether 概形。考虑一个整的闭子概形
$i : Z \to X$。考察形如 $i_*\mathcal{G}$ 的凝聚层往往很方便，
其中 $\mathcal{G}$ 是 $Z$ 上的凝聚层。我们尤其关注
$\mathcal{G}$ 为无挠秩 $1$ 层的情形。例如，$\mathcal{G}$
可以是 $Z$ 上的非零理想层；更特殊地，也可取
$\mathcal{G} = \mathcal{O}_Z$。

\medskip\noindent
本节始终会用到：凝聚层与有限型准凝聚层是同一回事，
凝聚层的准凝聚子商仍然凝聚；参见第
\ref{coherent-section-coherent-sheaves} 节。
凝聚层的支集是闭的；参见《模》引理
\ref{modules-lemma-support-finite-type-closed}。

\begin{lemma}
\label{coherent-lemma-prepare-filter-support}
设 $X$ 为 Noether 概形。
设 $\mathcal{F}$ 为 $X$ 上的凝聚层。
假设 $\text{Supp}(\mathcal{F}) = Z \cup Z'$，其中
$Z$、$Z'$ 闭。则存在凝聚层的短正合列
$$
0 \to \mathcal{G}' \to \mathcal{F} \to \mathcal{G} \to 0
$$
使得 $\text{Supp}(\mathcal{G}') \subset Z'$ 且
$\text{Supp}(\mathcal{G}) \subset Z$。
\end{lemma}

\begin{proof}
设 $\mathcal{I} \subset \mathcal{O}_X$ 为定义
$Z$ 上既约诱导闭子概形结构的理想层，参见《概形》引理
\ref{schemes-lemma-reduced-closed-subscheme}。
考虑子层
$\mathcal{G}'_n = \mathcal{I}^n\mathcal{F}$ 以及商层
$\mathcal{G}_n = \mathcal{F}/\mathcal{I}^n\mathcal{F}$。
对每个 $n$，都有短正合列
$$
0 \to \mathcal{G}'_n \to \mathcal{F} \to \mathcal{G}_n \to 0
$$
对每个点 $x$（其属于 $Z' \setminus Z$），有
$\mathcal{I}_x = \mathcal{O}_{X, x}$，从而
$\mathcal{G}_{n, x} = 0$。故
$\text{Supp}(\mathcal{G}_n) \subset Z$。注意
$X \setminus Z'$ 是 Noether 概形。因此由引理
\ref{coherent-lemma-power-ideal-kills-sheaf}，存在 $n$ 使得
$\mathcal{G}'_n|_{X \setminus Z'} =
\mathcal{I}^n\mathcal{F}|_{X \setminus Z'} = 0$。
对这样的 $n$，有
$\text{Supp}(\mathcal{G}'_n) \subset Z'$。
因此取 $\mathcal{G}' = \mathcal{G}'_n$ 且
$\mathcal{G} = \mathcal{G}_n$ 即可。
\end{proof}

\begin{lemma}
\label{coherent-lemma-prepare-filter-irreducible}
设 $X$ 为 Noether 概形。
设 $i : Z \to X$ 为整的闭子概形。
设 $\xi \in Z$ 为其泛点。
设 $\mathcal{F}$ 为 $X$ 上的凝聚层。
假设 $\mathcal{F}_\xi$ 被 $\mathfrak m_\xi$ 零化。
则存在整数 $r \geq 0$、凝聚理想层
$\mathcal{I} \subset \mathcal{O}_Z$ 以及凝聚层的单射
$$
i_*\left(\mathcal{I}^{\oplus r}\right) \to \mathcal{F}
$$
它在 $\xi$ 的一个邻域上是同构。
\end{lemma}

\begin{proof}
设 $\mathcal{J} \subset \mathcal{O}_X$ 为 $Z$ 的理想层。
设 $\mathcal{F}' \subset \mathcal{F}$ 为
$\mathcal{F}$ 中被 $\mathcal{J}$ 零化的局部截面所成的子层。
由《性质》引理
\ref{properties-lemma-sections-annihilated-by-ideal}，它是准凝聚层。
此外，$\mathcal{F}'_\xi = \mathcal{F}_\xi$，因为
$\mathcal{J}_\xi = \mathfrak m_\xi$，并可应用《性质》引理
\ref{properties-lemma-sections-annihilated-by-ideal} 的第 (3) 部分。
由引理 \ref{coherent-lemma-local-isomorphism} 可知，
$\mathcal{F}' \to \mathcal{F}$ 在 $\xi$ 的一个邻域上诱导同构。
因此可将 $\mathcal{F}$ 替换为 $\mathcal{F}'$，并假设
$\mathcal{F}$ 被 $\mathcal{J}$ 零化。

\medskip\noindent
假设 $\mathcal{J}\mathcal{F} = 0$。由引理
\ref{coherent-lemma-i-star-equivalence}，可写成
$\mathcal{F} = i_*\mathcal{G}$，其中 $\mathcal{G}$ 是
$Z$ 上的某个凝聚层。
假设可以找到态射
$\mathcal{I}^{\oplus r} \to \mathcal{G}$，它在泛点
$\xi$（即 $Z$ 的泛点）的一个邻域上是同构。于是应用 $i_*$
（它左正合）便得到引理结论。因此已经归结为
$X = Z$ 的情形。

\medskip\noindent
假设 $Z = X$ 是泛点为 $\xi$ 的整 Noether 概形。
注意，此时 $\mathcal{O}_{X, \xi} = \kappa(\xi)$ 是 $X$ 的函数域。
由于 $\mathcal{F}_\xi$ 是有限 $\mathcal{O}_\xi$-模，
可知 $r = \dim_{\kappa(\xi)} \mathcal{F}_\xi$ 有限。
因此层 $\mathcal{O}_X^{\oplus r}$ 与 $\mathcal{F}$
在 $\xi$ 处的层茎同构。
由引理 \ref{coherent-lemma-map-stalks-local-map}，存在非空开集
$U \subset X$ 及态射
$\psi : \mathcal{O}_X^{\oplus r}|_U \to \mathcal{F}|_U$，
它在 $\xi$ 处是同构，进而由引理
\ref{coherent-lemma-local-isomorphism} 可知，它在 $\xi$ 的一个邻域上是同构。
由《概形》引理
\ref{schemes-lemma-reduced-closed-subscheme}，存在准凝聚理想层
$\mathcal{I} \subset \mathcal{O}_X$，其相伴闭子概形
$Y \subset X$ 是 $U$ 的补集。
由引理 \ref{coherent-lemma-homs-over-open}，存在 $n \geq 0$ 及态射
$\mathcal{I}^n(\mathcal{O}_X^{\oplus r}) \to \mathcal{F}$，
它恢复 $\psi$（在 $U$ 上）。由于
$\mathcal{I}^n(\mathcal{O}_X^{\oplus r}) = (\mathcal{I}^n)^{\oplus r}$，
便得到引理所述的映射。该映射是单射，因为 $X$ 是整的，
而且它在 $X$ 的泛点处是单射（简单证明略）。
\end{proof}

\begin{lemma}
\label{coherent-lemma-coherent-filter}
设 $X$ 为 Noether 概形。
设 $\mathcal{F}$ 为 $X$ 上的凝聚层。
存在由凝聚子层组成的滤过
$$
0 = \mathcal{F}_0 \subset \mathcal{F}_1 \subset
\ldots \subset \mathcal{F}_m = \mathcal{F}
$$
使得对每个 $j = 1, \ldots, m$，都存在整的闭子概形
$Z_j \subset X$ 及非零凝聚理想层
$\mathcal{I}_j \subset \mathcal{O}_{Z_j}$，满足
$$
\mathcal{F}_j/\mathcal{F}_{j - 1}
\cong (Z_j \to X)_* \mathcal{I}_j
$$
\end{lemma}

\begin{proof}
考虑集合
$$
\mathcal{T} =
\left\{
\begin{matrix}
Z \subset X
\text{ closed such that there exists a coherent sheaf }
\mathcal{F} \\
\text{ 其中 }
\text{Supp}(\mathcal{F}) = Z
\text{ for which the lemma is wrong}
\end{matrix}
\right\}
$$
我们要证明 $\mathcal{T}$ 为空。若非如此，由于 $X$ 是 Noether 的，
可选取极小元 $Z \in \mathcal{T}$。这意味着存在
凝聚层 $\mathcal{F}$，它在 $X$ 上且支集为 $Z$，
而引理对它不成立。
显然 $Z \not = \emptyset$，因为支集为空的层只有零层，
而引理对零层成立（取 $m = 0$）。

\medskip\noindent
若 $Z$ 不可约，则可写成 $Z = Z_1 \cup Z_2$，
其中 $Z_1, Z_2$ 闭且严格小于 $Z$。
于是应用引理 \ref{coherent-lemma-prepare-filter-support}，得到凝聚层的短正合列
$$
0 \to
\mathcal{G}_1 \to
\mathcal{F} \to
\mathcal{G}_2 \to 0
$$
满足 $\text{Supp}(\mathcal{G}_i) \subset Z_i$。
由 $Z$ 的极小性，每个 $\mathcal{G}_i$
都有引理陈述中的滤过。考察 $\mathcal{F}$ 上的诱导滤过便得到矛盾。
因此 $Z$ 必不可约。

\medskip\noindent
假设 $Z$ 不可约。设 $\mathcal{J}$ 为定义 $Z$
的既约诱导闭子概形结构的理想层，参见《概形》引理
\ref{schemes-lemma-reduced-closed-subscheme}。
由引理 \ref{coherent-lemma-power-ideal-kills-sheaf} 可知，存在
$n \geq 0$ 使得 $\mathcal{J}^n\mathcal{F} = 0$。
因而得到滤过
$$
0 = \mathcal{J}^n\mathcal{F} \subset \mathcal{J}^{n - 1}\mathcal{F}
\subset \ldots \subset \mathcal{J}\mathcal{F} \subset \mathcal{F}
$$
其每个相继子商都被 $\mathcal{J}$ 零化。
因此，若这些子商各自都有引理陈述中的滤过，
则 $\mathcal{F}$ 也有。换言之，可以假设
$\mathcal{J}$ 确实零化 $\mathcal{F}$。

\medskip\noindent
在 $Z$ 可约且 $\mathcal{J}\mathcal{F} = 0$ 的情形，
可应用引理 \ref{coherent-lemma-prepare-filter-irreducible}。
这给出短正合列
$$
0 \to
i_*(\mathcal{I}^{\oplus r}) \to
\mathcal{F} \to
\mathcal{Q} \to 0
$$
其中 $\mathcal{Q}$ 定义为该商。
由刚引用的引理，$\mathcal{Q}$ 在 $\xi$ 的一个邻域上为零，
故 $\mathcal{Q}$ 的支集严格小于 $Z$。于是
$\mathcal{Q}$ 由 $Z$ 的极小性具有所需类型的滤过。
但这样 $\mathcal{F}$ 显然也有，得到最终矛盾。
\end{proof}

\begin{lemma}
\label{coherent-lemma-property-initial}
设 $X$ 为 Noether 概形。
设 $\mathcal{P}$ 为 $X$ 上凝聚层的一种性质。假设：
\begin{enumerate}
\item 对凝聚层的任意短正合列
$$
0 \to \mathcal{F}_1 \to \mathcal{F} \to \mathcal{F}_2 \to 0
$$
若 $\mathcal{F}_i$（$i = 1, 2$）具有性质
$\mathcal{P}$，则 $\mathcal{F}$ 也具有该性质。
\item 对每个整的闭子概形 $Z \subset X$ 以及每个准凝聚理想层
$\mathcal{I} \subset \mathcal{O}_Z$，性质 $\mathcal{P}$
对 $i_*\mathcal{I}$ 成立。
\end{enumerate}
则性质 $\mathcal{P}$ 对 $X$ 上的每个凝聚层均成立。
\end{lemma}

\begin{proof}
先注意，若 $\mathcal{F}$ 是凝聚层，并有凝聚子层的滤过
$$
0 = \mathcal{F}_0 \subset \mathcal{F}_1 \subset
\ldots \subset \mathcal{F}_m = \mathcal{F}
$$
使每个 $\mathcal{F}_i/\mathcal{F}_{i - 1}$
都具有性质 $\mathcal{P}$，则 $\mathcal{F}$ 也具有该性质。
这由 $\mathcal{P}$ 的性质 (1) 推出。
另一方面，由引理 \ref{coherent-lemma-coherent-filter}，
可将任意 $\mathcal{F}$ 滤过，使相继子商具有 (2) 中的形式。
故引理成立。
\end{proof}

\begin{lemma}
\label{coherent-lemma-property-irreducible}
设 $X$ 为 Noether 概形。设 $Z_0 \subset X$ 为泛点
$\xi$ 的不可约闭子集。设 $\mathcal{P}$ 为
$X$ 上支集包含于 $Z_0$ 的凝聚层的一种性质，并且：
\begin{enumerate}
\item 对凝聚层的任意短正合列，若三项中有两项具有性质
$\mathcal{P}$，则第三项也具有该性质。
\item 对每个整的闭子概形 $Z \subset Z_0 \subset X$，
$Z \not = Z_0$，以及每个准凝聚理想层
$\mathcal{I} \subset \mathcal{O}_Z$，性质
$\mathcal{P}$ 对 $(Z \to X)_*\mathcal{I}$ 成立。
\item 存在某个凝聚层 $\mathcal{G}$，它在 $X$ 上且满足：
\begin{enumerate}
\item $\text{Supp}(\mathcal{G}) = Z_0$；
\item $\mathcal{G}_\xi$ 被 $\mathfrak m_\xi$ 零化；
\item $\dim_{\kappa(\xi)} \mathcal{G}_\xi = 1$；并且
\item 性质 $\mathcal{P}$ 对 $\mathcal{G}$ 成立。
\end{enumerate}
\end{enumerate}
则性质 $\mathcal{P}$ 对每个凝聚层 $\mathcal{F}$ 均成立，
只要该层在 $X$ 上且其支集包含于 $Z_0$。
\end{lemma}

\begin{proof}
先注意，若 $\mathcal{F}$ 是支集包含于 $Z_0$ 的凝聚层，
并有凝聚子层的滤过
$$
0 = \mathcal{F}_0 \subset \mathcal{F}_1 \subset
\ldots \subset \mathcal{F}_m = \mathcal{F}
$$
使每个 $\mathcal{F}_i/\mathcal{F}_{i - 1}$
都具有性质 $\mathcal{P}$，则 $\mathcal{F}$ 也具有该性质。
或者，若 $\mathcal{F}$ 具有性质 $\mathcal{P}$，且除一个之外的所有
$\mathcal{F}_i/\mathcal{F}_{i - 1}$ 都具有性质
$\mathcal{P}$，则最后一个也具有该性质。这由假设 (1) 推出。

\medskip\noindent
作为第一个应用，可知支集严格包含于 $Z_0$ 的任意凝聚层
都具有性质 $\mathcal{P}$。事实上，这样的层有一个滤过
（参见引理 \ref{coherent-lemma-coherent-filter}），其子商依 (2)
均具有性质 $\mathcal{P}$。

\medskip\noindent
设 $\mathcal{G}$ 如 (3) 所述。由引理
\ref{coherent-lemma-prepare-filter-irreducible}，存在理想层
$\mathcal{I}$（定义在 $Z_0$ 上）、整数 $r \geq 1$ 以及短正合列
$$
0 \to
\left((Z_0 \to X)_*\mathcal{I}\right)^{\oplus r} \to
\mathcal{G} \to
\mathcal{Q} \to 0
$$
其中 $\mathcal{Q}$ 的支集严格包含于 $Z_0$。
由 (3)(c) 可知 $r = 1$。由于 $\mathcal{Q}$
也具有性质 $\mathcal{P}$，可知
$(Z_0 \to X)_*\mathcal{I}$ 具有性质 $\mathcal{P}$。

\medskip\noindent
接着，设 $\mathcal{I}' \not = 0$ 是 $Z_0$
上的另一个准凝聚理想层。考虑交
$\mathcal{I}'' = \mathcal{I}' \cap \mathcal{I}$，得到两个短正合列
$$
0 \to
(Z_0 \to X)_*\mathcal{I}'' \to
(Z_0 \to X)_*\mathcal{I} \to
\mathcal{Q} \to 0
$$
和
$$
0 \to
(Z_0 \to X)_*\mathcal{I}'' \to
(Z_0 \to X)_*\mathcal{I}' \to
\mathcal{Q}' \to 0.
$$
注意，凝聚层 $\mathcal{Q}$ 与
$\mathcal{Q}'$ 的支集都严格包含于 $Z_0$。
因此 $\mathcal{Q}$ 与 $\mathcal{Q}'$ 均具有性质
$\mathcal{P}$（见上）。于是利用 (1) 可知
$(Z_0 \to X)_*\mathcal{I}''$ 与
$(Z_0 \to X)_*\mathcal{I}'$ 也都具有 $\mathcal{P}$。

\medskip\noindent
证明的最后一步是注意到，任意凝聚层 $\mathcal{F}$
只要在 $X$ 上且支集包含于 $Z_0$，便都有一个滤过（再次参见引理
\ref{coherent-lemma-coherent-filter}），依刚才所述，其子商均具有性质
$\mathcal{P}$。
\end{proof}

\begin{lemma}
\label{coherent-lemma-property}
设 $X$ 为 Noether 概形。
设 $\mathcal{P}$ 为 $X$ 上凝聚层的一种性质，并且：
\begin{enumerate}
\item 对凝聚层的任意短正合列，若三项中有两项具有性质
$\mathcal{P}$，则第三项也具有该性质。
\item 对每个整的闭子概形 $Z \subset X$（其泛点为
$\xi$），存在某个凝聚层 $\mathcal{G}$，满足：
\begin{enumerate}
\item $\text{Supp}(\mathcal{G}) = Z$；
\item $\mathcal{G}_\xi$ 被 $\mathfrak m_\xi$ 零化；
\item $\dim_{\kappa(\xi)} \mathcal{G}_\xi = 1$；并且
\item 性质 $\mathcal{P}$ 对 $\mathcal{G}$ 成立。
\end{enumerate}
\end{enumerate}
则性质 $\mathcal{P}$ 对 $X$ 上的每个凝聚层均成立。
\end{lemma}

\begin{proof}
依引理 \ref{coherent-lemma-property-initial}，只需证明：
对所有整的闭子概形 $Z \subset X$ 及所有准凝聚理想层
$\mathcal{I} \subset \mathcal{O}_Z$，性质 $\mathcal{P}$
对 $(Z \to X)_*\mathcal{I}$ 成立。若非如此，则由于
$X$ 是 Noether 的，存在极小的整闭子概形
$Z_0 \subset X$，使得性质 $\mathcal{P}$ 对
$(Z_0 \to X)_*\mathcal{I}_0$ 不成立，其中
$\mathcal{I}_0 \subset \mathcal{O}_{Z_0}$ 是某个准凝聚理想层；
但性质 $\mathcal{P}$ 对 $(Z \to X)_*\mathcal{I}$
成立，只要 $Z \subset Z_0$ 是整的闭子概形、
$Z \not = Z_0$ 且
$\mathcal{I} \subset \mathcal{O}_Z$ 是准凝聚理想层。
由第 (2) 部分，存在 $\mathcal{G}$（对应于 $Z_0$）；
依引理 \ref{coherent-lemma-property-irreducible}，这种情况不可能发生。
\end{proof}

\begin{lemma}
\label{coherent-lemma-property-irreducible-higher-rank-cohomological}
设 $X$ 为 Noether 概形。设 $Z_0 \subset X$ 为泛点
$\xi$ 的不可约闭子集。设 $\mathcal{P}$ 为 $X$
上凝聚层的一种性质，并且：
\begin{enumerate}
\item 对凝聚层的任意短正合列
$$
0 \to \mathcal{F}_1 \to \mathcal{F} \to \mathcal{F}_2 \to 0
$$
若 $\mathcal{F}_i$（$i = 1, 2$）具有性质
$\mathcal{P}$，则 $\mathcal{F}$ 也具有该性质。
\item 若性质 $\mathcal{P}$ 对 $\mathcal{F}^{\oplus r}$
成立，其中 $r \geq 1$，则它也对 $\mathcal{F}$ 成立。
\item 对每个整的闭子概形 $Z \subset Z_0 \subset X$、
$Z \not = Z_0$ 及每个准凝聚理想层
$\mathcal{I} \subset \mathcal{O}_Z$，性质
$\mathcal{P}$ 对 $(Z \to X)_*\mathcal{I}$ 成立。
\item 存在某个凝聚层 $\mathcal{G}$，满足：
\begin{enumerate}
\item $\text{Supp}(\mathcal{G}) = Z_0$；
\item $\mathcal{G}_\xi$ 被 $\mathfrak m_\xi$ 零化；并且
\item 对每个准凝聚理想层
$\mathcal{J} \subset \mathcal{O}_X$，若
$\mathcal{J}_\xi = \mathcal{O}_{X, \xi}$，则存在准凝聚子层
$\mathcal{G}' \subset \mathcal{J}\mathcal{G}$，使得
$\mathcal{G}'_\xi = \mathcal{G}_\xi$，且性质
$\mathcal{P}$ 对 $\mathcal{G}'$ 成立。
\end{enumerate}
\end{enumerate}
则性质 $\mathcal{P}$ 对每个凝聚层 $\mathcal{F}$ 均成立，
只要该层在 $X$ 上且其支集包含于 $Z_0$。
\end{lemma}

\begin{proof}
注意，若 $\mathcal{F}$ 是凝聚层，并有凝聚子层的滤过
$$
0 = \mathcal{F}_0 \subset \mathcal{F}_1 \subset
\ldots \subset \mathcal{F}_m = \mathcal{F}
$$
使每个 $\mathcal{F}_i/\mathcal{F}_{i - 1}$
都具有性质 $\mathcal{P}$，则 $\mathcal{F}$ 也具有该性质。
这由假设 (1) 推出。

\medskip\noindent
作为第一个应用，可知支集严格包含于 $Z_0$ 的任意凝聚层
都具有性质 $\mathcal{P}$。事实上，这样的层有一个滤过
（参见引理 \ref{coherent-lemma-coherent-filter}），其子商依 (3)
均具有性质 $\mathcal{P}$。

\medskip\noindent
以 $i : Z_0 \to X$ 表示该闭浸入。
考虑 (4) 中的凝聚层 $\mathcal{G}$。
由引理 \ref{coherent-lemma-prepare-filter-irreducible}，
存在理想层 $\mathcal{I}$（定义在 $Z_0$ 上）以及短正合列
$$
0 \to
i_*\mathcal{I}^{\oplus r} \to
\mathcal{G} \to
\mathcal{Q} \to 0
$$
其中 $\mathcal{Q}$ 的支集严格包含于 $Z_0$。
特别地，$r > 0$ 且 $\mathcal{I}$ 非零，因为
$\mathcal{G}$ 的支集等于 $Z_0$。
任取非零准凝聚理想层 $\mathcal{I}' \subset \mathcal{I}$；
它定义在 $Z_0$ 上且包含于 $\mathcal{I}$。于是也得到短正合列
$$
0 \to
i_*(\mathcal{I}')^{\oplus r} \to
\mathcal{G} \to
\mathcal{Q}' \to 0
$$
其中 $\mathcal{Q}'$ 的支集真包含于 $Z_0$。
设 $\mathcal{J} \subset \mathcal{O}_X$ 为定义
$\mathcal{Q}'$ 的支集的准凝聚理想层（例如，对应于
$\mathcal{Q}'$ 的支集上的既约诱导闭子概形结构的理想）。
则 $\mathcal{J}_\xi = \mathcal{O}_{X, \xi}$。由引理
\ref{coherent-lemma-power-ideal-kills-sheaf} 可知，
$\mathcal{J}^n\mathcal{Q}' = 0$ 对某个 $n$ 成立。
因此 $\mathcal{J}^n\mathcal{G} \subset i_*(\mathcal{I}')^{\oplus r}$。
由引理的假设 (4)(c)，可知存在准凝聚子层
$\mathcal{G}' \subset \mathcal{J}^n\mathcal{G}$，满足
$\mathcal{G}'_\xi = \mathcal{G}_\xi$，且性质
$\mathcal{P}$ 成立。
从而得到短正合列
$$
0 \to \mathcal{G}' \to
i_*(\mathcal{I}')^{\oplus r} \to
\mathcal{Q}'' \to 0
$$
其中 $\mathcal{Q}''$ 的支集真包含于 $Z_0$。
因此由开头的说明和引理的性质 (1)，可知
$i_*(\mathcal{I}')^{\oplus r}$ 满足 $\mathcal{P}$。
再由 (2) 可知 $i_*\mathcal{I}'$ 满足 $\mathcal{P}$。
最后，对任意准凝聚理想层
$\mathcal{I}'' \subset \mathcal{O}_{Z_0}$，令
$\mathcal{I}' = \mathcal{I}'' \cap \mathcal{I}$，便得到短正合列
$$
0 \to
i_*(\mathcal{I}') \to
i_*(\mathcal{I}'') \to
\mathcal{Q}''' \to 0
$$
其中 $\mathcal{Q}'''$ 的支集真包含于 $Z_0$。
故性质 $\mathcal{P}$ 对 $i_*\mathcal{I}''$ 也成立。

\medskip\noindent
证明的最后一步是注意到，任意凝聚层 $\mathcal{F}$
只要在 $X$ 上且支集包含于 $Z_0$，便都有一个滤过（再次参见引理
\ref{coherent-lemma-coherent-filter}），依刚才所述，其子商均具有性质
$\mathcal{P}$。
\end{proof}

\begin{lemma}
\label{coherent-lemma-property-higher-rank-cohomological}
设 $X$ 为 Noether 概形。
设 $\mathcal{P}$ 为 $X$ 上凝聚层的一种性质，并且：
\begin{enumerate}
\item 对凝聚层的任意短正合列
$$
0 \to \mathcal{F}_1 \to \mathcal{F} \to \mathcal{F}_2 \to 0
$$
若 $\mathcal{F}_i$（$i = 1, 2$）具有性质
$\mathcal{P}$，则 $\mathcal{F}$ 也具有该性质。
\item 若性质 $\mathcal{P}$ 对 $\mathcal{F}^{\oplus r}$
成立，其中 $r \geq 1$，则它也对 $\mathcal{F}$ 成立。
\item 对每个整的闭子概形 $Z \subset X$（其泛点为
$\xi$），存在某个凝聚层 $\mathcal{G}$，满足：
\begin{enumerate}
\item $\text{Supp}(\mathcal{G}) = Z$；
\item $\mathcal{G}_\xi$ 被 $\mathfrak m_\xi$ 零化；并且
\item 对每个准凝聚理想层
$\mathcal{J} \subset \mathcal{O}_X$，若
$\mathcal{J}_\xi = \mathcal{O}_{X, \xi}$，则存在准凝聚子层
$\mathcal{G}' \subset \mathcal{J}\mathcal{G}$，使得
$\mathcal{G}'_\xi = \mathcal{G}_\xi$，且性质
$\mathcal{P}$ 对 $\mathcal{G}'$ 成立。
\end{enumerate}
\end{enumerate}
则性质 $\mathcal{P}$ 对 $X$ 上的每个凝聚层均成立。
\end{lemma}

\begin{proof}
这由引理
\ref{coherent-lemma-property-irreducible-higher-rank-cohomological}
推出，其方式与引理 \ref{coherent-lemma-property} 由引理
\ref{coherent-lemma-property-irreducible} 推出的方式完全相同。
\end{proof}

\section{有限态射与仿射性}
\label{coherent-section-finite-affine}

\noindent
本节利用前几节的结果证明：Noether 仿射概形在有限态射下的像
是仿射的。稍后将看到，该结果在更一般的情形下成立
（参见《极限》引理 \ref{limits-lemma-affine} 及命题
\ref{limits-proposition-affine}）。

\begin{lemma}
\label{coherent-lemma-finite-morphism-Noetherian}
设 $f : Y \to X$ 为概形态射。
假设 $f$ 有限且满射，并且 $X$ 局部 Noether。
设 $Z \subset X$ 为泛点 $\xi$ 的整闭子概形。
则存在凝聚层 $\mathcal{F}$，它在 $Y$ 上，并使得
$f_*\mathcal{F}$ 的支集等于 $Z$，且
$(f_*\mathcal{F})_\xi$ 被 $\mathfrak m_\xi$ 零化。
\end{lemma}

\begin{proof}
注意，由《态射》引理
\ref{morphisms-lemma-finite-type-noetherian}，
$Y$ 局部 Noether。
由于 $f$ 满射，纤维 $Y_\xi$ 非空。
选取 $\xi' \in Y$，使其映到 $\xi$。令
$Z' = \overline{\{\xi'\}}$。
可将 $Z' \subset Y$ 视为既约闭子概形，参见《概形》引理
\ref{schemes-lemma-reduced-closed-subscheme}。
因此层 $\mathcal{F} = (Z' \to Y)_*\mathcal{O}_{Z'}$
是 $Y$ 上的凝聚层（参见引理
\ref{coherent-lemma-finite-pushforward-coherent}）。
考察交换图
$$
\xymatrix{
Z' \ar[r]_{i'} \ar[d]_{f'} &
Y \ar[d]^f \\
Z \ar[r]^i &
X
}
$$
可见 $f_*\mathcal{F} = i_*f'_*\mathcal{O}_{Z'}$。
因此 $f_*\mathcal{F}$ 在 $\xi$ 处的层茎就是
$f'_*\mathcal{O}_{Z'}$ 在 $\xi$ 处的层茎。注意，由于
$Z'$ 是泛点为 $\xi'$ 的整概形，$\xi'$ 是 $Z'$
中位于 $\xi$ 上方的唯一一点，参见《代数》引理
\ref{algebra-lemma-finite-is-integral} 与
\ref{algebra-lemma-integral-no-inclusion}。
故 $f'_*\mathcal{O}_{Z'}$ 在 $\xi$ 处的层茎等于
$\mathcal{O}_{Z', \xi'} = \kappa(\xi')$。特别地，
$f_*\mathcal{F}$ 在 $\xi$ 处的层茎非零。
结合 $f_*\mathcal{F}$ 具有
$i_*f'_*(\text{something})$ 的形式这一事实，即得引理。
\end{proof}

\begin{lemma}
\label{coherent-lemma-affine-morphism-projection-ideal}
设 $f : Y \to X$ 为概形态射。
设 $\mathcal{F}$ 为 $Y$ 上的准凝聚层。
设 $\mathcal{I}$ 为 $X$ 上的准凝聚理想层。
若态射 $f$ 仿射，则
$\mathcal{I}f_*\mathcal{F} = f_*(f^{-1}\mathcal{I}\mathcal{F})$。
\end{lemma}

\begin{proof}
该记号的含义如下。由于 $f^{-1}$ 是正合函子，
$f^{-1}\mathcal{I}$ 是 $f^{-1}\mathcal{O}_X$ 的理想层。
通过映射
$f^\sharp : f^{-1}\mathcal{O}_X \to \mathcal{O}_Y$，
它作用在 $\mathcal{F}$ 上。于是
$f^{-1}\mathcal{I}\mathcal{F}$ 是由形如 $as$
的局部截面的和所生成的子层，其中 $a$ 是
$f^{-1}\mathcal{I}$ 的局部截面，而 $s$ 是
$\mathcal{F}$ 的局部截面。它是 $\mathcal{O}_Y$-模
$\mathcal{F}$ 的准凝聚子模，因为它也是自然映射
$f^*\mathcal{I} \otimes_{\mathcal{O}_Y} \mathcal{F} \to \mathcal{F}$
的像。

\medskip\noindent
说明这些之后，证明是直接的。问题是局部的，故可假设
$X$ 仿射。由于 $f$ 仿射，可知 $Y$ 也仿射。
因此可写成
$Y = \Spec(B)$、$X = \Spec(A)$、
$\mathcal{F} = \widetilde{M}$ 以及
$\mathcal{I} = \widetilde{I}$。在此情形，引理断言归结为
$$
I(M_A) = ((IB)M)_A
$$
其中 $M_A$ 表示相应的 $A$-模；它由 $B$-模 $M$ 得到。
\end{proof}

\begin{lemma}
\label{coherent-lemma-image-affine-finite-morphism-affine-Noetherian}
设 $f : Y \to X$ 为概形态射。假设：
\begin{enumerate}
\item $f$ 有限；
\item $f$ 满射；
\item $Y$ 仿射；并且
\item $X$ Noether。
\end{enumerate}
则 $X$ 仿射。
\end{lemma}

\begin{proof}
我们将证明，在引理的假设下，对任意凝聚
$\mathcal{O}_X$-模 $\mathcal{F}$ 都有
$H^1(X, \mathcal{F}) = 0$。特别地，这蕴含对
$H^1(X, \mathcal{I}) = 0$ 对
$\mathcal{O}_X$ 的每个准凝聚理想层成立。于是由引理
\ref{coherent-lemma-quasi-compact-h1-zero-covering} 或引理
\ref{coherent-lemma-quasi-separated-h1-zero-covering} 可知 $X$ 仿射。

\medskip\noindent
设 $\mathcal{P}$ 为凝聚层 $\mathcal{F}$（在 $X$ 上）
的如下性质：
$$
\mathcal{P}(\mathcal{F}) \Leftrightarrow H^1(X, \mathcal{F}) = 0.
$$
我们将应用引理
\ref{coherent-lemma-property-higher-rank-cohomological}。
因此必须对 $\mathcal{P}$ 验证该引理的 (1)、(2)、(3)。
性质 (1) 由层的短正合列所伴随的上同调长正合列推出。
性质 (2) 由 $H^1(X, -)$ 是加性函子推出。
为验证 (3)，设 $Z \subset X$ 为泛点 $\xi$ 的整闭子概形。
设 $\mathcal{F}$ 为 $Y$ 上的凝聚层，使得
$f_*\mathcal{F}$ 的支集等于 $Z$，且
$(f_*\mathcal{F})_\xi$ 被 $\mathfrak m_\xi$ 零化；
参见引理 \ref{coherent-lemma-finite-morphism-Noetherian}。
我们断言取 $\mathcal{G} = f_*\mathcal{F}$ 即可。
只需验证引理
\ref{coherent-lemma-property-higher-rank-cohomological} 的 (3)(c)。
因此假设 $\mathcal{J} \subset \mathcal{O}_X$
是满足
$\mathcal{J}_\xi = \mathcal{O}_{X, \xi}$ 的准凝聚理想层。
有限态射是仿射的，故由引理
\ref{coherent-lemma-affine-morphism-projection-ideal} 可知
$\mathcal{J}\mathcal{G} = f_*(f^{-1}\mathcal{J}\mathcal{F})$。
此外，如引理
\ref{coherent-lemma-affine-morphism-projection-ideal} 的证明中所指出，
层 $f^{-1}\mathcal{J}\mathcal{F}$ 是准凝聚
$\mathcal{O}_Y$-模。
由于 $Y$ 仿射，有
$H^1(Y, f^{-1}\mathcal{J}\mathcal{F}) = 0$，
参见引理 \ref{coherent-lemma-quasi-coherent-affine-cohomology-zero}。
由于 $f$ 有限，因而仿射，所以
$$
H^1(X, \mathcal{J}\mathcal{G}) =
H^1(X, f_*(f^{-1}\mathcal{J}\mathcal{F})) =
H^1(Y, f^{-1}\mathcal{J}\mathcal{F}) = 0
$$
这由引理 \ref{coherent-lemma-relative-affine-cohomology} 得到。
因此准凝聚子层
$\mathcal{G}' = \mathcal{J}\mathcal{G}$ 满足
$\mathcal{P}$。这如所要求验证了引理
\ref{coherent-lemma-property-higher-rank-cohomological} 的性质 (3)(c)。
\end{proof}

\section{Proj 上的凝聚层，I}
\label{coherent-section-coherent-proj}

\noindent
本节讨论 $\text{Proj}(A)$ 上的凝聚层，其中 $A$
是由 $A_1$ 在 $A_0$ 上生成的 Noether 分次环。
下一节将讨论 $A$ 不由次数 $1$ 的元素生成时的情形。
首先，我们为 Noether 环上的射影空间表述一个汇总性结果。

\begin{lemma}
\label{coherent-lemma-coherent-projective}
设 $R$ 为 Noether 环。
设 $n \geq 0$ 为整数。
对每个凝聚层 $\mathcal{F}$（在 $\mathbf{P}^n_R$ 上），
有下列结论：
\begin{enumerate}
\item 存在 $r \geq 0$、
$d_1, \ldots, d_r \in \mathbf{Z}$ 以及满射
$$
\bigoplus\nolimits_{j = 1, \ldots, r}
\mathcal{O}_{\mathbf{P}^n_R}(d_j)
\longrightarrow
\mathcal{F}.
$$
\item 有 $H^i(\mathbf{P}^n_R, \mathcal{F}) = 0$，
除非 $0 \leq i \leq n$。
\item 对任意 $i$，上同调群
$H^i(\mathbf{P}^n_R, \mathcal{F})$ 是有限 $R$-模。
\item 若 $i > 0$，则
$H^i(\mathbf{P}^n_R, \mathcal{F}(d)) = 0$
对所有充分大的 $d$ 成立。
\item 对任意 $k \in \mathbf{Z}$，分次
$R[T_0, \ldots, T_n]$-模
$$
\bigoplus\nolimits_{d \geq k} H^0(\mathbf{P}^n_R, \mathcal{F}(d))
$$
是有限 $R[T_0, \ldots, T_n]$-模。
\end{enumerate}
\end{lemma}

\begin{proof}
我们将用到 $\mathcal{O}_{\mathbf{P}^n_R}(1)$ 是概形
$\mathbf{P}^n_R$ 上的丰沛可逆层。由于
$\mathbf{P}^n_R$ 被标准仿射开集 $D_{+}(T_i)$ 覆盖，
这直接由定义推出。因此由《性质》命题
\ref{properties-proposition-characterize-ample}，
每个有限型准凝聚 $\mathcal{O}_{\mathbf{P}^n_R}$-模
都是 $\mathcal{O}_{\mathbf{P}^n_R}(1)$ 的张量幂的
某个有限直和的商。另一方面，由引理
\ref{coherent-lemma-coherent-Noetherian}，在 $R$ 上的射影空间中，
凝聚层与有限型准凝聚层是同一回事。故得到 (1)。

\medskip\noindent
射影 $n$-空间 $\mathbf{P}^n_R$ 由 $n + 1$ 个仿射开集覆盖，
即标准开集 $D_{+}(T_i)$（$i = 0, \ldots, n$），参见《构造》引理
\ref{constructions-lemma-standard-covering-projective-space}。
因此，对任意准凝聚层 $\mathcal{F}$（在
$\mathbf{P}^n_R$ 上），有
$H^i(\mathbf{P}^n_R, \mathcal{F}) = 0$（当
$i \geq n + 1$ 时），参见引理
\ref{coherent-lemma-vanishing-nr-affines}。故 (2) 成立。

\medskip\noindent
我们同时为 $\mathbf{P}^n_R$ 上的所有凝聚层证明
(3) 与 (4)，并对 $i$ 作降归纳。由 (2)，结论对
$i \geq n + 1$ 显然成立。假设已知 $i + 1$ 的结论，
现在证明 $i$ 的结论。（若 $i = 0$，则 (4) 为空断言。）
设 $\mathcal{F}$ 为 $\mathbf{P}^n_R$ 上的凝聚层。
选取 (1) 中的满射，并以 $\mathcal{G}$ 表示其核，从而有短正合列
$$
0 \to \mathcal{G} \to
\bigoplus\nolimits_{j = 1, \ldots, r}
\mathcal{O}_{\mathbf{P}^n_R}(d_j)
\to
\mathcal{F} \to 0
$$
由引理 \ref{coherent-lemma-coherent-abelian-Noetherian} 可知
$\mathcal{G}$ 凝聚。上同调长正合列给出正合列
$$
H^i(\mathbf{P}^n_R, \bigoplus\nolimits_{j = 1, \ldots, r}
\mathcal{O}_{\mathbf{P}^n_R}(d_j))
\to
H^i(\mathbf{P}^n_R, \mathcal{F})
\to
H^{i + 1}(\mathbf{P}^n_R, \mathcal{G}).
$$
由归纳假设，右侧 $R$-模有限；由引理
\ref{coherent-lemma-cohomology-projective-space-over-ring}，
左侧 $R$-模有限。由于 $R$ Noether，立即可知
$H^i(\mathbf{P}^n_R, \mathcal{F})$ 是有限 $R$-模。
这证明了断言 (3) 的归纳步骤。
由于 $\mathcal{O}_{\mathbf{P}^n_R}(d)$ 可逆，
可知在 $\mathbf{P}^n_R$ 上作扭转是正合函子
（因为它由与可逆层张量得到，参见《构造》定义
\ref{constructions-definition-twist}）。
这意味着对所有 $d \in \mathbf{Z}$，序列
$$
0 \to \mathcal{G}(d) \to
\bigoplus\nolimits_{j = 1, \ldots, r}
\mathcal{O}_{\mathbf{P}^n_R}(d_j + d)
\to
\mathcal{F}(d) \to 0
$$
短正合。所得上同调序列为
$$
H^i(\mathbf{P}^n_R, \bigoplus\nolimits_{j = 1, \ldots, r}
\mathcal{O}_{\mathbf{P}^n_R}(d_j + d))
\to
H^i(\mathbf{P}^n_R, \mathcal{F}(d))
\to
H^{i + 1}(\mathbf{P}^n_R, \mathcal{G}(d)).
$$
由归纳假设，$d \gg 0$ 时右侧的模为零；由引理
\ref{coherent-lemma-cohomology-projective-space-over-ring} 中的计算，
一旦 $d \geq -\min\{d_j\}$ 且 $i \geq 1$，左侧的模即为零。
故得到断言 (4) 的归纳步骤。
这完成 (3) 与 (4) 的证明。

\medskip\noindent
为证明 (5)，注意对所有充分大的 $d$，映射
$$
H^0(\mathbf{P}^n_R, \bigoplus\nolimits_{j = 1, \ldots, r}
\mathcal{O}_{\mathbf{P}^n_R}(d_j + d))
\to
H^0(\mathbf{P}^n_R, \mathcal{F}(d))
$$
是满射，这是由刚刚证明的
$H^1(\mathbf{P}^n_R, \mathcal{G}(d))$ 的消失得到的。
换言之，当 $k$ 充分大时，模
$$
M_k
=
\bigoplus\nolimits_{d \geq k} H^0(\mathbf{P}^n_R, \mathcal{F}(d))
$$
是相应模
$$
N_k
=
\bigoplus\nolimits_{d \geq k} H^0(\mathbf{P}^n_R,
\bigoplus\nolimits_{j = 1, \ldots, r}
\mathcal{O}_{\mathbf{P}^n_R}(d_j + d)
)
$$
的商。当 $k$ 充分小时（例如有
$k < -d_j$ 对所有 $j$ 成立），由第
\ref{coherent-section-cohomology-projective-space} 节中的计算，有
$$
N_k = \bigoplus\nolimits_{j = 1, \ldots, r}
R[T_0, \ldots, T_n](d_j)
$$
特别地，它有限生成。
现在设 $k \in \mathbf{Z}$ 任意。
选取 $k_{-} \ll k \ll k_{+}$。
考虑交换图
$$
\xymatrix{
N_{k_{-}} & N_{k_{+}} \ar[d] \ar[l] \\
M_k & M_{k_{+}} \ar[l]
}
$$
其中竖直箭头是上述满射，而水平箭头是显然的包含映射。
由以上说明，$N_{k_{-}}$ 是有限生成
$R[T_0, \ldots, T_n]$-模。因此 $N_{k_{+}}$
也是有限生成 $R[T_0, \ldots, T_n]$-模，
因为它是有限生成模的子模，而环
$R[T_0, \ldots, T_n]$ 是 Noether 的。
由于竖直箭头满射，$M_{k_{+}}$ 是有限生成
$R[T_0, \ldots, T_n]$-模。商
$M_k/M_{k_{+}}$ 作为 $R$-模是有限的，因为它是有限个有限
$R$-模 $H^0(\mathbf{P}^n_R, \mathcal{F}(d))$
的直和，其中 $k \leq d < k_{+}$。
注意，此处使用了 (3) 中 $i = 0$ 的情形。因此
$M_k/M_{k_{+}}$ 更是有限
$R[T_0, \ldots, T_n]$-模。
换言之，我们把 $M_k$ 夹在两个有限
$R[T_0, \ldots, T_n]$-模之间，结论得证。
\end{proof}

\begin{lemma}
\label{coherent-lemma-coherent-on-proj}
设 $A$ 为分次环，使得 $A_0$ Noether，且 $A$
由 $A_1$ 的有限多个元素在 $A_0$ 上生成。
令 $X = \text{Proj}(A)$。则 $X$ 是 Noether 概形。
设 $\mathcal{F}$ 为凝聚 $\mathcal{O}_X$-模。
\begin{enumerate}
\item 存在 $r \geq 0$、
$d_1, \ldots, d_r \in \mathbf{Z}$ 以及满射
$$
\bigoplus\nolimits_{j = 1, \ldots, r} \mathcal{O}_X(d_j)
\longrightarrow \mathcal{F}.
$$
\item 对任意 $p$，上同调群 $H^p(X, \mathcal{F})$
是有限 $A_0$-模。
\item 若 $p > 0$，则
$H^p(X, \mathcal{F}(d)) = 0$ 对所有充分大的 $d$ 成立。
\item 对任意 $k \in \mathbf{Z}$，分次 $A$-模
$$
\bigoplus\nolimits_{d \geq k} H^0(X, \mathcal{F}(d))
$$
是有限 $A$-模。
\end{enumerate}
\end{lemma}

\begin{proof}
由假设，存在分次 $A_0$-代数的满射
$$
A_0[T_0, \ldots, T_n] \longrightarrow A
$$
其中 $\deg(T_j) = 1$ 对 $j = 0, \ldots, n$ 成立。
由《构造》引理
\ref{constructions-lemma-surjective-graded-rings-generated-degree-1-map-proj}，
这定义一个闭浸入 $i : X \to \mathbf{P}^n_{A_0}$，满足
$i^*\mathcal{O}_{\mathbf{P}^n_{A_0}}(1) = \mathcal{O}_X(1)$。
特别地，$X$ 作为 Noether 概形
$\mathbf{P}^n_{A_0}$ 的闭子概形，是 Noether 的。
我们断言，引理关于 $\mathcal{F}$ 的各结论可由引理
\ref{coherent-lemma-coherent-projective} 关于凝聚层
$i_*\mathcal{F}$（引理 \ref{coherent-lemma-i-star-equivalence}）在
$\mathbf{P}^n_{A_0}$ 上的相应结论推出。例如，由该引理存在满射
$$
\bigoplus\nolimits_{j = 1, \ldots, r} \mathcal{O}_{\mathbf{P}^n_{A_0}}(d_j)
\longrightarrow i_*\mathcal{F}.
$$
由伴随，这对应于映射
$\bigoplus_{j = 1, \ldots, r} \mathcal{O}_X(d_j) \longrightarrow \mathcal{F}$，
它也满射。关于上同调的断言由等式
$H^p(X, \mathcal{F}(d)) = H^p(\mathbf{P}^n_{A_0}, i_*\mathcal{F}(d))$
及引理 \ref{coherent-lemma-relative-affine-cohomology} 推出。
\end{proof}

\begin{lemma}
\label{coherent-lemma-recover-tail-graded-module}
设 $A$ 为分次环，使得 $A_0$ Noether，且 $A$
由 $A_1$ 的有限多个元素在 $A_0$ 上生成。设 $M$
为有限分次 $A$-模。令 $X = \text{Proj}(A)$，并令
$\widetilde{M}$ 为准凝聚 $\mathcal{O}_X$-模；它在
$X$ 上且与 $M$ 相伴。
《构造》引理 \ref{constructions-lemma-apply-modules} 中的映射
$$
M_n \longrightarrow \Gamma(X, \widetilde{M}(n))
$$
对所有充分大的 $n$ 都是同构。
\end{lemma}

\begin{proof}
由于 $M$ 是有限 $A$-模，可知 $\widetilde{M}$
是有限型 $\mathcal{O}_X$-模，即凝聚
$\mathcal{O}_X$-模。
令 $N = \bigoplus_{n \in \mathbf{Z}} \Gamma(X, \widetilde{M}(n))$。
必须证明映射 $M \to N$（分次 $A$-模的映射）
在所有充分大的次数上是同构。
由《性质》引理 \ref{properties-lemma-proj-quasi-coherent}，
存在典范同构 $\widetilde{N} \to \widetilde{M}$，使得诱导映射
$N_n \to N_n = \Gamma(X, \widetilde{M}(n))$
为恒等映射。因此有映射
$\widetilde{M} \to \widetilde{N} \to \widetilde{M}$，
使得对所有 $n$，交换图
$$
\xymatrix{
M_n \ar[d] \ar[r] & N_n \ar[d] \ar@{=}[rd] \\
\Gamma(X, \widetilde{M}(n)) \ar[r] &
\Gamma(X, \widetilde{N}(n)) \ar[r]^{\cong} &
\Gamma(X, \widetilde{M}(n))
}
$$
交换。这意味着复合
$$
M_n \to \Gamma(X, \widetilde{M}(n))
\to \Gamma(X, \widetilde{N}(n)) \to \Gamma(X, \widetilde{M}(n))
$$
等于典范映射
$M_n \to \Gamma(X, \widetilde{M}(n))$。
显然，这蕴含复合
$\widetilde{M} \to \widetilde{N} \to \widetilde{M}$
为恒等映射。因此 $\widetilde{M} \to \widetilde{N}$
是同构。令 $K = \Ker(M \to N)$ 且
$Q = \Coker(M \to N)$。
回忆函子 $M \mapsto \widetilde{M}$ 正合，参见《构造》引理
\ref{constructions-lemma-proj-sheaves}。
因此 $\widetilde{K} = 0$ 且 $\widetilde{Q} = 0$。
回忆 $A$ 是 Noether 环，$M$ 是有限生成 $A$-模，
而 $N$ 是分次 $A$-模，并且由引理
\ref{coherent-lemma-coherent-on-proj} 的最后一部分，
$N' = \bigoplus_{n \geq 0} N_n$ 有限生成。
因此 $K' = \bigoplus_{n \geq 0} K_n$ 及
$Q' = \bigoplus_{n \geq 0} Q_n$ 是有限 $A$-模。
注意 $\widetilde{Q} = \widetilde{Q'}$ 且
$\widetilde{K} = \widetilde{K'}$。
故为完成证明，只需说明：若有限 $A$-模 $K$ 满足
$\widetilde{K} = 0$，则它只有有限多个非零齐次分量
$K_d$，其中 $d \geq 0$。
为此，设 $x_1, \ldots, x_r \in K$ 为齐次生成元，
其次数分别为 $d_1, \ldots, d_r$。
设 $f_1, \ldots, f_n \in A_1$ 为生成
$A$ 的元素（在 $A_0$ 上）。对每个 $i$ 与 $j$，存在
$n_{ij} \geq 0$，使得
$f_i^{n_{ij}} x_j = 0$ 在 $K_{d_j + n_{ij}}$ 中成立：
否则 $x_i/f_i^{d_i} \in K_{(f_i)}$ 将不为零，即
$\widetilde{K}$ 不为零。
于是 $K_d$ 在 $d > \max_j(d_j + \sum_i n_{ij})$ 时
为零，因为 $K_d$ 的每个元素都是若干项之和，
每一项都是 $f_i$ 中的一个单项式乘以某个 $x_j$，
且总次数为 $d$。
\end{proof}

\noindent
设 $A$ 为分次环，使得 $A_0$ Noether，且 $A$
由 $A_1$ 的有限多个元素在 $A_0$ 上生成。回忆
$A_+ = \bigoplus_{n > 0} A_n$ 是无关理想。
设 $M$ 为分次 $A$-模。回忆，称 $M$ 是
$A_+$-幂挠模，是指对所有 $x \in M$ 都存在
$n \geq 1$ 使得 $(A_+)^n x = 0$；参见《代数进阶》定义
\ref{more-algebra-definition-f-power-torsion}。
若 $M$ 有限生成，则这等价于
$M_n = 0$ 对 $n \gg 0$ 成立。有时把 $A_+$-幂挠模简称为挠模。
有时称分次 $A$-模 $M$ 无挠，是指只要
$x \in M$ 满足 $(A_+)^n x = 0$ 且 $n > 0$，
便有 $x = 0$。
以 $\text{Mod}_A$ 表示分次 $A$-模范畴，
以 $\text{Mod}^{fg}_A$ 表示其中有限生成对象组成的满子范畴，
并以 $\text{Mod}^{fg}_{A, torsion}$ 表示由模
$M$ 组成的满子范畴；这些模满足
$M_n = 0$（当 $n \gg 0$）。

\begin{proposition}
\label{coherent-proposition-coherent-modules-on-proj}
设 $A$ 为分次环，使得 $A_0$ Noether，且 $A$
由 $A_1$ 的有限多个元素在 $A_0$ 上生成。
令 $X = \text{Proj}(A)$。函子 $M \mapsto \widetilde M$
诱导等价
$$
\text{Mod}^{fg}_A/\text{Mod}^{fg}_{A, torsion}
\longrightarrow
\textit{Coh}(\mathcal{O}_X)
$$
其拟逆由
$\mathcal{F} \longmapsto \bigoplus_{n \geq 0} \Gamma(X, \mathcal{F}(n))$
给出。
\end{proposition}

\begin{proof}
子范畴 $\text{Mod}^{fg}_{A, torsion}$ 是
$\text{Mod}^{fg}_A$ 的 Serre 子范畴，参见《同调》定义
\ref{homology-definition-serre-subcategory}。
这由上述对象描述显然可见，也可由《代数进阶》引理
\ref{more-algebra-lemma-I-power-torsion} 推出。
因此箭头左侧的商范畴已在《同调》引理
\ref{homology-lemma-serre-subcategory-is-kernel} 中定义。
为定义命题中的函子，只需说明函子
$M \mapsto \widetilde M$ 把挠模送到 $0$。
这是显然的，因为对任意齐次 $f \in A_+$，
模 $M_f$ 为零，因而值 $M_{(f)}$（$\widetilde M$
在 $D_+(f)$ 上的值）也为零。

\medskip\noindent
由引理 \ref{coherent-lemma-coherent-on-proj}，所提出的拟逆有意义。
具体而言，该引理表明
$\mathcal{F} \longmapsto \bigoplus_{n \geq 0} \Gamma(X, \mathcal{F}(n))$
是函子
$\textit{Coh}(\mathcal{O}_X) \to \text{Mod}^{fg}_A$；
可将它与商函子
$\text{Mod}^{fg}_A \to \text{Mod}^{fg}_A/\text{Mod}^{fg}_{A, torsion}$
复合。

\medskip\noindent
由引理 \ref{coherent-lemma-recover-tail-graded-module}，
从左到右再到左的复合同构于恒等函子。

\medskip\noindent
最后，设 $\mathcal{F}$ 为凝聚 $\mathcal{O}_X$-模。
令 $M = \bigoplus_{n \in \mathbf{Z}} \Gamma(X, \mathcal{F}(n))$，
将其视为分次 $A$-模，从而我们的函子把 $\mathcal{F}$ 送到
$M_{\geq 0} = \bigoplus_{n \geq 0} M_n$。
由《性质》引理 \ref{properties-lemma-proj-quasi-coherent}，
典范映射 $\widetilde M \to \mathcal{F}$ 是同构。
由于包含映射
$M_{\geq 0} \to M$ 定义同构
$\widetilde{M_{\geq 0}} \to \widetilde M$，
可知从右到左再到右的复合也同构于恒等函子。
\end{proof}





\section{Proj 上的凝聚层，II}
\label{coherent-section-coherent-proj-general}

\noindent
在本节中，我们讨论 $\text{Proj}(A)$ 上的凝聚层，
其中 $A$ 是 Noether 分次环。多数结果将借助一个技巧，
由前一节关于 $A$ 由次数 $1$ 元生成时的相应结果推出。

\begin{lemma}
\label{coherent-lemma-coherent-on-proj-general}
设 $A$ 为 Noether 分次环。令 $X = \text{Proj}(A)$。则 $X$
是 Noether 概形。设 $\mathcal{F}$ 为凝聚 $\mathcal{O}_X$-模。
\begin{enumerate}
\item 存在 $r \geq 0$、
$d_1, \ldots, d_r \in \mathbf{Z}$ 以及满射
$$
\bigoplus\nolimits_{j = 1, \ldots, r} \mathcal{O}_X(d_j)
\longrightarrow \mathcal{F}.
$$
\item 对任意 $p$，上同调群 $H^p(X, \mathcal{F})$ 是有限
$A_0$-模。
\item 若 $p > 0$，则 $H^p(X, \mathcal{F}(d)) = 0$
对所有充分大的 $d$ 成立。
\item 对任意 $k \in \mathbf{Z}$，分次 $A$-模
$$
\bigoplus\nolimits_{d \geq k} H^0(X, \mathcal{F}(d))
$$
是有限 $A$-模。
\end{enumerate}
\end{lemma}

\begin{proof}
我们把该断言归约到引理
\ref{coherent-lemma-coherent-on-proj} 来证明。
由《代数》引理 \ref{algebra-lemma-graded-Noetherian} 与
\ref{algebra-lemma-S-plus-generated}，环 $A_0$ 是 Noether 的，
并且 $A$ 在 $A_0$ 上由有限多个正次数齐次元
$f_1, \ldots, f_r$ 生成。
令 $d$ 为充分可除的整数。按《代数》第
\ref{algebra-section-graded} 节的记号，令 $A' = A^{(d)}$。
于是 $A'$ 在 $A'_0 = A_0$ 上由次数 $1$ 元生成；参见
《代数》引理 \ref{algebra-lemma-uple-generated-degree-1}。
因此引理 \ref{coherent-lemma-coherent-on-proj} 适用于
$X' = \text{Proj}(A')$。

\medskip\noindent
由《构造》引理 \ref{constructions-lemma-d-uple}，存在概形同构
$i : X \to X'$ 以及同构
$\mathcal{O}_X(nd) \to i^*\mathcal{O}_{X'}(n)$；
它们与映射 $A' \to A$ 以及映射
$A_n \to H^0(X, \mathcal{O}_X(n))$ 和 $A'_n \to H^0(X', \mathcal{O}_{X'}(n))$
相容。因此，引理 \ref{coherent-lemma-coherent-on-proj} 蕴含 $X$
为 Noether 概形，并且 (1) 与 (2) 成立。为证明 (3) 与 (4)，
可以利用如下事实：对任意固定的 $k$、$p$ 与 $q$，都有
$$
\bigoplus\nolimits_{dn + q \geq k} H^p(X, \mathcal{F}(dn + q)) =
\bigoplus\nolimits_{dn + q \geq k} H^p(X', (i_*\mathcal{F}(q))(n)
$$
这是由上述相容性得到的。若 $p > 0$，则由引理
\ref{coherent-lemma-coherent-on-proj}，对充分大的 $k$
（它可依赖于 $q$），右端为零。因为整数模 $d$
只有有限多个同余类，可知 (3) 对 $\mathcal{F}$ 在 $X$
上成立。若 $p = 0$，
则由引理 \ref{coherent-lemma-coherent-on-proj}，右端是有限
$A'$-模。再次利用同余类的有限性，可知
$\bigoplus_{n \geq k} H^0(X, \mathcal{F}(n))$ 也是有限
$A'$-模。由于该 $A'$-模结构来自 $A$-模结构
（由上述相容性），可知它也是有限 $A$-模。
\end{proof}

\begin{lemma}
\label{coherent-lemma-recover-tail-graded-module-general}
设 $A$ 为 Noether 分次环，并令 $d$ 为 $A$ 在 $A_0$
上生成元的最小公倍数。设 $M$ 为有限分次 $A$-模。
令 $X = \text{Proj}(A)$，并令 $\widetilde{M}$ 为准凝聚
$\mathcal{O}_X$-模；它在 $X$ 上且与 $M$ 相伴。
设 $k \in \mathbf{Z}$。
\begin{enumerate}
\item $N' = \bigoplus_{n \geq k} H^0(X, \widetilde{M(n)})$
是有限 $A$-模；
\item $N = \bigoplus_{n \geq k} H^0(X, \widetilde{M}(n))$
是有限 $A$-模；
\item 存在典范映射 $N \to N'$；
\item 若 $k$ 充分小，则存在典范映射 $M \to N'$；
\item 映射 $M_n \to N'_n$ 对 $n \gg 0$ 是同构；
\item $N_n \to N'_n$ 对 $d | n$ 是同构。
\end{enumerate}
\end{lemma}

\begin{proof}
(3) 中的映射 $N \to N'$ 来自《构造》等式
(\ref{constructions-equation-multiply-more-generally})，
取整体截面即可得到。

\medskip\noindent
由《构造》等式
(\ref{constructions-equation-global-sections-more-generally})，
存在分次 $A$-模的映射
$M \to \bigoplus_{n \in \mathbf{Z}} H^0(X, \widetilde{M(n)})$。
若 $M$ 的生成元所在次数均为 $\geq k$，则像包含在子模
$N' \subset \bigoplus_{n \in \mathbf{Z}} H^0(X, \widetilde{M(n)})$
中，从而得到 (4) 中的映射。

\medskip\noindent
由《代数》引理 \ref{algebra-lemma-graded-Noetherian} 与
\ref{algebra-lemma-S-plus-generated}，环 $A_0$ 是 Noether 的，
并且 $A$ 在 $A_0$ 上由有限多个正次数齐次元
$f_1, \ldots, f_r$ 生成。
令 $d = \text{lcm}(\deg(f_i))$。例如，由《构造》引理
\ref{constructions-lemma-where-invertible} 可知 (6) 成立。

\medskip\noindent
因为 $M$ 是有限 $A$-模，可知 $\widetilde{M}$ 是有限型
$\mathcal{O}_X$-模，即凝聚 $\mathcal{O}_X$-模。
因此 (2) 由引理 \ref{coherent-lemma-coherent-on-proj-general} 推出。

\medskip\noindent
我们用一个技巧由 (2) 推出 (1)。对
$q \in \{0, \ldots, d - 1\}$，记
$$
{}^qN = \bigoplus\nolimits_{n + q \geq k} H^0(X, \widetilde{M(q)}(n))
$$
由 (2)，这些都是有限 $A$-模。Noether 环 $A$
在 $A^{(d)} = \bigoplus_{n \geq 0} A_{dn}$ 上有限，
因为它由 $f_i$ 在 $A^{(d)}$ 上生成，且 $f_i^d \in A^{(d)}$。
因此 ${}^qN$ 是有限 $A^{(d)}$-模。
此外，$A^{(d)}$ 是 Noether 的（这由《代数》引理
\ref{algebra-lemma-dehomogenize-finite-type} 推出）。
于是 $A^{(d)}$-子模
${}^qN^{(d)} = \bigoplus_{n \in \mathbf{Z}} {}^qN_{dn}$
是有限 $A^{(d)}$-模。利用同构
$\widetilde{M(dn + q)} = \widetilde{M(q)}(dn)$，可写成
$$
N' =
\bigoplus\nolimits_{q \in \{0, \ldots, d - 1\}}
\bigoplus\nolimits_{dn + q \geq k}
H^0(X, \widetilde{M(q)}(dn)) =
\bigoplus\nolimits_{q \in \{0, \ldots, d - 1\}} {}^qN^{(d)}
$$
因此 $N'$ 在 $A^{(d)}$ 上有限，因而当然也在 $A$ 上有限。
故 (1) 成立。

\medskip\noindent
设 $K$ 为有限 $A$-模，且满足 $\widetilde{K} = 0$。
我们断言，$K_n = 0$ 对 $d|n$ 且 $n \gg 0$ 成立。
如上论证可知，$K^{(d)}$ 是有限 $A^{(d)}$-模。设
$x_1, \ldots, x_m \in K$ 是 $K^{(d)}$ 在 $A^{(d)}$
上的齐次生成元，其所在次数分别为
$d_1, \ldots, d_m$，并且 $d | d_j$。
对每个 $i$ 与 $j$，存在 $n_{ij} \geq 0$，使得
$f_i^{n_{ij}} x_j = 0$ 在 $K_{d_j + n_{ij}}$ 中成立；否则
$x_j/f_i^{d_i/\deg(f_i)} \in K_{(f_i)}$ 将不为零，即
$\widetilde{K}$ 不为零。这里使用了
$\deg(f_i) | d | d_j$ 对所有 $i, j$ 成立。
由此可知，$K_n$ 为零，其中 $n$ 满足 $d | n$ 且
$n > \max_j (d_j + \sum_i n_{ij} \deg(f_i))$，
因为 $K_n$ 的每个元素都是若干项之和，
其中每一项都是 $f_i$ 中的一个单项式乘以某个 $x_j$，
且总次数为 $n$。

\medskip\noindent
为完成证明，必须说明 $M \to N'$ 在所有充分大的次数上
是同构。映射 $N \to N'$ 诱导同构
$\widetilde{N} \to \widetilde{N'}$，因为在仿射开集
$D_+(f_i) = D_+(f_i^d)$ 上，相应的模由性质 (6) 给出同构
$N_{(f_i)} \cong N_{(f_i^d)} \cong N'_{(f_i^d)} \cong N'_{(f_i)}$。
由《性质》引理 \ref{properties-lemma-proj-quasi-coherent}，
存在典范同构 $\widetilde{N} \to \widetilde{M}$。
复合 $\widetilde{N} \to \widetilde{M} \to \widetilde{N'}$
就是上述同构（证明从略；提示：在标准仿射开集上检查）。
因此映射 $M \to N'$ 诱导同构
$\widetilde{M} \to \widetilde{N'}$。
令 $K = \Ker(M \to N')$ 且 $Q = \Coker(M \to N')$。
回忆函子 $M \mapsto \widetilde{M}$ 正合；参见《构造》引理
\ref{constructions-lemma-proj-sheaves}。
因此 $\widetilde{K} = 0$ 且 $\widetilde{Q} = 0$。
由上一段的结果可知，$K_n = 0$ 且 $Q_n = 0$
对 $d | n$ 且 $n \gg 0$ 成立。此时终于可见使用 $N'$ 而不是
$N$ 的优势：函子 $M \leadsto N'$ 与扭转相容
（由构造立即可见）。因此，把 $M$ 替换为 $M(q)$
并重复整个论证，可知 $K_n = 0$ 且 $Q_n = 0$
对 $n \equiv q \bmod d$ 且 $n \gg 0$ 成立。
因为模 $n$ 只有有限多个同余类，证明完成。
\end{proof}

\noindent
设 $A$ 为 Noether 分次环。回忆
$A_+ = \bigoplus_{n > 0} A_n$ 是无关理想。
由《代数》引理 \ref{algebra-lemma-graded-Noetherian} 与
\ref{algebra-lemma-S-plus-generated}，环 $A_0$ 是 Noether 的，
并且 $A$ 在 $A_0$ 上由有限多个正次数齐次元
$f_1, \ldots, f_r$ 生成。
令 $d = \text{lcm}(\deg(f_i))$。设 $M$ 为分次 $A$-模。
在这种情形下，若齐次元 $x \in M$ 满足
$$
(A_+ x)_{nd} = 0\text{ 对所有 }n \gg 0
$$
则称它是{\it 无关的}\footnote{这是非标准记号。}。
若 $x \in M$ 齐次且无关，而 $f \in A$ 齐次，
则 $fx$ 也无关。因此，无关元素的集合生成一个分次子模
$M_{irrelevant} \subset M$。若 $M$ 的每个齐次元都无关，
则称 $M$ 是{\it 无关的}，即 $M_{irrelevant} = M$。
若 $M$ 有限生成，则这等价于
$M_{nd} = 0$ 对 $n \gg 0$ 成立。
以 $\text{Mod}_A$ 表示分次 $A$-模范畴，
以 $\text{Mod}^{fg}_A$ 表示有限生成对象组成的满子范畴，
并以 $\text{Mod}^{fg}_{A, irrelevant}$ 表示无关模组成的满子范畴。

\begin{proposition}
\label{coherent-proposition-coherent-modules-on-proj-general}
设 $A$ 为 Noether 分次环。令 $X = \text{Proj}(A)$。
函子 $M \mapsto \widetilde M$ 诱导等价
$$
\text{Mod}^{fg}_A/\text{Mod}^{fg}_{A, irrelevant}
\longrightarrow
\textit{Coh}(\mathcal{O}_X)
$$
其拟逆由
$\mathcal{F} \longmapsto \bigoplus_{n \geq 0} \Gamma(X, \mathcal{F}(n))$
给出。
\end{proposition}

\begin{proof}
我们敦促读者先阅读 $A$ 由次数 $1$ 元生成的情形；
参见命题 \ref{coherent-proposition-coherent-modules-on-proj}。
设 $f_1, \ldots, f_r \in A$ 为 $A$ 的正次数齐次元，
它们在 $A_0$ 上生成该环。令 $d$ 为这些次数 $d_i$
的最小公倍数；这些是诸 $f_i$ 的次数。设 $M$ 为有限 $A$-模。
我们来说明，$\widetilde{M}$ 为零，当且仅当 $M$
是无关分次 $A$-模（按命题陈述上文的定义）。
为此，设 $x \in M$ 为齐次元。
选取充分小的 $k \in \mathbf{Z}$，并令
$N \to N'$ 与 $M \to N'$ 如引理
\ref{coherent-lemma-recover-tail-graded-module-general} 中所述。
还可以选取充分大的 $l$，使得
$M_n \to N_n$ 对 $n \geq l$ 为同构。
若 $\widetilde{M}$ 为零，则 $N = 0$。
因此，对任意齐次元 $f \in A_+$，若它满足
$\deg(f) + \deg(x) = nd$ 且 $nd > l$，则 $fx$ 为零，因为
$N_{nd} \to N'_{nd}$ 与 $M_{nd} \to N'_{nd}$ 都是同构。
故 $x$ 无关。
反之，设 $M$ 无关。则 $M_{nd}$ 对 $n \gg 0$ 为零
（参见命题前的讨论）。这显然蕴含
$M_{(f_i)} = M_{(f_i^{d/\deg(f_i)})} = 0$，
从而由构造可知 $\widetilde{M} = 0$。

\medskip\noindent
由此可知，子范畴 $\text{Mod}^{fg}_{A, irrelevant}$ 是
$\text{Mod}^{fg}_A$ 的 Serre 子范畴，因为它是正合函子
$M \mapsto \widetilde M$ 的核；参见《同调》引理
\ref{homology-lemma-kernel-exact-functor} 与《构造》引理
\ref{constructions-lemma-proj-sheaves}。
因此，箭头左侧的商范畴已在《同调》引理
\ref{homology-lemma-serre-subcategory-is-kernel} 中定义。
为定义命题中的函子，只需说明函子
$M \mapsto \widetilde M$ 把无关模送到 $0$；
这已在上文证明。

\medskip\noindent
由引理 \ref{coherent-lemma-coherent-on-proj-general}，所提出的拟逆有意义。
具体而言，该引理表明
$\mathcal{F} \longmapsto \bigoplus_{n \geq 0} \Gamma(X, \mathcal{F}(n))$
是函子 $\textit{Coh}(\mathcal{O}_X) \to \text{Mod}^{fg}_A$；
可将它与商函子
$\text{Mod}^{fg}_A \to \text{Mod}^{fg}_A/\text{Mod}^{fg}_{A, irrelevant}$
复合。

\medskip\noindent
由引理 \ref{coherent-lemma-recover-tail-graded-module-general}，
从左到右再到左的复合同构于恒等函子。
具体而言，设 $M$ 为有限分次 $A$-模，并取充分小的
$k \in \mathbf{Z}$，再令 $N \to N'$ 与 $M \to N'$
如引理 \ref{coherent-lemma-recover-tail-graded-module-general} 中所述。
于是 $M \to N'$ 的核与余核只在有限多个次数上非零，
因而都是无关的。此外，映射 $N \to N'$ 的核与余核
在所有充分大的、可被 $d$ 整除的次数上为零，
故它们也都是无关模。因此，$M \to N'$ 与 $N \to N'$
在商范畴中均为同构，正如所需。

\medskip\noindent
最后，设 $\mathcal{F}$ 为凝聚 $\mathcal{O}_X$-模。
令 $M = \bigoplus_{n \in \mathbf{Z}} \Gamma(X, \mathcal{F}(n))$，
将它视为分次 $A$-模，从而我们的函子把 $\mathcal{F}$ 送到
$M_{\geq 0} = \bigoplus_{n \geq 0} M_n$。
由《性质》引理 \ref{properties-lemma-proj-quasi-coherent}，
典范映射 $\widetilde M \to \mathcal{F}$ 是同构。
因为包含映射 $M_{\geq 0} \to M$ 定义同构
$\widetilde{M_{\geq 0}} \to \widetilde M$，可知
从右到左再到右的复合也同构于恒等函子。
\end{proof}










\section{沿射影态射的高阶直像}
\label{coherent-section-projective-pushforward}

\noindent
我们先陈述并证明基底仿射时的一个结果，
随后由此推出关于射影态射的一些结果。

\begin{lemma}
\label{coherent-lemma-coherent-proper-ample}
设 $R$ 为 Noether 环。设 $X \to \Spec(R)$ 为固有态射。
设 $\mathcal{L}$ 为 $X$ 上的丰沛可逆层。设 $\mathcal{F}$
为凝聚 $\mathcal{O}_X$-模。
\begin{enumerate}
\item 分次环
$A = \bigoplus_{d \geq 0} H^0(X, \mathcal{L}^{\otimes d})$
是有限生成 $R$-代数。
\item 存在 $r \geq 0$、
$d_1, \ldots, d_r \in \mathbf{Z}$ 以及满射
$$
\bigoplus\nolimits_{j = 1, \ldots, r} \mathcal{L}^{\otimes d_j}
\longrightarrow \mathcal{F}.
$$
\item 对任意 $p$，上同调群 $H^p(X, \mathcal{F})$ 是有限
$R$-模。
\item 若 $p > 0$，则
$H^p(X, \mathcal{F} \otimes_{\mathcal{O}_X} \mathcal{L}^{\otimes d}) = 0$
对所有充分大的 $d$ 成立。
\item 对任意 $k \in \mathbf{Z}$，分次 $A$-模
$$
\bigoplus\nolimits_{d \geq k}
H^0(X, \mathcal{F} \otimes_{\mathcal{O}_X} \mathcal{L}^{\otimes d})
$$
是有限 $A$-模。
\end{enumerate}
\end{lemma}

\begin{proof}
由《态射》引理
\ref{morphisms-lemma-finite-type-over-affine-ample-very-ample}，
存在 $d > 0$ 及浸入 $i : X \to \mathbf{P}^n_R$，使得
$\mathcal{L}^{\otimes d} \cong i^*\mathcal{O}_{\mathbf{P}^n_R}(1)$。
由于 $X$ 在 $R$ 上固有，态射 $i$ 是闭浸入
（《态射》引理 \ref{morphisms-lemma-image-proper-scheme-closed}）。
因此有
$H^i(X, \mathcal{G}) = H^i(\mathbf{P}^n_R, i_*\mathcal{G})$，
其中 $\mathcal{G}$ 是 $X$ 上任意准凝聚层
（由引理 \ref{coherent-lemma-relative-affine-cohomology} 以及闭浸入为仿射态射这一事实；
参见《态射》引理 \ref{morphisms-lemma-closed-immersion-affine}）。
此外，若 $\mathcal{G}$ 凝聚，则 $i_*\mathcal{G}$ 也凝聚
（引理 \ref{coherent-lemma-i-star-equivalence}）。
以下将不再特别提及这些事实。

\medskip\noindent
(1) 的证明。令 $S = R[T_0, \ldots, T_n]$，从而
$\mathbf{P}^n_R = \text{Proj}(S)$。
注意，$A$ 是 $S$-代数（但环映射 $S \to A$
不是分次环同态，因为 $S_n$ 映入 $A_{dn}$）。
由投影公式
（《上同调》引理 \ref{cohomology-lemma-projection-formula}），有
$$
i_*(\mathcal{L}^{\otimes nd + q}) =
i_*(\mathcal{L}^{\otimes q})
\otimes_{\mathcal{O}_{\mathbf{P}^n_R}}
\mathcal{O}_{\mathbf{P}^n_R}(n)
$$
对所有 $n \in \mathbf{Z}$ 成立。由引理
\ref{coherent-lemma-coherent-projective}，可知
$\bigoplus_{n \geq 0} A_{nd + q}$ 是有限分次 $S$-模。
由于
$A = \bigoplus_{q \in \{0, \ldots, d - 1} \bigoplus_{n \geq 0} A_{nd + q}$，
可知 $A$ 作为 $S$-代数是有限的，因而 (1) 成立。

\medskip\noindent
(2) 的证明。这由《性质》命题
\ref{properties-proposition-characterize-ample} 推出。

\medskip\noindent
(3) 的证明。应用引理 \ref{coherent-lemma-coherent-projective}，并使用
$H^p(X, \mathcal{F}) = H^p(\mathbf{P}^n_R, i_*\mathcal{F})$。

\medskip\noindent
(4) 的证明。固定 $p > 0$。由投影公式，有
$$
i_*(\mathcal{F} \otimes_{\mathcal{O}_X} \mathcal{L}^{\otimes nd + q}) =
i_*(\mathcal{F} \otimes_{\mathcal{O}_X} \mathcal{L}^{\otimes q})
\otimes_{\mathcal{O}_{\mathbf{P}^n_R}}
\mathcal{O}_{\mathbf{P}^n_R}(n)
$$
对所有 $n \in \mathbf{Z}$ 成立。由引理
\ref{coherent-lemma-coherent-projective} 可知
$H^p(X, \mathcal{F} \otimes \mathcal{L}^{nd + q}) = 0$
对 $n \gg 0$ 成立。因为整数模 $d$ 只有有限多个同余类，
这证明了 (4)。

\medskip\noindent
(5) 的证明。固定整数 $k$。令
{\scriptsize $M = \bigoplus_{n \geq k} H^0(X, \mathcal{F} \otimes \mathcal{L}^{\otimes n})$}。
如上论证可知，$\bigoplus_{nd + q \geq k} M_{nd + q}$
是有限分次 $S$-模。由于
$M = \bigoplus_{q \in \{0, \ldots, d - 1\}}
\bigoplus_{nd + q \geq k} M_{nd + q}$，可知 $M$
是有限 $S$-模。由于 $S$-模结构经由环映射 $S \to A$
分解，可知 $M$ 是有限 $A$-模。
\end{proof}

\begin{lemma}
\label{coherent-lemma-kill-by-twisting}
设 $f : X \to S$ 为概形态射。
设 $\mathcal{F}$ 为准凝聚 $\mathcal{O}_X$-模。
设 $\mathcal{L}$ 为 $X$ 上的可逆层。假设
\begin{enumerate}
\item $S$ 是 Noether 的；
\item $f$ 是固有的；
\item $\mathcal{F}$ 是凝聚的；并且
\item $\mathcal{L}$ 在 $X/S$ 上相对丰沛。
\end{enumerate}
则存在 $n_0$，使得对所有 $n \geq n_0$，都有
$$
R^pf_*\left(\mathcal{F} \otimes_{\mathcal{O}_X} \mathcal{L}^{\otimes n}\right)
=
0
$$
对所有 $p > 0$ 成立。
\end{lemma}

\begin{proof}
选取有限仿射开覆盖 $S = \bigcup V_j$，并令
$X_j = f^{-1}(V_j)$。
显然，若能对有限多个系统
$(X_j \to V_j, \mathcal{L}|_{X_j}, \mathcal{F}|_{V_j})$
中的每一个解决该问题，结果即告成立。因此可设 $S$ 仿射。
在这种情形下，
$R^pf_*(\mathcal{F} \otimes \mathcal{L}^{\otimes n})$
为零等价于
$H^p(X, \mathcal{F} \otimes \mathcal{L}^{\otimes n})$ 为零；参见引理
\ref{coherent-lemma-quasi-coherence-higher-direct-images-application}。
因此，所需消失由引理 \ref{coherent-lemma-coherent-proper-ample} 推出
（该引理适用，因为由《态射》引理
\ref{morphisms-lemma-finite-type-over-affine-ample-very-ample}，
$\mathcal{L}$ 在 $X$ 上丰沛）。
\end{proof}

\begin{lemma}
\label{coherent-lemma-locally-projective-pushforward}
设 $S$ 为局部 Noether 概形。
设 $f : X \to S$ 为局部射影态射。
设 $\mathcal{F}$ 为凝聚 $\mathcal{O}_X$-模。
则 $R^if_*\mathcal{F}$ 是凝聚 $\mathcal{O}_S$-模，
这对所有 $i \geq 0$ 成立。
\end{lemma}

\begin{proof}
首先注意，局部射影态射是固有的
（《态射》引理 \ref{morphisms-lemma-locally-projective-proper}），
因而是有限型的。特别地，$X$ 是局部 Noether 的
（《态射》引理 \ref{morphisms-lemma-finite-type-noetherian}），
故该陈述有意义。此外，由引理
\ref{coherent-lemma-quasi-coherence-higher-direct-images}，
诸层 $R^pf_*\mathcal{F}$ 都准凝聚。

\medskip\noindent
说明这些之后，该陈述在 $S$ 上是局部的（例如由《上同调》引理
\ref{cohomology-lemma-localize-higher-direct-images}）。
因此可设 $S = \Spec(R)$ 是 Noether 环的谱，且
$X$ 是 $\mathbf{P}^n_R$ 的闭子概形，其中 $n$ 取某个值；参见《态射》引理
\ref{morphisms-lemma-characterize-locally-projective}。
在这种情形下，诸层 $R^pf_*\mathcal{F}$ 是与 $R$-模
$H^p(X, \mathcal{F})$ 相伴的准凝聚层；参见引理
\ref{coherent-lemma-quasi-coherence-higher-direct-images-application}。
因此，只需说明 $R$-模 $H^p(X, \mathcal{F})$
是有限 $R$-模（引理 \ref{coherent-lemma-coherent-Noetherian}）。
这由引理 \ref{coherent-lemma-coherent-proper-ample} 推出
（因为 $\mathcal{O}_{\mathbf{P}^n_R}(1)$ 限制到 $X$
后在 $X$ 上丰沛）。
\end{proof}








\section{丰沛可逆层与上同调}
\label{coherent-section-ample-cohomology}

\noindent
下面给出仿射基底上的固有概形之丰沛性判据，
它以扭转后的上同调消失来表述。

\begin{lemma}
\label{coherent-lemma-vanshing-gives-ample}
\begin{reference}
\cite[III Proposition 2.6.1]{EGA}
\end{reference}
设 $R$ 为 Noether 环。设 $f : X \to \Spec(R)$ 为固有态射。
设 $\mathcal{L}$ 为可逆 $\mathcal{O}_X$-模。
下列条件等价：
\begin{enumerate}
\item $\mathcal{L}$ 在 $X$ 上丰沛（这与许多其他条件等价；参见
《性质》命题 \ref{properties-proposition-characterize-ample} 与
《态射》引理
\ref{morphisms-lemma-finite-type-over-affine-ample-very-ample}）；
\item 对每个凝聚 $\mathcal{O}_X$-模 $\mathcal{F}$，存在
$n_0 \geq 0$，使得
$H^p(X, \mathcal{F} \otimes \mathcal{L}^{\otimes n}) = 0$
对所有 $n \geq n_0$ 及 $p > 0$ 成立；并且
\item 对每个准凝聚理想层
$\mathcal{I} \subset \mathcal{O}_X$，存在 $n \geq 1$，
使得 $H^1(X, \mathcal{I} \otimes \mathcal{L}^{\otimes n}) = 0$。
\end{enumerate}
\end{lemma}

\begin{proof}
蕴含关系 (1) $\Rightarrow$ (2) 由引理
\ref{coherent-lemma-coherent-proper-ample} 推出。
蕴含关系 (2) $\Rightarrow$ (3) 是平凡的。
蕴含关系 (3) $\Rightarrow$ (1) 是引理
\ref{coherent-lemma-quasi-compact-h1-zero-invertible}。
\end{proof}

\begin{lemma}
\label{coherent-lemma-surjective-finite-morphism-ample}
设 $R$ 为 Noether 环。设 $f : Y \to X$ 为在 $R$
上固有的概形之间的态射。设 $\mathcal{L}$ 为可逆
$\mathcal{O}_X$-模。假设 $f$ 有限且满射。
则 $\mathcal{L}$ 丰沛，当且仅当 $f^*\mathcal{L}$ 丰沛。
\end{lemma}

\begin{proof}
丰沛可逆层经拟仿射态射拉回后仍丰沛；参见《态射》引理
\ref{morphisms-lemma-pullback-ample-tensor-relatively-ample}。
这证明了其中一个蕴含关系，因为有限态射按定义是仿射的。

\medskip\noindent
假设 $f^*\mathcal{L}$ 丰沛。令 $P$ 为凝聚
$\mathcal{O}_X$-模 $\mathcal{F}$ 的如下性质：存在 $n_0$，
使得 $H^p(X, \mathcal{F} \otimes \mathcal{L}^{\otimes n}) = 0$
对所有 $n \geq n_0$ 及 $p > 0$ 成立。我们将证明 $P$
对任意凝聚 $\mathcal{O}_X$-模 $\mathcal{F}$ 成立；由引理
\ref{coherent-lemma-vanshing-gives-ample}，这蕴含 $\mathcal{L}$ 丰沛。
我们要应用引理 \ref{coherent-lemma-property-higher-rank-cohomological}。
因此必须对 $P$ 验证该引理的 (1)、(2) 与 (3)。
性质 (1) 由与层的短正合列相伴的上同调长正合列，以及
与可逆层张量是正合函子这一事实推出。性质 (2) 由
$H^p(X, -)$ 是加性函子推出。为证明 (3)，设 $Z \subset X$
为泛点是 $\xi$ 的整闭子概形。
设 $\mathcal{F}$ 为 $Y$ 上的凝聚层，使得
$f_*\mathcal{F}$ 的支撑等于 $Z$，且
$(f_*\mathcal{F})_\xi$ 被 $\mathfrak m_\xi$ 零化；
参见引理 \ref{coherent-lemma-finite-morphism-Noetherian}。我们断言，
取 $\mathcal{G} = f_*\mathcal{F}$ 即可。只需验证引理
\ref{coherent-lemma-property-higher-rank-cohomological} 的 (3)(c)。
故假设 $\mathcal{J} \subset \mathcal{O}_X$ 为准凝聚理想层，
并满足 $\mathcal{J}_\xi = \mathcal{O}_{X, \xi}$。
有限态射是仿射的，故由引理
\ref{coherent-lemma-affine-morphism-projection-ideal} 可知
$\mathcal{J}\mathcal{G} = f_*(f^{-1}\mathcal{J}\mathcal{F})$。
此外，正如引理 \ref{coherent-lemma-affine-morphism-projection-ideal}
的证明所指出，层 $f^{-1}\mathcal{J}\mathcal{F}$ 是凝聚
$\mathcal{O}_Y$-模。由于 $\mathcal{L}$ 丰沛，由引理
\ref{coherent-lemma-vanshing-gives-ample} 可知存在 $n_0$，使得
$$
H^p(Y, f^{-1}\mathcal{J}\mathcal{F}
\otimes_{\mathcal{O}_Y} f^*\mathcal{L}^{\otimes n}) = 0,
$$
对 $n \geq n_0$ 及 $p > 0$ 成立。因为 $f$ 有限，因而仿射，故
\begin{align*}
H^p(X, \mathcal{J}\mathcal{G} \otimes_{\mathcal{O}_X}
\mathcal{L}^{\otimes n})
& =
H^p(X, f_*(f^{-1}\mathcal{J}\mathcal{F}) \otimes_{\mathcal{O}_X}
\mathcal{L}^{\otimes n}) \\
& =
H^p(X, f_*(f^{-1}\mathcal{J}\mathcal{F} \otimes_{\mathcal{O}_Y}
f^*\mathcal{L}^{\otimes n})) \\
& =
H^p(Y, f^{-1}\mathcal{J}\mathcal{F} \otimes_{\mathcal{O}_Y}
f^*\mathcal{L}^{\otimes n}) = 0
\end{align*}
这里使用了投影公式
（《上同调》引理 \ref{cohomology-lemma-projection-formula}）及引理
\ref{coherent-lemma-relative-affine-cohomology}。
因此，准凝聚子层 $\mathcal{G}' = \mathcal{J}\mathcal{G}$
满足 $P$。这正如所需地验证了引理
\ref{coherent-lemma-property-higher-rank-cohomological} 的性质 (3)(c)。
\end{proof}

\noindent
上同调具有函子性。特别地，给定环化空间 $X$、可逆
$\mathcal{O}_X$-模 $\mathcal{L}$ 以及截面
$s \in \Gamma(X, \mathcal{L})$，可得到映射
$$
H^p(X, \mathcal{F})
\longrightarrow
H^p(X, \mathcal{F} \otimes_{\mathcal{O}_X} \mathcal{L}), \quad
\xi \longmapsto s\xi
$$
它由映射
$\mathcal{F} \to \mathcal{F} \otimes_{\mathcal{O}_X} \mathcal{L}$
诱导；该映射是乘以 $s$。令
$\Gamma_*(X, \mathcal{L}) =
\bigoplus_{n \geq 0} \Gamma(X, \mathcal{L}^{\otimes n})$
为分次环；参见《模》定义 \ref{modules-definition-gamma-star}。
给定 $\mathcal{O}_X$-模层 $\mathcal{F}$ 及整数
$p \geq 0$，令
$$
H^p_*(X, \mathcal{L}, \mathcal{F}) =
\bigoplus\nolimits_{n \in \mathbf{Z}} H^p(X,
\mathcal{F} \otimes_{\mathcal{O}_X} \mathcal{L}^{\otimes n})
$$
由上述乘法，它是分次 $\Gamma_*(X, \mathcal{L})$-模。
警告：记号 $H^p_*(X, \mathcal{L}, \mathcal{F})$ 并非标准记号。

\begin{lemma}
\label{coherent-lemma-invert-s-cohomology}
设 $X$ 为概形。设 $\mathcal{L}$ 为 $X$ 上的可逆层。
设 $s \in \Gamma(X, \mathcal{L})$。
设 $\mathcal{F}$ 为准凝聚 $\mathcal{O}_X$-模。
若 $X$ 拟紧且拟分离，则典范映射
$$
H^p_*(X, \mathcal{L}, \mathcal{F})_{(s)}
\longrightarrow
H^p(X_s, \mathcal{F})
$$
把 $\xi/s^n$ 映到 $s^{-n}\xi$，并且是同构。
\end{lemma}

\begin{proof}
注意，当 $p = 0$ 时，这就是《性质》引理
\ref{properties-lemma-invert-s-sections}。
我们用归纳原理（引理 \ref{coherent-lemma-induction-principle}）证明该陈述；
对拟紧开集 $U \subset X$，令 $P(U)$ 为如下性质：
对所有 $p \geq 0$，映射
$$
H^p_*(U, \mathcal{L}, \mathcal{F})_{(s)}
\longrightarrow
H^p(U_s, \mathcal{F})
$$
是同构。

\medskip\noindent
若 $U$ 仿射，则由引理
\ref{coherent-lemma-quasi-coherent-affine-cohomology-zero} 与《性质》引理
\ref{properties-lemma-affine-cap-s-open}，
上述箭头两端对 $p > 0$ 均为零，故陈述成立。
若 $P$ 对 $U$、$V$ 及 $U \cap V$ 成立，
则可使用 Mayer--Vietoris 序列
（《上同调》引理 \ref{cohomology-lemma-mayer-vietoris}），
得到长正合列之间的映射
\begingroup\small
$$
\xymatrix{
H^{p - 1}_*(U \cap V, \mathcal{L}, \mathcal{F})_{(s)} \ar[r] \ar[d] &
H^p_*(U \cup V, \mathcal{L}, \mathcal{F})_{(s)} \ar[r] \ar[d] &
H^p_*(U, \mathcal{L}, \mathcal{F})_{(s)}
\oplus
H^p_*(V, \mathcal{L}, \mathcal{F})_{(s)} \ar[d] \\
H^{p - 1}(U_s \cap V_s, \mathcal{F}) \ar[r]&
H^p(U_s \cup V_s, \mathcal{F}) \ar[r] &
H^p(U_s, \mathcal{F})
\oplus
H^p(V_s, \mathcal{F})
}
$$
\endgroup
（这里只展示一个片段）。注意，$U_s \cap V_s = (U \cap V)_s$，
并且 $U_s \cup V_s = (U \cup V)_s$。因此左右两个竖直映射
都是同构（图中未画出的右侧再一个映射和左侧再一个映射也如此）。
由五引理（《同调》引理 \ref{homology-lemma-five-lemma}），
可知 $P(U \cup V)$ 成立。这就完成了证明。
\end{proof}

\begin{lemma}
\label{coherent-lemma-section-affine-open-kills-classes}
设 $X$ 为概形。
设 $\mathcal{L}$ 为可逆 $\mathcal{O}_X$-模。
设 $s \in \Gamma(X, \mathcal{L})$ 为截面。假设
\begin{enumerate}
\item $X$ 拟紧且拟分离；并且
\item $X_s$ 仿射。
\end{enumerate}
则对每个准凝聚 $\mathcal{O}_X$-模 $\mathcal{F}$、
每个 $p > 0$ 以及所有 $\xi \in H^p(X, \mathcal{F})$，
都存在 $n \geq 0$，使得 $s^n\xi = 0$ 在
$H^p(X, \mathcal{F} \otimes_{\mathcal{O}_X} \mathcal{L}^{\otimes n})$
中成立。
\end{lemma}

\begin{proof}
回忆，$H^p(X_s, \mathcal{G})$
对每个准凝聚模 $\mathcal{G}$ 都为零；这是引理
\ref{coherent-lemma-quasi-coherent-affine-cohomology-zero}。
因此本引理由引理 \ref{coherent-lemma-invert-s-cohomology} 推出。
\end{proof}

\noindent
下述引理的更一般版本见《极限》引理
\ref{limits-lemma-ample-on-reduction}。

\begin{lemma}
\label{coherent-lemma-ample-on-reduction}
设 $i : Z \to X$ 为 Noether 概形的闭浸入，
并诱导底拓扑空间的同胚。
设 $\mathcal{L}$ 为 $X$ 上的可逆层。
则 $i^*\mathcal{L}$ 在 $Z$ 上丰沛，当且仅当
$\mathcal{L}$ 在 $X$ 上丰沛。
\end{lemma}

\begin{proof}
若 $\mathcal{L}$ 丰沛，则例如由《态射》引理
\ref{morphisms-lemma-pullback-ample-tensor-relatively-ample}，
$i^*\mathcal{L}$ 丰沛。假设 $i^*\mathcal{L}$ 丰沛。
必须证明 $\mathcal{L}$ 在 $X$ 上丰沛。
设 $\mathcal{I} \subset \mathcal{O}_X$ 为定义闭子概形
$Z$ 的凝聚理想层。由于在集合意义下 $i(Z) = X$，
由引理 \ref{coherent-lemma-power-ideal-kills-sheaf} 可知，
$\mathcal{I}^n = 0$ 对某个 $n$ 成立。
考虑由闭子概形组成的序列
$$
X = Z_n \supset Z_{n - 1} \supset Z_{n - 2} \supset \ldots \supset Z_1 = Z
$$
它们分别由
$0 = \mathcal{I}^n \subset \mathcal{I}^{n - 1} \subset \ldots \subset
\mathcal{I}$ 定义。于是每个闭浸入 $Z_{i - 1} \to Z_{i}$
都由平方为零的凝聚理想层定义。这样便归约到
$\mathcal{I}^2 = 0$ 的情形。

\medskip\noindent
考虑准凝聚 $\mathcal{O}_X$-模的短正合列
$$
0 \to \mathcal{I} \to \mathcal{O}_X \to i_*\mathcal{O}_Z \to 0
$$
与 $\mathcal{L}^{\otimes n}$ 张量后，得到短正合列
\begin{equation}
\label{coherent-equation-ses}
0 \to \mathcal{I} \otimes_{\mathcal{O}_X} \mathcal{L}^{\otimes n}
\to \mathcal{L}^{\otimes n} \to i_*i^*\mathcal{L}^{\otimes n} \to 0
\end{equation}
因为 $\mathcal{I}^2 = 0$，可使用《态射》引理
\ref{morphisms-lemma-i-star-equivalence}，把 $\mathcal{I}$
看作准凝聚 $\mathcal{O}_Z$-模；于是
$\mathcal{I} \otimes_{\mathcal{O}_X} \mathcal{L}^{\otimes n} =
\mathcal{I} \otimes_{\mathcal{O}_Z} i^*\mathcal{L}^{\otimes n}$，
这里作了显然的记号滥用。此外，该层在 $Z$ 上的上同调
典范地等同于它在 $X$ 上的上同调（因为 $i$ 是同胚）。

\medskip\noindent
设 $x \in X$ 为一点，并以 $z \in Z$ 表示相应的点。
因为 $i^*\mathcal{L}$ 丰沛，存在 $n$ 及截面
$s \in \Gamma(Z, i^*\mathcal{L}^{\otimes n})$，使得
$z \in Z_s$ 且 $Z_s$ 仿射。把 $s$ 提升为
$\mathcal{L}^{\otimes n}$ 在 $X$ 上的截面之障碍，是来自层的短正合列
(\ref{coherent-equation-ses}) 的连接像
$$
\xi = \partial s \in
H^1(X, \mathcal{I} \otimes_{\mathcal{O}_X} \mathcal{L}^{\otimes n}) =
H^1(Z, \mathcal{I} \otimes_{\mathcal{O}_Z} i^*\mathcal{L}^{\otimes n})
$$
若把 $s$ 替换为 $s^{e + 1}$，则 $\xi$ 被替换为
$\partial(s^{e + 1}) = (e + 1) s^e \xi$；其所在群为
$H^1(Z, \mathcal{I} \otimes_{\mathcal{O}_Z} i^*\mathcal{L}^{\otimes (e + 1)n})$，
因为
$$
0 \to
\bigoplus\nolimits_{m \geq 0}
\mathcal{I} \otimes_{\mathcal{O}_X} \mathcal{L}^{\otimes m} \to
\bigoplus\nolimits_{m \geq 0}
\mathcal{L}^{\otimes m} \to
\bigoplus\nolimits_{m \geq 0}
i_*i^*\mathcal{L}^{\otimes m} \to 0
$$
的连接映射由《上同调》引理
\ref{cohomology-lemma-boundary-derivation-over-cup-product} 是导子。
由引理 \ref{coherent-lemma-section-affine-open-kills-classes}，
可知 $s^e \xi$ 对充分大的 $e$ 为零。
因此，把 $s$ 替换为某个幂后，可以假设 $s$ 是截面
$s' \in \Gamma(X, \mathcal{L}^{\otimes n})$ 的像。
于是 $X_{s'}$ 是开子概形，而 $Z_s \to X_{s'}$
是 Noether 概形之间的满射闭浸入，且 $Z_s$ 仿射。
因此，由引理 \ref{coherent-lemma-image-affine-finite-morphism-affine-Noetherian}，
$X_s$ 仿射，于是可知 $\mathcal{L}$ 丰沛。
\end{proof}

\noindent
下述引理的更一般版本见《极限》引理
\ref{limits-lemma-thickening-quasi-affine}。

\begin{lemma}
\label{coherent-lemma-thickening-quasi-affine}
设 $i : Z \to X$ 为 Noether 概形的闭浸入，
并诱导底拓扑空间的同胚。
则 $X$ 拟仿射，当且仅当 $Z$ 拟仿射。
\end{lemma}

\begin{proof}
回忆，概形拟仿射，当且仅当其结构层丰沛；参见《性质》引理
\ref{properties-lemma-quasi-affine-O-ample}。
因此，若 $Z$ 拟仿射，则 $\mathcal{O}_Z$ 丰沛；
由引理 \ref{coherent-lemma-ample-on-reduction}，$\mathcal{O}_X$ 丰沛；
从而 $X$ 拟仿射。反向阅读刚才给出的论证，
便得到逆命题的一个证明；它也可以用初等方法看出。
\end{proof}

\begin{lemma}
\label{coherent-lemma-affine-in-presence-ample}
设 $X$ 为概形。设 $\mathcal{L}$ 为丰沛可逆
$\mathcal{O}_X$-模。设 $n_0$ 为整数。
若 $H^p(X, \mathcal{L}^{\otimes -n}) = 0$
对 $n \geq n_0$ 及 $p > 0$ 成立，则 $X$ 仿射。
\end{lemma}

\begin{proof}
我们断言，$H^p(X, \mathcal{F}) = 0$ 对每个准凝聚
$\mathcal{O}_X$-模及 $p > 0$ 成立。由《性质》定义
\ref{properties-definition-ample}，$X$ 拟紧，
故该断言由引理 \ref{coherent-lemma-quasi-compact-h1-zero-covering}
完成证明。由《性质》引理
\ref{properties-lemma-ample-separated}，概形 $X$ 分离。
设 $X$ 被 $e + 1$ 个仿射开集覆盖。于是
$H^p(X, \mathcal{F}) = 0$ 对 $p > e$ 成立；参见引理
\ref{coherent-lemma-vanishing-nr-affines}。因此可对 $p$ 作降归纳
来证明该断言。把 $\mathcal{F}$ 写成有限型准凝聚模的滤过余极限
（《性质》引理
\ref{properties-lemma-quasi-coherent-colimit-finite-type}），
并使用《上同调》引理
\ref{cohomology-lemma-quasi-separated-cohomology-colimit}，
可设 $\mathcal{F}$ 为有限型。于是可以选取 $n > n_0$，
使得 $\mathcal{F} \otimes_{\mathcal{O}_X} \mathcal{L}^{\otimes n}$
由整体截面生成；参见《性质》命题
\ref{properties-proposition-characterize-ample}。
这意味着存在短正合列
$$
0 \to \mathcal{F}' \to
\bigoplus\nolimits_{i \in I} \mathcal{L}^{\otimes -n}
\to \mathcal{F} \to 0
$$
其中 $I$ 是某个集合（事实上可选取有限的 $I$）。
由归纳假设，$H^{p + 1}(X, \mathcal{F}') = 0$；
而由假设（结合已经使用过的上同调与余极限的交换性），有
$H^p(X, \bigoplus_{i \in I} \mathcal{L}^{\otimes -n}) = 0$。
由上同调长正合列，可知
$H^p(X, \mathcal{F}) = 0$，正如所需。
\end{proof}

\begin{lemma}
\label{coherent-lemma-affine-if-quasi-affine}
设 $X$ 为拟仿射概形。
若 $H^p(X, \mathcal{O}_X) = 0$ 对 $p > 0$ 成立，
则 $X$ 仿射。
\end{lemma}

\begin{proof}
由《性质》引理 \ref{properties-lemma-quasi-affine-O-ample}，
$\mathcal{O}_X$ 丰沛，故结论由引理
\ref{coherent-lemma-affine-in-presence-ample} 推出。
\end{proof}







\section{Chow 引理}
\label{coherent-section-chows-lemma}

\noindent
本节证明 Noether 情形的 Chow 引理
（引理 \ref{coherent-lemma-chow-Noetherian}）。
在《极限》第 \ref{limits-section-chows-lemma} 节中，
我们证明非 Noether 情形的一些变体。

\begin{lemma}
\label{coherent-lemma-chow-Noetherian}
\begin{reference}
\cite[II Theorem 5.6.1(a)]{EGA}
\end{reference}
设 $S$ 为 Noether 概形。
设 $f : X \to S$ 为有限型分离态射。
则存在 $n \geq 0$ 及交换图
$$
\xymatrix{
X \ar[rd] & X' \ar[d] \ar[l]^\pi \ar[r] & \mathbf{P}^n_S \ar[dl] \\
& S &
}
$$
其中 $X' \to \mathbf{P}^n_S$ 是浸入，且
$\pi : X' \to X$ 固有并满射。此外，还可以作如此安排：
存在稠密开子概形 $U \subset X$，使得
$\pi^{-1}(U) \to U$ 为同构。
\end{lemma}

\begin{proof}
在余下证明中遇到的所有概形，都将在 Noether 概形 $S$
上有限型，因而是 Noether 的
（参见《态射》引理 \ref{morphisms-lemma-finite-type-noetherian}）。
它们之间的所有态射都自动拟紧、局部有限型且拟分离；参见
《态射》引理 \ref{morphisms-lemma-permanence-finite-type} 及《性质》引理
\ref{properties-lemma-locally-Noetherian-quasi-separated} 与
\ref{properties-lemma-morphism-Noetherian-schemes-quasi-compact}。

\medskip\noindent
概形 $X$ 只有有限多个不可约分支
（《性质》引理 \ref{properties-lemma-Noetherian-irreducible-components}）。
设 $X = X_1 \cup \ldots \cup X_r$ 是 $X$
分解为不可约分支的表达式。
令 $\eta_i \in X_i$ 为泛点。对每个点 $x \in X$，
存在仿射开集 $U_x \subset X$，它包含 $x$ 以及所有泛点
$\eta_i$；参见《性质》引理
\ref{properties-lemma-point-and-maximal-points-affine}。
因为 $X$ 拟紧，可以找到有限仿射开覆盖
$X = U_1 \cup \ldots \cup U_m$，使得每个 $U_i$
都包含 $\eta_1, \ldots, \eta_r$。特别地，可知开集
$U = U_1 \cap \ldots \cap U_m \subset X$ 稠密。
以下只会用到这一点，以及诸 $U_i$ 是覆盖 $X$ 的仿射开集这一事实。

\medskip\noindent
令 $X^* \subset X$ 为 $U \to X$ 的概形论闭包；参见《态射》定义
\ref{morphisms-definition-scheme-theoretic-image}。
令 $U_i^* = X^* \cap U_i$。注意，$U_i^*$ 是 $U_i$
的闭子概形，因而 $U_i^*$ 仿射。由于 $U$ 在 $X$ 中稠密，
态射 $X^* \to X$ 是满射闭浸入。它在 $U$ 上是同构。
因此可以把 $X$ 替换为 $X^*$，把 $U_i$ 替换为 $U_i^*$，
并假设 $U$ 在 $X$ 中概形论稠密；参见《态射》定义
\ref{morphisms-definition-scheme-theoretically-dense}。

\medskip\noindent
由《态射》引理 \ref{morphisms-lemma-quasi-projective-finite-type-over-S}，
可以找到浸入 $j_i : U_i \to \mathbf{P}_S^{n_i}$，
这对每个 $i$ 成立。
由《态射》引理 \ref{morphisms-lemma-quasi-compact-immersion}，
可以找到闭子概形 $Z_i \subset \mathbf{P}_S^{n_i}$，
使得 $j_i : U_i \to Z_i$ 是概形论稠密开浸入。
注意，$Z_i \to S$ 固有；参见《态射》引理
\ref{morphisms-lemma-locally-projective-proper}。考虑态射
$$
j = (j_1|_U, \ldots, j_m|_U) : U \longrightarrow
\mathbf{P}_S^{n_1} \times_S \ldots \times_S \mathbf{P}_S^{n_m}.
$$
由上述引理，可以找到闭子概形 $Z$，它位于
$\mathbf{P}_S^{n_1} \times_S \ldots \times_S \mathbf{P}_S^{n_m}$
中，并且 $j : U \to Z$ 是开浸入，
并且 $U$ 在 $Z$ 中概形论稠密。态射 $Z \to S$ 固有。
考虑第 $i$ 个投影
$$
\text{pr}_i|_Z : Z \longrightarrow \mathbf{P}^{n_i}_S.
$$
该态射经由 $Z_i$ 分解（参见《态射》引理
\ref{morphisms-lemma-factor-factor}）。以 $p_i : Z \to Z_i$
表示诱导态射。例如由《态射》引理
\ref{morphisms-lemma-image-proper-scheme-closed}，这是固有态射。
此时已有 $U \subset U_i \subset Z_i$ 是概形论稠密开浸入。
此外，可把 $Z$ 看作“对角”态射
$U \to Z_1 \times_S \ldots \times_S Z_m$ 的概形论像。

\medskip\noindent
令 $V_i = p_i^{-1}(U_i)$。注意，
$p_i|_{V_i} : V_i \to U_i$ 固有。令
$X' = V_1 \cup \ldots \cup V_m$。由构造，$X'$ 可浸入概形
$\mathbf{P}^{n_1}_S \times_S \ldots \times_S \mathbf{P}^{n_m}_S$。
因此，由 Segre 嵌入（参见《构造》引理
\ref{constructions-lemma-segre-embedding}），可知 $X'$
可浸入 $S$ 上的某个射影空间。

\medskip\noindent
我们断言，诸态射 $p_i|_{V_i}: V_i \to U_i$
可粘合为态射 $X' \to X$。事实上，显然 $p_i|_U$
是从 $U$ 到 $U$ 的恒等映射。由构造，$U \subset X'$
概形论稠密，因而它在开子概形 $V_i \cap V_j$ 中也概形论稠密。
于是
$p_i|_{V_i \cap V_j} = p_j|_{V_i \cap V_j}$，
这里把两端视为到分离 $S$-概形 $X$ 的态射；参见《态射》引理
\ref{morphisms-lemma-equality-of-morphisms}。
把所得态射记为 $\pi : X' \to X$。

\medskip\noindent
我们断言 $\pi^{-1}(U_i) = V_i$。
由于 $\pi|_{V_i} = p_i|_{V_i}$，可知
$V_i \subset \pi^{-1}(U_i)$。考虑交换图
$$
\xymatrix{
V_i \ar[r] \ar[rd]_{p_i|_{V_i}} & \pi^{-1}(U_i) \ar[d] \\
& U_i
}
$$
因为 $V_i \to U_i$ 固有，由《态射》引理
\ref{morphisms-lemma-image-proper-scheme-closed}，
可知水平箭头的像闭。由于
$V_i \subset \pi^{-1}(U_i)$ 概形论稠密（因为它包含 $U$），
可知 $V_i = \pi^{-1}(U_i)$，正如所断言。

\medskip\noindent
这表明 $\pi^{-1}(U_i) \to U_i$ 可与固有态射
$p_i|_{V_i} : V_i \to U_i$ 识别。因此可知，$X$
有有限仿射覆盖 $X = \bigcup U_i$，使得 $\pi$
在该覆盖每个成员上的限制都固有。
于是由《态射》引理 \ref{morphisms-lemma-proper-local-on-the-base}，
可知 $\pi$ 固有。

\medskip\noindent
最后必须说明 $\pi^{-1}(U) = U$。为此，使用交换图
$$
\xymatrix{
U \ar[r] \ar[rd] & \pi^{-1}(U) \ar[d] \\
& U
}
$$
并使用 $\text{id}_U : U \to U$ 固有且 $U$
在 $\pi^{-1}(U)$ 中概形论稠密，便可如上作同样论证。
\end{proof}

\begin{remark}
\label{coherent-remark-chow-Noetherian}
在 Chow 引理 \ref{coherent-lemma-chow-Noetherian} 的情形中：
\begin{enumerate}
\item 态射 $\pi$ 实际上是 H-射影的（因而射影；参见《态射》引理
\ref{morphisms-lemma-H-projective}），因为态射
$X' \to \mathbf{P}^n_S \times_S X = \mathbf{P}^n_X$
是闭浸入（使用 $\pi$ 固有这一事实；参见《态射》引理
\ref{morphisms-lemma-image-proper-scheme-closed}）。
\item 可以假设 $\pi^{-1}(U)$ 在 $X'$ 中概形论稠密。
事实上，只需把 $X'$ 替换为 $\pi^{-1}(U)$ 的概形论闭包。
在这种情形下，可把 $U$ 看作 $X'$ 的概形论稠密开子概形。
参见《态射》第 \ref{morphisms-section-scheme-theoretic-image} 节。
\item 若 $X$ 既约，则可选取既约的 $X'$。这由 (2) 显然可见。
\end{enumerate}
\end{remark}






\section{凝聚层的高阶直像}
\label{coherent-section-proper-pushforward}

\noindent
本节证明一个基本事实：凝聚层沿固有态射的高阶直像仍凝聚。

\begin{proposition}
\label{coherent-proposition-proper-pushforward-coherent}
\begin{reference}
\cite[III Theorem 3.2.1]{EGA}
\end{reference}
设 $S$ 为局部 Noether 概形。
设 $f : X \to S$ 为固有态射。
设 $\mathcal{F}$ 为凝聚 $\mathcal{O}_X$-模。
则 $R^if_*\mathcal{F}$ 是凝聚 $\mathcal{O}_S$-模，
这对所有 $i \geq 0$ 成立。
\end{proposition}

\begin{proof}
由于该问题在 $S$ 上是局部的，可以假设 $S$ 是 Noether 概形。
由于固有态射是有限型的，可知在这种情形下 $X$ 也是 Noether 概形。
考虑性质 $\mathcal{P}$；它定义在 $X$ 上的凝聚层上，规则如下：
$$
\mathcal{P}(\mathcal{F}) \Leftrightarrow
R^pf_*\mathcal{F}\text{ is coherent for all }p \geq 0
$$
我们将用引理 \ref{coherent-lemma-property} 的结论证明
$\mathcal{P}$ 对 $X$ 上每个凝聚层都成立。

\medskip\noindent
设
$$
0 \to \mathcal{F}_1 \to \mathcal{F}_2 \to \mathcal{F}_3 \to 0
$$
为 $X$ 上凝聚层的短正合列。考虑高阶直像的长正合列
$$
R^{p - 1}f_*\mathcal{F}_3 \to
R^pf_*\mathcal{F}_1 \to
R^pf_*\mathcal{F}_2 \to
R^pf_*\mathcal{F}_3 \to
R^{p + 1}f_*\mathcal{F}_1
$$
显然，若诸层 $\mathcal{F}_i$ 中三项取二具有性质
$\mathcal{P}$，则第三个层的高阶直像在这个正合复形中夹在
两个凝聚层之间。因此，由引理
\ref{coherent-lemma-coherent-abelian-Noetherian} 与
\ref{coherent-lemma-coherent-Noetherian-quasi-coherent-sub-quotient}，
这些高阶直像也凝聚。故性质 $\mathcal{P}$ 对第三个层也成立。

\medskip\noindent
设 $Z \subset X$ 为整闭子概形。必须找到凝聚层
$\mathcal{F}$（在 $X$ 上），其支撑包含在 $Z$ 中；其在泛点
$\xi$（即 $Z$ 的泛点）处的茎是 $1$ 维向量空间，
标量域为 $\kappa(\xi)$，
并且 $\mathcal{P}$ 对 $\mathcal{F}$ 成立。
以 $g = f|_Z : Z \to S$ 表示 $f$ 的限制。
假设能找到凝聚层 $\mathcal{G}$（在 $Z$ 上），使得
(a) $\mathcal{G}_\xi$ 是 $1$ 维向量空间，其标量域为 $\kappa(\xi)$，
(b) $R^pg_*\mathcal{G} = 0$ 对 $p > 0$ 成立，并且
(c) $g_*\mathcal{G}$ 凝聚。于是可考虑
$\mathcal{F} = (Z \to X)_*\mathcal{G}$。因为 $Z \to X$ 是闭浸入，
可知 $(Z \to X)_*\mathcal{G}$ 在 $X$ 上凝聚，并且
$R^p(Z \to X)_*\mathcal{G} = 0$ 对 $p > 0$ 成立
（引理 \ref{coherent-lemma-finite-pushforward-coherent}）。
因此，由相对 Leray 谱序列
（《上同调》引理 \ref{cohomology-lemma-relative-Leray}），
将有 $R^pf_*\mathcal{F} = R^pg_*\mathcal{G} = 0$ 对 $p > 0$ 成立，
且 $f_*\mathcal{F} = g_*\mathcal{G}$ 凝聚。
最后，$\mathcal{F}_\xi = ((Z \to X)_*\mathcal{G})_\xi = \mathcal{G}_\xi$，
这验证了关于 $\xi$ 处之茎的条件。
因此，一切归结为寻找具有 (a)、(b)、(c) 的凝聚层
$\mathcal{G}$（在 $Z$ 上）。

\medskip\noindent
可以把 Chow 引理 \ref{coherent-lemma-chow-Noetherian}
应用于态射 $Z \to S$。由此得到 Chow 引理陈述中的交换图
$$
\xymatrix{
Z \ar[rd]_g & Z' \ar[d]^-{g'} \ar[l]^\pi \ar[r]_i & \mathbf{P}^m_S \ar[dl] \\
& S &
}
$$
此外，令 $U \subset Z$ 为稠密开子概形，使得
$\pi^{-1}(U) \to U$ 为同构。由评注
\ref{coherent-remark-chow-Noetherian} 中的讨论，可知
$i' = (i, \pi) : Z' \to \mathbf{P}^m_Z$ 是闭浸入。因此
$$
\mathcal{L} = i^*\mathcal{O}_{\mathbf{P}^m_S}(1) \cong
(i')^*\mathcal{O}_{\mathbf{P}^m_Z}(1)
$$
既 $g'$-相对丰沛，也 $\pi$-相对丰沛（例如由《态射》引理
\ref{morphisms-lemma-characterize-ample-on-finite-type}）。
因此，由引理 \ref{coherent-lemma-kill-by-twisting}，存在 $n \geq 0$，
使得 $R^p\pi_*\mathcal{L}^{\otimes n} = 0$ 对所有 $p > 0$ 成立，
且 $R^p(g')_*\mathcal{L}^{\otimes n} = 0$ 对所有 $p > 0$ 成立。
令 $\mathcal{G} = \pi_*\mathcal{L}^{\otimes n}$。
性质 (a) 成立，因为 $\pi_*\mathcal{L}^{\otimes n}|_U$
是可逆层（因为 $\pi^{-1}(U) \to U$ 是同构）。
性质 (b) 与 (c) 成立，因为由相对 Leray 谱序列
（《上同调》引理 \ref{cohomology-lemma-relative-Leray}），有
$$
E_2^{p, q} = R^pg_* R^q\pi_*\mathcal{L}^{\otimes n}
\Rightarrow
R^{p + q}(g')_*\mathcal{L}^{\otimes n}
$$
而由 $n$ 的选取，$E_2^{p, q}$ 中仅 $q = 0$ 的项可能非零，
且 $R^{p + q}(g')_*\mathcal{L}^{\otimes n}$ 中仅
$p = q = 0$ 的项可能非零。这蕴含
$R^pg_*\mathcal{G} = 0$ 对 $p > 0$ 成立，并且
$g_*\mathcal{G} = (g')_*\mathcal{L}^{\otimes n}$。
最后，应用前面的引理 \ref{coherent-lemma-locally-projective-pushforward}，
可知 $g_*\mathcal{G} = (g')_*\mathcal{L}^{\otimes n}$
如所需地凝聚。
\end{proof}

\begin{lemma}
\label{coherent-lemma-proper-over-affine-cohomology-finite}
设 $S = \Spec(A)$，其中 $A$ 为 Noether 环。
设 $f : X \to S$ 为固有态射。
设 $\mathcal{F}$ 为凝聚 $\mathcal{O}_X$-模。
则 $H^i(X, \mathcal{F})$ 是有限 $A$-模，
这对所有 $i \geq 0$ 成立。
\end{lemma}

\begin{proof}
这正是命题 \ref{coherent-proposition-proper-pushforward-coherent} 的仿射情形。
具体而言，由引理 \ref{coherent-lemma-quasi-coherence-higher-direct-images} 与
\ref{coherent-lemma-quasi-coherence-higher-direct-images-application}，已知
$R^if_*\mathcal{F}$ 是与 $A$-模 $H^i(X, \mathcal{F})$
相伴的准凝聚层；而由引理 \ref{coherent-lemma-coherent-Noetherian}，
该层凝聚，当且仅当 $H^i(X, \mathcal{F})$
是有限型 $A$-模。
\end{proof}

\begin{lemma}
\label{coherent-lemma-graded-finiteness}
设 $A$ 为 Noether 环。
设 $B$ 为有限生成分次 $A$-代数。
设 $f : X \to \Spec(A)$ 为固有态射。
令 $\mathcal{B} = f^*\widetilde B$。
设 $\mathcal{F}$ 为有限型准凝聚分次 $\mathcal{B}$-模。
\begin{enumerate}
\item 对每个 $p \geq 0$，分次 $B$-模 $H^p(X, \mathcal{F})$
是有限 $B$-模。
\item 若 $\mathcal{L}$ 为丰沛可逆 $\mathcal{O}_X$-模，
则存在整数 $d_0$，使得
$H^p(X, \mathcal{F} \otimes \mathcal{L}^{\otimes d}) = 0$
对所有 $p > 0$ 及 $d \geq d_0$ 成立。
\end{enumerate}
\end{lemma}

\begin{proof}
为证明该引理，考虑纤维积图
$$
\xymatrix{
X' = \Spec(B) \times_{\Spec(A)} X
\ar[r]_-\pi \ar[d]_{f'} &
X \ar[d]^f \\
\Spec(B) \ar[r] &
\Spec(A)
}
$$
注意，$f'$ 是固有态射；参见《态射》引理
\ref{morphisms-lemma-base-change-proper}。
此外，$B$ 是有限生成 $A$-代数，因而是 Noether 的
（《代数》引理 \ref{algebra-lemma-Noetherian-permanence}）。
这蕴含 $X'$ 是 Noether 概形
（《态射》引理 \ref{morphisms-lemma-finite-type-noetherian}）。
由《构造》引理 \ref{constructions-lemma-spec-properties}，
注意 $X'$ 是准凝聚 $\mathcal{O}_X$-代数 $\mathcal{B}$ 的相对谱。
因为 $\mathcal{F}$ 是准凝聚 $\mathcal{B}$-模，
可知存在唯一的准凝聚 $\mathcal{O}_{X'}$-模 $\mathcal{F}'$，
使得 $\pi_*\mathcal{F}' = \mathcal{F}$；参见《态射》引理
\ref{morphisms-lemma-affine-equivalence-modules}
由于 $\mathcal{F}$ 作为 $\mathcal{B}$-模是有限型的，
可知 $\mathcal{F}'$ 是有限型 $\mathcal{O}_{X'}$-模
（细节从略）。换言之，$\mathcal{F}'$ 是凝聚
$\mathcal{O}_{X'}$-模（引理 \ref{coherent-lemma-coherent-Noetherian}）。
由于态射 $\pi : X' \to X$ 仿射，由引理
\ref{coherent-lemma-relative-affine-cohomology} 有
$$
H^p(X, \mathcal{F}) = H^p(X', \mathcal{F}')
$$
因此 (1) 由引理 \ref{coherent-lemma-proper-over-affine-cohomology-finite} 推出。
给定 (2) 中的 $\mathcal{L}$，令
$\mathcal{L}' = \pi^*\mathcal{L}$。注意，由《态射》引理
\ref{morphisms-lemma-pullback-ample-tensor-relatively-ample}，
$\mathcal{L}'$ 在 $X'$ 上丰沛。由投影公式
（《上同调》引理 \ref{cohomology-lemma-projection-formula}），有
$\pi_*(\mathcal{F}' \otimes \mathcal{L}') = \mathcal{F} \otimes \mathcal{L}$。
因此，按上述相同论证，(2) 由引理
\ref{coherent-lemma-kill-by-twisting} 推出。
\end{proof}

\section{形式函数定理}
\label{coherent-section-theorem-formal-functions}

\noindent
本节研究形如 $\{I^n\mathcal{F}\}_{n \geq 0}$ 或
$\{\mathcal{F}/I^n\mathcal{F}\}_{n \geq 0}$ 的层序列在
$n$ 变化时其上同调的行为。

\medskip\noindent
这里及下文采用如下记号。给定概形态射 $f : X \to Y$、
准凝聚层 $\mathcal{F}$（在 $X$ 上），以及准凝聚理想层
$\mathcal{I} \subset \mathcal{O}_Y$，以
$\mathcal{I}^n\mathcal{F}$ 表示由 $f^{-1}(\mathcal{I}^n)$
的局部截面与 $\mathcal{F}$ 的局部截面之乘积生成的准凝聚子层。用公式表示为
$$
\mathcal{I}^n\mathcal{F}
=
\Im\left(
f^*(\mathcal{I}^n) \otimes_{\mathcal{O}_X} \mathcal{F}
\longrightarrow
\mathcal{F}
\right).
$$
注意有自然映射
$$
f^{-1}(\mathcal{I}^n) \otimes_{f^{-1}\mathcal{O}_Y} \mathcal{I}^m\mathcal{F}
\longrightarrow
f^*(\mathcal{I}^n) \otimes_{\mathcal{O}_X} \mathcal{I}^m\mathcal{F}
\longrightarrow
\mathcal{I}^{n + m}\mathcal{F}
$$
因此，$\mathcal{I}^n$ 的一个截面借助高阶直像的函子性给出映射
{\small $R^pf_*(\mathcal{I}^m\mathcal{F}) \to
R^pf_*(\mathcal{I}^{n + m}\mathcal{F})$}。
先局部化再层化，可见存在 $\mathcal{O}_Y$-模映射
$$
\mathcal{I}^n \otimes_{\mathcal{O}_Y} R^pf_*(\mathcal{I}^m\mathcal{F})
\longrightarrow
R^pf_*(\mathcal{I}^{n + m}\mathcal{F}).
$$
换言之，$\bigoplus_{n \geq 0} R^pf_*(\mathcal{I}^n\mathcal{F})$
是分次 $\bigoplus_{n \geq 0} \mathcal{I}^n$-模。

\medskip\noindent
若 $Y = \Spec(A)$ 且 $\mathcal{I} = \widetilde{I}$，则将
$\mathcal{I}^n\mathcal{F}$ 简记为 $I^n\mathcal{F}$。上面引入的映射赋予
$M = \bigoplus H^p(X, I^n\mathcal{F})$ 分次
$S = \bigoplus I^n$-模的结构。若 $f$ 固有，$A$ 是 Noether 环且
$\mathcal{F}$ 凝聚，则该模事实上是有限型的。

\begin{lemma}
\label{coherent-lemma-cohomology-powers-ideal-times-F}
\begin{reference}
\cite[III Cor 3.3.2]{EGA}
\end{reference}
设 $A$ 为 Noether 环，$I \subset A$ 为理想，并令
$B = \bigoplus_{n \geq 0} I^n$。设
$f : X \to \Spec(A)$ 为固有态射，并设凝聚层
$\mathcal{F}$（在 $X$ 上）。
则对每个 $p \geq 0$，分次 $B$-模
$\bigoplus_{n \geq 0} H^p(X, I^n\mathcal{F})$ 是有限 $B$-模。
\end{lemma}

\begin{proof}
令 $\mathcal{B} = f^*\widetilde{B}$。由于 $\mathcal{B}$ 满射到
$\bigoplus I^n\mathcal{O}_X$，显然
$\bigoplus I^n\mathcal{F}$ 是有限型分次 $\mathcal{B}$-模。
故结论由引理 \ref{coherent-lemma-graded-finiteness} 的 (1) 推出。
\end{proof}

\begin{lemma}
\label{coherent-lemma-cohomology-powers-ideal-times-sheaf}
给定概形态射 $f : X \to Y$、准凝聚层
$\mathcal{F}$（在 $X$ 上），以及准凝聚理想层
$\mathcal{I} \subset \mathcal{O}_Y$。假设 $Y$ 局部 Noether，
$f$ 固有且 $\mathcal{F}$ 凝聚。则
$$
\mathcal{M} =
\bigoplus\nolimits_{n \geq 0} R^pf_*(\mathcal{I}^n\mathcal{F})
$$
是准凝聚有限型分次
$\mathcal{A} = \bigoplus_{n \geq 0} \mathcal{I}^n$-模。
\end{lemma}

\begin{proof}
该陈述在 $Y$ 上是局部的，故可归约到 $Y$ 仿射的情形。
在仿射情形下，结论由引理
\ref{coherent-lemma-cohomology-powers-ideal-times-F} 推出。细节从略。
\end{proof}

\begin{lemma}
\label{coherent-lemma-cohomology-powers-ideal-application}
设 $A$ 为 Noether 环，$I \subset A$ 为理想，
$f : X \to \Spec(A)$ 为固有态射，并设
凝聚层 $\mathcal{F}$（在 $X$ 上）。则对每个 $p \geq 0$，
存在整数 $c \geq 0$，使得：
\begin{enumerate}
\item 乘法映射
$I^{n - c} \otimes H^p(X, I^c\mathcal{F}) \to H^p(X, I^n\mathcal{F})$
对所有 $n \geq c$ 都是满射；
\item 映射
$H^p(X, I^{n + m}\mathcal{F}) \to H^p(X, I^n\mathcal{F})$
之像包含于子模 $I^{m - e} H^p(X, I^n\mathcal{F})$，其中
$e = \max(0, c - n)$；这里 $n + m \geq c$ 且 $n, m \geq 0$；
\item 对 $n \geq c$，有
$$
\Ker(H^p(X, I^n\mathcal{F}) \to H^p(X, \mathcal{F})) =
\Ker(H^p(X, I^n\mathcal{F}) \to H^p(X, I^{n - c}\mathcal{F}))
$$
\item 存在映射
$I^nH^p(X, \mathcal{F}) \to H^p(X, I^{n - c}\mathcal{F})$（对 $n \geq c$），
使得当 $n \geq 2c$ 时，复合映射
$$
H^p(X, I^n\mathcal{F}) \to 
I^{n - c}H^p(X, \mathcal{F}) \to
H^p(X, I^{n - 2c}\mathcal{F})
$$
与
$$
I^nH^p(X, \mathcal{F}) \to
H^p(X, I^{n - c}\mathcal{F}) \to
I^{n - 2c}H^p(X, \mathcal{F})
$$
均为典范映射；并且
\item 逆向系统 $(H^p(X, I^n\mathcal{F}))$ 与
$(I^nH^p(X, \mathcal{F}))$ 是 pro-同构的。
\end{enumerate}
\end{lemma}

\begin{proof}
记 $M_n = H^p(X, I^n\mathcal{F})$（对 $n \geq 1$），并令
$M_0 = H^p(X, \mathcal{F})$，从而有映射
$\ldots \to M_3 \to M_2 \to M_1 \to M_0$。令
$B = \bigoplus_{n \geq 0} I^n$，则
$M = \bigoplus_{n \geq 0} M_n$ 是有限分次 $B$-模；参见引理
\ref{coherent-lemma-cohomology-powers-ideal-times-F}。
注意，乘积
$B_n \otimes M_m \to M_{m + n}$，$a \otimes m \mapsto a \cdot m$
与该逆向系统中的映射相容；更确切地说，对
$n, m' \geq 0$ 且 $m \geq m'$，交换图
$$
\xymatrix{
B_n \otimes_A M_m \ar[r] \ar[d] & M_{n + m} \ar[d] \\
B_n \otimes_A M_{m'} \ar[r] & M_{n + m'}
}
$$
交换。

\medskip\noindent
证明 (1)。选取 $d_1, \ldots, d_t \geq 0$ 及 $x_i \in M_{d_i}$，
使得 $M$ 由 $x_1, \ldots, x_t$ 在 $B$ 上生成。
对任意 $c \geq \max\{d_i\}$，可得
$B_{n - c} \cdot M_c = M_n$ 对 $n \geq c$ 成立，因而 (1) 成立。

\medskip\noindent
证明 (2)。令 $c$ 如 (1) 的证明中所取，并设 $n + m \geq c$。
有 $M_{n + m} = B_{n + m - c} \cdot M_c$。
若 $c > n$，则利用 $M_c \to M_n$ 以及上面指出的乘积与转移映射的相容性，
可知 $M_{n + m} \to M_n$ 之像包含于 $I^{n + m - c}M_n$。
若 $c \leq n$，则写成
$M_{n + m} = B_m \cdot B_{n - c} \cdot M_c =
B_m \cdot M_n$，从而看出其像包含于 $I^m M_n$。这证明了 (2)。

\medskip\noindent
令 $K_n \subset M_n$ 为映射 $M_n \to M_0$ 的核。
上面指出的乘积与转移映射的相容性表明
$K = \bigoplus K_n \subset M$ 是分次 $B$-子模。
由于 $B$ 是 Noether 的，且 $M$ 是有限生成分次 $B$-模，
$K$ 是有限生成分次 $B$-模。选取
$d'_1, \ldots, d'_{t'} \geq 0$ 及 $y_i \in K_{d'_i}$，
使得 $K$ 由 $y_1, \ldots, y_{t'}$ 在 $B$ 上生成。
令 $c = \max(d'_i, d'_j)$。由于
$y_i \in \Ker(M_{d'_i} \to M_0)$，可知
$B_n \cdot y_i \subset \Ker(M_{n + d'_i} \to M_n)$。
由此可见，$K_n = \Ker(M_n \to M_{n - c})$
对 $n \geq c$ 成立。这证明了 (3)。

\medskip\noindent
考虑下列实线部分交换的图
$$
\xymatrix{
I^n \otimes_A M_0 \ar[r] \ar[d] &
I^nM_0 \ar[r] \ar@{..>}[d] &
M_0 \ar[d] \\
M_n \ar[r] &
M_{n - c} \ar[r] &
M_0
}
$$
由上可知，满射 $I^n \otimes_A M_0 \to I^nM_0$ 的核映入
$K_n$，故得到虚线箭头，并且复合映射
$I^nM_0 \to M_{n - c} \to M_0$ 是典范映射。
于是，复合映射
$I^nM_0 \to M_{n - c} \to I^{n - 2c}M_0$
对 $n \geq 2c$ 显然是典范映射。
考虑复合映射 $M_n \to I^{n - c}M_0 \to M_{n - 2c}$。
第一映射把形如 $a \cdot m$ 的元素（其中
$a \in I^{n - c}$ 且 $m \in M_c$）送到 $a m'$，这里
$m'$ 是 $m$ 在 $M_0$ 中的像。第二映射再把它送到
$a \cdot m'$（位于 $M_{n - 2c}$ 中），故 (4) 成立。

\medskip\noindent
(5) 是 (4) 与 pro-对象态射之定义的直接推论。
\end{proof}

\noindent
在引理 \ref{coherent-lemma-cohomology-powers-ideal-times-F} 与
\ref{coherent-lemma-cohomology-powers-ideal-application} 的情形下，考虑逆向系统
$$
\mathcal{F}/I\mathcal{F} \leftarrow
\mathcal{F}/I^2\mathcal{F} \leftarrow
\mathcal{F}/I^3\mathcal{F} \leftarrow \ldots
$$
我们希望知道上同调群会发生什么。下面是第一个结论。

\begin{lemma}
\label{coherent-lemma-ML-cohomology-powers-ideal}
设 $A$ 为 Noether 环，$I \subset A$ 为理想，
$f : X \to \Spec(A)$ 为固有态射，并设凝聚层
$\mathcal{F}$（在 $X$ 上）。固定 $p \geq 0$。
存在 $c \geq 0$，使得：
\begin{enumerate}
\item 对所有 $n \geq c$，有
$$
\Ker(H^p(X, \mathcal{F}) \to H^p(X, \mathcal{F}/I^n\mathcal{F})) \subset
I^{n - c}H^p(X, \mathcal{F}).
$$
\item 逆向系统
$$
\left(H^p(X, \mathcal{F}/I^n\mathcal{F})\right)_{n \in \mathbf{N}}
$$
满足 Mittag--Leffler 条件（参见《同调》定义
\ref{homology-definition-Mittag-Leffler}）；并且
\item 有
$$
\Im(H^p(X, \mathcal{F}/I^k\mathcal{F})
\to H^p(X, \mathcal{F}/I^n\mathcal{F}))
=
\Im(H^p(X, \mathcal{F})
\to H^p(X, \mathcal{F}/I^n\mathcal{F}))
$$
这对所有 $k \geq n + c$ 成立。
\end{enumerate}
\end{lemma}

\begin{proof}
令 $c = \max\{c_p, c_{p + 1}\}$，其中 $c_p, c_{p + 1}$ 是引理
\ref{coherent-lemma-cohomology-powers-ideal-application} 分别对
$H^p$ 与 $H^{p + 1}$ 给出的整数。

\medskip\noindent
先证明 (1)。考虑短正合列
$$
0 \to I^n\mathcal{F} \to \mathcal{F} \to \mathcal{F}/I^n\mathcal{F} \to 0
$$
由上同调长正合列可知
$$
\Ker(
H^p(X, \mathcal{F}) \to H^p(X, \mathcal{F}/I^n\mathcal{F})
)
=
\Im(
H^p(X, I^n\mathcal{F}) \to H^p(X, \mathcal{F})
)
$$
因此，由引理 \ref{coherent-lemma-cohomology-powers-ideal-application} 的 (2)，
它包含于 $I^{n - c}H^p(X, \mathcal{F})$，这对 $n \geq c$ 成立。

\medskip\noindent
注意，(3) 按 Mittag--Leffler 系统的定义蕴含 (2)。

\medskip\noindent
下面证明 (3)。固定 $n$。考虑交换图
$$
\xymatrix{
0 \ar[r] &
I^n\mathcal{F} \ar[r] &
\mathcal{F} \ar[r] &
\mathcal{F}/I^n\mathcal{F} \ar[r] &
0 \\
0 \ar[r] &
I^{n + m}\mathcal{F} \ar[r] \ar[u] &
\mathcal{F} \ar[r] \ar[u] &
\mathcal{F}/I^{n + m}\mathcal{F} \ar[r] \ar[u] &
0
}
$$
它诱导出如下交换图
$$
\xymatrix{
H^p(X, \mathcal{F}) \ar[r] &
H^p(X, \mathcal{F}/I^n\mathcal{F}) \ar[r]_\delta &
H^{p + 1}(X, I^n\mathcal{F}) \ar[r] &
H^{p + 1}(X, \mathcal{F}) \\
H^p(X, \mathcal{F}) \ar[r] \ar[u]^1 &
H^p(X, \mathcal{F}/I^{n + m}\mathcal{F}) \ar[r] \ar[u]^\gamma &
H^{p + 1}(X, I^{n + m}\mathcal{F}) \ar[u]^\alpha \ar[r]^-\beta &
H^{p + 1}(X, \mathcal{F}) \ar[u]_1
}
$$
其两行正合。由引理
\ref{coherent-lemma-cohomology-powers-ideal-application} 的 (4)，$\beta$ 的核等于
$\alpha$ 的核，这对 $m \geq c$ 成立。
追图可知 $\gamma$ 之像包含于 $\delta$ 之核，故 (3) 成立
（令 $k = n + m$ 即得）。
\end{proof}

\begin{theorem}[形式函数定理]
\label{coherent-theorem-formal-functions}
设 $A$ 为 Noether 环，$I \subset A$ 为理想，
$f : X \to \Spec(A)$ 为固有态射，并设凝聚层
$\mathcal{F}$（在 $X$ 上）。固定 $p \geq 0$。
映射系
$$
H^p(X, \mathcal{F})/I^nH^p(X, \mathcal{F})
\longrightarrow
H^p(X, \mathcal{F}/I^n\mathcal{F})
$$
定义极限之间的同构
$$
H^p(X, \mathcal{F})^\wedge
\longrightarrow
\lim_n H^p(X, \mathcal{F}/I^n\mathcal{F})
$$
其中左端是 $A$-模 $H^p(X, \mathcal{F})$ 关于理想 $I$ 的完备化；
参见《代数》第 \ref{algebra-section-completion} 节。
而且，对极限拓扑而言，这事实上是同胚。
\end{theorem}

\begin{proof}
这可由引理 \ref{coherent-lemma-ML-cohomology-powers-ideal} 推出如下。
令 $M = H^p(X, \mathcal{F})$、
$M_n = H^p(X, \mathcal{F}/I^n\mathcal{F})$，并记
$N_n = \Im(M \to M_n)$。由引理
\ref{coherent-lemma-ML-cohomology-powers-ideal} 的 (2) 与 (3)，可知
$(M_n)$ 是 Mittag--Leffler 系统，并且 $N_n \subset M_n$
等于 $M_k$ 的像，这对所有 $k \gg n$ 成立。因此，作为具有极限拓扑的拓扑模，
$\lim M_n = \lim N_n$。另一方面，诸 $N_n$ 构成模 $M$
之商模的逆向系统，故 $\lim N_n$ 是 $M$ 关于诸核
$K_n = \Ker(M \to N_n)$ 所给拓扑的完备化。
由引理 \ref{coherent-lemma-ML-cohomology-powers-ideal} 的 (1)，有
$K_n \subset I^{n - c}M$；又由于 $N_n \subset M_n$ 被 $I^n$ 零化，
故 $I^n M \subset K_n$。因此，以子模 $K_n$ 作为 $0$ 的开邻域基本系
所定义的拓扑与 $I$-进拓扑相同，并得到诱导映射
$M^\wedge = \lim M/I^nM \to \lim N_n = \lim M_n$
是拓扑模的同构\footnote{
为确保无误，$M^\wedge$ 上的极限拓扑与其 $I$-进拓扑相同；
这由《代数》引理 \ref{algebra-lemma-hathat-finitely-generated} 推出。
另见《代数进阶》第 \ref{more-algebra-section-topological-ring} 节。}。
\end{proof}

\begin{lemma}
\label{coherent-lemma-spell-out-theorem-formal-functions}
设 $A$ 为环，$I \subset A$ 为理想。假设 $A$ 是 Noether 的且关于
$I$ 完备。设 $f : X \to \Spec(A)$ 为固有态射，
并设凝聚层 $\mathcal{F}$（在 $X$ 上）。则
$$
H^p(X, \mathcal{F}) = \lim_n H^p(X, \mathcal{F}/I^n\mathcal{F})
$$
对所有 $p \geq 0$ 成立。
\end{lemma}

\begin{proof}
这是形式函数定理（定理 \ref{coherent-theorem-formal-functions}）在完备 Noether
基环情形下的重新表述。事实上，在此情形下，$A$-模
$H^p(X, \mathcal{F})$ 是有限的
（引理 \ref{coherent-lemma-proper-over-affine-cohomology-finite}），因而是
$I$-进完备的（《代数》引理 \ref{algebra-lemma-completion-tensor}）；
由此可见左端不必再取完备化。
\end{proof}

\begin{lemma}
\label{coherent-lemma-formal-functions-stalk}
给定概形态射 $f : X \to Y$ 以及准凝聚层
$\mathcal{F}$（在 $X$ 上）。假设
\begin{enumerate}
\item $Y$ 局部 Noether；
\item $f$ 固有；并且
\item $\mathcal{F}$ 凝聚。
\end{enumerate}
设 $y \in Y$ 为一点。考虑纤维 $X_1 = X_y$ 的无穷小邻域
$$
\xymatrix{
X_n =
\Spec(\mathcal{O}_{Y, y}/\mathfrak m_y^n) \times_Y X
\ar[r]_-{i_n} \ar[d]_{f_n} &
X \ar[d]^f \\
\Spec(\mathcal{O}_{Y, y}/\mathfrak m_y^n) \ar[r]^-{c_n} & Y
}
$$
并令 $\mathcal{F}_n = i_n^*\mathcal{F}$。则有
$$
\left(R^pf_*\mathcal{F}\right)_y^\wedge
\cong
\lim_n H^p(X_n, \mathcal{F}_n)
$$
这是 $\mathcal{O}_{Y, y}^\wedge$-模的同构。
\end{lemma}

\begin{proof}
这只是形式函数定理，即定理 \ref{coherent-theorem-formal-functions}，
一个特殊情形的重新表述。下面详述之。注意
$\mathcal{O}_{Y, y}$ 是 Noether 局部环。考虑典范态射
$c : \Spec(\mathcal{O}_{Y, y}) \to Y$；参见《概形》公式
(\ref{schemes-equation-canonical-morphism})。该态射是平坦的，
因为它给出局部环的等同。暂记
$f' : X' \to \Spec(\mathcal{O}_{Y, y})$ 为 $f$ 到这个局部环的基变换。
由引理 \ref{coherent-lemma-flat-base-change-cohomology} 可知
$c^*R^pf_*\mathcal{F} = R^pf'_*\mathcal{F}'$。
此外，纤维 $X_y$ 与 $X'_y$ 的无穷小邻域彼此等同
（验证从略；提示：态射 $c_n$ 经由 $c$ 分解）。

\medskip\noindent
因此可以假设 $Y = \Spec(A)$ 是 Noether 局部环 $A$ 的谱，
其极大理想为 $\mathfrak m$，并且 $y \in Y$ 对应于闭点
（即对应于 $\mathfrak m$）。特别地，随即有
$$
\left(R^pf_*\mathcal{F}\right)_y =
\Gamma(Y, R^pf_*\mathcal{F}) =
H^p(X, \mathcal{F}).
$$

\medskip\noindent
在这种情形下，诸态射 $c_n$ 也全都是闭浸入，
故其基变换 $i_n$ 也都是闭浸入。注意
$i_{n, *}\mathcal{F}_n = i_{n, *}i_n^*\mathcal{F}
= \mathcal{F}/\mathfrak m^n\mathcal{F}$。由 $i_n$ 的 Leray 谱序列及引理
\ref{coherent-lemma-finite-pushforward-coherent}，可知
$$
H^p(X_n, \mathcal{F}_n) =
H^p(X, i_{n, *}\mathcal{F}_n) =
H^p(X, \mathcal{F}/\mathfrak m^n\mathcal{F})
$$
因此确可应用形式函数定理计算该引理陈述中的极限，结论得证。
\end{proof}

\noindent
下面这个引理稍后会推广到维数 $ > 0$ 的纤维，即再下一个引理。

\begin{lemma}
\label{coherent-lemma-higher-direct-images-zero-finite-fibre}
设 $f : X \to Y$ 为概形态射，并设 $y \in Y$。假设
\begin{enumerate}
\item $Y$ 局部 Noether；
\item $f$ 固有；并且
\item $f^{-1}(\{y\})$ 有限。
\end{enumerate}
则对任意凝聚层 $\mathcal{F}$（在 $X$ 上），都有
$(R^pf_*\mathcal{F})_y = 0$ 对所有 $p > 0$ 成立。
\end{lemma}

\begin{proof}
纤维 $X_y$ 是有限的，并且由《态射》引理
\ref{morphisms-lemma-finite-fibre}，它是有限离散空间。
而且，每个无穷小邻域 $X_n$ 的底拓扑空间都与之相同。
因此，由《概形》引理
\ref{schemes-lemma-scheme-finite-discrete-affine}，每个概形 $X_n$ 都仿射。
所以 $H^p(X_n, \mathcal{F}_n) = 0$ 对所有 $p > 0$ 成立。
于是由引理 \ref{coherent-lemma-formal-functions-stalk} 可知
$(R^pf_*\mathcal{F})_y^\wedge = 0$。注意，由命题
\ref{coherent-proposition-proper-pushforward-coherent}，
$R^pf_*\mathcal{F}$ 是凝聚的，因而 $R^pf_*\mathcal{F}_y$ 是有限
$\mathcal{O}_{Y, y}$-模。由 Nakayama 引理
（《代数》引理 \ref{algebra-lemma-NAK}），若局部环上有限模的完备化为零，
则该模为零。因此 $(R^pf_*\mathcal{F})_y = 0$。
\end{proof}


\begin{lemma}
\label{coherent-lemma-higher-direct-images-zero-above-dimension-fibre}
设 $f : X \to Y$ 为概形态射，并设 $y \in Y$。假设
\begin{enumerate}
\item $Y$ 局部 Noether；
\item $f$ 固有；并且
\item $\dim(X_y) = d$。
\end{enumerate}
则对任意凝聚层 $\mathcal{F}$（在 $X$ 上），都有
$(R^pf_*\mathcal{F})_y = 0$ 对所有 $p > d$ 成立。
\end{lemma}

\begin{proof}
纤维 $X_y$ 在 $\Spec(\kappa(y))$ 上是有限型的。
故由《态射》引理 \ref{morphisms-lemma-finite-type-noetherian}，
$X_y$ 是 Noether 概形。因此，由《性质》引理
\ref{properties-lemma-Noetherian-topology}，$X_y$ 的底拓扑空间是 Noether 的。
而且，每个无穷小邻域 $X_n$ 的底拓扑空间与 $X_y$ 的相同。
故由《上同调》命题 \ref{cohomology-proposition-vanishing-Noetherian}，
$H^p(X_n, \mathcal{F}_n) = 0$ 对所有 $p > d$ 成立。
于是可知 $(R^pf_*\mathcal{F})_y^\wedge = 0$；
由引理 \ref{coherent-lemma-formal-functions-stalk}，这对 $p > d$ 成立。注意，由命题
\ref{coherent-proposition-proper-pushforward-coherent}，
$R^pf_*\mathcal{F}$ 是凝聚的，因而 $R^pf_*\mathcal{F}_y$ 是有限
$\mathcal{O}_{Y, y}$-模。由 Nakayama 引理
（《代数》引理 \ref{algebra-lemma-NAK}），若局部环上有限模的完备化为零，
则该模为零。因此 $(R^pf_*\mathcal{F})_y = 0$。
\end{proof}

\section{形式函数定理的应用}
\label{coherent-section-applications-formal-functions}


\noindent
我们会按需在这里增补更多内容。目前需要 Noether 情形下有限态射的如下刻画。

\begin{lemma}
\label{coherent-lemma-characterize-finite}
（更一般的版本见《态射进阶》引理
\ref{more-morphisms-lemma-characterize-finite}。）
设 $f : X \to S$ 为概形态射，并假设 $S$ 局部 Noether。
则下列条件等价：
\begin{enumerate}
\item $f$ 有限；
\item $f$ 固有且纤维有限。
\end{enumerate}
\end{lemma}

\begin{proof}
由《态射》引理 \ref{morphisms-lemma-finite-proper}，有限态射固有。
由《态射》引理 \ref{morphisms-lemma-finite-quasi-finite}，有限态射拟有限。
拟有限态射有有限纤维；参见《态射》引理
\ref{morphisms-lemma-quasi-finite}。故有限态射固有且有有限纤维。

\medskip\noindent
假设 $f$ 固有且有有限纤维。要证明 $f$ 有限。
事实上，只须证明 $f$ 仿射。确实，若 $f$ 仿射，则由《态射》引理
\ref{morphisms-lemma-integral-universally-closed} 可知 $f$ 整；
继而由《态射》引理 \ref{morphisms-lemma-finite-integral} 可知 $f$ 有限。

\medskip\noindent
为证明 $f$ 仿射，可以假设 $S$ 仿射；目标是证明 $X$ 也仿射。
由于 $f$ 固有，可知 $X$ 分离且拟紧。因此可以用引理
\ref{coherent-lemma-quasi-separated-h1-zero-covering} 的判据证明 $X$ 仿射。
为此，设 $\mathcal{I} \subset \mathcal{O}_X$ 为有限型理想层。
特别地，$\mathcal{I}$ 是 $X$ 上的凝聚层。由引理
\ref{coherent-lemma-higher-direct-images-zero-finite-fibre}，可得
$R^1f_*\mathcal{I}_s = 0$ 对所有 $s \in S$ 成立。
换言之，$R^1f_*\mathcal{I} = 0$。于是由 $f$ 的 Leray 谱序列可见
$H^1(X , \mathcal{I}) = H^1(S, f_*\mathcal{I})$。
由于 $S$ 仿射，且 $f_*\mathcal{I}$ 准凝聚
（《概形》引理 \ref{schemes-lemma-push-forward-quasi-coherent}），
由引理 \ref{coherent-lemma-quasi-coherent-affine-cohomology-zero} 可得所需的
$H^1(S, f_*\mathcal{I}) = 0$。故如所需地，
$H^1(X, \mathcal{I}) = 0$。
\end{proof}

\noindent
作为推论，有如下有用结论。

\begin{lemma}
\label{coherent-lemma-proper-finite-fibre-finite-in-neighbourhood}
\begin{slogan}
固有态射在有限纤维的一个邻域上有限。
\end{slogan}
（更一般的版本见《态射进阶》引理
\ref{more-morphisms-lemma-proper-finite-fibre-finite-in-neighbourhood}。）
设 $f : X \to S$ 为概形态射，并设 $s \in S$。假设
\begin{enumerate}
\item $S$ 局部 Noether；
\item $f$ 固有；并且
\item $f^{-1}(\{s\})$ 是有限集。
\end{enumerate}
则存在开邻域 $V \subset S$（为 $s$ 的邻域），使得
$f|_{f^{-1}(V)} : f^{-1}(V) \to V$ 有限。
\end{lemma}

\begin{proof}
由《态射》引理 \ref{morphisms-lemma-finite-fibre}，态射 $f$
在 $f^{-1}(\{s\})$ 的所有点处拟有限。由《态射》引理
\ref{morphisms-lemma-quasi-finite-points-open}，$f$ 拟有限的点构成开集
$U \subset X$。令 $Z = X \setminus U$。则 $s \not \in f(Z)$。
由于 $f$ 固有，集合 $f(Z) \subset S$ 闭。选取任意开邻域
$V \subset S$（为 $s$ 的邻域），使得 $Z \cap V = \emptyset$。于是
$f^{-1}(V) \to V$ 局部拟有限且固有，故其拟有限
（《态射》引理 \ref{morphisms-lemma-quasi-finite-locally-quasi-compact}），
从而有有限纤维（《态射》引理 \ref{morphisms-lemma-quasi-finite}），
继而由引理 \ref{coherent-lemma-characterize-finite} 可知它有限。
\end{proof}

\begin{lemma}
\label{coherent-lemma-ample-on-fibre}
设 $f : X \to Y$ 为概形的固有态射，其中 $Y$ Noether。
设 $\mathcal{L}$ 为可逆 $\mathcal{O}_X$-模，
$\mathcal{F}$ 为凝聚 $\mathcal{O}_X$-模。
设 $y \in Y$ 为一点，使得 $\mathcal{L}_y$ 在 $X_y$ 上丰沛。
则存在 $d_0$，使得对所有 $d \geq d_0$，都有
$$
R^pf_*(\mathcal{F} \otimes_{\mathcal{O}_X} \mathcal{L}^{\otimes d})_y = 0
\text{ 对 }p > 0
$$
且映射
$$
f_*(\mathcal{F} \otimes_{\mathcal{O}_X} \mathcal{L}^{\otimes d})_y
\longrightarrow
H^0(X_y, \mathcal{F}_y \otimes_{\mathcal{O}_{X_y}} \mathcal{L}_y^{\otimes d})
$$
为满射。
\end{lemma}

\begin{proof}
注意，$\mathcal{O}_{Y, y}$ 是 Noether 局部环。考虑典范态射
$c : \Spec(\mathcal{O}_{Y, y}) \to Y$；参见《概形》公式
(\ref{schemes-equation-canonical-morphism})。该态射是平坦的，
因为它给出局部环的等同。暂记
$f' : X' \to \Spec(\mathcal{O}_{Y, y})$ 为 $f$ 到这个局部环的基变换。
由引理 \ref{coherent-lemma-flat-base-change-cohomology} 可知
$c^*R^pf_*\mathcal{F} = R^pf'_*\mathcal{F}'$。
此外，纤维 $X_y$ 与 $X'_y$ 彼此等同。
因此可以假设 $Y = \Spec(A)$ 是 Noether 局部环
$(A, \mathfrak m, \kappa)$ 的谱，且 $y \in Y$ 对应于
$\mathfrak m$。在这种情形下，
$R^pf_*(\mathcal{F} \otimes_{\mathcal{O}_X} \mathcal{L}^{\otimes d})_y =
H^p(X, \mathcal{F} \otimes_{\mathcal{O}_X} \mathcal{L}^{\otimes d})$
对所有 $p \geq 0$ 成立。记投影为
$f_y : X_y \to \Spec(\kappa)$。

\medskip\noindent
令 $B = \text{Gr}_\mathfrak m(A) =
\bigoplus_{n \geq 0} \mathfrak m^n/\mathfrak m^{n + 1}$。
考虑层 $\mathcal{B} = f_y^*\widetilde{B}$；它是准凝聚分次
$\mathcal{O}_{X_y}$-代数层。采用第
\ref{coherent-section-theorem-formal-functions} 节中的记号，并以
$I$ 由 $\mathfrak m$ 代替。由于 $X_y$ 是 $X$ 中由
$\mathfrak m\mathcal{O}_X$ 截出的闭子概形，可以把
$\mathfrak m^n\mathcal{F}/\mathfrak m^{n + 1}\mathcal{F}$
视为凝聚 $\mathcal{O}_{X_y}$-模；参见引理
\ref{coherent-lemma-i-star-equivalence}。于是
$\bigoplus_{n \geq 0} \mathfrak m^n\mathcal{F}/\mathfrak m^{n + 1}\mathcal{F}$
是有限型准凝聚分次 $\mathcal{B}$-模；这是因为它在
$\mathcal{B}$ 上由零次部分生成，并且其零次部分是
$\mathcal{F}_y = \mathcal{F}/\mathfrak m \mathcal{F}$，
这是凝聚 $\mathcal{O}_{X_y}$-模。故由引理
\ref{coherent-lemma-graded-finiteness} 的 (2)，可知
$$
H^p(X_y, \mathfrak m^n \mathcal{F}/ \mathfrak m^{n + 1}\mathcal{F}
\otimes_{\mathcal{O}_{X_y}} \mathcal{L}_y^{\otimes d}) = 0
$$
对所有 $p > 0$、$d \geq d_0$ 及 $n \geq 0$ 成立。
由引理 \ref{coherent-lemma-relative-affine-cohomology}，这等价于说
{\small
$
H^p(X, \mathfrak m^n \mathcal{F}/ \mathfrak m^{n + 1}\mathcal{F}
\otimes_{\mathcal{O}_X} \mathcal{L}^{\otimes d}) = 0
$}
对所有 $p > 0$、$d \geq d_0$ 及 $n \geq 0$ 成立。

\medskip\noindent
考虑凝聚 $\mathcal{O}_X$-模的短正合列
$$
0 \to \mathfrak m^n\mathcal{F}/\mathfrak m^{n + 1} \mathcal{F}
\to \mathcal{F}/\mathfrak m^{n + 1} \mathcal{F}
\to \mathcal{F}/\mathfrak m^n \mathcal{F} \to 0
$$
与 $\mathcal{L}^{\otimes d}$ 张量是正合函子，因而得到短正合列
$$
0 \to
\mathfrak m^n\mathcal{F}/\mathfrak m^{n + 1} \mathcal{F}
\otimes_{\mathcal{O}_X} \mathcal{L}^{\otimes d}
\to \mathcal{F}/\mathfrak m^{n + 1} \mathcal{F}
\otimes_{\mathcal{O}_X} \mathcal{L}^{\otimes d}
\to \mathcal{F}/\mathfrak m^n \mathcal{F}
\otimes_{\mathcal{O}_X} \mathcal{L}^{\otimes d} \to 0
$$
利用上同调长正合列和上面的消失，通过归纳可得：
\begin{enumerate}
\item $H^p(X, \mathcal{F}/\mathfrak m^n \mathcal{F}
\otimes_{\mathcal{O}_X} \mathcal{L}^{\otimes d}) = 0$
对所有 $p > 0$、$d \geq d_0$ 及 $n \geq 0$ 成立；并且
\item $H^0(X, \mathcal{F}/\mathfrak m^n \mathcal{F}
\otimes_{\mathcal{O}_X} \mathcal{L}^{\otimes d}) \to
H^0(X_y, \mathcal{F}_y \otimes_{\mathcal{O}_{X_y}} \mathcal{L}_y^{\otimes d})$
对所有 $d \geq d_0$ 及 $n \geq 1$ 都是满射。
\end{enumerate}
由形式函数定理（定理 \ref{coherent-theorem-formal-functions}），可知
$\mathfrak m$-进完备化（对
$H^p(X, \mathcal{F} \otimes_{\mathcal{O}_X} \mathcal{L}^{\otimes d})$）
在所有 $d \geq d_0$ 与 $p > 0$ 时都为零。
由于 $H^p(X, \mathcal{F} \otimes_{\mathcal{O}_X} \mathcal{L}^{\otimes d})$
是有限 $A$-模（引理 \ref{coherent-lemma-proper-over-affine-cohomology-finite}），
由 Nakayama 引理（《代数》引理 \ref{algebra-lemma-NAK}）可知
$H^p(X, \mathcal{F} \otimes_{\mathcal{O}_X} \mathcal{L}^{\otimes d})$
对所有 $d \geq d_0$ 与 $p > 0$ 都为零。
对 $p = 0$，由引理 \ref{coherent-lemma-ML-cohomology-powers-ideal} 的 (3) 可得
$H^0(X, \mathcal{F} \otimes_{\mathcal{O}_X} \mathcal{L}^{\otimes d}) \to
H^0(X_y, \mathcal{F}_y \otimes_{\mathcal{O}_{X_y}} \mathcal{L}_y^{\otimes d})$
为满射，这就给出了引理的最后一项陈述。
\end{proof}

\begin{lemma}
\label{coherent-lemma-ample-in-neighbourhood}
（更一般的版本见《态射进阶》引理
\ref{more-morphisms-lemma-ample-in-neighbourhood}。）
设 $f : X \to Y$ 为概形的固有态射，其中 $Y$ Noether。
设 $\mathcal{L}$ 为可逆 $\mathcal{O}_X$-模。
设 $y \in Y$ 为一点，使得 $\mathcal{L}_y$ 在 $X_y$ 上丰沛。
则存在开邻域 $V \subset Y$（为 $y$ 的邻域），使得
$\mathcal{L}|_{f^{-1}(V)}$ 在 $f^{-1}(V)/V$ 上丰沛。
\end{lemma}

\begin{proof}
选取 $d_0$，如引理 \ref{coherent-lemma-ample-on-fibre} 在
$\mathcal{F} = \mathcal{O}_X$ 时所给。选取 $d \geq d_0$，
使得可以找到 $r \geq 0$ 及截面
$s_{y, 0}, \ldots, s_{y, r} \in H^0(X_y, \mathcal{L}_y^{\otimes d})$，
它们定义闭浸入
$$
\varphi_y =
\varphi_{\mathcal{L}_y^{\otimes d}, (s_{y, 0}, \ldots, s_{y, r})} :
X_y \to \mathbf{P}^r_{\kappa(y)}.
$$
这可由《态射》引理
\ref{morphisms-lemma-finite-type-over-affine-ample-very-ample} 做到；
还要用《态射》引理
\ref{morphisms-lemma-image-proper-scheme-closed} 看出 $\varphi_y$
是闭浸入，并用《构造》第
\ref{constructions-section-projective-space} 节关于以可逆层与截面描述
到射影空间之态射的结论。由 $d_0$ 的选取，以 $Y$ 的一个开邻域
（为 $y$ 的邻域）替换后，可以选取
$s_0, \ldots, s_r \in H^0(X, \mathcal{L}^{\otimes d})$，
它们分别映到 $s_{y, 0}, \ldots, s_{y, r}$。
令 $X_{s_i} \subset X$ 为 $s_i$ 生成
$\mathcal{L}^{\otimes d}$ 的开子集。由于诸
$s_{y, i}$ 生成 $\mathcal{L}_y^{\otimes d}$，可见
$X_y \subset U = \bigcup X_{s_i}$。由于 $X \to Y$ 是闭映射，
可知存在开邻域 $y \in V \subset Y$，使得
$f^{-1}(V) \subset U$。将 $Y$ 替换为 $V$ 后，可以假设诸
$s_i$ 生成 $\mathcal{L}^{\otimes d}$。于是得到态射
$$
\varphi = \varphi_{\mathcal{L}^{\otimes d}, (s_0, \ldots, s_r)} :
X \longrightarrow \mathbf{P}^r_Y
$$
它满足 $\mathcal{L}^{\otimes d} \cong \varphi^*\mathcal{O}_{\mathbf{P}^r_Y}(1)$，
且其到 $y$ 的基变换给出 $\varphi_y$。

\medskip\noindent
我们以一个小技巧完成证明；“正确的”证明是直接说明，$\varphi$
在基变换到 $y$ 的某个开邻域后是闭浸入。具体而言，由引理
\ref{coherent-lemma-proper-finite-fibre-finite-in-neighbourhood}，可知 $\varphi$
在纤维 $\mathbf{P}^r_{\kappa(y)}$ 的一个开邻域上有限；这里该纤维是
$\mathbf{P}^r_Y \to Y$ 位于 $y$ 上方的纤维。
利用 $\mathbf{P}^r_Y \to Y$ 是闭映射，缩小 $Y$ 后，
可以假设 $\varphi$ 有限。于是
$\mathcal{L}^{\otimes d} \cong \varphi^*\mathcal{O}_{\mathbf{P}^r_Y}(1)$
由非常一般的《态射》引理
\ref{morphisms-lemma-pullback-ample-tensor-relatively-ample} 可知是丰沛的。
\end{proof}

\section{上同调与基变换（三）}
\label{coherent-section-cohomology-and-base-change-perfect}

\noindent
本节证明一个十分一般的现象之最简单情形；该现象将在《概形的导出范畴》第
\ref{perfect-section-cohomology-and-base-change-perfect} 节讨论。
关于把下列引理翻译成代数语言的表述，见评注
\ref{coherent-remark-explain-perfect-direct-image}。

\begin{lemma}
\label{coherent-lemma-perfect-direct-image}
设 $A$ 为 Noether 环，并令 $S = \Spec(A)$。设
$f : X \to S$ 为概形的固有态射。设 $\mathcal{F}$ 为凝聚
$\mathcal{O}_X$-模，并且在 $S$ 上平坦。则
\begin{enumerate}
\item $R\Gamma(X, \mathcal{F})$ 是 $D(A)$ 的完美对象；并且
\item 对任意环同态 $A \to A'$，基变换映射
$$
R\Gamma(X, \mathcal{F}) \otimes_A^{\mathbf{L}} A'
\longrightarrow
R\Gamma(X_{A'}, \mathcal{F}_{A'})
$$
是同构。
\end{enumerate}
\end{lemma}

\begin{proof}
选取有限仿射开覆盖 $X = \bigcup_{i = 1, \ldots, n} U_i$。
由引理 \ref{coherent-lemma-separated-case-relative-cech} 与
\ref{coherent-lemma-base-change-complex}，{\v C}ech 复形
$K^\bullet = {\check C}^\bullet(\mathcal{U}, \mathcal{F})$ 满足
$$
K^\bullet \otimes_A A' = R\Gamma(X_{A'}, \mathcal{F}_{A'})
$$
这对所有环同态 $A \to A'$ 成立。令
$K_{alt}^\bullet = {\check C}_{alt}^\bullet(\mathcal{U}, \mathcal{F})$
为交错 {\v C}ech 复形。由《上同调》引理
\ref{cohomology-lemma-alternating-usual}，存在同伦等价
$K_{alt}^\bullet \to K^\bullet$（在 $A$-模范畴中）。特别地，同样有
$$
K_{alt}^\bullet \otimes_A A' = R\Gamma(X_{A'}, \mathcal{F}_{A'})
$$
由于 $\mathcal{F}$ 在 $A$ 上平坦，可见每个 $K_{alt}^n$
在 $A$ 上平坦（参见《态射》引理
\ref{morphisms-lemma-flat-module-characterize}）。而且，
$K_{alt}^\bullet$ 有上界（这正是改用交错 {\v C}ech 复形的原因），故
$K_{alt}^\bullet \otimes_A A' = K_{alt}^\bullet \otimes_A^{\mathbf{L}} A'$；
这是导出张量积定义的结果（参见《代数进阶》第
\ref{more-algebra-section-derived-tensor-product} 节）。由引理
\ref{coherent-lemma-proper-over-affine-cohomology-finite}，上同调群
$H^i(K_{alt}^\bullet)$ 是有限 $A$-模。由于
$K_{alt}^\bullet$ 有界，可得 $K_{alt}^\bullet$ 是伪凝聚的；
参见《代数进阶》引理
\ref{more-algebra-lemma-Noetherian-pseudo-coherent}。
给定任意 $A$-模 $M$，令 $A' = A \oplus M$，其中 $M$ 是平方零理想，
即 $(a, m) \cdot (a', m') = (aa', am' + a'm)$。
由上可见，$K_{alt}^\bullet \otimes_A^\mathbf{L} A'$
仅在次数 $0, \ldots, n$ 有上同调。因此，
$K_{alt}^\bullet \otimes_A^\mathbf{L} M$ 仅在次数
$0, \ldots, n$ 有上同调。故 $K_{alt}^\bullet$ 有有限 Tor 维数；
参见《代数进阶》定义 \ref{more-algebra-definition-tor-amplitude}。
最后应用《代数进阶》引理 \ref{more-algebra-lemma-perfect} 即得结论。
\end{proof}

\begin{remark}
\label{coherent-remark-explain-perfect-direct-image}
引理 \ref{coherent-lemma-perfect-direct-image} 的一个推论是：存在由有限射影
$A$-模组成的有限复形 $M^\bullet$，使得
$$
H^i(X_{A'}, \mathcal{F}_{A'}) = H^i(M^\bullet \otimes_A A')
$$
关于 $A'$ 函子性地成立。$\mathcal{F}$ 在 $A$ 上平坦这一条件是本质的；
参见 \cite{Hartshorne}。
\end{remark}

\section{凝聚形式模}
\label{coherent-section-coherent-formal}

\noindent
由于尚无形式概形理论可供使用，我们发展一些语言来代替
“Noether 进形式概形上的凝聚模”这一概念。

\medskip\noindent
设 $X$ 为 Noether 概形，并设 $\mathcal{I} \subset \mathcal{O}_X$
为准凝聚理想层。我们考虑逆向系统 $(\mathcal{F}_n)$，
其各项是凝聚 $\mathcal{O}_X$-模，并满足：
\begin{enumerate}
\item $\mathcal{F}_n$ 被 $\mathcal{I}^n$ 零化；并且
\item 转移映射诱导同构
$\mathcal{F}_{n + 1}/\mathcal{I}^n\mathcal{F}_{n + 1} \to \mathcal{F}_n$。
\end{enumerate}
这类逆向系统的态射按通常方式定义。以
$\textit{Coh}(X, \mathcal{I})$ 表示这些逆向系统所成的范畴。
下面将证明一系列关于该范畴中对象的引理。事实上，后续大多数引理
都是下述仿射情形范畴描述的直接推论。

\begin{lemma}
\label{coherent-lemma-inverse-systems-affine}
若 $X = \Spec(A)$ 是 Noether 环的谱，且 $\mathcal{I}$ 是与理想
$I \subset A$ 相伴的准凝聚理想层，则
$\textit{Coh}(X, \mathcal{I})$ 等价于有限 $A^\wedge$-模范畴，
其中 $A^\wedge$ 是 $A$ 关于 $I$ 的完备化。
\end{lemma}

\begin{proof}
以 $\text{Mod}^{fg}_{A, I}$ 表示逆向系统 $(M_n)$ 所成的范畴，
其中各项是有限 $A$-模；这些系统满足：(1) $M_n$ 被 $I^n$ 零化；
(2) $M_{n + 1}/I^nM_{n + 1} = M_n$。由 $X$ 上凝聚层与有限
$A$-模之间的对应（引理 \ref{coherent-lemma-coherent-Noetherian}），只须证明
$\text{Mod}^{fg}_{A, I}$ 等价于有限 $A^\wedge$-模范畴。
为此，只须证明：给定对象 $(M_n)$（属于
$\text{Mod}^{fg}_{A, I}$），模
$$
M = \lim M_n
$$
是有限 $A^\wedge$-模，且 $M/I^nM = M_n$。
由于转移映射都是满射，可见 $M \to M_1$ 是满射。
选取 $x_1, \ldots, x_t \in M$，使其像生成 $M_1$。
这诱导系统间的映射 $(A/I^n)^{\oplus t} \to M_n$。
由 Nakayama 引理（《代数》引理 \ref{algebra-lemma-NAK}），
这些映射均为满射。令 $K_n \subset (A/I^n)^{\oplus t}$ 为其核。
性质 (2) 蕴含 $K_{n + 1} \to K_n$ 是满射；特别地，系统
$(K_n)$ 满足 Mittag--Leffler 条件。由《同调》引理
\ref{homology-lemma-Mittag-Leffler}，得到正合列
$0 \to K \to (A^\wedge)^{\oplus t} \to M \to 0$，
其中 $K = \lim K_n$。故 $M$ 是有限 $A^\wedge$-模。
由于 $K \to K_n$ 为满射，得到
$$
M/I^nM = \Coker(K \to (A/I^n)^{\oplus t}) =
(A/I^n)^{\oplus t}/K_n = M_n
$$
如所需。
\end{proof}

\begin{lemma}
\label{coherent-lemma-inverse-systems-abelian}
设 $X$ 为 Noether 概形，并设 $\mathcal{I} \subset \mathcal{O}_X$
为准凝聚理想层。
\begin{enumerate}
\item 范畴 $\textit{Coh}(X, \mathcal{I})$ 是阿贝尔范畴。
\item 对开集 $U \subset X$，限制函子
$\textit{Coh}(X, \mathcal{I}) \to \textit{Coh}(U, \mathcal{I}|_U)$
是正合的。
\item $\textit{Coh}(X, \mathcal{I})$ 中的正合性可通过限制到
$X$ 的一个开覆盖之各成员上来检验。
\end{enumerate}
\end{lemma}

\begin{proof}
设 $\alpha =(\alpha_n) : (\mathcal{F}_n) \to (\mathcal{G}_n)$ 为
$\textit{Coh}(X, \mathcal{I})$ 的态射。$\alpha$ 的余核是逆向系统
$(\Coker(\alpha_n))$（细节从略）。为描述其核，对 $l \geq m$ 令
$$
\mathcal{K}'_{l, m} = \Im(\Ker(\alpha_l) \to \mathcal{F}_m)
$$
我们断言：
\begin{enumerate}
\item[(a)] 逆向系统 $(\mathcal{K}'_{l, m})_{l \geq m}$ 最终常值，
设其最终值为 $\mathcal{K}'_m$；
\item[(b)] 系统
$(\mathcal{K}'_m/\mathcal{I}^n\mathcal{K}'_m)_{m \geq n}$ 最终常值，
设其最终值为 $\mathcal{K}_n$；
\item[(c)] 系统 $(\mathcal{K}_n)$ 构成
$\textit{Coh}(X, \mathcal{I})$ 的一个对象；并且
\item[(d)] 该对象是 $\alpha$ 的核。
\end{enumerate}
为证明 (a)、(b)、(c)，可以在仿射局部工作；设
$X = \Spec(A)$，且 $\mathcal{I}$ 对应于理想 $I \subset A$。
由引理 \ref{coherent-lemma-inverse-systems-affine}，$\alpha$ 对应于映射
$f : M \to N$（在有限 $A^\wedge$-模之间）。记
$K = \Ker(f)$。注意 $A^\wedge$ 是 Noether 环
（《代数》引理 \ref{algebra-lemma-completion-Noetherian-Noetherian}）。
选取整数 $c \geq 0$，使得
$K \cap I^n M \subset I^{n - c}K$ 对 $n \geq c$ 成立
（《代数》引理 \ref{algebra-lemma-Artin-Rees}），并且使它对映射
$f$ 与理想 $I^\wedge = IA^\wedge$ 满足《代数》引理
\ref{algebra-lemma-map-AR}。于是，$\mathcal{K}'_{l, m}$ 对应于 $A$-模
$$
K'_{l, m} = \frac{a^{-1}(I^lN) + I^mM}{I^mM} =
\frac{K + I^{l - c}f^{-1}(I^cN) + I^mM}{I^mM} =
\frac{K + I^mM}{I^mM}
$$
其中最后一个等号在 $l \geq m + c$ 时成立。因此
$\mathcal{K}'_m$ 对应于 $A$-模 $K/K \cap I^mM$，且
$\mathcal{K}'_m/\mathcal{I}^n\mathcal{K}'_m$ 对应于
$$
\frac{K}{K \cap I^mM + I^nK} = \frac{K}{I^nK}
$$
当 $m \geq n + c$ 时，这由上面对 $c$ 的选取成立。
故 $\mathcal{K}_n$ 对应于 $K/I^nK$。

\medskip\noindent
下面证明 (d)。由上面的仿射描述，复合映射
$(\mathcal{K}_n) \to (\mathcal{F}_n) \to (\mathcal{G}_n)$
显然为零。设 $\beta : (\mathcal{H}_n) \to (\mathcal{F}_n)$
为满足 $\alpha \circ \beta = 0$ 的态射。则
$\mathcal{H}_l \to \mathcal{F}_l$ 的像包含于 $\Ker(\alpha_l)$。
由于 $\mathcal{H}_m = \mathcal{H}_l/\mathcal{I}^m\mathcal{H}_l$
对 $l \geq m$ 成立，
得到一族映射 $\mathcal{H}_m \to \mathcal{K}'_{l, m}$，
从而得到映射 $\mathcal{H}_m \to \mathcal{K}_m'$。
又由于
$\mathcal{H}_n = \mathcal{H}_m/\mathcal{I}^n\mathcal{H}_m$，
得到一族映射
$\mathcal{H}_n \to \mathcal{K}'_m/\mathcal{I}^n\mathcal{K}'_m$，
因而如所需地得到映射 $\mathcal{H}_n \to \mathcal{K}_n$。

\medskip\noindent
为完成 (1) 的证明，还须证明 $\Coim = \Im$
（在 $\textit{Coh}(X, \mathcal{I})$ 中）。上面已经看出，
在仿射情形下，取核与余核和把
$\textit{Coh}(X, \mathcal{I})$ 描述为模范畴相交换。
由于模范畴中 $\Im = \Coim$，这给出
$\Coim = \Im$（在 $\textit{Coh}(X, \mathcal{I})$ 中）。
引理的 (2)、(3) 由我们对核与余核的构造立即推出。
\end{proof}

\begin{lemma}
\label{coherent-lemma-inverse-systems-surjective}
设 $X$ 为 Noether 概形，并设 $\mathcal{I} \subset \mathcal{O}_X$
为准凝聚理想层。映射
$(\mathcal{F}_n) \to (\mathcal{G}_n)$ 在
$\textit{Coh}(X, \mathcal{I})$ 中为满射，当且仅当
$\mathcal{F}_1 \to \mathcal{G}_1$ 为满射。
\end{lemma}

\begin{proof}
从略。提示：在仿射开集上考察，使用引理
\ref{coherent-lemma-inverse-systems-affine} 及《代数》引理
\ref{algebra-lemma-NAK}。
\end{proof}

\noindent
设 $X$ 为 Noether 概形，并设 $\mathcal{I} \subset \mathcal{O}_X$
为准凝聚理想层。有函子
\begin{equation}
\label{coherent-equation-completion-functor}
\textit{Coh}(\mathcal{O}_X) \longrightarrow \textit{Coh}(X, \mathcal{I}), \quad
\mathcal{F} \longmapsto \mathcal{F}^\wedge
\end{equation}
它把凝聚 $\mathcal{O}_X$-模 $\mathcal{F}$ 送到对象
$\mathcal{F}^\wedge = (\mathcal{F}/\mathcal{I}^n\mathcal{F})$；
该对象属于 $\textit{Coh}(X, \mathcal{I})$。

\begin{lemma}
\label{coherent-lemma-exact}
函子 (\ref{coherent-equation-completion-functor}) 是正合的。
\end{lemma}

\begin{proof}
只须在 $X$ 上局部检验。因此可以假设 $X$ 仿射，
即处于引理 \ref{coherent-lemma-inverse-systems-affine} 的情形。
该函子就是函子
$\text{Mod}^{fg}_A \to \text{Mod}^{fg}_{A^\wedge}$，
它把有限 $A$-模 $M$ 送到完备化 $M^\wedge$。
故结论由《代数》引理 \ref{algebra-lemma-completion-flat} 推出。
\end{proof}

\begin{lemma}
\label{coherent-lemma-completion-internal-hom}
设 $X$ 为 Noether 概形，并设 $\mathcal{I} \subset \mathcal{O}_X$
为准凝聚理想层。设 $\mathcal{F}$、$\mathcal{G}$ 为凝聚
$\mathcal{O}_X$-模。令
$\mathcal{H} = \SheafHom_{\mathcal{O}_X}(\mathcal{G}, \mathcal{F})$。
则
$$
\lim H^0(X, \mathcal{H}/\mathcal{I}^n\mathcal{H}) =
\Mor_{\textit{Coh}(X, \mathcal{I})}
(\mathcal{G}^\wedge, \mathcal{F}^\wedge).
$$
\end{lemma}

\begin{proof}
为证明该结论，可以在 $X$ 上仿射局部地工作。
故可假设 $X = \Spec(A)$，并且 $\mathcal{F}$、$\mathcal{G}$
由有限 $A$-模 $M$、$N$ 给出。于是 $\mathcal{H}$ 对应于有限
$A$-模 $H = \Hom_A(M, N)$。经由引理
\ref{coherent-lemma-inverse-systems-affine} 的等价，引理的陈述化为
$$
H^\wedge = \Hom_{A^\wedge}(M^\wedge, N^\wedge)
$$
由《代数》引理 \ref{algebra-lemma-completion-flat}（使用三次），有
$$
H^\wedge = \Hom_A(M, N) \otimes_A A^\wedge =
\Hom_{A^\wedge}(M \otimes_A A^\wedge, N \otimes_A A^\wedge) =
\Hom_{A^\wedge}(M^\wedge, N^\wedge)
$$
其中第二个等号使用了 $A^\wedge$ 在 $A$ 上平坦这一事实
（参见《代数进阶》引理
\ref{more-algebra-lemma-pseudo-coherence-and-base-change-ext}）。
引理得证。
\end{proof}

\noindent
设 $X$ 为 Noether 概形。设 $\mathcal{I} \subset \mathcal{O}_X$
为准凝聚理想层。称对象 $(\mathcal{F}_n)$（属于
$\textit{Coh}(X, \mathcal{I})$）是 {\it $\mathcal{I}$-幂挠的}，或称它
{\it 被 $\mathcal{I}$ 的某个幂零化}，若存在 $c \geq 1$，使得
$\mathcal{F}_n = \mathcal{F}_c$ 对所有 $n \geq c$ 成立。
若是如此，就称 $(\mathcal{F}_n)$ {\it 被 $\mathcal{I}^c$ 零化}。
若 $X = \Spec(A)$ 仿射，则经由引理
\ref{coherent-lemma-inverse-systems-affine} 的等价，这些对象恰对应于有限
$A$-模；这些模被 $I$ 的某个幂或被 $I^c$ 零化。

\begin{lemma}
\label{coherent-lemma-existence-easy}
设 $X$ 为 Noether 概形，并设 $\mathcal{I} \subset \mathcal{O}_X$
为准凝聚理想层。设 $\mathcal{G}$ 为凝聚 $\mathcal{O}_X$-模。
设 $(\mathcal{F}_n)$ 为 $\textit{Coh}(X, \mathcal{I})$ 的对象。
\begin{enumerate}
\item 若 $\alpha : (\mathcal{F}_n) \to \mathcal{G}^\wedge$ 是一个映射，
其核与余核均被 $\mathcal{I}$ 的某个幂零化，则存在三元组
$(\mathcal{F}, a, \beta)$，且它在唯一同构意义下唯一；其中
\begin{enumerate}
\item $\mathcal{F}$ 是凝聚 $\mathcal{O}_X$-模；
\item $a : \mathcal{F} \to \mathcal{G}$ 是 $\mathcal{O}_X$-模映射，
其核与余核均被 $\mathcal{I}$ 的某个幂零化；
\item $\beta : (\mathcal{F}_n) \to \mathcal{F}^\wedge$ 是同构；并且
\item $\alpha = a^\wedge \circ \beta$。
\end{enumerate}
\item 若 $\alpha : \mathcal{G}^\wedge \to (\mathcal{F}_n)$ 是一个映射，
其核与余核均被 $\mathcal{I}$ 的某个幂零化，则存在三元组
$(\mathcal{F}, a, \beta)$，且它在唯一同构意义下唯一；其中
\begin{enumerate}
\item $\mathcal{F}$ 是凝聚 $\mathcal{O}_X$-模；
\item $a : \mathcal{G} \to \mathcal{F}$ 是 $\mathcal{O}_X$-模映射，
其核与余核均被 $\mathcal{I}$ 的某个幂零化；
\item $\beta : \mathcal{F}^\wedge \to (\mathcal{F}_n)$ 是同构；并且
\item $\alpha = \beta \circ a^\wedge$。
\end{enumerate}
\end{enumerate}
\end{lemma}

\begin{proof}
证明 (1)。唯一性表明，只须构造 $(\mathcal{F}, a, \beta)$，
并可在 $X$ 上 Zariski 局部地完成构造。因此可以假设 $X = \Spec(A)$，
且 $\mathcal{I}$ 对应于理想 $I \subset A$。在此情形下可应用引理
\ref{coherent-lemma-inverse-systems-affine}。设 $M'$ 为有限 $A^\wedge$-模，
它对应于 $(\mathcal{F}_n)$。设 $N$ 为有限 $A$-模，
它对应于 $\mathcal{G}$。于是 $\alpha$ 对应于映射
$$
\varphi : M' \longrightarrow N^\wedge
$$
其核与余核被某个 $I^t$ 零化，其中 $t$ 为某个整数。回忆
$N^\wedge = N \otimes_A A^\wedge$
（《代数》引理 \ref{algebra-lemma-completion-tensor}）。
由《代数进阶》引理
\ref{more-algebra-lemma-application-formal-glueing}，存在 $A$-模映射
$\psi : M \to N$，其核与余核是 $I$-幂挠的，并且存在同构
$M \otimes_A A^\wedge = M'$，它与 $\varphi$ 相容。
由于 $N$ 与 $M'$ 都是有限模，可知 $M$ 是有限 $A$-模；
参见《代数进阶》评注
\ref{more-algebra-remark-formal-glueing-algebras}。故
$M \otimes_A A^\wedge = M^\wedge$。我们略去验证由此得到的三元组
$(M, N \to M, M^\wedge \to M')$ 在唯一同构意义下唯一。

\medskip\noindent
(2) 的证明完全相同，从略。
\end{proof}

\begin{lemma}
\label{coherent-lemma-torsion-hom-ext}
设 $X$ 为 Noether 概形，并设 $\mathcal{I} \subset \mathcal{O}_X$
为准凝聚理想层。$\textit{Coh}(X, \mathcal{I})$ 中任何被
$\mathcal{I}$ 的某个幂零化的对象，都在
(\ref{coherent-equation-completion-functor}) 的本质像中。
此外，若 $\mathcal{F}$、$\mathcal{G}$ 属于
$\textit{Coh}(\mathcal{O}_X)$，且 $\mathcal{F}$ 或 $\mathcal{G}$
之一被 $\mathcal{I}$ 的某个幂零化，则映射
$$
\xymatrix{
\Hom_X(\mathcal{F}, \mathcal{G}) \ar[d] &
\Ext_X(\mathcal{F}, \mathcal{G}) \ar[d] \\
\Hom_{\textit{Coh}(X, \mathcal{I})}(\mathcal{F}^\wedge, \mathcal{G}^\wedge) &
\Ext_{\textit{Coh}(X, \mathcal{I})}(\mathcal{F}^\wedge, \mathcal{G}^\wedge)
}
$$
均为同构。
\end{lemma}

\begin{proof}
设 $(\mathcal{F}_n)$ 是 $\textit{Coh}(X, \mathcal{I})$ 的对象，
且被 $\mathcal{I}^c$ 零化，其中 $c \geq 1$。于是
$\mathcal{F}_n \to \mathcal{F}_c$ 对 $n \geq c$ 是同构。
故若令 $\mathcal{F} = \mathcal{F}_c$，则可见
$\mathcal{F}^\wedge \cong (\mathcal{F}_n)$。这证明了第一项断言。

\medskip\noindent
设 $\mathcal{F}$、$\mathcal{G}$ 为 $\textit{Coh}(\mathcal{O}_X)$
的对象，使得 $\mathcal{F}$ 或 $\mathcal{G}$ 之一被
$\mathcal{I}^c$ 零化，其中 $c \geq 1$。则
$\mathcal{H} = \SheafHom_{\mathcal{O}_X}(\mathcal{G}, \mathcal{F})$
是凝聚 $\mathcal{O}_X$-模，并被 $\mathcal{I}^c$ 零化。因此
$$
\Hom_X(\mathcal{G}, \mathcal{F}) =
H^0(X, \mathcal{H}) =
\lim H^0(X, \mathcal{H}/\mathcal{I}^n\mathcal{H}) =
\Mor_{\textit{Coh}(X, \mathcal{I})}
(\mathcal{G}^\wedge, \mathcal{F}^\wedge).
$$
参见引理 \ref{coherent-lemma-completion-internal-hom}。
这证明了关于同态的陈述。

\medskip\noindent
记号 $\Ext$ 指《同调》第 \ref{homology-section-extensions} 节定义的扩张。
$\Ext$ 上映射的单射性立即由 $\Hom$ 上映射的双射性推出。
为证明满射性，假设 $\mathcal{F}$ 被 $I$ 的某个幂零化。
引理 \ref{coherent-lemma-existence-easy} 的 (1) 表明：给定
$\textit{Coh}(U, I\mathcal{O}_U)$ 中的扩张
$$
0 \to \mathcal{G}^\wedge \to (\mathcal{E}_n) \to \mathcal{F}^\wedge \to 0
$$
态射 $\mathcal{G}^\wedge \to (\mathcal{E}_n)$ 同构于
$\mathcal{G} \to \mathcal{E}^\wedge$，其中
$\mathcal{G} \to \mathcal{E}$ 是
$\textit{Coh}(\mathcal{O}_U)$ 中的某个态射。另一种情形同理。
\end{proof}

\begin{lemma}
\label{coherent-lemma-finite-over-rees-algebra}
设 $X$ 为 Noether 概形，并设 $\mathcal{I} \subset \mathcal{O}_X$
为准凝聚理想层。若 $(\mathcal{F}_n)$ 是
$\textit{Coh}(X, \mathcal{I})$ 的对象，则
$\bigoplus \Ker(\mathcal{F}_{n + 1} \to \mathcal{F}_n)$ 是有限型分次
准凝聚 $\bigoplus \mathcal{I}^n/\mathcal{I}^{n + 1}$-模。
\end{lemma}

\begin{proof}
该问题在 $X$ 上是局部的，故可假设 $X$ 仿射，
即处于引理 \ref{coherent-lemma-inverse-systems-affine} 的情形。
在这种情形下，若 $(\mathcal{F}_n)$ 对应于有限 $A^\wedge$-模
$M$，则 $\bigoplus \Ker(\mathcal{F}_{n + 1} \to \mathcal{F}_n)$
对应于 $\bigoplus I^nM/I^{n + 1}M$；后者显然是
$\bigoplus I^n/I^{n + 1}$ 上的有限模。
\end{proof}

\begin{lemma}
\label{coherent-lemma-inverse-systems-pullback}
设 $f : X \to Y$ 为 Noether 概形之间的态射。
设 $\mathcal{J} \subset \mathcal{O}_Y$ 为准凝聚理想层，并令
$\mathcal{I} = f^{-1}\mathcal{J} \mathcal{O}_X$。则有右正合函子
$$
f^* : \textit{Coh}(Y, \mathcal{J}) \longrightarrow \textit{Coh}(X, \mathcal{I})
$$
它把 $(\mathcal{G}_n)$ 送到 $(f^*\mathcal{G}_n)$。若 $f$ 平坦，
则 $f^*$ 是正合函子。
\end{lemma}

\begin{proof}
由于 $f^* : \textit{Coh}(\mathcal{O}_Y) \to \textit{Coh}(\mathcal{O}_X)$
右正合，有
$$
f^*\mathcal{G}_n =
f^*(\mathcal{G}_{n + 1}/\mathcal{I}^n\mathcal{G}_{n + 1}) =
f^*\mathcal{G}_{n + 1}/f^{-1}\mathcal{I}^nf^*\mathcal{G}_{n + 1} =
f^*\mathcal{G}_{n + 1}/\mathcal{J}^nf^*\mathcal{G}_{n + 1}
$$
故一个系统的拉回仍是系统。引理
\ref{coherent-lemma-inverse-systems-abelian} 的证明中关于余核的构造表明，
$f^* : \textit{Coh}(Y, \mathcal{J}) \to \textit{Coh}(X, \mathcal{I})$
总是右正合的。若 $f$ 平坦，则
$f^* : \textit{Coh}(\mathcal{O}_Y) \to \textit{Coh}(\mathcal{O}_X)$
是正合函子。由引理 \ref{coherent-lemma-inverse-systems-abelian}
的证明中对核的构造可知，在这种情形下，
$f^* : \textit{Coh}(Y, \mathcal{J}) \to \textit{Coh}(X, \mathcal{I})$
也把核变成核。
\end{proof}

\begin{lemma}
\label{coherent-lemma-inverse-systems-pullback-equivalence}
设 $f : X' \to X$ 为 Noether 概形之间的态射。设 $Z \subset X$
为闭子概形，并以 $Z' = f^{-1}Z$ 表示其概形论逆像。
设 $\mathcal{I} \subset \mathcal{O}_X$、
$\mathcal{I}' \subset \mathcal{O}_{X'}$ 为对应的准凝聚理想层。
若 $f$ 平坦且诱导态射 $Z' \to Z$ 是同构，则拉回函子
$f^* : \textit{Coh}(X, \mathcal{I}) \to \textit{Coh}(X', \mathcal{I}')$
（引理 \ref{coherent-lemma-inverse-systems-pullback}）是等价。
\end{lemma}

\begin{proof}
若 $X$ 与 $X'$ 仿射，则结论立即由《代数进阶》引理
\ref{more-algebra-lemma-neighbourhood-equivalence} 推出。
为在一般情形下证明它，令 $Z_n \subset X$、$Z'_n \subset X'$
为第 $n$ 个无穷小邻域，分别对应于 $Z$、$Z'$。
诱导态射 $Z_n \to Z'_n$ 在底拓扑空间上是同胚。
另一方面，若 $z' \in Z'$ 映到 $z \in Z$，则环同态
$\mathcal{O}_{X, z} \to \mathcal{O}_{X', z'}$ 平坦，且诱导同构
$\mathcal{O}_{X, z}/\mathcal{I}_z \to \mathcal{O}_{X', z'}/\mathcal{I}'_{z'}$。
因此，它诱导同构
$\mathcal{O}_{X, z}/\mathcal{I}_z^n \to
\mathcal{O}_{X', z'}/(\mathcal{I}'_{z'})^n$，
这对所有 $n \geq 1$ 成立；例如见《代数进阶》引理
\ref{more-algebra-lemma-neighbourhood-isomorphism}。
故 $Z'_n \to Z_n$ 是概形的同构。
因此，$f^*$ 在凝聚 $\mathcal{O}_X$-模范畴（其中对象被
$\mathcal{I}^n$ 零化）与凝聚 $\mathcal{O}_{X'}$-模范畴（其中对象被
$(\mathcal{I}')^n$ 零化）之间诱导等价；参见引理
\ref{coherent-lemma-i-star-equivalence}。这显然蕴含本引理。
\end{proof}

\begin{lemma}
\label{coherent-lemma-inverse-systems-ideals-equivalence}
设 $X$ 为 Noether 概形。设
$\mathcal{I}, \mathcal{J} \subset \mathcal{O}_X$ 为准凝聚理想层。
若 $V(\mathcal{I}) = V(\mathcal{J})$ 是 $X$ 的同一个闭子集，
则 $\textit{Coh}(X, \mathcal{I})$ 与
$\textit{Coh}(X, \mathcal{J})$ 等价。
\end{lemma}

\begin{proof}
首先假设 $X = \Spec(A)$ 仿射。设 $I, J \subset A$ 为与
$\mathcal{I}, \mathcal{J}$ 对应的理想。则
$V(I) = V(J)$ 蕴含 $I^c \subset J$ 且 $J^d \subset I$，
其中 $c, d \geq 1$ 为某些整数；这是 Zariski 拓扑的初等性质
（参见《代数》第 \ref{algebra-section-spectrum-ring} 节及引理
\ref{algebra-lemma-Noetherian-power}）。故 $I$-进与 $J$-进完备化
（对 $A$）相同；参见《代数》引理
\ref{algebra-lemma-change-ideal-completion}。因此，在此情形下，
所求等价由引理 \ref{coherent-lemma-inverse-systems-affine} 推出。

\medskip\noindent
一般而言，利用上面的论述及 $X$ 拟紧这一事实，选取
$c, d \geq 1$，使得 $\mathcal{I}^c \subset \mathcal{J}$ 且
$\mathcal{J}^d \subset \mathcal{I}$。给定对象 $(\mathcal{F}_n)$
（属于 $\textit{Coh}(X, \mathcal{I})$），
断言逆向系统
$$
(\mathcal{F}_{cn}/\mathcal{J}^n\mathcal{F}_{cn})
$$
属于 $\textit{Coh}(X, \mathcal{J})$。这可在一个仿射覆盖的各成员上检验；
细节从略。以同样方式，可以构造 $\textit{Coh}(X, \mathcal{I})$
的对象，其起点是 $\textit{Coh}(X, \mathcal{J})$ 的一个对象。
我们略去验证这些构造定义互为拟逆的函子。
\end{proof}

\section{Grothendieck 存在定理（一）}
\label{coherent-section-existence}

\noindent
本节讨论射影情形下的 Grothendieck 存在定理。我们将使用第
\ref{coherent-section-coherent-formal} 节发展的凝聚形式模概念。
建议熟悉形式概形的读者阅读 \cite{EGA} 中该定理的陈述与证明。

\begin{lemma}
\label{coherent-lemma-fully-faithful}
设 $A$ 为关于理想 $I$ 完备的 Noether 环。
设 $f : X \to \Spec(A)$ 为固有态射，并令
$\mathcal{I} = I\mathcal{O}_X$。则函子
(\ref{coherent-equation-completion-functor}) 是全忠实的。
\end{lemma}

\begin{proof}
设 $\mathcal{F}$、$\mathcal{G}$ 为凝聚 $\mathcal{O}_X$-模。
则 $\mathcal{H} = \SheafHom_{\mathcal{O}_X}(\mathcal{G}, \mathcal{F})$
是凝聚 $\mathcal{O}_X$-模；参见《模》引理
\ref{modules-lemma-internal-hom-locally-kernel-direct-sum}。
由引理 \ref{coherent-lemma-completion-internal-hom}，映射
$$
\lim_n H^0(X, \mathcal{H}/\mathcal{I}^n\mathcal{H})
\to
\Mor_{\textit{Coh}(X, \mathcal{I})}
(\mathcal{G}^\wedge, \mathcal{F}^\wedge)
$$
是双射。因此，函子 (\ref{coherent-equation-completion-functor}) 的全忠实性
由应用于凝聚层 $\mathcal{H}$ 的形式函数定理
（引理 \ref{coherent-lemma-spell-out-theorem-formal-functions}）推出。
\end{proof}

\begin{lemma}
\label{coherent-lemma-vanishing-projective}
设 $A$ 为 Noether 环，$I \subset A$ 为理想。
设 $f : X \to \Spec(A)$ 为固有态射，并设 $\mathcal{L}$
为 $f$-丰沛可逆层。令 $\mathcal{I} = I\mathcal{O}_X$。
设 $(\mathcal{F}_n)$ 为 $\textit{Coh}(X, \mathcal{I})$ 的对象。
则存在整数 $d_0$，使得
$$
H^1(X, \Ker(\mathcal{F}_{n + 1} \to \mathcal{F}_n)
\otimes \mathcal{L}^{\otimes d} )
= 0
$$
对所有 $n \geq 0$ 及所有 $d \geq d_0$ 成立。
\end{lemma}

\begin{proof}
令 $B = \bigoplus I^n/I^{n + 1}$，并令
$\mathcal{B} = \bigoplus \mathcal{I}^n/\mathcal{I}^{n + 1} = f^*\widetilde{B}$。
由引理 \ref{coherent-lemma-finite-over-rees-algebra}，分次准凝聚
$\mathcal{B}$-模
$\mathcal{G} = \bigoplus \Ker(\mathcal{F}_{n + 1} \to \mathcal{F}_n)$
是有限型的。故本引理由引理 \ref{coherent-lemma-graded-finiteness} 的 (2) 推出。
\end{proof}

\begin{lemma}
\label{coherent-lemma-existence-projective}
设 $A$ 为关于理想 $I$ 完备的 Noether 环。
设 $f : X \to \Spec(A)$ 为射影态射，并令
$\mathcal{I} = I\mathcal{O}_X$。则函子
(\ref{coherent-equation-completion-functor}) 是等价。
\end{lemma}

\begin{proof}
(\ref{coherent-equation-completion-functor}) 是全忠实的；这已在引理
\ref{coherent-lemma-fully-faithful} 中看到。因此只须证明该函子
本质满射。

\medskip\noindent
先证明每个对象 $(\mathcal{F}_n)$（属于
$\textit{Coh}(X, \mathcal{I})$）都是函子
(\ref{coherent-equation-completion-functor}) 之像中某个对象的商。
设 $\mathcal{L}$ 为 $f$-丰沛可逆层（在 $X$ 上）。
按引理 \ref{coherent-lemma-vanishing-projective} 选取
$d_0$。选取 $d \geq d_0$，使得
$\mathcal{F}_1 \otimes \mathcal{L}^{\otimes d}$
由某些整体截面 $s_{1, 1}, \ldots, s_{t, 1}$ 生成。
由于系统的转移映射
$$
H^0(X, \mathcal{F}_{n + 1} \otimes \mathcal{L}^{\otimes d})
\longrightarrow
H^0(X, \mathcal{F}_n \otimes \mathcal{L}^{\otimes d})
$$
因 $H^1$ 的消失而成为满射，可以把
$s_{1, 1}, \ldots, s_{t, 1}$ 提升为相容整体截面系
$s_{1, n}, \ldots, s_{t, n}$；这些是
$\mathcal{F}_n \otimes \mathcal{L}^{\otimes d}$ 的截面。
它们确定相容映射系
$$
(s_{1, n}, \ldots, s_{t, n}) :
(\mathcal{L}^{\otimes -d})^{\oplus t} \longrightarrow \mathcal{F}_n
$$
由引理 \ref{coherent-lemma-inverse-systems-surjective}，得到所需的满射
$$
\left((\mathcal{L}^{\otimes -d})^{\oplus t}\right)^\wedge
\longrightarrow
(\mathcal{F}_n)
$$

\medskip\noindent
上一段的结论以及 $\textit{Coh}(X, \mathcal{I})$ 是阿贝尔范畴这一事实
（引理 \ref{coherent-lemma-inverse-systems-abelian}）蕴含：
$\textit{Coh}(X, \mathcal{I})$ 的每个对象，都是来自
$\textit{Coh}(\mathcal{O}_X)$ 的两个对象之间某个映射的余核。
由于 (\ref{coherent-equation-completion-functor}) 全忠实且正合
（引理 \ref{coherent-lemma-fully-faithful} 与 \ref{coherent-lemma-exact}），结论成立。
\end{proof}

\section{Grothendieck 存在定理（二）}
\label{coherent-section-existence-proper}

\noindent
本节讨论固有情形下的 Grothendieck 存在定理。在给出其陈述与证明前，
还需要进一步发展第 \ref{coherent-section-coherent-formal} 节引入的凝聚形式模范畴
$\textit{Coh}(X, \mathcal{I})$ 的一些理论。

\begin{remark}
\label{coherent-remark-inverse-systems-kernel-cokernel-annihilated-by}
设 $X$ 为 Noether 概形，并设
$\mathcal{I}, \mathcal{K} \subset \mathcal{O}_X$ 为准凝聚理想层。
设 $\alpha : (\mathcal{F}_n) \to (\mathcal{G}_n)$ 为
$\textit{Coh}(X, \mathcal{I})$ 的态射。
给定仿射开集 $\Spec(A) = U \subset X$，其中
$\mathcal{I}|_U, \mathcal{K}|_U$ 对应于理想 $I, K \subset A$，
以 $\alpha_U : M \to N$ 表示有限 $A^\wedge$-模之间的映射，
它经由引理 \ref{coherent-lemma-inverse-systems-affine} 对应于 $\alpha|_U$。
我们断言下列条件等价：
\begin{enumerate}
\item 存在整数 $t \geq 1$，使得
$\Ker(\alpha_n)$ 与 $\Coker(\alpha_n)$ 被 $\mathcal{K}^t$ 零化，
这对所有 $n \geq 1$ 成立；
\item 对任意上述仿射开集 $\Spec(A) = U \subset X$，模
$\Ker(\alpha_U)$ 与 $\Coker(\alpha_U)$ 被某个 $K^t$ 零化，
其中 $t \geq 1$ 为某个整数；并且
\item 存在有限仿射开覆盖 $X = \bigcup U_i$，使得 (2) 的结论
对 $\alpha_{U_i}$ 成立。
\end{enumerate}
若这些等价条件成立，就称 $\alpha$ 为
{\it 核与余核被 $\mathcal{K}$ 的某个幂零化的映射}。
为看出其等价性，使用如下交换代数事实：假设给定 $A$-模的正合列
$$
0 \to T \to M \to N \to Q \to 0
$$
其中 $T$ 与 $Q$ 被 $K^t$ 零化，这里 $K \subset A$ 为某个理想。
则对每个 $f, g \in K^t$，存在典范映射
$"fg": N \to M$，使得 $M \to N \to M$ 等于乘以 $fg$。
事实上，对 $y \in N$，可以选取 $x \in M$，其像为
$fy$（在 $N$ 中），再令 $"fg"(y) = gx$。由此显然可见，
$\Ker(M/JM \to N/JN)$ 与 $\Coker(M/JM \to N/JN)$
均被 $K^{2t}$ 零化，这对任意理想 $J \subset A$ 成立。

\medskip\noindent
把该交换代数事实应用于 $\alpha_{U_i}$ 与 $J = I^n$，
可见 (3) 蕴含 (1)。反之，假设 (1) 成立且
$M \to N$ 等于 $\alpha_U$。则存在 $t \geq 1$，使得
$\Ker(M/I^nM \to N/I^nN)$ 与 $\Coker(M/I^nM \to N/I^nN)$
被 $K^t$ 零化，这对所有 $n$ 成立。得到映射
$"fg" : N/I^nN \to M/I^nM$，它们在极限中诱导映射
$N \to M$；这里 $N$ 与 $M$ 都是 $I$-进完备的。因为与
$N \to M \to N$ 的复合是乘以 $fg$，可知 $fg$ 零化 $T$ 与 $Q$。
换言之，$T$ 与 $Q$ 如所需地被 $K^{2t}$ 零化。
\end{remark}

\begin{lemma}
\label{coherent-lemma-existence-tricky}
设 $X$ 为 Noether 概形。设
$\mathcal{I}, \mathcal{K} \subset \mathcal{O}_X$ 为准凝聚理想层。
设 $X_e \subset X$ 为由 $\mathcal{K}^e$ 截出的闭子概形。
令 $\mathcal{I}_e = \mathcal{I}\mathcal{O}_{X_e}$。
设 $(\mathcal{F}_n)$ 为 $\textit{Coh}(X, \mathcal{I})$ 的对象。
假设：
\begin{enumerate}
\item 函子
$\textit{Coh}(\mathcal{O}_{X_e}) \to \textit{Coh}(X_e, \mathcal{I}_e)$
对所有 $e \geq 1$ 都是等价；并且
\item 存在凝聚层 $\mathcal{H}$（在 $X$ 上）及映射
$\alpha : (\mathcal{F}_n) \to \mathcal{H}^\wedge$，其核与余核
被 $\mathcal{K}$ 的某个幂零化。
\end{enumerate}
则 $(\mathcal{F}_n)$ 在 (\ref{coherent-equation-completion-functor}) 的本质像中。
\end{lemma}

\begin{proof}
在本证明中，将不再特别说明地使用如下事实：对闭浸入
$i : Z \to X$，函子 $i_*$ 在 $Z$ 上凝聚模范畴与 $X$ 上被
$Z$ 的理想层零化的凝聚模范畴之间给出等价；参见引理
\ref{coherent-lemma-i-star-equivalence}。特别地，可以把
$\textit{Coh}(\mathcal{O}_{X_e})$ 等同于凝聚 $\mathcal{O}_X$-模范畴
（其中对象被 $\mathcal{K}^e$ 零化），并把
$\textit{Coh}(X_e, \mathcal{I}_e)$
等同于 $\textit{Coh}(X, \mathcal{I})$ 中由被 $\mathcal{K}^e$
零化的对象组成的满子范畴。而且，(1) 告诉我们，这两个范畴在完备化函子
(\ref{coherent-equation-completion-functor}) 下等价。

\medskip\noindent
应用该等价，得到凝聚 $\mathcal{O}_X$-模 $\mathcal{G}_e$，
它被 $\mathcal{K}^e$ 零化，并对应于系统
$(\mathcal{F}_n/\mathcal{K}^e\mathcal{F}_n)$；该系统属于
$\textit{Coh}(X, \mathcal{I})$。映射
$\mathcal{F}_n/\mathcal{K}^{e + 1}\mathcal{F}_n \to
\mathcal{F}_n/\mathcal{K}^e\mathcal{F}_n$ 对应于典范映射
$\mathcal{G}_{e + 1} \to \mathcal{G}_e$，它们诱导同构
$\mathcal{G}_{e + 1}/\mathcal{K}^e\mathcal{G}_{e + 1} \to \mathcal{G}_e$。
故 $(\mathcal{G}_e)$ 是 $\textit{Coh}(X, \mathcal{K})$ 的对象。
映射 $\alpha$ 诱导映射系
$$
\mathcal{F}_n/\mathcal{K}^e\mathcal{F}_n
\longrightarrow
\mathcal{H}/(\mathcal{I}^n + \mathcal{K}^e)\mathcal{H}
$$
从而得到映射
$\mathcal{G}_e \to \mathcal{H}/\mathcal{K}^e\mathcal{H}$
（再次使用范畴等价）。由假设 (2)，存在整数 $t \geq 1$，
使得 $\mathcal{K}^t$ 零化所有映射
$\mathcal{F}_n \to \mathcal{H}/\mathcal{I}^n\mathcal{H}$ 的核与余核。
于是，$\mathcal{K}^{2t}$ 零化映射
$\mathcal{F}_n/\mathcal{K}^e\mathcal{F}_n \to
\mathcal{H}/(\mathcal{I}^n + \mathcal{K}^e)\mathcal{H}$ 的核与余核；
参见评注 \ref{coherent-remark-inverse-systems-kernel-cokernel-annihilated-by}。
继而可知 $\mathcal{K}^{4t}$ 零化下列映射的核与余核：
$$
\mathcal{G}_e
\longrightarrow
\mathcal{H}/\mathcal{K}^e\mathcal{H},
$$
参见评注 \ref{coherent-remark-inverse-systems-kernel-cokernel-annihilated-by}。
应用引理 \ref{coherent-lemma-existence-easy}，得到凝聚 $\mathcal{O}_X$-模
$\mathcal{F}$、映射 $a : \mathcal{F} \to \mathcal{H}$，以及同构
$\beta : (\mathcal{G}_e) \to (\mathcal{F}/\mathcal{K}^e\mathcal{F})$
（在 $\textit{Coh}(X, \mathcal{K})$ 中）。
反向追溯，对给定的 $n$，三元组
$(\mathcal{F}/\mathcal{I}^n\mathcal{F}, a \bmod \mathcal{I}^n, \beta
\bmod \mathcal{I}^n)$ 是该引理中对应于态射
$\alpha_n \bmod \mathcal{K}^e :
(\mathcal{F}_n/\mathcal{K}^e\mathcal{F}_n) \to
(\mathcal{H}/(\mathcal{I}^n + \mathcal{K}^e)\mathcal{H})$
（作为 $\textit{Coh}(X, \mathcal{K})$ 的态射）的三元组。
因此，引理 \ref{coherent-lemma-existence-easy} 的唯一性给出典范同构
$\mathcal{F}/\mathcal{I}^n\mathcal{F} \to \mathcal{F}_n$，
并与涉及的所有态射相容。这就完成了引理的证明。
\end{proof}

\begin{lemma}
\label{coherent-lemma-inverse-systems-push-pull}
设 $Y$ 为 Noether 概形。设
$\mathcal{J}, \mathcal{K} \subset \mathcal{O}_Y$ 为准凝聚理想层。
设 $f : X \to Y$ 为固有态射，并且它在
$V = Y \setminus V(\mathcal{K})$ 上为同构。
令 $\mathcal{I} = f^{-1}\mathcal{J} \mathcal{O}_X$。
设 $(\mathcal{G}_n)$ 为 $\textit{Coh}(Y, \mathcal{J})$ 的对象，
$\mathcal{F}$ 为凝聚 $\mathcal{O}_X$-模，并设
$\beta : (f^*\mathcal{G}_n)  \to \mathcal{F}^\wedge$ 为
$\textit{Coh}(X, \mathcal{I})$ 中的同构。则存在映射
$$
\alpha :
(\mathcal{G}_n)
\longrightarrow
(f_*\mathcal{F})^\wedge
$$
（在 $\textit{Coh}(Y, \mathcal{J})$ 中），其核与余核被
$\mathcal{K}$ 的某个幂零化。
\end{lemma}

\begin{proof}
由于 $f$ 是固有态射，可知 $f_*\mathcal{F}$ 是凝聚
$\mathcal{O}_Y$-模（命题
\ref{coherent-proposition-proper-pushforward-coherent}）。故引理的陈述有意义。
考虑复合映射
$$
\gamma_n : \mathcal{G}_n \to
f_*f^*\mathcal{G}_n \to
f_*(\mathcal{F}/\mathcal{I}^n\mathcal{F}).
$$
这里第一个映射是伴随映射，第二个是 $f_*\beta_n$。
我们断言：存在唯一的引理中所述 $\alpha$，使得复合映射
$$
\mathcal{G}_n \xrightarrow{\alpha_n}
f_*\mathcal{F}/\mathcal{J}^nf_*\mathcal{F} \to
f_*(\mathcal{F}/\mathcal{I}^n\mathcal{F})
$$
都等于 $\gamma_n$，这对所有 $n$ 成立。由于唯一性，可以假设
$Y = \Spec(B)$ 仿射。令 $J \subset B$ 为与理想层
$\mathcal{J}$ 对应的理想。令
$$
M_n = H^0(X, \mathcal{F}/\mathcal{I}^n\mathcal{F})
\quad\text{且}\quad
M = H^0(X, \mathcal{F})
$$
由引理 \ref{coherent-lemma-ML-cohomology-powers-ideal} 与定理
\ref{coherent-theorem-formal-functions}，诸模 $M_n$ 的逆极限等于完备化
$M^\wedge = \lim M/J^nM$。令
$N_n = H^0(Y, \mathcal{G}_n)$ 且 $N = \lim N_n$。
经由引理 \ref{coherent-lemma-inverse-systems-affine} 的范畴等价，有限
$B^\wedge$-模 $N$ 与 $M^\wedge$ 分别对应于
$(\mathcal{G}_n)$ 与 $f_*\mathcal{F}^\wedge$。
由此，$\alpha$ 必须是 $\textit{Coh}(Y, \mathcal{J})$ 的态射，
它对应于同态
$$
\lim \gamma_n : N = \lim_n N_n \longrightarrow \lim M_n = M^\wedge
$$
（在有限 $B^\wedge$-模之间）。

\medskip\noindent
还须证明 $\alpha$ 的核与余核被 $\mathcal{K}$ 的某个幂零化。
令 $Y' = \Spec(B^\wedge)$，并令 $X' = Y' \times_Y X$。
令 $\mathcal{K}'$、$\mathcal{J}'$、$\mathcal{G}'_n$ 以及
$\mathcal{I}'$、$\mathcal{F}'$ 分别为 $\mathcal{K}$、$\mathcal{J}$、
$\mathcal{G}_n$ 以及 $\mathcal{I}$、$\mathcal{F}$ 到
$Y'$ 与 $X'$ 的拉回。投影态射 $f' : X' \to Y'$ 是 $f$
沿 $Y' \to Y$ 的基变换。注意 $Y' \to Y$ 是平坦概形态射，
因为由《代数》引理 \ref{algebra-lemma-completion-flat}，
$B \to B^\wedge$ 平坦。因此，由引理
\ref{coherent-lemma-flat-base-change-cohomology}，$f'_*\mathcal{F}'$ 以及
$f'_*(f')^*\mathcal{G}_n'$ 分别是 $f_*\mathcal{F}$ 以及
$f_*f^*\mathcal{G}_n$ 到 $Y'$ 的拉回。
构造的唯一性表明，$\alpha$ 到 $Y'$ 的拉回就是对应映射
$\alpha'$，它为 $Y'$ 上情形构造。而且，为检验 $\alpha$ 的核与余核被
$\mathcal{K}^t$ 零化，只须检验 $\alpha'$ 的核与余核被
$(\mathcal{K}')^t$ 零化。确实，为此须检验映射
$\alpha_n$ 与 $\alpha'_n$ 的核与余核（见评注
\ref{coherent-remark-inverse-systems-kernel-cokernel-annihilated-by}），
而环同态 $B \to B^\wedge$ 在被 $J^n$ 与被 $(J')^n$ 零化的模
之间诱导范畴等价；参见《代数进阶》引理
\ref{more-algebra-lemma-neighbourhood-equivalence}。
故可以假设 $B$ 关于 $J$ 完备。

\medskip\noindent
假设 $Y = \Spec(B)$ 仿射，$\mathcal{J}$ 对应于理想
$J \subset B$，且 $B$ 关于 $J$ 完备。
在这种情形下，$(\mathcal{G}_n)$ 位于函子
$\textit{Coh}(\mathcal{O}_Y) \to \textit{Coh}(Y, \mathcal{J})$
的本质像中。设 $\mathcal{G}$ 是凝聚 $\mathcal{O}_Y$-模，满足
$(\mathcal{G}_n) = \mathcal{G}^\wedge$。注意
$f^*(\mathcal{G}^\wedge) = (f^*\mathcal{G})^\wedge$。
因此，引理 \ref{coherent-lemma-fully-faithful} 告诉我们，$\beta$ 来自同构
$b : f^*\mathcal{G} \to \mathcal{F}$，且 $\alpha$ 是把完备化函子
应用于下列映射所得的态射：
$$
\mathcal{G} \to f_*f^*\mathcal{G} \cong f_*\mathcal{F}
$$
故要验证伴随映射 $c : \mathcal{G} \to f_*f^*\mathcal{G}$ 的核与余核
被 $\mathcal{K}$ 的某个幂零化。然而，由于限制
$f|_{f^{-1}(V)} : f^{-1}(V) \to V$ 是同构，可知 $c|_V$ 是同构。
因此，凝聚层 $\Ker(c)$ 与 $\Coker(c)$ 支撑在
$V(\mathcal{K})$ 上，故被 $\mathcal{K}$ 的某个幂零化
（引理 \ref{coherent-lemma-power-ideal-kills-sheaf}），如所需。
\end{proof}

\noindent
下述命题是实践中最常用的 Grothendieck 存在定理形式。

\begin{proposition}
\label{coherent-proposition-existence-proper}
设 $A$ 为关于理想 $I$ 完备的 Noether 环。
设 $f : X \to \Spec(A)$ 为固有概形态射。
令 $\mathcal{I} = I\mathcal{O}_X$。则函子
(\ref{coherent-equation-completion-functor}) 是等价。
\end{proposition}

\begin{proof}
已经看到，函子 (\ref{coherent-equation-completion-functor}) 是全忠实的；
见引理 \ref{coherent-lemma-fully-faithful}。因此只须证明该函子本质满射。

\medskip\noindent
考虑集合 $\Xi$，其元素是准凝聚理想层
$\mathcal{K} \subset \mathcal{O}_X$，且满足：每个对象
$(\mathcal{F}_n)$（被 $\mathcal{K}$ 零化）均位于本质像中。
要证明 $(0)$ 属于 $\Xi$。
若不然，由于 $X$ 是 Noether 概形，存在极大准凝聚理想层
$\mathcal{K}$，它不属于 $\Xi$；参见引理 \ref{coherent-lemma-acc-coherent}。
将 $X$ 替换为 $X$ 中与 $\mathcal{K}$ 对应的闭子概形后，
可以假设
每个非零 $\mathcal{K}$ 均属于 $\Xi$。（这里使用了被
$\mathcal{K}$ 零化的凝聚模与 $\mathcal{K}$ 所对应闭子概形上的
凝聚模之间的对应；参见引理 \ref{coherent-lemma-i-star-equivalence}。）
设 $(\mathcal{F}_n)$ 为 $\textit{Coh}(X, \mathcal{I})$ 的对象。
将证明该对象位于函子 (\ref{coherent-equation-completion-functor}) 的本质像中，
从而完成命题的证明。

\medskip\noindent
应用 Chow 引理（引理 \ref{coherent-lemma-chow-Noetherian}），得到固有满射
$f : X' \to X$，它在稠密开集 $U \subset X$ 上为同构，
且 $X'$ 在 $A$ 上射影。设 $\mathcal{K}$ 为截出约化补集
$X \setminus U$ 的准凝聚理想层。由射影情形的 Grothendieck 存在定理
（引理 \ref{coherent-lemma-existence-projective}），存在凝聚模
$\mathcal{F}'$（在 $X'$ 上），使得
$(\mathcal{F}')^\wedge \cong (f^*\mathcal{F}_n)$。
由命题 \ref{coherent-proposition-proper-pushforward-coherent}，
$\mathcal{O}_X$-模 $\mathcal{H} = f_*\mathcal{F}'$ 凝聚；
由引理 \ref{coherent-lemma-inverse-systems-push-pull}，存在态射
$(\mathcal{F}_n) \to \mathcal{H}^\wedge$
（在 $\textit{Coh}(X, \mathcal{I})$ 中），其核与余核被
$\mathcal{K}$ 的某个幂零化。诸幂 $\mathcal{K}^e$ 都属于
$\Xi$，因而 (\ref{coherent-equation-completion-functor}) 对闭子概形
$X_e = V(\mathcal{K}^e)$ 是等价。最后应用引理
\ref{coherent-lemma-existence-tricky} 即得结论。
\end{proof}

\section{在基底上为固有}
\label{coherent-section-proper-over-base}

\noindent
本节很短，只是指出在基底上固有的闭子集，以及支撑在基底上固有的
有限型准凝聚模所具有的一些有用性质。

\begin{lemma}
\label{coherent-lemma-closed-proper-over-base}
设 $f : X \to S$ 为局部有限型概形态射，并设 $Z \subset X$ 为闭子集。
下列条件等价：
\begin{enumerate}
\item 态射 $Z \to S$ 在赋予 $Z$ 诱导的约化闭子概形结构时固有
（《概形》定义 \ref{schemes-definition-reduced-induced-scheme}）；
\item 对 $Z$ 的某个闭子概形结构，态射 $Z \to S$ 固有；
\item 对 $Z$ 的任意闭子概形结构，态射 $Z \to S$ 固有。
\end{enumerate}
\end{lemma}

\begin{proof}
(3) $\Rightarrow$ (1) 与 (1) $\Rightarrow$ (2) 都是立即的，
故只须证明 (2) 蕴含 (3)。我们建议读者自行找出这个事实的证明。
设 $Z'$ 与 $Z''$ 是 $Z$ 上的闭子概形结构，且 $Z' \to S$ 固有。
须证明 $Z'' \to S$ 固有。令 $Z''' = Z' \cup Z''$ 为概形论并；
参见《态射》定义
\ref{morphisms-definition-scheme-theoretic-intersection-union}。
则 $Z'''$ 是 $Z$ 上的另一个闭子概形结构。例如，这可由
《态射》引理 \ref{morphisms-lemma-scheme-theoretic-union}
对概形论并的描述推出。由于 $Z'' \to Z'''$ 是闭浸入，
只须证明 $Z''' \to S$ 固有（见《态射》引理
\ref{morphisms-lemma-closed-immersion-proper} 与
\ref{morphisms-lemma-composition-proper}）。
态射 $Z' \to Z'''$ 是双射闭浸入，特别地，它满且万有闭。
于是，$Z' \to S$ 分离蕴含 $Z''' \to S$ 分离；见《态射》引理
\ref{morphisms-lemma-image-universally-closed-separated}。
此外，$Z''' \to S$ 局部有限型，因为 $X \to S$ 局部有限型
（《态射》引理 \ref{morphisms-lemma-immersion-locally-finite-type}
与 \ref{morphisms-lemma-composition-finite-type}）。
由于 $Z' \to S$ 拟紧，且 $Z' \to Z'''$ 是同胚，可知
$Z''' \to S$ 拟紧。最后，由于 $Z' \to S$ 万有闭，
《态射》引理 \ref{morphisms-lemma-image-proper-is-proper}
表明 $Z''' \to S$ 也万有闭。这就完成了证明。
\end{proof}

\begin{definition}
\label{coherent-definition-proper-over-base}
设 $f : X \to S$ 为局部有限型概形态射，并设 $Z \subset X$ 为闭子集。
若引理 \ref{coherent-lemma-closed-proper-over-base} 的等价条件成立，
则称{\it $Z$ 在 $S$ 上固有}。
\end{definition}

\noindent
如果态射 $f : X \to S$ 不是局部有限型的，则上述定义所用的引理不成立。
因此，当 $f$ 不是局部有限型时，我们建议读者不要使用这一术语。

\begin{lemma}
\label{coherent-lemma-closed-closed-proper-over-base}
设 $f : X \to S$ 为局部有限型概形态射，并设
$Y \subset Z \subset X$ 为闭子集。若 $Z$ 在 $S$ 上固有，
则 $Y$ 也具有同一性质。
\end{lemma}

\begin{proof}
略。
\end{proof}

\begin{lemma}
\label{coherent-lemma-base-change-closed-proper-over-base}
考虑概形的笛卡尔图
$$
\xymatrix{
X' \ar[d]_{f'} \ar[r]_{g'} & X \ar[d]^f \\
S' \ar[r]^g & S
}
$$
其中 $f$ 局部有限型。若 $Z$ 是 $X$ 的闭子集并在 $S$ 上固有，
则 $(g')^{-1}(Z)$ 是 $X'$ 的闭子集并在 $S'$ 上固有。
\end{lemma}

\begin{proof}
由《态射》引理 \ref{morphisms-lemma-base-change-finite-type}，
$f'$ 局部有限型，故陈述有意义。赋予 $Z$ 诱导的约化闭子概形结构。
以 $Z' = (g')^{-1}(Z)$ 表示概形论逆像
（《概形》定义 \ref{schemes-definition-inverse-image-closed-subscheme}）。
则
$Z' = X' \times_X Z = (S' \times_S X) \times_X Z = S' \times_S Z$
在 $S'$ 上固有，因为它是 $Z$（在 $S$ 上）的基变换
（《态射》引理 \ref{morphisms-lemma-base-change-proper}）。
\end{proof}

\begin{lemma}
\label{coherent-lemma-functoriality-closed-proper-over-base}
设 $S$ 为概形。设 $f : X \to Y$ 为两个在 $S$ 上局部有限型的
概形之间的态射。
\begin{enumerate}
\item 若 $Y$ 在 $S$ 上分离，且 $Z \subset X$ 是在 $S$ 上固有的
闭子集，则 $f(Z)$ 是 $Y$ 的闭子集并在 $S$ 上固有。
\item 若 $f$ 万有闭，且 $Z \subset X$ 是在 $S$ 上固有的闭子集，
则 $f(Z)$ 是 $Y$ 的闭子集并在 $S$ 上固有。
\item 若 $f$ 固有，且 $Z \subset Y$ 是在 $S$ 上固有的闭子集，
则 $f^{-1}(Z)$ 是 $X$ 的闭子集并在 $S$ 上固有。
\end{enumerate}
\end{lemma}

\begin{proof}
证明 (1)。假设 $Y$ 在 $S$ 上分离，且 $Z \subset X$
是一个在 $S$ 上固有的闭子集。赋予 $Z$ 诱导的约化闭子概形结构，
并把《态射》引理
\ref{morphisms-lemma-scheme-theoretic-image-is-proper}
应用于态射 $Z \to Y$（在 $S$ 上），即得结论。

\medskip\noindent
证明 (2)。假设 $f$ 万有闭，且 $Z \subset X$ 是在 $S$ 上固有的
闭子集。赋予 $Z$ 与 $Z' = f(Z)$ 各自诱导的约化闭子概形结构，
从而得到诱导态射 $Z \to Z'$。
以 $Z'' = f^{-1}(Z')$ 表示概形论逆像
（《概形》定义 \ref{schemes-definition-inverse-image-closed-subscheme}）。
则 $Z'' \to Z'$ 作为 $f$ 的基变换是万有闭的
（《态射》引理 \ref{morphisms-lemma-base-change-proper}）。
于是，$Z \to Z'$ 作为闭浸入 $Z \to Z''$ 与 $Z'' \to Z'$ 的复合
也是万有闭的（《态射》引理
\ref{morphisms-lemma-closed-immersion-proper} 与
\ref{morphisms-lemma-composition-proper}）。
由《态射》引理 \ref{morphisms-lemma-image-universally-closed-separated}，
$Z' \to S$ 分离。由于 $Z \to S$ 拟紧且 $Z \to Z'$ 满，
可知 $Z' \to S$ 拟紧。由于 $Z' \to S$ 是 $Z' \to Y$ 与
$Y \to S$ 的复合，可知 $Z' \to S$ 局部有限型
（《态射》引理 \ref{morphisms-lemma-immersion-locally-finite-type}
与 \ref{morphisms-lemma-composition-finite-type}）。
最后，由于 $Z \to S$ 万有闭，《态射》引理
\ref{morphisms-lemma-image-proper-is-proper}
表明 $Z' \to S$ 也万有闭。这就完成了证明。

\medskip\noindent
证明 (3)。假设 $f$ 固有，且 $Z \subset Y$ 是在 $S$ 上固有的闭子集。
赋予 $Z$ 诱导的约化闭子概形结构。
以 $Z' = f^{-1}(Z)$ 表示概形论逆像
（《概形》定义 \ref{schemes-definition-inverse-image-closed-subscheme}）。
则 $Z' \to Z$ 作为 $f$ 的基变换是固有的
（《态射》引理 \ref{morphisms-lemma-base-change-proper}）。
因此，$Z' \to S$ 作为 $Z' \to Z$ 与 $Z \to S$ 的复合是固有的
（《态射》引理 \ref{morphisms-lemma-composition-proper}）。
这就完成了证明。
\end{proof}

\begin{lemma}
\label{coherent-lemma-union-closed-proper-over-base}
设 $f : X \to S$ 为局部有限型概形态射。设
$Z_i \subset X$（$i = 1, \ldots, n$）为闭子集。
若 $Z_i$（$i = 1, \ldots, n$）都在 $S$ 上固有，则
$Z_1 \cup \ldots \cup Z_n$ 也具有同一性质。
\end{lemma}

\begin{proof}
赋予 $Z_i$ 各自诱导的约化闭子概形结构。态射
$$
Z_1 \amalg \ldots \amalg Z_n \longrightarrow X
$$
由《态射》引理 \ref{morphisms-lemma-closed-immersion-finite}
与 \ref{morphisms-lemma-finite-union-finite} 是有限的。
由于有限态射万有闭（《态射》引理
\ref{morphisms-lemma-finite-proper}），并且
$Z_1 \amalg \ldots \amalg Z_n$ 在 $S$ 上固有，
由引理 \ref{coherent-lemma-functoriality-closed-proper-over-base} 的 (2)
可知其像 $Z_1 \cup \ldots \cup Z_n$ 在 $S$ 上固有。
\end{proof}

\noindent
设 $f : X \to S$ 为局部有限型概形态射，并设 $\mathcal{F}$ 为有限型
准凝聚 $\mathcal{O}_X$-模。则支撑
$\text{Supp}(\mathcal{F})$（即 $\mathcal{F}$ 的支撑）是 $X$ 的闭子集；
见《态射》引理
\ref{morphisms-lemma-support-finite-type}。因此，说
“$\mathcal{F}$ 的支撑在 $S$ 上固有”是有意义的。

\begin{lemma}
\label{coherent-lemma-module-support-proper-over-base}
设 $f : X \to S$ 为局部有限型概形态射，并设 $\mathcal{F}$ 为有限型
准凝聚 $\mathcal{O}_X$-模。下列条件等价：
\begin{enumerate}
\item $\mathcal{F}$ 的支撑在 $S$ 上固有；
\item $\mathcal{F}$ 的概形论支撑
（《态射》定义 \ref{morphisms-definition-scheme-theoretic-support}）
在 $S$ 上固有；
\item 存在闭子概形 $Z \subset X$ 及有限型准凝聚
$\mathcal{O}_Z$-模 $\mathcal{G}$，使得 (a) $Z \to S$ 固有，
且 (b) $(Z \to X)_*\mathcal{G} = \mathcal{F}$。
\end{enumerate}
\end{lemma}

\begin{proof}
$\text{Supp}(\mathcal{F})$（即 $\mathcal{F}$ 的支撑）是 $X$ 的闭子集；
见《态射》引理 \ref{morphisms-lemma-support-finite-type}。
故可应用定义 \ref{coherent-definition-proper-over-base}。
由于 $\mathcal{F}$ 的概形论支撑是一个闭子概形，其底层闭子集为
$\text{Supp}(\mathcal{F})$，定义 \ref{coherent-definition-proper-over-base}
表明 (1) 与 (2) 等价。(2) 蕴含 (3) 是清楚的。
反之，若 (3) 成立，则
$\text{Supp}(\mathcal{F}) \subset Z$
（这是 $X$ 的闭子集之间的包含），因而例如由引理
\ref{coherent-lemma-closed-closed-proper-over-base}，
$\text{Supp}(\mathcal{F})$ 在 $S$ 上固有。
\end{proof}

\begin{lemma}
\label{coherent-lemma-base-change-module-support-proper-over-base}
考虑概形的笛卡尔图
$$
\xymatrix{
X' \ar[d]_{f'} \ar[r]_{g'} & X \ar[d]^f \\
S' \ar[r]^g & S
}
$$
其中 $f$ 局部有限型。设 $\mathcal{F}$ 为有限型准凝聚
$\mathcal{O}_X$-模。若 $\mathcal{F}$ 的支撑在 $S$ 上固有，
则 $(g')^*\mathcal{F}$ 的支撑在 $S'$ 上固有。
\end{lemma}

\begin{proof}
陈述有意义，因为由《模》引理
\ref{modules-lemma-pullback-finite-type}，
$(g')*\mathcal{F}$ 是有限型的。由《态射》引理
\ref{morphisms-lemma-support-finite-type}，
有 $\text{Supp}((g')^*\mathcal{F}) = (g')^{-1}(\text{Supp}(\mathcal{F}))$。
故本引理由引理 \ref{coherent-lemma-base-change-closed-proper-over-base} 推出。
\end{proof}

\begin{lemma}
\label{coherent-lemma-cat-module-support-proper-over-base}
设 $f : X \to S$ 为局部有限型概形态射。设 $\mathcal{F}$、$\mathcal{G}$
为有限型准凝聚 $\mathcal{O}_X$-模。
\begin{enumerate}
\item 若 $\mathcal{F}$、$\mathcal{G}$ 的支撑都在 $S$ 上固有，
则 $\mathcal{F} \oplus \mathcal{G}$、$\mathcal{G}$ 被
$\mathcal{F}$ 的任意扩张、$\Im(u)$ 与 $\Coker(u)$
（这里所取的是任意 $\mathcal{O}_X$-模映射
$u : \mathcal{F} \to \mathcal{G}$），以及 $\mathcal{F}$ 或 $\mathcal{G}$
的任意准凝聚商也具有同一性质。
\item 若 $S$ 局部 Noether，则凝聚 $\mathcal{O}_X$-模中支撑在
$S$ 上固有者所成范畴，是凝聚 $\mathcal{O}_X$-模所成
阿贝尔范畴的 Serre 子范畴（《同调》定义
\ref{homology-definition-serre-subcategory}）。
\end{enumerate}
\end{lemma}

\begin{proof}
证明 (1)。设 $Z$、$Z'$ 分别为 $\mathcal{F}$ 与 $\mathcal{G}$ 的支撑。
则 (1) 中提到的所有层，其支撑都包含于 $Z \cup Z'$。
只要检验这些层为有限型且准凝聚，引理
\ref{coherent-lemma-closed-closed-proper-over-base} 与
\ref{coherent-lemma-union-closed-proper-over-base}
便表明断言成立。关于准凝聚性，参见《概形》第
\ref{schemes-section-quasi-coherent} 节。关于“有限型”，
建议读者参阅《模》第 \ref{modules-section-finite-type} 节。

\medskip\noindent
证明 (2)。证明与 (1) 的证明相同。注意这些断言有意义：
由《态射》引理 \ref{morphisms-lemma-finite-type-noetherian}，
$X$ 局部 Noether；还要使用第 \ref{coherent-section-coherent-sheaves} 节
对凝聚模范畴的描述。
\end{proof}

\begin{lemma}
\label{coherent-lemma-support-proper-over-base-pushforward}
设 $S$ 为局部 Noether 概形。设 $f : X \to S$ 为局部有限型概形态射。
设 $\mathcal{F}$ 为凝聚 $\mathcal{O}_X$-模，其支撑在 $S$ 上固有。
则 $R^pf_*\mathcal{F}$ 是凝聚 $\mathcal{O}_S$-模，
这对所有 $p \geq 0$ 成立。
\end{lemma}

\begin{proof}
由引理 \ref{coherent-lemma-module-support-proper-over-base}，存在闭浸入
$i : Z \to X$ 及有限型准凝聚 $\mathcal{O}_Z$-模 $\mathcal{G}$，
使得 (a) $g = f \circ i : Z \to S$ 固有，且 (b)
$i_*\mathcal{G} = \mathcal{F}$。由命题
\ref{coherent-proposition-proper-pushforward-coherent}，
$R^pg_*\mathcal{G}$ 在 $S$ 上凝聚。另一方面，
$R^qi_*\mathcal{G} = 0$ 对 $q > 0$ 成立（引理
\ref{coherent-lemma-finite-pushforward-coherent}）。
由《上同调》引理 \ref{cohomology-lemma-relative-Leray}，
$R^pf_*\mathcal{F} = R^pg_*\mathcal{G}$，故结论成立。
\end{proof}

\begin{lemma}
\label{coherent-lemma-systems-with-proper-support}
设 $S$ 为 Noether 概形，$f : X \to S$ 为有限型态射，并设
$\mathcal{I} \subset \mathcal{O}_X$ 为准凝聚理想层。下列范畴是
$\textit{Coh}(X, \mathcal{I})$ 的 Serre 子范畴：
\begin{enumerate}
\item $\textit{Coh}(X, \mathcal{I})$ 的满子范畴，其中的对象
$(\mathcal{F}_n)$ 满足 $\mathcal{F}_1$ 的支撑在 $S$ 上固有；
\item $\textit{Coh}(X, \mathcal{I})$ 的满子范畴，其中的对象
$(\mathcal{F}_n)$ 满足：存在闭子概形 $Z \subset X$，
它在 $S$ 上固有，并且
$\mathcal{I}_Z \mathcal{F}_n = 0$ 对所有 $n \geq 1$ 成立。
\end{enumerate}
\end{lemma}

\begin{proof}
将使用《同调》引理 \ref{homology-lemma-characterize-serre-subcategory}
的判据。此外，还将使用如下事实：若
$0 \to (\mathcal{G}_n) \to (\mathcal{F}_n) \to (\mathcal{H}_n) \to 0$
是 $\textit{Coh}(X, \mathcal{I})$ 的短正合列，则
(a) $\mathcal{G}_n \to \mathcal{F}_n \to \mathcal{H}_n \to 0$
对所有 $n \geq 1$ 正合；并且 (b) $\mathcal{G}_n$ 是
$\Ker(\mathcal{F}_m \to \mathcal{H}_m)$ 的商，其中某个 $m \geq n$。
见引理 \ref{coherent-lemma-inverse-systems-abelian} 的证明。

\medskip\noindent
证明 (1)。设 $(\mathcal{F}_n)$ 为
$\textit{Coh}(X, \mathcal{I})$ 的对象。则
$\text{Supp}(\mathcal{F}_n) = \text{Supp}(\mathcal{F}_1)$
对所有 $n \geq 1$ 成立。
因此，由上述 (a) 与 (b)，对
任意短正合列
$0 \to (\mathcal{G}_n) \to (\mathcal{F}_n) \to (\mathcal{H}_n) \to 0$
（在 $\textit{Coh}(X, \mathcal{I})$ 中），
都有
$\text{Supp}(\mathcal{G}_1) \cup \text{Supp}(\mathcal{H}_1) =
\text{Supp}(\mathcal{F}_1)$。
这证明 (1) 所定义的范畴是
$\textit{Coh}(X, \mathcal{I})$ 的 Serre 子范畴。

\medskip\noindent
证明 (2)。这里作同样的论证。设
$0 \to (\mathcal{G}_n) \to (\mathcal{F}_n) \to (\mathcal{H}_n) \to 0$
为 $\textit{Coh}(X, \mathcal{I})$ 的短正合列。
若 $Z \subset X$ 为闭子概形，且 $\mathcal{I}_Z$
零化 $\mathcal{F}_n$（对所有 $n$），则由上述 (a) 与 (b)，
$\mathcal{I}_Z$ 零化 $\mathcal{G}_n$ 与
$\mathcal{H}_n$（对所有 $n$）。因此，若 $Z \to S$ 固有，便可知 (2)
所定义的范畴在 $\textit{Coh}(X, \mathcal{I})$ 中对取子对象和
商对象封闭。最后，设 $Z \subset X$ 与 $Y \subset X$ 是在 $S$
上固有的闭子概形，并且
$\mathcal{I}_Z \mathcal{G}_n = 0$ 且
$\mathcal{I}_Y \mathcal{H}_n = 0$，这对所有 $n \geq 1$ 成立。
由上述 (a) 可知
$\mathcal{I}_{Z \cup Y} = \mathcal{I}_Z \cdot \mathcal{I}_Y$
零化 $\mathcal{F}_n$（对所有 $n$）。由引理
\ref{coherent-lemma-union-closed-proper-over-base}
（以及定义 \ref{coherent-definition-proper-over-base}，它表明在该并集上
可以任选概形结构），$Z \cup Y \to S$ 固有，证明完成。
\end{proof}

\section{Grothendieck 存在定理（三）}
\label{coherent-section-existence-proper-support}

\noindent
为陈述 Grothendieck 存在定理的一般形式，先引入一些记号。
设 $A$ 为关于理想 $I$ 完备的 Noether 环。设
$f : X \to \Spec(A)$ 为分离有限型概形态射，并令
$\mathcal{I} = I\mathcal{O}_X$。在此情形下，以
$$
\textit{Coh}_{\text{support proper over } A}(\mathcal{O}_X)
$$
表示 $\textit{Coh}(\mathcal{O}_X)$ 的满子范畴，其中的对象是凝聚
$\mathcal{O}_X$-模，且其支撑在 $\Spec(A)$ 上固有。由引理
\ref{coherent-lemma-cat-module-support-proper-over-base}，这是
$\textit{Coh}(\mathcal{O}_X)$ 的 Serre 子范畴。类似地，以
$$
\textit{Coh}_{\text{support proper over } A}(X, \mathcal{I})
$$
表示 $\textit{Coh}(X, \mathcal{I})$ 的满子范畴，其中的对象
$(\mathcal{F}_n)$ 满足：$\mathcal{F}_1$ 的支撑在
$\Spec(A)$ 上固有。由引理
\ref{coherent-lemma-systems-with-proper-support} 的 (1)，这是
$\textit{Coh}(X, \mathcal{I})$ 的 Serre 子范畴。
由于商模的支撑包含在原模的支撑中，
(\ref{coherent-equation-completion-functor}) 诱导函子
\begin{equation}
\label{coherent-equation-completion-functor-proper-over-A}
\textit{Coh}_{\text{support proper over }A}(\mathcal{O}_X)
\longrightarrow
\textit{Coh}_{\text{support proper over }A}(X, \mathcal{I})
\end{equation}
现在可以陈述本节的主要定理。

\begin{theorem}[Grothendieck 存在定理]
\label{coherent-theorem-grothendieck-existence}
\begin{reference}
\cite[III Theorem 5.1.5]{EGA}
\end{reference}
设 $A$ 为关于理想 $I$ 完备的 Noether 环。设 $X$ 为 $A$ 上分离的
有限型概形。则函子
(\ref{coherent-equation-completion-functor-proper-over-A})
$$
\textit{Coh}_{\text{support proper over }A}(\mathcal{O}_X)
\longrightarrow
\textit{Coh}_{\text{support proper over }A}(X, \mathcal{I})
$$
是等价。
\end{theorem}

\begin{proof}
将不再特别说明地使用引理 \ref{coherent-lemma-i-star-equivalence}
给出的范畴等价。在本证明中，对闭子概形
$Z \subset X$（它在 $A$ 上固有），若 $X$ 上的凝聚模被 $Z$ 的理想层零化，
或者等价地，它是 $Z$ 上凝聚模的直像，就称该模“支撑在 $Z$ 上”。
由命题 \ref{coherent-proposition-existence-proper}，已知凝聚模与支撑在
$Z$ 上的凝聚模系统之间的函子具有所述结论。因此，只须证明
$\textit{Coh}_{\text{support proper over }A}(\mathcal{O}_X)$ 的每个对象
以及 $\textit{Coh}_{\text{support proper over }A}(X, \mathcal{I})$
的每个对象，都支撑在某个闭子概形
$Z \subset X$ 上，该子概形在 $A$ 上固有。对
$\textit{Coh}_{\text{support proper over }A}(\mathcal{O}_X)$ 的对象，
这由定义成立。下面将依照命题 \ref{coherent-proposition-existence-proper}
的证明方法，对
$\textit{Coh}_{\text{support proper over }A}(X, \mathcal{I})$
的对象证明该断言。建议读者先阅读该命题的证明。

\medskip\noindent
考虑集合 $\Xi$，其元素为准凝聚理想层
$\mathcal{K} \subset \mathcal{O}_X$，并具有如下性质：对每个对象
$(\mathcal{F}_n)$（属于
$\textit{Coh}_{\text{support proper over }A}(X, \mathcal{I})$）
若它被 $\mathcal{K}$ 零化，
上述断言均成立。要证明 $(0)$ 属于 $\Xi$。
若不然，由于 $X$ 是 Noether 概形，存在极大准凝聚理想层
$\mathcal{K}$，它不属于 $\Xi$；见引理 \ref{coherent-lemma-acc-coherent}。
将 $X$ 替换为 $X$ 中与 $\mathcal{K}$ 对应的闭子概形后，
可以假设每个非零 $\mathcal{K}$ 均属于 $\Xi$。
设 $(\mathcal{F}_n)$ 为
$\textit{Coh}_{\text{support proper over }A}(X, \mathcal{I})$ 的对象。
将证明该对象支撑在某个闭子概形
$Z \subset X$ 上，该子概形在 $A$ 上固有，从而完成定理的证明。

\medskip\noindent
应用 Chow 引理（引理 \ref{coherent-lemma-chow-Noetherian}），得到固有满射
$f : Y \to X$，它在稠密开集 $U \subset X$ 上为同构，且 $Y$
在 $A$ 上 H-拟射影。选取开浸入 $j : Y \to Y'$，其中 $Y'$
在 $A$ 上射影；见《态射》引理
\ref{morphisms-lemma-H-quasi-projective-open-H-projective}。
注意
$$
\text{Supp}(f^*\mathcal{F}_n) = f^{-1}\text{Supp}(\mathcal{F}_n) =
f^{-1}\text{Supp}(\mathcal{F}_1)
$$
第一个等式来自《态射》引理
\ref{morphisms-lemma-support-finite-type}。由假设以及引理
\ref{coherent-lemma-functoriality-closed-proper-over-base} 的 (3)，
$f^{-1}\text{Supp}(\mathcal{F}_1)$ 在 $A$ 上固有。
因此，由引理 \ref{coherent-lemma-functoriality-closed-proper-over-base} 的 (1)，
$f^{-1}\text{Supp}(\mathcal{F}_1)$ 在 $j$ 下的像是 $Y'$ 中的闭集。
故由引理 \ref{coherent-lemma-pushforward-coherent-on-open}，
$\mathcal{F}'_n = j_*f^*\mathcal{F}_n$ 在 $Y'$ 上凝聚。
于是 $(\mathcal{F}_n')$ 是
$\textit{Coh}(Y', I\mathcal{O}_{Y'})$ 的对象。
由射影情形的 Grothendieck 存在定理
（引理 \ref{coherent-lemma-existence-projective}），存在凝聚
$\mathcal{O}_{Y'}$-模 $\mathcal{F}'$ 及同构
$(\mathcal{F}')^\wedge \cong (\mathcal{F}'_n)$
（在 $\textit{Coh}(Y', I\mathcal{O}_{Y'})$ 中）。
由于 $\mathcal{F}'/I\mathcal{F}' = \mathcal{F}'_1$，可知
$$
\text{Supp}(\mathcal{F}') \cap V(I\mathcal{O}_{Y'}) =
\text{Supp}(\mathcal{F}'_1) = j(f^{-1}\text{Supp}(\mathcal{F}_1))
$$
结构态射 $p' : Y' \to \Spec(A)$ 固有，故
$p'(\text{Supp}(\mathcal{F}') \setminus j(Y))$
是 $\Spec(A)$ 中的闭集。$\Spec(A)$ 的每个非空闭子集都包含
$V(I)$ 的一点，因为由《代数》引理
\ref{algebra-lemma-radical-completion}，$I$ 包含于 $A$ 的 Jacobson 根。
上述等式表明
$\text{Supp}(\mathcal{F}') \cap (p')^{-1}V(I) \subset j(Y)$，
从而 $\text{Supp}(\mathcal{F}') \subset j(Y)$。
因此，$\mathcal{F}'|_Y = j^*\mathcal{F}'$
支撑在某个闭子概形 $Z'$ 上，该子概形位于 $Y$ 中且在 $A$ 上固有，并且
$(\mathcal{F}'|_Y)^\wedge = (f^*\mathcal{F}_n)$。

\medskip\noindent
设 $\mathcal{K}$ 为截出约化补集 $X \setminus U$ 的准凝聚理想层。
由命题 \ref{coherent-proposition-proper-pushforward-coherent}，
$\mathcal{O}_X$-模 $\mathcal{H} = f_*(\mathcal{F}'|_Y)$ 凝聚；
由引理 \ref{coherent-lemma-inverse-systems-push-pull}，存在态射
$\alpha : (\mathcal{F}_n) \to \mathcal{H}^\wedge$
（在 $\textit{Coh}(X, \mathcal{I})$ 中），其核与余核被幂
$\mathcal{K}^t$ 零化，其中该幂来自 $\mathcal{K}$。得到正合列
$$
0 \to \Ker(\alpha) \to (\mathcal{F}_n) \to
\mathcal{H}^\wedge \to \Coker(\alpha) \to 0
$$
（在 $\textit{Coh}(X, \mathcal{I})$ 中）。
若 $Z_0 \subset X$ 是 $\mathcal{H}$ 的概形论支撑，则在集合意义下，
$Z_0 \subset f(Z')$ 是清楚的。于是，由引理
\ref{coherent-lemma-closed-closed-proper-over-base} 与引理
\ref{coherent-lemma-functoriality-closed-proper-over-base} 的 (2)，
$Z_0$ 在 $A$ 上固有。故 $\mathcal{H}^\wedge$
属于引理 \ref{coherent-lemma-systems-with-proper-support} 的 (2)
所定义的子范畴，从而尤其属于
$\textit{Coh}_{\text{support proper over }A}(X, \mathcal{I})$。
由引理 \ref{coherent-lemma-systems-with-proper-support} 的 (1)，
可知 $\Ker(\alpha)$ 与 $\Coker(\alpha)$ 属于
$\textit{Coh}_{\text{support proper over }A}(X, \mathcal{I})$。
由归纳假设，更确切地说，因为 $\mathcal{K}^t$ 属于 $\Xi$，
$\Ker(\alpha)$ 与 $\Coker(\alpha)$ 属于引理
\ref{coherent-lemma-systems-with-proper-support} 的 (2) 所定义的子范畴。
由于按该引理这是 Serre 子范畴，$(\mathcal{F}_n)$ 也属于其中，
这正是所要证明的。
\end{proof}

\begin{remark}[展开 Grothendieck 存在定理]
\label{coherent-remark-reformulate-existence-theorem}
设 $A$ 为关于理想 $I$ 完备的 Noether 环。记
$S = \Spec(A)$ 与 $S_n = \Spec(A/I^n)$。
设 $X \to S$ 为分离有限型态射。对 $n \geq 1$，令
$X_n = X \times_S S_n$。图示如下：
$$
\xymatrix{
X_1 \ar[r]_{i_1} \ar[d] & X_2 \ar[r]_{i_2} \ar[d] & X_3 \ar[r] \ar[d] &
\ldots & X \ar[d] \\
S_1 \ar[r] & S_2 \ar[r] & S_3 \ar[r] & \ldots & S
}
$$
在此情形下，考虑系统 $(\mathcal{F}_n, \varphi_n)$，其中
\begin{enumerate}
\item $\mathcal{F}_n$ 是凝聚 $\mathcal{O}_{X_n}$-模；
\item $\varphi_n : i_n^*\mathcal{F}_{n + 1} \to \mathcal{F}_n$
是同构；并且
\item $\text{Supp}(\mathcal{F}_1)$ 在 $S_1$ 上固有。
\end{enumerate}
定理 \ref{coherent-theorem-grothendieck-existence} 断言，完备化函子
$$
\begin{matrix}
\text{coherent }\mathcal{O}_X\text{-模 }\mathcal{F} \\
\text{with support proper over }A
\end{matrix}
\quad
\longrightarrow
\quad
\begin{matrix}
\text{systems }(\mathcal{F}_n) \\
\text{如上}
\end{matrix}
$$
是范畴等价。在 $X$ 在 $A$ 上固有这一特殊情形下，可以省略支撑条件。
\end{remark}

\section{Grothendieck 代数化定理}
\label{coherent-section-algebraization}

\noindent
第一个结论是用闭子概形和有限态射来表述 Grothendieck 存在定理。

\begin{lemma}
\label{coherent-lemma-algebraize-formal-closed-subscheme}
设 $A$ 为关于理想 $I$ 完备的 Noether 环。记
$S = \Spec(A)$ 与 $S_n = \Spec(A/I^n)$。
设 $X \to S$ 为分离有限型态射。对 $n \geq 1$，令
$X_n = X \times_S S_n$。假设给定交换图
$$
\xymatrix{
Z_1 \ar[r] \ar[d] & Z_2 \ar[r] \ar[d] & Z_3 \ar[r] \ar[d] & \ldots \\
X_1 \ar[r]^{i_1} & X_2 \ar[r]^{i_2} & X_3 \ar[r] & \ldots
}
$$
（诸方块均为笛卡尔方块），并假设：
\begin{enumerate}
\item $Z_1 \to X_1$ 是闭浸入；并且
\item $Z_1 \to S_1$ 固有。
\end{enumerate}
则存在概形的闭浸入 $Z \to X$，使得
$Z_n = Z \times_S S_n$。而且，$Z$ 在 $S$ 上固有。
\end{lemma}

\begin{proof}
以 $j_n : Z_n \to X_n$ 表示竖直态射。由于陈述中的方块是笛卡尔的，
$j_n$ 到 $X_1$ 的基变换就是 $j_1$。因此，《态射》引理
\ref{morphisms-lemma-check-closed-infinitesimally} 表明 $j_n$ 是闭浸入。
令 $\mathcal{F}_n = j_{n, *}\mathcal{O}_{Z_n}$，从而
$j_n^\sharp$ 是满射 $\mathcal{O}_{X_n} \to \mathcal{F}_n$。
再次利用方块的笛卡尔性可知，$\mathcal{F}_{n + 1}$ 到 $X_n$
的拉回就是 $\mathcal{F}_n$。故按评注
\ref{coherent-remark-reformulate-existence-theorem} 中的重新表述，
Grothendieck 存在定理告诉我们，存在映射
$\mathcal{O}_X \to \mathcal{F}$（在凝聚 $\mathcal{O}_X$-模之间），
其到 $X_n$ 的限制恢复
$\mathcal{O}_{X_n} \to \mathcal{F}_n$。而且，$\mathcal{F}$
的支撑在 $S$ 上固有。由于完备化函子正合（引理 \ref{coherent-lemma-exact}），
余核 $\mathcal{Q}$（对应于 $\mathcal{O}_X \to \mathcal{F}$）
具有零完备化。
由于 $\mathcal{F}$ 的支撑在 $S$ 上固有，$\mathcal{Q}$ 也如此，
这蕴含 $\mathcal{Q} = 0$；例如，这是因为
(\ref{coherent-equation-completion-functor-proper-over-A}) 按
Grothendieck 存在定理是等价。因此，有
$\mathcal{F} = \mathcal{O}_X/\mathcal{J}$，其中 $\mathcal{J}$
是某个准凝聚理想层。
令 $Z = V(\mathcal{J})$ 即完成证明。
\end{proof}

\noindent
在下述引理中，实际上只须假设 $Y_1 \to X_1$ 有限，因为这将蕴含
$Y_n \to X_n$ 也有限（见《态射进阶》引理
\ref{more-morphisms-lemma-thicken-property-morphisms-cartesian}）。

\begin{lemma}
\label{coherent-lemma-algebraize-formal-scheme-finite-over-proper}
设 $A$ 为关于理想 $I$ 完备的 Noether 环。记
$S = \Spec(A)$ 与 $S_n = \Spec(A/I^n)$。
设 $X \to S$ 为分离有限型态射。对 $n \geq 1$，令
$X_n = X \times_S S_n$。假设给定交换图
$$
\xymatrix{
Y_1 \ar[r] \ar[d] & Y_2 \ar[r] \ar[d] & Y_3 \ar[r] \ar[d] & \ldots \\
X_1 \ar[r]^{i_1} & X_2 \ar[r]^{i_2} & X_3 \ar[r] & \ldots
}
$$
（诸方块均为笛卡尔方块），并假设：
\begin{enumerate}
\item $Y_n \to X_n$ 是有限态射；并且
\item $Y_1 \to S_1$ 固有。
\end{enumerate}
则存在概形的有限态射 $Y \to X$，使得
$Y_n = Y \times_S S_n$。而且，$Y$ 在 $S$ 上固有。
\end{lemma}

\begin{proof}
以 $f_n : Y_n \to X_n$ 表示竖直态射。令
$\mathcal{F}_n = f_{n, *}\mathcal{O}_{Y_n}$。这是凝聚
$\mathcal{O}_{X_n}$-模，因为 $f_n$ 有限（引理
\ref{coherent-lemma-finite-pushforward-coherent}）。
利用方块的笛卡尔性可知，$\mathcal{F}_{n + 1}$ 到 $X_n$
的拉回就是 $\mathcal{F}_n$。故按评注
\ref{coherent-remark-reformulate-existence-theorem} 中的重新表述，
Grothendieck 存在定理告诉我们，存在凝聚 $\mathcal{O}_X$-模
$\mathcal{F}$，其到 $X_n$ 的限制恢复 $\mathcal{F}_n$。
而且，$\mathcal{F}$ 的支撑在 $S$ 上固有。
由于完备化函子全忠实（定理
\ref{coherent-theorem-grothendieck-existence}），诸乘法映射
$\mathcal{F}_n \otimes_{\mathcal{O}_{X_n}} \mathcal{F}_n \to
\mathcal{F}_n$ 相容地给出 $\mathcal{F}$ 上的代数结构。
令 $Y = \underline{\Spec}_X(\mathcal{F})$ 即完成证明。
\end{proof}

\begin{lemma}
\label{coherent-lemma-algebraize-morphism}
设 $A$ 为关于理想 $I$ 完备的 Noether 环。记
$S = \Spec(A)$ 与 $S_n = \Spec(A/I^n)$。设 $X$、$Y$ 为 $S$
上的概形。对 $n \geq 1$，令 $X_n = X \times_S S_n$ 及
$Y_n = Y \times_S S_n$。假设给定相容的交换图系统
$$
\xymatrix{
& & X_{n + 1} \ar[rd] \ar[rr]_{g_{n + 1}} & & Y_{n + 1} \ar[ld] \\
X_n \ar[rru] \ar[rd] \ar[rr]_{g_n} & & Y_n \ar[rru] \ar[ld] & S_{n + 1} \\
& S_n \ar[rru]
}
$$
并假设：
\begin{enumerate}
\item $X \to S$ 固有；并且
\item $Y \to S$ 分离且有限型。
\end{enumerate}
则存在唯一的概形态射 $g : X \to Y$（在 $S$ 上），使得
$g_n$ 是 $g$ 到 $S_n$ 的基变换。
\end{lemma}

\begin{proof}
态射 $(1, g_n) : X_n \to X_n \times_S Y_n$ 是闭浸入，因为
$Y_n \to S_n$ 分离（《概形》引理
\ref{schemes-lemma-section-immersion}）。因此，由引理
\ref{coherent-lemma-algebraize-formal-closed-subscheme}，存在闭子概形
$Z \subset X \times_S Y$，它在 $S$ 上固有，且到 $S_n$ 的基变换恢复
$X_n \subset X_n \times_S Y_n$。第一投影 $p : Z \to X$ 是固有态射
（因为 $Z$ 在 $S$ 上固有；见《态射》引理
\ref{morphisms-lemma-image-proper-scheme-closed}），其到 $S_n$ 的基变换
是同构，这对所有 $n$ 成立。特别地，由引理
\ref{coherent-lemma-proper-finite-fibre-finite-in-neighbourhood}，
$p : Z \to X$ 在 $X_0$ 的某个开邻域上有限。由于 $X$ 在 $S$ 上固有，
该开邻域就是整个 $X$，从而 $p : Z \to X$ 有限。
应用命题 \ref{coherent-proposition-existence-proper} 的等价可知
$p_*\mathcal{O}_Z = \mathcal{O}_X$，因为模 $I^n$ 后这一点对所有
$n$ 成立。因此 $p$ 是同构，并得到态射
$g$，它是复合 $X \cong Z \to Y$。唯一性的证明从略。
\end{proof}

\noindent
为了证明“抽象”的代数化定理，必须假设有一个丰沛可逆层；
没有这一假设，结论不成立。

\begin{theorem}[Grothendieck 代数化定理]
\label{coherent-theorem-algebraization}
设 $A$ 为关于理想 $I$ 完备的 Noether 环。令
$S = \Spec(A)$ 且 $S_n = \Spec(A/I^n)$。考虑交换图
$$
\xymatrix{
X_1 \ar[r]_{i_1} \ar[d] & X_2 \ar[r]_{i_2} \ar[d] & X_3 \ar[r] \ar[d] &
\ldots \\
S_1 \ar[r] & S_2 \ar[r] & S_3 \ar[r] & \ldots
}
$$
（诸方块均为笛卡尔方块）。假设给定 $(\mathcal{L}_n, \varphi_n)$，
其中每个 $\mathcal{L}_n$ 都是 $X_n$ 上的可逆层，并且
$\varphi_n : i_n^*\mathcal{L}_{n + 1} \to \mathcal{L}_n$
是同构。若
\begin{enumerate}
\item $X_1 \to S_1$ 固有；并且
\item $\mathcal{L}_1$ 在 $X_1$ 上丰沛，
\end{enumerate}
则存在概形的固有态射 $X \to S$、丰沛可逆
$\mathcal{O}_X$-模 $\mathcal{L}$，以及同构
$X_n \cong X \times_S S_n$ 与
$\mathcal{L}_n \cong \mathcal{L}|_{X_n}$，它们与态射
$i_n$ 和 $\varphi_n$ 相容。
\end{theorem}

\begin{proof}
由于图中的方块是笛卡尔的，且态射 $S_n \to S_{n + 1}$ 是闭浸入，
可知态射 $i_n$ 也是闭浸入。特别地，可以把
$X_m$ 看作 $X_n$ 的闭子概形（当 $m < n$ 时）。事实上，$X_m$ 是由准凝聚理想层
$I^m\mathcal{O}_{X_n}$ 截出的闭子概形。而且，概形
$X_1, X_2, X_3, \ldots$ 的底层拓扑空间都被等同，故可以（且确实）
把层 $\mathcal{O}_{X_n}$ 看作位于同一个底层拓扑空间上；
凝聚 $\mathcal{O}_{X_n}$-模亦然。令
$$
\mathcal{F}_n =
\Ker(\mathcal{O}_{X_{n + 1}} \to \mathcal{O}_{X_n})
$$
从而得到短正合列
$$
0 \to \mathcal{F}_n \to \mathcal{O}_{X_{n + 1}} \to \mathcal{O}_{X_n} \to 0
$$
由上可知 $\mathcal{F}_n = I^n\mathcal{O}_{X_{n + 1}}$。
于是 $\mathcal{F}_n$ 是 $X_{n + 1}$ 上被 $I$ 零化的凝聚层，
故可以（且确实）把它看作凝聚 $\mathcal{O}_{X_1}$-模。
注意，对 $m > n$，层
$$
I^n\mathcal{O}_{X_m}/I^{n + 1}\mathcal{O}_{X_m}
$$
同构地映到 $\mathcal{F}_n$，所用映射为
$\mathcal{O}_{X_m} \to \mathcal{O}_{X_{n + 1}}$。因此，给定
$n_1, n_2 \geq 0$，
可以选取 $m > n_1 + n_2$ 并考虑乘法映射
$$
I^{n_1}\mathcal{O}_{X_m} \times I^{n_2}\mathcal{O}_{X_m}
\longrightarrow
I^{n_1 + n_2}\mathcal{O}_{X_m} \to \mathcal{F}_{n_1 + n_2}
$$
它诱导 $\mathcal{O}_{X_1}$-双线性映射
$$
\mathcal{F}_{n_1} \times \mathcal{F}_{n_2} \longrightarrow
\mathcal{F}_{n_1 + n_2}
$$
继而定义分次 $\mathcal{O}_{X_1}$-代数结构，其底层分次模为
$\mathcal{F} = \bigoplus_{n \geq 0} \mathcal{F}_n$。

\medskip\noindent
令 $B = \bigoplus I^n/I^{n + 1}$；这是有限生成分次
$A/I$-代数。令 $\mathcal{B} = (X_1 \to S_1)^*\widetilde{B}$。
上述讨论给出分次 $\mathcal{O}_{X_1}$-代数的典范满射
$$
\mathcal{B} \longrightarrow \mathcal{F}
$$
特别地，$\mathcal{F}$ 是有限型准凝聚分次 $\mathcal{B}$-模。
由引理 \ref{coherent-lemma-graded-finiteness}，可找到整数 $d_0$，使得
$H^1(X_1, \mathcal{F} \otimes \mathcal{L}^{\otimes d}) = 0$
对所有 $d \geq d_0$ 成立。选取 $d \geq d_0$，使得存在截面
$s_{0, 1}, \ldots, s_{N, 1} \in \Gamma(X_1, \mathcal{L}_1^{\otimes d})$，
它们诱导浸入
$$
\psi_1 : X_1 \to \mathbf{P}^N_{S_1}
$$
（在 $S_1$ 上）；见《态射》引理
\ref{morphisms-lemma-finite-type-over-affine-ample-very-ample}。
由于 $X_1$ 在 $S_1$ 上固有，可知 $\psi_1$ 是闭浸入；见
《态射》引理 \ref{morphisms-lemma-image-proper-scheme-closed}
及《概形》引理 \ref{schemes-lemma-immersion-when-closed}。
将把 $\psi_1$“提升”为 $X_n$ 到 $\mathbf{P}^N$ 的相容闭浸入系统。

\medskip\noindent
把证明第一段中的短正合列与
$\mathcal{L}_{n + 1}^{\otimes d}$ 张量，得到短正合列
$$
0 \to \mathcal{F}_n \otimes \mathcal{L}_{n + 1}^{\otimes d} \to
\mathcal{L}_{n + 1}^{\otimes d} \to \mathcal{L}_{n + 1}^{\otimes d} \to 0
$$
使用同构 $\varphi_n$，得到同构
$\mathcal{L}_{n + 1} \otimes \mathcal{O}_{X_l} = \mathcal{L}_l$
（对 $l \leq n$）。因此上述序列变为
$$
0 \to \mathcal{F}_n \otimes \mathcal{L}_1^{\otimes d} \to
\mathcal{L}_{n + 1}^{\otimes d} \to \mathcal{L}_n^{\otimes d} \to 0
$$
$H^1(X, \mathcal{F}_n \otimes \mathcal{L}_1^{\otimes d})$ 的消失
使我们能够归纳地把
$s_{0, 1}, \ldots, s_{N, 1} \in \Gamma(X_1, \mathcal{L}_1^{\otimes d})$
提升为截面
$s_{0, n}, \ldots, s_{N, n} \in \Gamma(X_n, \mathcal{L}_n^{\otimes d})$。
从而得到交换图
$$
\xymatrix{
X_1 \ar[r]_{i_1} \ar[d]_{\psi_1} &
X_2 \ar[r]_{i_2} \ar[d]_{\psi_2} &
X_3 \ar[r] \ar[d]_{\psi_3} &
\ldots \\
\mathbf{P}^N_{S_1} \ar[r] &
\mathbf{P}^N_{S_2} \ar[r] &
\mathbf{P}^N_{S_3} \ar[r] & \ldots
}
$$
其中，按《构造》第 \ref{constructions-section-projective-space} 节的记号，
$\psi_n = \varphi_{(\mathcal{L}_n, (s_{0, n}, \ldots, s_{N, n}))}$。
由于定理陈述中的方块是笛卡尔的，上图中的方块也是笛卡尔的。
应用引理 \ref{coherent-lemma-algebraize-formal-closed-subscheme} 即得结论。
\end{proof}
